\input{preamble}

% OK, start here.
%
\begin{document}

\title{Compléments sur les morphismes d'espaces algébriques}


\maketitle

\phantomsection
\label{section-phantom}

\tableofcontents

\section{Introduction}
\label{section-introduction}

\noindent
Dans ce chapitre, nous poursuivons notre étude des propriétés des morphismes d'espaces algébriques.
Une référence fondamentale est \cite{Kn}.





\section{Conventions}
\label{section-conventions}

\noindent
Nous supposons une fois pour toutes que tous les schémas appartiennent à
un grand site fppf $\Sch_{fppf}$. De plus, tous les anneaux $A$ considérés
ont la propriété que $\Spec(A)$ est (isomorphe) à un
objet de ce grand site.

\medskip\noindent
Soit $S$ un schéma et soit $X$ un espace algébrique sur $S$.
Dans ce chapitre et les suivants, nous noterons $X \times_S X$
le produit de $X$ par lui-même (dans la catégorie des espaces algébriques
sur $S$), au lieu de $X \times X$.








\section{Morphismes radiciels}
\label{section-radicial}

\noindent
Il s'avère qu'un morphisme radiciel n'est pas la même chose qu'un morphisme
universellement injectif, contrairement à ce qui se passe pour les
morphismes de schémas. En fait, il est un peu plus fort.

\begin{definition}
\label{definition-radicial}
Soit $S$ un schéma. Soit $f : X \to Y$ un morphisme d'espaces algébriques
au-dessus de $S$. On dit que $f$ est {\it radiciel} si, pour tout morphisme
$\Spec(K) \to Y$ où $K$ est un corps, la réduction
$(\Spec(K) \times_Y X)_{red}$ est soit vide, soit
est représentable par le spectre d'une extension de corps purement inséparable de $K$.
\end{definition}

\begin{lemma}
\label{lemma-radicial-implies-universally-injective}
Un morphisme radiciel d'espaces algébriques est universellement injectif.
\end{lemma}

\begin{proof}
Soit $S$ un schéma. Soit $f : X \to Y$ un morphisme radiciel
d'espaces algébriques au-dessus de $S$.
Il est clair d'après la définition que, pour un morphisme
$\Spec(K) \to Y$, il existe au plus un relèvement de ce morphisme
en un morphisme vers $X$. Nous concluons donc que $f$ est universellement
injectif par
Morphismes d'espaces,
lemme \ref{spaces-morphisms-lemma-universally-injective}.
\end{proof}

\begin{example}
\label{example-universally-injective-not-radicial}
Il n'est plus vrai que « universellement injectif » soit équivalent à « radiciel ».
Par exemple, le morphisme
$$
X = [\Spec(\overline{\mathbf{Q}})/
\text{Gal}(\overline{\mathbf{Q}}/\mathbf{Q})]
\longrightarrow
S = \Spec(\mathbf{Q})
$$
d'Espaces,
exemple \ref{spaces-example-Qbar}
est universellement injectif, mais n'est pas radiciel au sens ci-dessus.
\end{example}

\noindent
Néanmoins, il arrive souvent que l'implication réciproque soit vraie.

\begin{lemma}
\label{lemma-when-universally-injective-radicial}
Soit $S$ un schéma. Soit $f : X \to Y$ un morphisme universellement injectif
d'espaces algébriques au-dessus de $S$.
\begin{enumerate}
\item Si $f$ est décent, alors $f$ est radiciel.
\item Si $f$ est quasi-séparé, alors $f$ est radiciel.
\item Si $f$ est localement séparé, alors $f$ est radiciel.
\end{enumerate}
\end{lemma}

\begin{proof}
Soit $\mathcal{P}$ une propriété des morphismes d'espaces algébriques
qui est stable par changement de base et composition et qui est satisfaite pour
les immersions fermées. Supposons que $f : X \to Y$ possède $\mathcal{P}$ et
soit universellement injectif. Alors, dans la situation de la
définition \ref{definition-radicial},
le morphisme $(\Spec(K) \times_Y X)_{red} \to \Spec(K)$
est universellement injectif et possède $\mathcal{P}$. Cela ramène le
problème de démontrer
$$
\mathcal{P} + \text{universellement injectif}
\Rightarrow
\text{radicial}
$$
au problème de démontrer que tout espace algébrique réduit non vide $X$
sur un corps, dont le morphisme structural $X \to \Spec(K)$ est universellement
injectif et possède $\mathcal{P}$, est représentable par le spectre d'un corps.
En effet, alors $X \to \Spec(K)$ sera un morphisme de schémas et
nous concluons par l'équivalence entre radiciel et universellement injectif pour
les morphismes de schémas, voir
Morphismes, lemme \ref{morphisms-lemma-universally-injective}.

\medskip\noindent
Démontrons (1). Supposons que $f$ soit décent et universellement injectif. D'après
Espaces décents,
les lemmes \ref{decent-spaces-lemma-base-change-relative-conditions},
\ref{decent-spaces-lemma-composition-relative-conditions} et
\ref{decent-spaces-lemma-properties-trivial-implications}
(pour voir qu'une immersion est décente), nous voyons que la discussion du
premier paragraphe s'applique.
Soit $X$ un espace algébrique réduit décent non vide
universellement injectif sur un corps $K$. En particulier, nous voyons que $|X|$
est un singleton. D'après
Espaces décents, lemme \ref{decent-spaces-lemma-when-field},
nous concluons que $X \cong \Spec(L)$ pour une extension
$K \subset L$, comme souhaité.

\medskip\noindent
Un morphisme quasi-séparé est décent, voir
Espaces décents,
lemme \ref{decent-spaces-lemma-properties-trivial-implications}.
Ainsi, (1) implique (2).

\medskip\noindent
Démontrons (3).
Rappelons que les axiomes de séparation sont stables par changement de base
et composition et que les immersions fermées sont séparées, voir
Morphismes d'espaces,
les lemmes \ref{spaces-morphisms-lemma-base-change-separated},
\ref{spaces-morphisms-lemma-composition-separated} et
\ref{spaces-morphisms-lemma-immersions-monomorphisms}.
La discussion du premier paragraphe de la démonstration s'applique donc.
Soit $X$ un espace algébrique réduit, universellement injectif et
localement séparé sur un corps $K$.
En particulier, $|X|$ est un singleton, donc $X$ est quasi-compact, voir
Propriétés des espaces, lemme \ref{spaces-properties-lemma-quasi-compact-space}.
Nous pouvons trouver un morphisme \'etale surjectif $U \to X$ avec $U$ affine, voir
Propriétés des espaces,
lemme \ref{spaces-properties-lemma-quasi-compact-affine-cover}.
Considérons le morphisme de schémas
$$
j :
U \times_X U
\longrightarrow
U \times_{\Spec(K)} U
$$
Comme $X \to \Spec(K)$ est universellement injectif, $j$ est surjectif,
et comme $X \to \Spec(K)$ est localement séparé, $j$ est une immersion.
Une immersion surjective est une immersion fermée, voir
Schémas, lemme \ref{schemes-lemma-immersion-when-closed}.
Ainsi $R = U \times_X U$ est affine comme sous-schéma fermé d'un schéma affine.
En particulier, $R$ est quasi-compact.
Il s'ensuit que $X = U/R$ est quasi-séparé, et le résultat découle de (2).
\end{proof}

\begin{remark}
\label{remark-weakly-radicial}
Soit $X \to Y$ un morphisme d'espaces algébriques.
Pour certaines applications (des morphismes radiciels),
il suffit d'exiger que, pour tout
$\Spec(K) \to Y$ où $K$ est un corps,
\begin{enumerate}
\item l'espace $|\Spec(K) \times_Y X|$ soit un singleton,
\item il existe un monomorphisme
$\Spec(L) \to \Spec(K) \times_Y X$, et
\item $K \subset L$ soit purement inséparable.
\end{enumerate}
Si nécessaire, nous appellerons ultérieurement un tel morphisme {\it faiblement radiciel}.
Par exemple, si $X \to Y$ est un morphisme faiblement radiciel surjectif,
alors $X(k) \to Y(k)$ est surjectif pour tout corps algébriquement clos $k$.
Remarquons que le changement de base
$X_{\overline{\mathbf{Q}}} \to \Spec(\overline{\mathbf{Q}})$
du morphisme de
l'exemple \ref{example-universally-injective-not-radicial}
est faiblement radiciel, mais pas radiciel. L'analogue du
lemme \ref{lemma-when-universally-injective-radicial}
est que, si $X \to Y$ possède la propriété ($\beta$) et est universellement
injectif, alors il est faiblement radiciel (démonstration omise).
\end{remark}

\begin{lemma}
\label{lemma-check-universally-injective}
Soit $S$ un schéma. Soit $f : X \to Y$ un morphisme d'espaces algébriques
au-dessus de $S$. Supposons que
\begin{enumerate}
\item $f$ soit localement de type fini,
\item pour tout morphisme \'etale $V \to Y$, l'application $|X \times_Y V| \to |V|$
soit injective.
\end{enumerate}
Alors $f$ est universellement injectif.
\end{lemma}

\begin{proof}
La question est locale pour la topologie \'etale sur $Y$, d'après
Morphismes d'espaces, lemme
\ref{spaces-morphisms-lemma-universally-injective-local}.
Nous pouvons donc supposer que $Y$ est un schéma.
Alors $Y$ est en particulier décent et, d'après Espaces décents, lemme
\ref{decent-spaces-lemma-conditions-on-point-in-fibre-and-qf}
nous voyons que $f$ est localement quasi-fini.
Soit $y \in Y$ un point et soit $X_y$ la fibre schématique.
Supposons que $X_y$ ne soit pas vide. D'après Espaces sur les corps, lemme
\ref{spaces-over-fields-lemma-locally-quasi-finite-over-field}
nous voyons que $X_y$ est un schéma localement quasi-fini sur
$\kappa(y)$. Comme $|X_y| \subset |X|$ est la fibre de $|X| \to |Y|$
au-dessus de $y$, nous voyons que $X_y$ possède un unique point $x$. Il en est de même
pour $X_y \times_{\Spec(\kappa(y))} \Spec(k)$ pour toute
extension finie séparable $k/\kappa(y)$,
car nous pouvons réaliser $k$ comme corps résiduel en un point
au-dessus de $y$ dans un schéma \'etale au-dessus de $Y$,
voir Compléments sur les morphismes, lemme
\ref{more-morphisms-lemma-realize-prescribed-residue-field-extension-etale}.
Ainsi $X_y$ est géométriquement connexe, voir
Variétés, lemme \ref{varieties-lemma-characterize-geometrically-disconnected}.
Cela implique que l'extension finie $\kappa(x)/\kappa(y)$
est purement inséparable.

\medskip\noindent
Nous concluons (dans le cas où $Y$ est un schéma)
que, pour tout $y \in Y$, la fibre $X_y$ est soit vide,
soit $(X_y)_{red} = \Spec(\kappa(x))$ avec
$\kappa(y) \subset \kappa(x)$ purement inséparable.
Ainsi $f$ est radiciel (certains détails sont omis), donc universellement injectif par
le lemme \ref{lemma-radicial-implies-universally-injective}.
\end{proof}




\section{Monomorphismes}
\label{section-monomorphisms}

\noindent
Cette section fait suite à
Morphismes d'espaces, section \ref{spaces-morphisms-section-monomorphisms}.
Nous voudrions savoir si tout monomorphisme d'espaces algébriques est représentable.
Si vous pouvez démontrer que c'est vrai ou fournir un
contre-exemple, veuillez envoyer un courriel à
\href{mailto:stacks.project@gmail.com}{stacks.project@gmail.com}.
Pour le moment, on connaît les cas suivants
\begin{enumerate}
\item pour les monomorphismes qui sont localement de type fini
(plus généralement, tout morphisme séparé et localement quasi-fini
est représentable par Morphismes d'espaces, lemme
\ref{spaces-morphisms-lemma-locally-quasi-finite-separated-representable}
et un monomorphisme localement de type fini est
localement quasi-fini par Morphismes d'espaces, lemme
\ref{spaces-morphisms-lemma-monomorphism-loc-finite-type-loc-quasi-finite}),
\item si le but est une union disjointe de spectres d'anneaux locaux de dimension zéro
(Espaces décents, lemme
\ref{decent-spaces-lemma-monomorphism-toward-disjoint-union-dim-0-rings}), et
\item pour les monomorphismes plats (voir ci-dessous).
\end{enumerate}

\begin{lemma}[David Rydh]
\label{lemma-flat-case}
Un monomorphisme plat d'espaces algébriques est représentable par des schémas.
\end{lemma}

\begin{proof}
Soit $f : X \to Y$ un monomorphisme plat d'espaces algébriques.
Pour démontrer que $f$ est représentable, il faut montrer que
$X \times_Y V$ est un schéma pour tout schéma $V$ muni d'un morphisme vers $Y$.
Comme la propriété d'être un schéma est locale (Propriétés des espaces,
lemme \ref{spaces-properties-lemma-subscheme}), nous pouvons
supposer que $V$ est affine. Nous pouvons donc supposer que $Y = \Spec(B)$
est un schéma affine. Ensuite, nous pouvons supposer que $X$ est quasi-compact
en remplaçant $X$ par un ouvert quasi-compact. L'espace $X$ est
séparé puisque $X \to X \times_{\Spec(B)} X$ est un isomorphisme.
En appliquant Limites d'espaces, lemme \ref{spaces-limits-lemma-enough-local},
nous nous ramenons au cas où $B$ est local, où $X \to \Spec(B)$ est un
monomorphisme plat et où
il existe un point $x \in X$ s'envoyant sur le point fermé de $\Spec(B)$.
Alors $X \to \Spec(B)$ est surjectif puisque les générisations
se relèvent le long des morphismes plats d'espaces algébriques séparés, voir
Espaces décents, lemme \ref{decent-spaces-lemma-generalizations-lift-flat}.
Ainsi, $\{X \to \Spec(B)\}$ est un recouvrement fpqc.
Alors $X \to \Spec(B)$ est un morphisme qui devient un isomorphisme
après changement de base par $X \to \Spec(B)$. Il est donc un isomorphisme par
descente fpqc, voir Descente sur les espaces, lemme
\ref{spaces-descent-lemma-descending-property-isomorphism}.
\end{proof}

\noindent
Le résultat suivant est (en un certain sens) une variante du lemme ci-dessus.

\begin{lemma}
\label{lemma-ui-case}
Soit $S$ un schéma. Soit $f : X \to Y$ un monomorphisme quasi-compact
d'espaces algébriques tel que, pour tout $T \to Y$, l'application
$$
\mathcal{O}_T \to f_{T,*}\mathcal{O}_{X \times_Y T}
$$
soit injective. Alors $f$ est un isomorphisme (et donc représentable par des schémas).
\end{lemma}

\begin{proof}
La question est locale pour la topologie \'etale sur $Y$; nous pouvons donc supposer que $Y = \Spec(A)$
est affine. Alors $X$ est quasi-compact et nous pouvons choisir un schéma affine
$U = \Spec(B)$ et un morphisme \'etale surjectif $U \to X$
(Propriétés des espaces, lemme
\ref{spaces-properties-lemma-quasi-compact-affine-cover}).
Remarquons que $U \times_X U = \Spec(B \otimes_A B)$. Ainsi, la catégorie des
faisceaux quasi-cohérents de $\mathcal{O}_X$-modules est équivalente à la
catégorie $DD_{B/A}$ des données de descente sur les modules pour $A \to B$.
Voir Propriétés des espaces, proposition
\ref{spaces-properties-proposition-quasi-coherent},
Descente, définition \ref{descent-definition-descent-datum-modules}, et
Descente, sous-section \ref{descent-subsection-descent-modules-morphisms}.
D'autre part,
$$
A \to B
$$
est une application d'anneaux universellement injective. En effet, étant donné un
$A$-module $M$, nous voyons que $A \oplus M \to B \otimes_A (A \oplus M)$
est injective par l'hypothèse du lemme. Ainsi,
$DD_{B/A}$ est équivalente à la catégorie des $A$-modules par
Descente, théorème \ref{descent-theorem-descent}. Ainsi, le tiré en arrière le long de
$f : X \to \Spec(A)$ détermine une équivalence de catégories des
modules quasi-cohérents. En particulier, $f^*$ est exact sur les
modules quasi-cohérents et nous voyons que $f$ est plat
(un petit détail est omis). De plus, il est clair que $f$ est surjectif
(par exemple parce que $\Spec(B) \to \Spec(A)$ est surjectif).
Ainsi, $\{X \to \Spec(A)\}$ est un recouvrement fpqc.
Alors $X \to \Spec(A)$ est un morphisme qui devient un isomorphisme
après changement de base par $X \to \Spec(A)$. Il est donc un isomorphisme par
descente fpqc, voir Descente sur les espaces, lemme
\ref{spaces-descent-lemma-descending-property-isomorphism}.
\end{proof}

\begin{lemma}
\label{lemma-flat-surjective-monomorphism}
Un monomorphisme quasi-compact, plat et surjectif d'espaces algébriques
est un isomorphisme.
\end{lemma}

\begin{proof}
Un tel morphisme satisfait les hypothèses du lemme \ref{lemma-ui-case}.
\end{proof}







\section{Faisceau conormal d'une immersion}
\label{section-conormal-sheaf}

\noindent
Soit $S$ un schéma. Soit $i : Z \to X$ une immersion fermée d'espaces algébriques
au-dessus de $S$. Soit $\mathcal{I} \subset \mathcal{O}_X$ le
faisceau quasi-cohérent d'idéaux correspondant, voir
Morphismes d'espaces,
lemme \ref{spaces-morphisms-lemma-closed-immersion-ideals}.
Considérons la suite exacte courte
$$
0 \to \mathcal{I}^2 \to \mathcal{I} \to \mathcal{I}/\mathcal{I}^2 \to 0
$$
de faisceaux quasi-cohérents sur $X$. Comme le faisceau $\mathcal{I}/\mathcal{I}^2$
est annulé par $\mathcal{I}$, il correspond à un faisceau sur $Z$ par
Morphismes d'espaces, lemme \ref{spaces-morphisms-lemma-i-star-equivalence}.
Ce $\mathcal{O}_Z$-module quasi-cohérent est le
{\it faisceau conormal de $Z$ dans $X$} et est souvent noté
$\mathcal{I}/\mathcal{I}^2$, par l'abus de notation mentionné dans
Morphismes d'espaces,
section \ref{spaces-morphisms-section-closed-immersions-quasi-coherent}.

\medskip\noindent
Dans le cas où $i : Z \to X$ est une immersion (localement fermée), nous définissons le
faisceau conormal de $i$ comme le faisceau conormal de l'immersion fermée
$i : Z \to X \setminus \partial Z$, voir
Morphismes d'espaces, remarque \ref{spaces-morphisms-remark-immersion}.
Il est souvent noté
$\mathcal{I}/\mathcal{I}^2$ où $\mathcal{I}$ est le faisceau d'idéaux
de l'immersion fermée $i : Z \to X \setminus \partial Z$.

\begin{definition}
\label{definition-conormal-sheaf}
Soit $i : Z \to X$ une immersion. Le {\it faisceau conormal
$\mathcal{C}_{Z/X}$ de $Z$ dans $X$} ou le {\it faisceau conormal de $i$}
est le $\mathcal{O}_Z$-module quasi-cohérent $\mathcal{I}/\mathcal{I}^2$
décrit ci-dessus.
\end{definition}

\noindent
Dans \cite[IV, définition 16.1.2]{EGA}, ce faisceau est noté
$\mathcal{N}_{Z/X}$. Nous ne suivrons pas cette convention, car nous
souhaitons réserver la notation $\mathcal{N}_{Z/X}$
au {\it faisceau normal de l'immersion}. Il est défini par
$$
\mathcal{N}_{Z/X} =
\SheafHom_{\mathcal{O}_Z}(\mathcal{C}_{Z/X}, \mathcal{O}_Z) =
\SheafHom_{\mathcal{O}_Z}(\mathcal{I}/\mathcal{I}^2, \mathcal{O}_Z)
$$
pourvu que le faisceau conormal soit de présentation finie (sinon le
faisceau normal peut même ne pas être quasi-cohérent). Nous reviendrons sur le
faisceau normal ultérieurement (insérer ici une référence future).

\begin{lemma}
\label{lemma-etale-conormal}
Soit $S$ un schéma. Soit $i : Z \to X$ une immersion.
Soit $\varphi : U \to X$ un morphisme \'etale où $U$ est un schéma.
Posons $Z_U = U \times_X Z$, qui est un sous-schéma localement fermé de $U$.
Alors
$$
\mathcal{C}_{Z/X}|_{Z_U} = \mathcal{C}_{Z_U/U}
$$
canoniquement et fonctoriellement en $U$.
\end{lemma}

\begin{proof}
Soit $T \subset X$ un sous-espace fermé tel que $i$ définisse une
immersion fermée dans $X \setminus T$.
Soit $\mathcal{I}$ le faisceau quasi-cohérent d'idéaux sur
$X \setminus T$ définissant $Z$. Le lemme affirme simplement que
$\mathcal{I}|_{U \setminus \varphi^{-1}(T)}$ est le faisceau d'idéaux de
l'immersion $Z_U \to U \setminus \varphi^{-1}(T)$.
C'est clair d'après la construction de $\mathcal{I}$ dans
Morphismes d'espaces, lemme \ref{spaces-morphisms-lemma-closed-immersion-ideals}.
\end{proof}

\begin{lemma}
\label{lemma-conormal-functorial}
Soit $S$ un schéma. Soit
$$
\xymatrix{
Z \ar[r]_i \ar[d]_f & X \ar[d]^g \\
Z' \ar[r]^{i'} & X'
}
$$
un diagramme commutatif d'espaces algébriques au-dessus de $S$.
Supposons que $i$ et $i'$ soient des immersions. Il existe une application canonique
de $\mathcal{O}_Z$-modules
$$
f^*\mathcal{C}_{Z'/X'}
\longrightarrow
\mathcal{C}_{Z/X}
$$
\end{lemma}

\begin{proof}
D'abord, trouvons des sous-espaces ouverts $U' \subset X'$ et $U \subset X$ tels que
$g(U) \subset U'$ et tels que $i(Z) \subset U$ et $i(Z') \subset U'$
soient fermés (l'existence est omise). En remplaçant $X$ par $U$ et $X'$ par
$U'$, nous pouvons supposer que $i$ et $i'$ sont des immersions fermées.
Soient $\mathcal{I}' \subset \mathcal{O}_{X'}$ et
$\mathcal{I} \subset \mathcal{O}_X$ les faisceaux quasi-cohérents d'idéaux
associés à $i'$ et $i$, voir
Morphismes d'espaces, lemme \ref{spaces-morphisms-lemma-closed-immersion-ideals}.
Considérons la composée
$$
g^{-1}\mathcal{I}' \to g^{-1}\mathcal{O}_{X'}
\xrightarrow{g^\sharp} \mathcal{O}_X \to
\mathcal{O}_X/\mathcal{I} = i_*\mathcal{O}_Z
$$
Comme $g(i(Z)) \subset Z'$, nous concluons que cette composée est nulle (voir
l'énoncé sur les factorisations dans
Morphismes d'espaces,
lemme \ref{spaces-morphisms-lemma-closed-immersion-ideals}).
Nous obtenons ainsi un diagramme commutatif
$$
\xymatrix{
0 \ar[r] &
\mathcal{I} \ar[r] &
\mathcal{O}_X \ar[r] &
i_*\mathcal{O}_Z \ar[r] &
0 \\
0 \ar[r] &
g^{-1}\mathcal{I}' \ar[r] \ar[u] &
g^{-1}\mathcal{O}_{X'} \ar[r] \ar[u] &
g^{-1}i'_*\mathcal{O}_{Z'} \ar[r] \ar[u] &
0
}
$$
La ligne inférieure est exacte puisque $g^{-1}$ est un foncteur exact.
Par exactitude, nous voyons aussi que
$(g^{-1}\mathcal{I}')^2 = g^{-1}((\mathcal{I}')^2)$.
Le diagramme induit donc une application
$g^{-1}(\mathcal{I}'/(\mathcal{I}')^2) \to \mathcal{I}/\mathcal{I}^2$.
En tirant en arrière (par exemple à l'aide de $i^{-1}$) vers $Z$, nous obtenons
$i^{-1}g^{-1}(\mathcal{I}'/(\mathcal{I}')^2) \to \mathcal{C}_{Z/X}$.
Comme $i^{-1}g^{-1} = f^{-1}(i')^{-1}$, cela donne une application
$f^{-1}\mathcal{C}_{Z'/X'} \to \mathcal{C}_{Z/X}$, qui induit
l'application cherchée.
\end{proof}

\begin{lemma}
\label{lemma-conormal-functorial-more}
Soit $S$ un schéma. Le faisceau conormal de la
définition \ref{definition-conormal-sheaf} et sa fonctorialité du
lemme \ref{lemma-conormal-functorial}
satisfont les propriétés suivantes :
\begin{enumerate}
\item Si $Z \to X$ est une immersion de schémas au-dessus de $S$, alors le faisceau conormal
coïncide avec celui de
Morphismes, définition \ref{morphisms-definition-conormal-sheaf}.
\item Si, dans le
lemme \ref{lemma-conormal-functorial},
tous les espaces sont des schémas, alors l'application
$f^*\mathcal{C}_{Z'/X'} \to \mathcal{C}_{Z/X}$ est la même
que celle construite dans
Morphismes, lemme \ref{morphisms-lemma-conormal-functorial}.
\item Étant donné un diagramme commutatif
$$
\xymatrix{
Z \ar[r]_i \ar[d]_f & X \ar[d]^g \\
Z' \ar[r]^{i'} \ar[d]_{f'} & X' \ar[d]^{g'} \\
Z'' \ar[r]^{i''} & X''
}
$$
alors l'application $(f' \circ f)^*\mathcal{C}_{Z''/X''} \to \mathcal{C}_{Z/X}$
est la même que la composée de
$f^*\mathcal{C}_{Z'/X'} \to \mathcal{C}_{Z/X}$
avec le tiré en arrière par $f$ de
$(f')^*\mathcal{C}_{Z''/X''} \to \mathcal{C}_{Z'/X'}$.
\end{enumerate}
\end{lemma}

\begin{proof}
Démonstration omise. Remarquons que la partie (1) est un cas particulier du
lemme \ref{lemma-etale-conormal}.
\end{proof}

\begin{lemma}
\label{lemma-conormal-functorial-flat}
Soit $S$ un schéma. Soit
$$
\xymatrix{
Z \ar[r]_i \ar[d]_f & X \ar[d]^g \\
Z' \ar[r]^{i'} & X'
}
$$
un diagramme de produit fibré d'espaces algébriques au-dessus de $S$. Supposons que
$i$ et $i'$ soient des immersions. Alors l'application canonique
$f^*\mathcal{C}_{Z'/X'} \to \mathcal{C}_{Z/X}$ du
lemme \ref{lemma-conormal-functorial}
est surjective. Si $g$ est plat, c'est un isomorphisme.
\end{lemma}

\begin{proof}
Choisissons un diagramme commutatif
$$
\xymatrix{
U \ar[r] \ar[d] & X \ar[d] \\
U' \ar[r] & X'
}
$$
où $U$ et $U'$ sont des schémas et où les flèches horizontales sont surjectives
et \'etales, voir
Espaces, lemme \ref{spaces-lemma-lift-morphism-presentations}.
En utilisant alors
les lemmes \ref{lemma-etale-conormal} et \ref{lemma-conormal-functorial-more},
nous voyons que la question se ramène au cas d'un morphisme de schémas.
Dans le cas des schémas, il s'agit du
lemme \ref{morphisms-lemma-conormal-functorial-flat} de Morphismes.
\end{proof}

\begin{lemma}
\label{lemma-transitivity-conormal}
Soit $S$ un schéma.
Soient $Z \to Y \to X$ des immersions d'espaces algébriques.
Alors il existe une suite exacte canonique
$$
i^*\mathcal{C}_{Y/X} \to
\mathcal{C}_{Z/X} \to
\mathcal{C}_{Z/Y} \to 0
$$
où les applications proviennent du
lemme \ref{lemma-conormal-functorial}
et où $i : Z \to Y$ est le premier morphisme.
\end{lemma}

\begin{proof}
Soit $U$ un schéma et soit $U \to X$ un morphisme \'etale surjectif. Par
les lemmes \ref{lemma-etale-conormal} et \ref{lemma-conormal-functorial-more},
l'exactitude de la suite se traduit immédiatement par
l'exactitude de la suite correspondante pour les immersions de schémas
$Z \times_X U \to Y \times_X U \to U$. Le lemme découle donc du
lemme \ref{morphisms-lemma-transitivity-conormal} de Morphismes.
\end{proof}








\section{Cône normal d'une immersion}
\label{section-normal-cone}

\noindent
Soit $S$ un schéma. Soit $i : Z \to X$ une immersion fermée d'espaces algébriques
au-dessus de $S$. Soit $\mathcal{I} \subset \mathcal{O}_X$ le
faisceau quasi-cohérent d'idéaux correspondant, voir
Morphismes d'espaces, lemme
\ref{spaces-morphisms-lemma-closed-immersion-ideals}.
Considérons le faisceau quasi-cohérent d'algèbres graduées de $\mathcal{O}_X$
$\bigoplus_{n \geq 0} \mathcal{I}^n/\mathcal{I}^{n + 1}$.
Comme les faisceaux $\mathcal{I}^n/\mathcal{I}^{n + 1}$
sont tous annulés par $\mathcal{I}$, cette algèbre graduée
correspond à un faisceau quasi-cohérent d'algèbres graduées de $\mathcal{O}_Z$ par
Morphismes d'espaces, lemme \ref{spaces-morphisms-lemma-i-star-equivalence}.
Cette algèbre graduée quasi-cohérente de $\mathcal{O}_Z$ est appelée
{\it algèbre conormale de $Z$ dans $X$} et est souvent simplement notée
$\bigoplus_{n \geq 0} \mathcal{I}^n/\mathcal{I}^{n + 1}$
par l'abus de notation mentionné dans
Morphismes d'espaces, section
\ref{spaces-morphisms-section-closed-immersions-quasi-coherent}.

\medskip\noindent
Dans le cas où $i : Z \to X$ est une immersion (localement fermée), nous définissons l'algèbre conormale
de $i$ comme l'algèbre conormale de l'immersion fermée
$i : Z \to X \setminus \partial Z$, voir Morphismes d'espaces, remarque
\ref{spaces-morphisms-remark-immersion}.
Elle est souvent notée
$\bigoplus_{n \geq 0} \mathcal{I}^n/\mathcal{I}^{n + 1}$
où $\mathcal{I}$ est le faisceau d'idéaux
de l'immersion fermée $i : Z \to X \setminus \partial Z$.

\begin{definition}
\label{definition-conormal-algebra}
Soit $i : Z \to X$ une immersion. L'{\it algèbre conormale
$\mathcal{C}_{Z/X, *}$ de $Z$ dans $X$} ou l'{\it algèbre conormale de $i$}
est le faisceau quasi-cohérent d'algèbres graduées de $\mathcal{O}_Z$
$\bigoplus_{n \geq 0} \mathcal{I}^n/\mathcal{I}^{n + 1}$ décrit ci-dessus.
\end{definition}

\noindent
Ainsi $\mathcal{C}_{Z/X, 1} = \mathcal{C}_{Z/X}$ est le faisceau conormal
de l'immersion. De plus, $\mathcal{C}_{Z/X, 0} = \mathcal{O}_Z$ et
$\mathcal{C}_{Z/X, n}$ est un $\mathcal{O}_Z$-module quasi-cohérent
caractérisé par la propriété
\begin{equation}
\label{equation-conormal-in-degree-n}
i_*\mathcal{C}_{Z/X, n} = \mathcal{I}^n/\mathcal{I}^{n + 1}
\end{equation}
où $i : Z \to X \setminus \partial Z$ et $\mathcal{I}$ est le faisceau
d'idéaux de $i$ comme ci-dessus. Enfin, remarquons qu'il existe une application canonique surjective
\begin{equation}
\label{equation-conormal-algebra-quotient}
\text{Sym}^*(\mathcal{C}_{Z/X}) \longrightarrow \mathcal{C}_{Z/X, *}
\end{equation}
d'algèbres graduées quasi-cohérentes de $\mathcal{O}_Z$, qui est un isomorphisme
en degrés $0$ et $1$.

\begin{lemma}
\label{lemma-etale-conormal-algebra}
Soit $S$ un schéma. Soit $i : Z \to X$ une immersion d'espaces algébriques
au-dessus de $S$. Soit $\varphi : U \to X$ un morphisme \'etale où $U$ est un
schéma. Posons $Z_U = U \times_X Z$, qui est un sous-schéma localement fermé de $U$.
Alors
$$
\mathcal{C}_{Z/X, *}|_{Z_U} = \mathcal{C}_{Z_U/U, *}
$$
canoniquement et fonctoriellement en $U$.
\end{lemma}

\begin{proof}
Soit $T \subset X$ un sous-espace fermé tel que $i$ définisse une immersion
fermée dans $X \setminus T$. Soit $\mathcal{I}$ le faisceau quasi-cohérent
d'idéaux sur $X \setminus T$ définissant $Z$. Le lemme découle
du fait que
$\mathcal{I}|_{U \setminus \varphi^{-1}(T)}$ est le faisceau d'idéaux de
l'immersion $Z_U \to U \setminus \varphi^{-1}(T)$.
C'est clair d'après la construction de $\mathcal{I}$ dans
Morphismes d'espaces, lemme
\ref{spaces-morphisms-lemma-closed-immersion-ideals}.
\end{proof}

\begin{lemma}
\label{lemma-conormal-algebra-functorial}
Soit $S$ un schéma. Soit
$$
\xymatrix{
Z \ar[r]_i \ar[d]_f & X \ar[d]^g \\
Z' \ar[r]^{i'} & X'
}
$$
un diagramme commutatif d'espaces algébriques au-dessus de $S$.
Supposons que $i$ et $i'$ soient des immersions. Il existe une application canonique
d'algèbres graduées de $\mathcal{O}_Z$
$$
f^*\mathcal{C}_{Z'/X', *}
\longrightarrow
\mathcal{C}_{Z/X, *}
$$
\end{lemma}

\begin{proof}
Commençons par trouver des sous-espaces ouverts $U' \subset X'$ et $U \subset X$ tels que
$g(U) \subset U'$ et tels que $i(Z) \subset U$ et $i(Z') \subset U'$
soient fermés (l'existence est omise). En remplaçant $X$ par $U$ et $X'$ par
$U'$, nous pouvons supposer que $i$ et $i'$ sont des immersions fermées.
Soient $\mathcal{I}' \subset \mathcal{O}_{X'}$ et
$\mathcal{I} \subset \mathcal{O}_X$ les faisceaux quasi-cohérents d'idéaux
associés à $i'$ et $i$, voir
Morphismes d'espaces, lemme \ref{spaces-morphisms-lemma-closed-immersion-ideals}.
Considérons la composée
$$
g^{-1}\mathcal{I}' \to g^{-1}\mathcal{O}_{X'}
\xrightarrow{g^\sharp} \mathcal{O}_X \to
\mathcal{O}_X/\mathcal{I} = i_*\mathcal{O}_Z
$$
Comme $g(i(Z)) \subset Z'$, nous concluons que cette composée est nulle (voir
l'énoncé sur les factorisations dans
Morphismes d'espaces,
lemme \ref{spaces-morphisms-lemma-closed-immersion-ideals}).
Nous obtenons ainsi un diagramme commutatif
$$
\xymatrix{
0 \ar[r] &
\mathcal{I} \ar[r] &
\mathcal{O}_X \ar[r] &
i_*\mathcal{O}_Z \ar[r] &
0 \\
0 \ar[r] &
g^{-1}\mathcal{I}' \ar[r] \ar[u] &
g^{-1}\mathcal{O}_{X'} \ar[r] \ar[u] &
g^{-1}i'_*\mathcal{O}_{Z'} \ar[r] \ar[u] &
0
}
$$
La ligne inférieure est exacte puisque $g^{-1}$ est un foncteur exact.
Par exactitude, nous voyons aussi que
$(g^{-1}\mathcal{I}')^n = g^{-1}((\mathcal{I}')^n)$ pour tout $n \geq 1$.
Le diagramme induit donc une application
$g^{-1}((\mathcal{I}')^n/(\mathcal{I}')^{n + 1}) \to
\mathcal{I}^n/\mathcal{I}^{n + 1}$.
En tirant en arrière (par exemple à l'aide de $i^{-1}$) vers $Z$, nous obtenons
$i^{-1}g^{-1}((\mathcal{I}')^n/(\mathcal{I}')^{n + 1}) \to
\mathcal{C}_{Z/X, n}$.
Comme $i^{-1}g^{-1} = f^{-1}(i')^{-1}$, cela donne des applications
$f^{-1}\mathcal{C}_{Z'/X', n} \to \mathcal{C}_{Z/X, n}$, qui induisent
l'application cherchée.
\end{proof}

\begin{lemma}
\label{lemma-conormal-algebra-functorial-flat}
Soit $S$ un schéma. Soit
$$
\xymatrix{
Z \ar[r]_i \ar[d]_f & X \ar[d]^g \\
Z' \ar[r]^{i'} & X'
}
$$
un carré cartésien d'espaces algébriques au-dessus de $S$ avec
$i$ et $i'$ immersions. Alors l'application canonique
$f^*\mathcal{C}_{Z'/X', *} \to \mathcal{C}_{Z/X, *}$ du
lemme \ref{lemma-conormal-algebra-functorial}
est surjective. Si $g$ est plat, c'est un isomorphisme.
\end{lemma}

\begin{proof}
Nous pouvons vérifier l'assertion après localisation \'etale de $X'$.
Dans ce cas, nous pouvons supposer que $X' \to X$ est un morphisme de schémas,
d'où $Z$ et $Z'$ sont des schémas, et le résultat découle
du cas des schémas, voir
Diviseurs, lemme \ref{divisors-lemma-conormal-algebra-functorial-flat}.
\end{proof}

\noindent
Nous utilisons les mêmes conventions pour les cônes et les fibrés vectoriels sur
les espaces algébriques que pour les schémas (où nous utilisons
les conventions de l'EGA), voir
Constructions, sections \ref{constructions-section-cone} et
\ref{constructions-section-vector-bundle}.
En particulier, un fibré vectoriel est un objet très général
(et n'est pas localement isomorphe à un fibré en espaces affines).

\begin{definition}
\label{definition-normal-cone}
Soit $S$ un schéma. Soit $i : Z \to X$ une immersion d'espaces algébriques
au-dessus de $S$. Le {\it cône normal $C_ZX$} de $Z$ dans $X$ est
$$
C_ZX = \underline{\Spec}_Z(\mathcal{C}_{Z/X, *})
$$
voir Morphismes d'espaces,
définition \ref{spaces-morphisms-definition-relative-spec}. Le
{\it fibré normal} de $Z$ dans $X$ est le fibré vectoriel
$$
N_ZX = \underline{\Spec}_Z(\text{Sym}(\mathcal{C}_{Z/X}))
$$
\end{definition}

\noindent
Ainsi $C_ZX \to Z$ est un cône au-dessus de $Z$ et $N_ZX \to Z$ est un fibré vectoriel
au-dessus de $Z$. De plus, la surjection canonique
(\ref{equation-conormal-algebra-quotient}) d'algèbres graduées
définit une immersion fermée canonique
\begin{equation}
\label{equation-normal-cone-in-normal-bundle}
C_ZX \longrightarrow N_ZX
\end{equation}
de cônes au-dessus de $Z$.










\section{Faisceau des différentielles d'un morphisme}
\label{section-sheaf-differentials}

\noindent
Nous invitons le lecteur à consulter la section correspondante
dans le chapitre d'algèbre commutative
(Algèbre, section \ref{algebra-section-differentials}),
la section correspondante dans le chapitre sur les morphismes de schémas
(Morphismes, section \ref{morphisms-section-sheaf-differentials})
ainsi que
Modules sur les sites, section \ref{sites-modules-section-differentials}.
Nous montrons d'abord que la notion de faisceau des différentielles pour un
morphisme de schémas coïncide avec la notion correspondante pour les
petits sites \'etale (annelés).

\medskip\noindent
Pour énoncer clairement le lemme suivant, nous revenons temporairement à la notation
où $\mathcal{F}^a$ désigne le faisceau de $\mathcal{O}_{X_\etale}$-modules
associé à un $\mathcal{O}_X$-module quasi-cohérent $\mathcal{F}$
sur le schéma $X$, voir
Descente, définition \ref{descent-definition-structure-sheaf}.

\begin{lemma}
\label{lemma-match-modules-differentials}
Soit $f : X \to Y$ un morphisme de schémas. Soit
$f_{small} : X_\etale \to Y_\etale$ le morphisme associé
des petits sites \'etale, voir
Descente, remarque \ref{descent-remark-change-topologies-ringed}.
Alors il existe un isomorphisme canonique
$$
(\Omega_{X/Y})^a = \Omega_{X_\etale/Y_\etale}
$$
compatible avec les dérivations universelles. Ici, le premier module
est le faisceau sur $X_\etale$ associé
au $\mathcal{O}_X$-module quasi-cohérent $\Omega_{X/Y}$, voir
Morphismes, définition \ref{morphisms-definition-sheaf-differentials},
et le second module est celui de
Modules sur les sites,
définition \ref{sites-modules-definition-module-differentials}.
\end{lemma}

\begin{proof}
Soit $h : U \to X$ un morphisme \'etale. Dans ce cas, l'application naturelle
$h^*\Omega_{X/Y} \to \Omega_{U/Y}$ est un isomorphisme, voir
Compléments sur les morphismes,
lemme \ref{more-morphisms-lemma-sheaf-differentials-etale-localization}.
Cela signifie qu'il existe une $\mathcal{O}_{Y_\etale}$-dérivation naturelle
$$
\text{d}^a : \mathcal{O}_{X_\etale} \longrightarrow (\Omega_{X/Y})^a
$$
puisque nous venons de voir que la valeur de $(\Omega_{X/Y})^a$ sur tout objet
$U$ de $X_\etale$ s'identifie canoniquement à
$\Gamma(U, \Omega_{U/Y})$. Par la propriété universelle de
$\text{d}_{X/Y} :
\mathcal{O}_{X_\etale}
\to
\Omega_{X_\etale/Y_\etale}$
il existe une unique application $\mathcal{O}_{X_\etale}$-linéaire
$c : \Omega_{X_\etale/Y_\etale} \to (\Omega_{X/Y})^a$
telle que
$\text{d}^a = c \circ \text{d}_{X/Y}$.

\medskip\noindent
Réciproquement, supposons que $\mathcal{F}$ soit un
$\mathcal{O}_{X_\etale}$-module
et que $D : \mathcal{O}_{X_\etale} \to \mathcal{F}$ soit une
$\mathcal{O}_{Y_\etale}$-dérivation. Nous pouvons simplement restreindre
$D$ au petit site de Zariski $X_{Zar}$ de $X$. Comme les faisceaux sur $X_{Zar}$
coïncident avec les faisceaux sur $X$, voir
Descente, remarque \ref{descent-remark-Zariski-site-space},
nous voyons que $D|_{X_{Zar}} : \mathcal{O}_X \to \mathcal{F}|_{X_{Zar}}$
est simplement une dérivation $Y$ « usuelle ». Nous obtenons donc une application
$\psi : \Omega_{X/Y} \longrightarrow \mathcal{F}|_{X_{Zar}}$
telle que $D|_{X_{Zar}} = \psi \circ \text{d}$. En particulier, si nous
appliquons ceci avec $\mathcal{F} = \Omega_{X_\etale/Y_\etale}$
nous obtenons une application
$$
c' :
\Omega_{X/Y}
\longrightarrow
\Omega_{X_\etale/Y_\etale}|_{X_{Zar}}
$$
Considérons le morphisme de sites annelés
$\text{id}_{small, \etale, Zar} : X_\etale \to X_{Zar}$
décrit dans
Descente, remarque \ref{descent-remark-change-topologies-ringed} et
lemme \ref{descent-lemma-compare-sites}.
Comme le foncteur de restriction $\mathcal{F} \mapsto \mathcal{F}|_{X_{Zar}}$
est égal à $\text{id}_{small, \etale, Zar, *}$, comme
$\text{id}_{small, \etale, Zar}^*$ est adjoint à gauche de
$\text{id}_{small, \etale, Zar, *}$ et comme
$(\Omega_{X/Y})^a = \text{id}_{small, \etale, Zar}^*\Omega_{X/Y}$,
nous voyons que $c'$ est adjoint à une application
$$
c'' :
(\Omega_{X/Y})^a
\longrightarrow
\Omega_{X_\etale/Y_\etale}.
$$
Nous affirmons que $c''$ et $c'$ sont inverses l'un de l'autre.
Cette assertion achève la démonstration du lemme.
Pour le voir, il suffit de montrer que $c''(\text{d}(f)) = \text{d}_{X/Y}(f)$
et $c(\text{d}_{X/Y}(f)) = \text{d}(f)$ lorsque $f$ est une section locale de
$\mathcal{O}_X$ sur un ouvert de $X$. Nous omettons la vérification.
\end{proof}

\noindent
Cette observation ouvre la voie à la définition suivante. Pour une autre possibilité, voir
la remarque \ref{remark-alternative}.

\begin{definition}
\label{definition-sheaf-differentials}
Soit $S$ un schéma. Soit $f : X \to Y$ un morphisme d'espaces algébriques
au-dessus de $S$. Le {\it faisceau des différentielles $\Omega_{X/Y}$ de $X$ sur $Y$}
est le faisceau des différentielles
(Modules sur les sites,
définition \ref{sites-modules-definition-sheaf-differentials})
pour le morphisme de topoi annelés
$$
(f_{small}, f^\sharp) :
(X_\etale, \mathcal{O}_X)
\to
(Y_\etale, \mathcal{O}_Y)
$$
de
Propriétés des espaces,
lemme \ref{spaces-properties-lemma-morphism-ringed-topoi}.
La {\it dérivation universelle $Y$} sera notée
$\text{d}_{X/Y} : \mathcal{O}_X \to \Omega_{X/Y}$.
\end{definition}

\noindent
D'après
le lemme \ref{lemma-match-modules-differentials},
cela n'entre pas en conflit avec la notion déjà existante
dans le cas où $X$ et $Y$ sont représentables. Désormais, si $X$ et $Y$
sont représentables, nous ne faisons plus de distinction entre le faisceau des différentielles
défini ci-dessus et celui défini dans
Morphismes, définition \ref{morphisms-definition-sheaf-differentials}.
Nous voulons relier ceci aux modules usuels des différentielles pour les
morphismes de schémas. Voici le lemme clé.

\begin{lemma}
\label{lemma-localize-differentials}
Soit $S$ un schéma. Soit $f : X \to Y$ un morphisme d'espaces algébriques
au-dessus de $S$. Considérons un diagramme commutatif quelconque
$$
\xymatrix{
U \ar[d]_a \ar[r]_\psi & V \ar[d]^b \\
X \ar[r]^f & Y
}
$$
où les flèches verticales sont des morphismes \'etales d'espaces algébriques. Alors
$$
\Omega_{X/Y}|_{U_\etale} = \Omega_{U/V}
$$
En particulier, si $U$ et $V$ sont des schémas, ceci est égal au
faisceau usuel des différentielles du morphisme de schémas $U \to V$.
\end{lemma}

\begin{proof}
D'après
Propriétés des espaces, lemme \ref{spaces-properties-lemma-etale-morphism-topoi}
et l'Équation (\ref{spaces-properties-equation-restrict}),
on peut considérer la restriction d'un faisceau sur $X_\etale$ à
$U_\etale$ comme l'image inverse par $a_{small}$. De même pour $b$. Par
Modules sur les sites, Lemme \ref{sites-modules-lemma-localize-differentials},
on a
$$
\Omega_{X/Y}|_{U_\etale} =
\Omega_{\mathcal{O}_{U_\etale}/
a_{small}^{-1}f_{small}^{-1}\mathcal{O}_{Y_\etale}}
$$
Comme $a_{small}^{-1}f_{small}^{-1}\mathcal{O}_{Y_\etale}
= \psi_{small}^{-1}b_{small}^{-1}\mathcal{O}_{Y_\etale}
= \psi_{small}^{-1}\mathcal{O}_{V_\etale}$, le lemme en résulte.
\end{proof}

\begin{lemma}
\label{lemma-module-differentials-quasi-coherent}
Soit $S$ un schéma. Soit $f : X \to Y$ un morphisme d'espaces algébriques
sur $S$. Alors $\Omega_{X/Y}$ est un $\mathcal{O}_X$-module quasi-cohérent.
\end{lemma}

\begin{proof}
Choisissons un diagramme comme dans le
Lemme \ref{lemma-localize-differentials},
avec $a$ et $b$ surjectifs et $U$ et $V$ des schémas.
On voit alors que $\Omega_{X/Y}|_U = \Omega_{U/V}$, qui est
quasi-cohérent (par exemple, d'après
Morphismes, Lemme \ref{morphisms-lemma-differentials-diagonal}).
Nous concluons donc que $\Omega_{X/Y}$ est quasi-cohérent par
Propriétés des espaces,
Lemme \ref{spaces-properties-lemma-characterize-quasi-coherent}.
\end{proof}

\begin{remark}
\label{remark-alternative}
Maintenant que nous savons que $\Omega_{X/Y}$ est quasi-cohérent, nous pouvons tenter
de le construire d'une autre manière. On peut par exemple utiliser le résultat de
Propriétés des espaces,
section \ref{spaces-properties-section-quasi-coherent-presentation},
pour construire le faisceau des différentielles par recollement.
Par exemple, si $Y$ est un schéma et si $U \to X$ est un morphisme étale surjectif
d'un schéma vers $X$, alors $\Omega_{U/Y}$ est
un $\mathcal{O}_U$-module quasi-cohérent et, comme $s, t : R \to U$
sont étales, on obtient un isomorphisme
$$
\alpha : s^*\Omega_{U/Y} \to \Omega_{R/Y} \to t^*\Omega_{U/Y}
$$
en utilisant
Morphismes, Lemme \ref{morphisms-lemma-triangle-differentials-smooth}.
On vérifie qu'il satisfait la condition de cocycle, et c'est terminé.
Si $Y$ n'est pas un schéma, on définit $\Omega_{U/Y}$ comme le conoyau
du morphisme $(U \to Y)^*\Omega_{Y/S} \to \Omega_{U/S}$, puis on procède comme
précédemment. Cette procédure en deux étapes est quelque peu disgracieuse. Une autre possibilité
consiste à recoller les faisceaux $\Omega_{U/V}$ pour tout diagramme comme dans le
Lemme \ref{lemma-localize-differentials},
mais cela n'est pas non plus très élégant. Les deux méthodes fonctionnent néanmoins et
donnent une construction légèrement plus élémentaire du faisceau des
différentielles.
\end{remark}

\begin{lemma}
\label{lemma-functoriality-differentials}
Soit $S$ un schéma. Soit
$$
\xymatrix{
X' \ar[d] \ar[r]_f & X \ar[d] \\
Y' \ar[r] & Y
}
$$
un diagramme commutatif d'espaces algébriques. Le morphisme
$f^\sharp : \mathcal{O}_X \to f_*\mathcal{O}_{X'}$, composé avec le morphisme
$f_*\text{d}_{X'/Y'} : f_*\mathcal{O}_{X'} \to f_*\Omega_{X'/Y'}$, est une
$Y$-dérivation. On obtient donc un morphisme canonique de $\mathcal{O}_X$-modules
$\Omega_{X/Y} \to f_*\Omega_{X'/Y'}$ et, par
l'adjonction entre $f_*$ et $f^*$, un
homomorphisme canonique de $\mathcal{O}_{X'}$-modules
$$
c_f : f^*\Omega_{X/Y} \longrightarrow \Omega_{X'/Y'}.
$$
Il est uniquement caractérisé par la propriété suivante :
$f^*\text{d}_{X/Y}(t)$ est envoyé sur $\text{d}_{X'/Y'}(f^* t)$
pour toute section locale $t$ de $\mathcal{O}_X$.
\end{lemma}

\begin{proof}
Il s'agit d'un cas particulier de
Modules sur les sites, Lemme
\ref{sites-modules-lemma-functoriality-differentials-ringed-topoi}.
\end{proof}

\begin{lemma}
\label{lemma-check-functoriality-differentials}
Soit $S$ un schéma. Soit
$$
\xymatrix{
X'' \ar[d] \ar[r]_g & X' \ar[d] \ar[r]_f & X \ar[d] \\
Y'' \ar[r] & Y' \ar[r] & Y
}
$$
un diagramme commutatif d'espaces algébriques sur $S$. Alors on a
$$
c_{f \circ g} = c_g \circ g^* c_f
$$
comme morphismes $(f \circ g)^*\Omega_{X/Y} \to \Omega_{X''/Y''}$.
\end{lemma}

\begin{proof}
Démonstration omise. Indication : utiliser la caractérisation de $c_f, c_g, c_{f \circ g}$
par l'effet de ces morphismes sur les sections locales.
\end{proof}

\begin{lemma}
\label{lemma-triangle-differentials}
Soit $S$ un schéma.
Soient $f : X \to Y$ et $g : Y \to B$ des morphismes d'espaces algébriques sur $S$.
Alors il existe une suite exacte canonique
$$
f^*\Omega_{Y/B} \to \Omega_{X/B} \to \Omega_{X/Y} \to 0
$$
où les morphismes proviennent d'applications du
Lemme \ref{lemma-functoriality-differentials}.
\end{lemma}

\begin{proof}
Cela découle de la version de ce résultat pour les schémas, voir
Morphismes, Lemme \ref{morphisms-lemma-triangle-differentials},
par localisation étale, voir le
Lemme \ref{lemma-localize-differentials}.
\end{proof}

\begin{lemma}
\label{lemma-immersion-differentials}
Soit $S$ un schéma. Si $X \to Y$ est une immersion
d'espaces algébriques sur $S$, alors $\Omega_{X/S}$ est nul.
\end{lemma}

\begin{proof}
Cela découle de la version de ce résultat pour les schémas, voir
Morphismes, Lemme \ref{morphisms-lemma-immersion-differentials},
par localisation étale, voir le
Lemme \ref{lemma-localize-differentials}.
\end{proof}

\begin{lemma}
\label{lemma-differentials-relative-immersion}
Soit $S$ un schéma. Soit $B$ un espace algébrique sur $S$.
Soit $i : Z \to X$ une immersion d'espaces algébriques sur $B$.
Il existe une suite exacte canonique
$$
\mathcal{C}_{Z/X} \to i^*\Omega_{X/B} \to \Omega_{Z/B} \to 0
$$
où la première flèche est induite par $\text{d}_{X/B}$
et où la seconde flèche provient du
Lemme \ref{lemma-functoriality-differentials}.
\end{lemma}

\begin{proof}
Il s'agit de la version pour les espaces algébriques de
Morphismes, Lemme \ref{morphisms-lemma-differentials-relative-immersion},
et elle découlera de ce lemme par
localisation étale, voir les
Lemmes \ref{lemma-localize-differentials} et
\ref{lemma-etale-conormal}.
Il faut toutefois s'assurer que l'on peut définir globalement la première flèche.
Expliquons donc ici le sens de « induite par $\text{d}_{X/B}$ ».
On peut en effet supposer que $i$ est une immersion fermée après avoir remplacé $X$
par un sous-espace ouvert. Soit $\mathcal{I} \subset \mathcal{O}_X$
le faisceau quasi-cohérent d'idéaux correspondant à $Z \subset X$.
Alors $\text{d}_{X/S} : \mathcal{I} \to \Omega_{X/S}$
envoie le sous-faisceau $\mathcal{I}^2 \subset \mathcal{I}$ dans
$\mathcal{I}\Omega_{X/S}$. Il induit donc un morphisme
$\mathcal{I}/\mathcal{I}^2 \to \Omega_{X/S}/\mathcal{I}\Omega_{X/S}$
qui est $\mathcal{O}_X/\mathcal{I}$-linéaire.
D'après
Morphismes d'espaces, Lemme \ref{spaces-morphisms-lemma-i-star-equivalence},
cela correspond au morphisme $\mathcal{C}_{Z/X} \to i^*\Omega_{X/S}$ voulu.
\end{proof}

\begin{lemma}
\label{lemma-differentials-relative-immersion-section}
Soit $S$ un schéma. Soit $B$ un espace algébrique sur $S$.
Soit $i : Z \to X$ une immersion d'espaces algébriques sur $B$, et
supposons que $i$ admette (localement pour la topologie étale) un inverse à gauche. Alors la
suite canonique
$$
0 \to \mathcal{C}_{Z/X} \to i^*\Omega_{X/B} \to \Omega_{Z/B} \to 0
$$
du
Lemme \ref{lemma-differentials-relative-immersion}
est (localement pour la topologie étale) exacte scindée.
\end{lemma}

\begin{proof}
Précision : nous affirmons que, si $g : X \to Z$ est un inverse à gauche de $i$
sur $B$, alors $i^*c_g$ est un inverse à droite du morphisme
$i^*\Omega_{X/B} \to \Omega_{Z/B}$.
Ceci étant dit, le résultat découle du résultat correspondant pour les
morphismes de schémas par localisation étale, voir les
Lemmes \ref{lemma-localize-differentials} et
\ref{lemma-etale-conormal}.
\end{proof}

\begin{lemma}
\label{lemma-base-change-differentials}
Soit $S$ un schéma.
Soit $X \to Y$ un morphisme d'espaces algébriques sur $S$.
Soit $g : Y' \to Y$ un morphisme d'espaces algébriques sur $S$.
Soit $X' = X_{Y'}$ le changement de base de $X$.
Notons $g' : X' \to X$ la projection.
Alors le morphisme
$$
(g')^*\Omega_{X/Y} \to \Omega_{X'/Y'}
$$
du
Lemme \ref{lemma-functoriality-differentials}
est un isomorphisme.
\end{lemma}

\begin{proof}
Cela découle de la version pour les schémas, voir
Morphismes, Lemme \ref{morphisms-lemma-base-change-differentials},
et de la localisation étale, voir le
Lemme \ref{lemma-localize-differentials}.
\end{proof}

\begin{lemma}
\label{lemma-differential-product}
Soit $S$ un schéma.
Soient $f : X \to B$ et $g : Y \to B$ des morphismes d'espaces algébriques
sur $S$ ayant le même but.
Soient $p : X \times_B Y \to X$ et $q : X \times_B Y \to Y$ les
morphismes de projection. Les morphismes du
Lemme \ref{lemma-functoriality-differentials}
$$
p^*\Omega_{X/B} \oplus q^*\Omega_{Y/B}
\longrightarrow
\Omega_{X \times_B Y/B}
$$
donnent un isomorphisme.
\end{lemma}

\begin{proof}
Cela découle de la version pour les schémas, voir
Morphismes, Lemme \ref{morphisms-lemma-differential-product},
et de la localisation étale, voir le
Lemme \ref{lemma-localize-differentials}.
\end{proof}

\begin{lemma}
\label{lemma-finite-type-differentials}
Soit $S$ un schéma.
Soit $f : X \to Y$ un morphisme d'espaces algébriques sur $S$.
Si $f$ est localement de type fini, alors $\Omega_{X/Y}$ est
un $\mathcal{O}_X$-module de type fini.
\end{lemma}

\begin{proof}
Cela découle de la version pour les schémas, voir
Morphismes, Lemme \ref{morphisms-lemma-finite-type-differentials},
et de la localisation étale, voir le
Lemme \ref{lemma-localize-differentials}.
\end{proof}

\begin{lemma}
\label{lemma-finite-presentation-differentials}
Soit $S$ un schéma.
Soit $f : X \to Y$ un morphisme d'espaces algébriques sur $S$.
Si $f$ est localement de présentation finie, alors $\Omega_{X/Y}$ est
un $\mathcal{O}_X$-module de présentation finie.
\end{lemma}

\begin{proof}
Cela découle de la version pour les schémas, voir
Morphismes, Lemme \ref{morphisms-lemma-finite-presentation-differentials},
et de la localisation étale, voir le
Lemme \ref{lemma-localize-differentials}.
\end{proof}

\begin{lemma}
\label{lemma-smooth-omega-finite-locally-free}
Soit $S$ un schéma.
Soit $f : X \to Y$ un morphisme lisse d'espaces algébriques sur $S$.
Alors le module des différentielles $\Omega_{X/Y}$
est localement libre de type fini.
\end{lemma}

\begin{proof}
L'énoncé est local pour la topologie étale sur $X$ et $Y$ par le
Lemme \ref{lemma-localize-differentials}.
Il découle donc du cas des schémas, voir
Morphismes, Lemme \ref{morphisms-lemma-smooth-omega-finite-locally-free}.
\end{proof}













\section{Invariance topologique du site étale}
\label{section-topological-invariance}

\noindent
Nous montrons que le site $X_{spaces, \etale}$ est un « invariant
topologique ». Il s'ensuit que $X_\etale$, qui
est formé des objets représentables de $X_{spaces, \etale}$,
est lui aussi un invariant topologique, voir le Lemme \ref{lemma-topological-invariance}.

\begin{theorem}
\label{theorem-topological-invariance}
Soit $S$ un schéma.
Soit $f : X \to Y$ un morphisme d'espaces algébriques sur $S$.
Supposons que $f$ soit entier, universellement injectif et surjectif.
Le foncteur
$$
V \longmapsto V_X = X \times_Y V
$$
définit une équivalence de catégories
$Y_{spaces, \etale} \to X_{spaces, \etale}$.
\end{theorem}

\begin{proof}
Le morphisme $f$ est représentable et constitue un homéomorphisme universel, voir
Morphismes d'espaces,
section \ref{spaces-morphisms-section-universal-homeomorphisms}.

\medskip\noindent
Nous montrons d'abord que le foncteur est fidèle.
Supposons que $V', V$ soient des objets de $Y_{spaces, \etale}$ et
que $a, b : V' \to V$ soient des morphismes distincts sur $Y$.
Comme $V', V$ sont étales sur $Y$, l'égalisateur
$$
E =  V' \times_{(a, b), V \times_Y V, \Delta_{V/Y}} V
$$
de $a, b$ est lui aussi étale sur $Y$. Ainsi, $E \to V'$ est un monomorphisme étale
(c'est-à-dire une immersion ouverte) qui est un isomorphisme si et seulement s'il est
surjectif. Comme $X \to Y$ est un homéomorphisme universel, nous voyons que c'est
le cas si et seulement si $E_X = V'_X$, c'est-à-dire si et seulement si $a_X = b_X$.

\medskip\noindent
Nous montrons ensuite que le foncteur est pleinement fidèle.
Supposons que $V', V$ soient des objets de $Y_{spaces, \etale}$ et
que $c : V'_X \to V_X$ soit un morphisme sur $X$. Nous voulons construire
un morphisme $a : V' \to V$ sur $Y$ tel que $a_X = c$.
Soit $a' : V'' \to V'$ un morphisme étale surjectif tel que $V''$ soit
un espace algébrique séparé. Si l'on peut construire un morphisme
$a'' : V'' \to V$ tel que $a''_X = c \circ a'_X$, alors les deux composées
$$
V'' \times_{V'} V'' \xrightarrow{\text{pr}_i} V'' \xrightarrow{a''} V
$$
sont égales par la fidélité du foncteur établie au premier
paragraphe. Ainsi, $a''$ se factorise par un unique morphisme
$a : V' \to V$, puisque $V'$ est (en tant que faisceau) le quotient de $V''$ par
la relation d'équivalence $V'' \times_{V'} V''$. Nous pouvons donc supposer que
$V'$ est séparé. Dans ce cas, le graphe
$$
\Gamma_c \subset (V' \times_Y V)_X
$$
est ouvert et fermé (détails omis). Comme $X \to Y$ est un homéomorphisme
universel, il existe un sous-espace ouvert et fermé
$\Gamma \subset V' \times_Y V$ tel que $\Gamma_X = \Gamma_c$.
La projection $\Gamma \to V'$ est un morphisme étale dont le changement
de base à $X$ est un isomorphisme. Ainsi, $\Gamma \to V'$ est étale,
universellement injectif et surjectif, donc est un isomorphisme d'après
Morphismes d'espaces,
Lemme \ref{spaces-morphisms-lemma-etale-universally-injective-open}.
Ainsi, $\Gamma$ est le graphe d'un morphisme $a : V' \to V$, comme souhaité.

\medskip\noindent
Enfin, nous montrons que le foncteur est essentiellement surjectif.
Supposons que $U$ soit un objet de $X_{spaces, \etale}$.
Il faut trouver un objet $V$ de $Y_{spaces, \etale}$
tel que $V_X \cong U$. Soit $U' \to U$ un morphisme étale surjectif
tel que $U' \cong V'_X$ et $U' \times_U U' \cong V''_X$
pour certains objets $V'', V'$ de $Y_{spaces, \etale}$.
Alors, par la pleine fidélité du foncteur, on obtient des morphismes
$s, t : V'' \to V'$ tels que $t_X = \text{pr}_0$ et $s_X = \text{pr}_1$
comme morphismes $U' \times_U U' \to U'$. Puisque
$(\text{pr}_0, \text{pr}_1) : U' \times_U U' \to U' \times_S U'$
est une relation d'équivalence étale et que $U' \to V'$ et
$U' \times_U U' \to V''$ sont universellement injectifs et surjectifs,
on en déduit que
$(t, s) : V'' \to V' \times_S V'$ est une relation d'équivalence étale.
Alors le quotient $V = V'/V''$ (voir
Espaces, Théorème \ref{spaces-theorem-presentation})
est l'espace algébrique $V$ sur $Y$. Il existe un morphisme
$V' \to V$ tel que $V'' = V' \times_V V'$. On obtient ainsi un morphisme
$V \to Y$ (voir
Descente sur les espaces, Lemme
\ref{spaces-descent-lemma-fpqc-universal-effective-epimorphisms}).
Après changement de base à $X$, on voit que l'on a un morphisme $U' \to V_X$
et un isomorphisme compatible $U' \times_{V_X} U' = U' \times_U U'$, ce qui
implique que $V_X \cong U$ (en appliquant encore une fois le lemme qui vient d'être cité).

\medskip\noindent
Choisissons un schéma $W$ et un morphisme étale surjectif $W \to Y$.
Choisissons un schéma $U'$ et un morphisme étale surjectif $U' \to U \times_X W_X$.
Notons que $U'$ et $U' \times_U U'$ sont des schémas étales sur $X$ dont le
morphisme structural vers $X$ se factorise par le schéma $W_X$.
Ainsi, d'après
Cohomologie étale,
Théorème \ref{etale-cohomology-theorem-topological-invariance},
il existe des schémas $V', V''$ étales sur $W$ dont les changements de base à
$W_X$ sont respectivement isomorphes à $U'$ et $U' \times_U U'$.
Cela achève la démonstration.
\end{proof}

\begin{lemma}
\label{lemma-topological-invariance}
Sous les hypothèses et avec les notations du
Théorème \ref{theorem-topological-invariance},
l'équivalence de catégories
$Y_{spaces, \etale} \to X_{spaces, \etale}$
induit par restriction des équivalences de catégories
$Y_\etale \to X_\etale$ et $Y_{affine, \etale} \to X_{affine, \etale}$.
\end{lemma}

\begin{proof}
Il s'agit simplement de dire que, pour tout objet
$V \in Y_{spaces, \etale}$, $V$ est un schéma (respectivement un schéma affine) si et
seulement si $V \times_Y X$ est un schéma (respectivement un schéma affine). Comme $V \times_Y X \to V$
est entier, universellement injectif et surjectif (en tant que changement
de base de $X \to Y$), cela
découle du Lemme
\ref{spaces-limits-lemma-integral-universally-bijective-scheme} de Limites d'espaces et de la
Proposition \ref{spaces-limits-proposition-affine}.
\end{proof}

\begin{remark}
\label{remark-topological-invariance-etale-site}
\begin{reference}
Courriel de Lenny Taelman daté du 1er mai 2016.
\end{reference}
Un homéomorphisme universel d'espaces algébriques n'est pas nécessairement représentable, voir
Morphismes d'espaces,
Exemple \ref{spaces-morphisms-example-universal-homeomorphism}.
En fait, le Théorème \ref{theorem-topological-invariance} n'est
pas valable pour tous les homéomorphismes universels. Pour le voir, soit $k$ un
corps algébriquement clos de caractéristique $0$ et soit
$$
\mathbf{A}^1 \to X \to \mathbf{A}^1
$$
comme dans Morphismes d'espaces,
Exemple \ref{spaces-morphisms-example-universal-homeomorphism}.
Rappelons que le premier morphisme est étale et identifie
$t$ à $-t$ pour $t \in \mathbf{A}^1_k \setminus \{0\}$
et que le second morphisme est notre homéomorphisme universel.
Comme $\mathbf{A}^1_k$ n'admet aucun
revêtement fini étale connexe non trivial
(car $k$ est algébriquement clos de caractéristique zéro ; détails omis),
il suffit de construire un revêtement fini étale connexe non trivial
$Y \to X$. Pour cela, soit $Y$ la droite affine
à origine dédoublée
(Schémas, Exemple \ref{schemes-example-affine-space-zero-doubled}).
Alors $Y = Y_1 \cup Y_2$, avec $Y_i = \mathbf{A}^1_k$, est obtenu par recollement
le long de $\mathbf{A}^1_k \setminus \{0\}$.
Pour définir le morphisme $Y \to X$, nous utilisons les morphismes
$$
Y_1 \xrightarrow{1} \mathbf{A}^1_k \to X
\quad\text{et}\quad
Y_2 \xrightarrow{-1} \mathbf{A}^1_k \to X.
$$
Ces morphismes se recollent sur $Y_1 \cap Y_2$ par construction de $X$ et
définissent donc un morphisme $Y \to X$. En fait, nous affirmons que
$$
\xymatrix{
Y \ar[d] & Y_1 \amalg Y_2 \ar[l] \ar[d] \\
X & \mathbf{A}^1_k \ar[l]
}
$$
est un carré cartésien. Nous omettons les détails ; on peut utiliser par exemple
Groupoïdes, Lemme \ref{groupoids-lemma-criterion-fibre-product}.
Comme $\mathbf{A}^1_k \to X$ est étale et
surjectif, cela montre que $Y \to X$
est fini étale de degré $2$, ce qui fournit l'exemple voulu.

\medskip\noindent
Plus simplement, on peut raisonner comme suit. Le schéma $Y$ est muni de l'action du groupe
$G = \{+1, -1\}$, qui opère librement, l'élément $-1$ agissant en échangeant
$Y_1$ et $Y_2$ et en changeant le signe de la coordonnée. Alors
$X = Y/G$ (voir Espaces, Définition \ref{spaces-definition-quotient})
et donc $Y \to X$ est fini étale. On peut aussi montrer directement
qu'il existe un homéomorphisme universel $X \to \mathbf{A}^1_k$
en utilisant $t \mapsto t^2$ sur les espaces affines. En fait, ce $X$ est
le même que le $X$ ci-dessus.
\end{remark}








\section{Épaississements}
\label{section-thickenings}

\noindent
La terminologie suivante n'est peut-être pas tout à fait standard, mais elle est commode.

\begin{definition}
\label{definition-thickening}
Épaississements. Soit $S$ un schéma.
\begin{enumerate}
\item Nous disons qu'un espace algébrique $X'$ est un {\it épaississement} d'un espace
algébrique $X$ si $X$ est un sous-espace fermé de $X'$ et si les espaces topologiques
associés sont égaux.
\item Nous disons que $X'$ est un {\it épaississement du premier ordre} de $X$ si
$X$ est un sous-espace fermé de $X'$ et si le faisceau quasi-cohérent d'idéaux
$\mathcal{I} \subset \mathcal{O}_{X'}$ qui définit $X$ est de carré nul.
\item Nous disons que $X'$ est un {\it épaississement d'ordre fini} de $X$
si $X$ est un sous-espace fermé de $X$ et si le faisceau quasi-cohérent d'idéaux
$\mathcal{I} \subset \mathcal{O}_{X'}$ qui définit $X$ est nilpotent, c'est-à-dire s'il
existe un entier $n \geq 0$ tel que $\mathcal{I}^{n + 1} = 0$.
\item Nous disons que $X'$ est un {\it épaississement d'ordre $n$} de $X$
si $X$ est un sous-espace fermé de $X$ et si $\mathcal{I}^{n + 1} = 0$,
où $\mathcal{I} \subset \mathcal{O}_{X'}$ est le
faisceau quasi-cohérent d'idéaux qui définit $X$.
\item Étant donnés deux épaississements $X \subset X'$ et $Y \subset Y'$, un
{\it morphisme d'épaississements} est un morphisme $f' : X' \to Y'$ tel que
$f(X) \subset Y$, c'est-à-dire tel que $f'|_X$ se factorise par le sous-espace
fermé $Y$. Dans cette situation, nous posons $f = f'|_X : X \to Y$ et disons
que $(f, f') : (X \subset X') \to (Y \subset Y')$ est un morphisme
d'épaississements.
\item Soit $B$ un espace algébrique. Nous définissons de même les
{\it épaississements sur $B$} et les
{\it morphismes d'épaississements sur $B$}. Cela signifie que les espaces
$X, X', Y, Y'$ ci-dessus sont des espaces algébriques munis d'un morphisme
structural vers $B$ et que les morphismes
$X \to X'$, $Y \to Y'$ et $f' : X' \to Y'$ sont des morphismes sur $B$.
\end{enumerate}
\end{definition}

\noindent
L'équivalence fondamentale.
Notons que si $X \subset X'$ est un épaississement, alors $X \to X'$
est entier et universellement bijectif. Il en résulte que
\begin{equation}
\label{equation-equivalence-etale-spaces}
X_{spaces, \etale} = X'_{spaces, \etale}
\end{equation}
par le foncteur d'image inverse, voir le
Théorème \ref{theorem-topological-invariance}.
Nous pouvons donc considérer $\mathcal{O}_{X'}$ comme un faisceau sur
$X_{spaces, \etale}$. Nous obtenons ainsi une équivalence canonique
de topos localement annelés
\begin{equation}
\label{equation-fundamental-equivalence}
(\Sh(X'_{spaces, \etale}), \mathcal{O}_{X'})
\cong
(\Sh(X_{spaces, \etale}), \mathcal{O}_{X'})
\end{equation}
Dans la suite, nous la combinerons fréquemment avec le résultat de pleine fidélité de
Propriétés des espaces, Théorème \ref{spaces-properties-theorem-fully-faithful}.
Par exemple, à l'immersion fermée $i_X : X \to X'$ correspond
le morphisme surjectif $i_X^\sharp : \mathcal{O}_{X'} \to \mathcal{O}_X$.

\medskip\noindent
Soit $S$ un schéma et soit $B$ un espace algébrique sur $S$.
Soit $(f, f') : (X \subset X') \to (Y \subset Y')$ un morphisme
d'épaississements sur $B$. Notons que le diagramme de foncteurs continus
$$
\xymatrix{
X_{spaces, \etale} &
Y_{spaces, \etale} \ar[l] \\
X'_{spaces, \etale} \ar[u] &
Y'_{spaces, \etale} \ar[u] \ar[l]
}
$$
est commutatif et que les flèches verticales sont des équivalences. Ainsi,
$f_{spaces, \etale}$, $f_{small}$,
$f'_{spaces, \etale}$, et $f'_{small}$
définissent tous le même morphisme de topos. Nous pouvons donc considérer
$$
(f')^\sharp :
f_{spaces, \etale}^{-1}\mathcal{O}_{Y'}
\longrightarrow
\mathcal{O}_{X'}
$$
comme un morphisme de faisceaux de $\mathcal{O}_B$-algèbres qui s'insère dans le
diagramme commutatif
$$
\xymatrix{
f_{spaces, \etale}^{-1}\mathcal{O}_Y \ar[r]_-{f^\sharp} \ar[r] &
\mathcal{O}_X \\
f_{spaces, \etale}^{-1}\mathcal{O}_{Y'} \ar[r]^-{(f')^\sharp}
\ar[u]^{i_Y^\sharp} &
\mathcal{O}_{X'} \ar[u]_{i_X^\sharp}
}
$$
Ici, $i_X : X \to X'$ et $i_Y : Y \to Y'$ désignent les
immersions fermées données.

\begin{lemma}
\label{lemma-first-order-thickening-maps}
Soit $S$ un schéma. Soit $B$ un espace algébrique sur $S$.
Soient $X \subset X'$ et $Y \subset Y'$ des épaississements
d'espaces algébriques sur $B$. Soit $f : X \to Y$ un morphisme d'espaces
algébriques sur $B$. Étant donné un morphisme de $\mathcal{O}_B$-algèbres
$$
\alpha : f_{spaces, \etale}^{-1}\mathcal{O}_{Y'} \to \mathcal{O}_{X'}
$$
tel que le diagramme
$$
\xymatrix{
f_{spaces, \etale}^{-1}\mathcal{O}_Y \ar[r]_-{f^\sharp} \ar[r] &
\mathcal{O}_X \\
f_{spaces, \etale}^{-1}\mathcal{O}_{Y'} \ar[r]^-\alpha
\ar[u]^{i_Y^\sharp} &
\mathcal{O}_{X'} \ar[u]_{i_X^\sharp}
}
$$
soit commutatif, il existe un unique morphisme $(f, f')$
d'épaississements sur $B$ tel que $\alpha = (f')^\sharp$.
\end{lemma}

\begin{proof}
Pour trouver $f'$, d'après
Propriétés des espaces, Théorème \ref{spaces-properties-theorem-fully-faithful},
il suffit de montrer que le morphisme de topos annelés
$$
(f_{spaces, \etale}, \alpha) :
(\Sh(X_{spaces, \etale}), \mathcal{O}_{X'})
\longrightarrow
(\Sh(Y_{spaces, \etale}), \mathcal{O}_{Y'})
$$
est un morphisme de topos localement annelés. Cela découle directement
de la définition des morphismes de topos localement annelés
(Modules sur les sites,
Définition \ref{sites-modules-definition-morphism-locally-ringed-topoi}),
du fait que $(f, f^\sharp)$ est un morphisme de topos localement annelés
(Propriétés des espaces,
Lemme \ref{spaces-properties-lemma-morphism-locally-ringed}),
du fait que $\alpha$ s'insère dans le diagramme commutatif donné et
du fait que les noyaux de $i_X^\sharp$ et de $i_Y^\sharp$ sont
localement nilpotents. Enfin, l'égalité $f' \circ i_X = i_Y \circ f$
découle de la commutativité du diagramme et d'une nouvelle application de
Propriétés des espaces, Théorème \ref{spaces-properties-theorem-fully-faithful}.
Nous omettons la vérification que $f'$ est un morphisme sur $B$.
\end{proof}

\begin{lemma}
\label{lemma-open-subspace-thickening}
Soit $S$ un schéma. Soit $X \subset X'$ un épaississement
d'espaces algébriques sur $S$. Pour tout sous-espace ouvert $U \subset X$, il
existe un unique sous-espace ouvert $U' \subset X'$ tel que
$U = X \times_{X'} U'$.
\end{lemma}

\begin{proof}
Soit $U' \to X'$ l'objet de $X'_{spaces, \etale}$
qui correspond à l'objet $U \to X$ de $X_{spaces, \etale}$
par (\ref{equation-equivalence-etale-spaces}). Le morphisme
$U' \to X'$ est étale et universellement injectif, donc une immersion ouverte, voir
Morphismes d'espaces,
Lemme \ref{spaces-morphisms-lemma-etale-universally-injective-open}.
\end{proof}

\noindent
Soit $i_X : X \to X'$ un épaississement d'espaces algébriques. Toute section locale
du noyau $\mathcal{I} = \Ker(i_X^\sharp) \subset \mathcal{O}_{X'}$ est
localement nilpotente. Cela donne un certain contrôle sur le faisceau structural
$\mathcal{O}_{X'}$, mais, techniquement, la classe des épaississements d'ordre fini
$X \subset X'$ est beaucoup plus facile à manier.
En effet, si $X'$ est un épaississement d'ordre $n$, nous avons une filtration
$$
0 \subset \mathcal{I}^n \subset \mathcal{I}^{n - 1} \subset
\ldots \subset \mathcal{I} \subset \mathcal{O}_{X'}
$$
et nous voyons que $X'$ est muni d'une filtration par des sous-espaces fermés
$$
X = X_0 \subset X_1 \subset \ldots \subset X_{n - 1} \subset X_{n + 1} = X'
$$
telle que chaque paire $X_i \subset X_{i + 1}$ est un épaississement du premier ordre
sur $B$. De simples arguments de récurrence permettent de reformuler nombre de résultats démontrés pour les
épaississements du premier ordre en des résultats sur les épaississements d'ordre fini.

\begin{lemma}
\label{lemma-first-order-thickening-surjective}
Soit $S$ un schéma. Soit $X \subset X'$ un épaississement
d'espaces algébriques sur $S$. Soit $U$ un objet affine de
$X_{spaces, \etale}$. Alors
$$
\Gamma(U, \mathcal{O}_{X'}) \to \Gamma(U, \mathcal{O}_X)
$$
est surjectif, où l'on considère $\mathcal{O}_{X'}$ comme un faisceau sur
$X_{spaces, \etale}$ au moyen de (\ref{equation-fundamental-equivalence}).
\end{lemma}

\begin{proof}
Soit $U' \to X'$ le morphisme étale d'espaces algébriques tel que
$U = X \times_{X'} U'$, voir le Théorème \ref{theorem-topological-invariance}.
D'après Limites d'espaces, Lemme \ref{spaces-limits-lemma-affine}, nous voyons
que $U'$ est un schéma affine. Par conséquent,
$\Gamma(U, \mathcal{O}_{X'}) = \Gamma(U', \mathcal{O}_{U'}) \to
\Gamma(U, \mathcal{O}_U)$
est surjectif, puisque $U \to U'$ est une immersion fermée de schémas affines.
Nous donnons ci-dessous une démonstration directe pour les épaississements d'ordre fini,
cas le plus utilisé en pratique.
\end{proof}

\begin{proof}[Démonstration pour les épaississements d'ordre fini]
Nous pouvons supposer que $X \subset X'$ est un épaississement du premier ordre, par le
principe expliqué ci-dessus. Notons $\mathcal{I}$ le noyau de la surjection
$\mathcal{O}_{X'} \to \mathcal{O}_X$. Comme $\mathcal{I}$ est un
$\mathcal{O}_{X'}$-module quasi-cohérent et que $\mathcal{I}^2 = 0$ par définition
d'un épaississement du premier ordre, nous pouvons appliquer
Morphismes d'espaces, Lemme
\ref{spaces-morphisms-lemma-i-star-equivalence}
pour voir que $\mathcal{I}$ est un $\mathcal{O}_X$-module quasi-cohérent.
Le lemme découle donc de la suite exacte longue de cohomologie
associée à la suite exacte courte
$$
0 \to \mathcal{I} \to \mathcal{O}_{X'} \to \mathcal{O}_X \to 0
$$
et du fait que $H^1_\etale(U, \mathcal{I}) = 0$, puisque
$\mathcal{I}$ est quasi-cohérent, voir
Descente, Proposition \ref{descent-proposition-same-cohomology-quasi-coherent},
et Cohomologie des schémas, Lemme
\ref{coherent-lemma-quasi-coherent-affine-cohomology-zero}.
\end{proof}

\begin{lemma}
\label{lemma-thickening-scheme}
Soit $S$ un schéma. Soit $X \subset X'$ un épaississement d'espaces algébriques
sur $S$. Si $X$ est (représentable par) un schéma, alors $X'$ l'est aussi.
\end{lemma}

\begin{proof}
Notons que $X'_{red} = X_{red}$. Par conséquent, si $X$ est un schéma, alors
$X'_{red}$ est un schéma. Le résultat découle donc de
Limites d'espaces, Lemme
\ref{spaces-limits-lemma-reduction-scheme}.
Nous donnons ci-dessous une démonstration directe pour les épaississements d'ordre fini,
cas le plus souvent utilisé en pratique.
\end{proof}

\begin{proof}[Démonstration pour les épaississements d'ordre fini]
Il suffit de le démontrer lorsque $X'$ est un épaississement du premier ordre de $X$. D'après
Propriétés des espaces, Lemme \ref{spaces-properties-lemma-subscheme},
il existe un plus grand sous-espace ouvert de $X'$ qui soit un schéma. Nous devons donc
montrer que tout point $x$ de $|X'| = |X|$ est contenu dans un sous-espace ouvert de
$X'$ qui soit un schéma. En utilisant le
Lemme \ref{lemma-open-subspace-thickening},
nous pouvons remplacer $X \subset X'$ par $U \subset U'$, avec $x \in U$ et $U$
un schéma affine. Nous pouvons donc supposer que $X$ est affine.
Nous sommes ainsi ramenés au cas traité dans le paragraphe suivant.

\medskip\noindent
Supposons que $X \subset X'$ soit un épaississement du premier ordre, où $X$ est un schéma
affine. Posons $A = \Gamma(X, \mathcal{O}_X)$ et
$A' = \Gamma(X', \mathcal{O}_{X'})$. D'après le
Lemme \ref{lemma-first-order-thickening-surjective},
le morphisme $A' \to A$ est surjectif. Son noyau $I$ est un idéal de carré nul. D'après
Propriétés des espaces,
Lemme \ref{spaces-properties-lemma-morphism-to-affine-scheme},
nous obtenons un morphisme canonique $f : X' \to \Spec(A')$ qui s'insère
dans le diagramme commutatif suivant
$$
\xymatrix{
X \ar@{=}[d] \ar[r] &  X' \ar[d]^f \\
\Spec(A) \ar[r] & \Spec(A')
}
$$
Comme les flèches horizontales sont des épaississements, il est clair que $f$ est
universellement injectif et surjectif. Il suffit donc de montrer que
$f$ est étale, car alors
Morphismes d'espaces,
Lemme \ref{spaces-morphisms-lemma-etale-universally-injective-open},
entraînera que $f$ est un isomorphisme.

\medskip\noindent
Pour démontrer que $f$ est étale, choisissons un schéma affine $U'$ et un
morphisme étale $U' \to X'$. Il suffit de montrer que
$U' \to X' \to \Spec(A')$ est étale, voir
Propriétés des espaces, Définition \ref{spaces-properties-definition-etale}.
Écrivons $U' = \Spec(B')$. Posons $U = X \times_{X'} U'$. Comme $U$
est un sous-espace fermé de $U'$, c'est un sous-schéma fermé, et donc
$U = \Spec(B)$ avec $B' \to B$ surjectif. Notons
$J = \Ker(B' \to B)$ et remarquons que $J = \Gamma(U, \mathcal{I})$,
où $\mathcal{I} = \Ker(\mathcal{O}_{X'} \to \mathcal{O}_X)$
sur $X_{spaces, \etale}$, comme dans la démonstration du
Lemme \ref{lemma-first-order-thickening-surjective}.
Le morphisme $U' \to X' \to \Spec(A')$ induit un diagramme
commutatif
$$
\xymatrix{
0 \ar[r] &
J \ar[r] &
B' \ar[r] &
B \ar[r] & 0 \\
0 \ar[r] &
I \ar[r] \ar[u] &
A' \ar[r] \ar[u] &
A \ar[r] \ar[u] & 0
}
$$
Maintenant, comme $\mathcal{I}$ est un $\mathcal{O}_X$-module quasi-cohérent,
nous avons $\mathcal{I} = (\widetilde I)^a$, voir
Descente, Définition \ref{descent-definition-structure-sheaf},
pour les notations, et
Descente, Proposition \ref{descent-proposition-equivalence-quasi-coherent},
pour la justification. Nous voyons donc que $J = I \otimes_A B$.
Enfin, notons que $A \to B$ est étale, car $U \to X$ est étale en tant que
changement de base du morphisme étale $U' \to X'$.
Nous concluons que $A' \to B'$ est étale par
Algèbre, Lemme \ref{algebra-lemma-lift-etale-infinitesimal}.
\end{proof}

\begin{lemma}
\label{lemma-thickening-equivalence}
Soit $S$ un schéma. Soit $X \subset X'$ un épaississement d'espaces algébriques
sur $S$. Le foncteur
$$
V' \longmapsto V = X \times_{X'} V'
$$
définit une équivalence de catégories
$X'_\etale \to X_\etale$.
\end{lemma}

\begin{proof}
Le foncteur $V' \mapsto V$ définit une équivalence de catégories
$X'_{spaces, \etale} \to X_{spaces, \etale}$, voir le
Théorème \ref{theorem-topological-invariance}.
Il suffit donc de montrer que $V$ est un schéma si et seulement si $V'$ est
un schéma. C'est le contenu du
Lemme \ref{lemma-thickening-scheme}.
\end{proof}

\noindent
Les épaississements du premier ordre se décrivent comme suit.

\begin{lemma}
\label{lemma-first-order-thickening}
Soit $S$ un schéma.
Soit $f : X \to B$ un morphisme d'espaces algébriques sur $S$.
Considérons une suite exacte courte
$$
0 \to \mathcal{I} \to \mathcal{A} \to \mathcal{O}_X \to 0
$$
de faisceaux sur $X_\etale$, où $\mathcal{A}$ est un faisceau de
$f^{-1}\mathcal{O}_B$-algèbres, $\mathcal{A} \to \mathcal{O}_X$ est une surjection
de faisceaux de $f^{-1}\mathcal{O}_B$-algèbres, et $\mathcal{I}$ son noyau.
Si
\begin{enumerate}
\item $\mathcal{I}$ est un idéal de carré nul dans $\mathcal{A}$, et
\item $\mathcal{I}$ est quasi-cohérent en tant que $\mathcal{O}_X$-module,
\end{enumerate}
alors il existe un épaississement du premier ordre
$X \subset X'$ sur $B$ et un isomorphisme
$\mathcal{O}_{X'} \to \mathcal{A}$ de $f^{-1}\mathcal{O}_B$-algèbres
compatible aux surjections vers $\mathcal{O}_X$.
\end{lemma}

\begin{proof}
Dans cette démonstration, nous reprenons certains des arguments employés dans les
démonstrations des
Lemmes \ref{lemma-first-order-thickening-surjective} et
\ref{lemma-thickening-scheme}.
Nous traitons d'abord le cas $B = S = \Spec(\mathbf{Z})$.
Soit $U$ un schéma affine et soit $U \to X$ étale.
Alors
$$
0 \to \mathcal{I}(U) \to \mathcal{A}(U) \to \mathcal{O}_X(U) \to 0
$$
est exacte, car $H^1(U_\etale, \mathcal{I}) = 0$ puisque
$\mathcal{I}$ est quasi-cohérent, voir
Descente, Proposition \ref{descent-proposition-same-cohomology-quasi-coherent},
et Cohomologie des schémas, Lemme
\ref{coherent-lemma-quasi-coherent-affine-cohomology-zero}.
Si $V \to U$ est un morphisme d'objets affines de $X_{spaces, \etale}$,
alors
$$
\mathcal{I}(V) = \mathcal{I}(U) \otimes_{\mathcal{O}_X(U)} \mathcal{O}_X(V)
$$
puisque $\mathcal{I}$ est un $\mathcal{O}_X$-module quasi-cohérent, voir
Descente, Proposition \ref{descent-proposition-equivalence-quasi-coherent}.
Ainsi $\mathcal{A}(U) \to \mathcal{A}(V)$ est un homomorphisme d'anneaux
étale, voir
Algèbre, Lemme \ref{algebra-lemma-lift-etale-infinitesimal}.
Nous voyons donc que
$$
U \longmapsto U' = \Spec(\mathcal{A}(U))
$$
est un foncteur de $X_{affine, \etale}$ vers la catégorie des schémas
affines et des morphismes étales. En fait, nous affirmons que ce foncteur peut
s'étendre en un foncteur $U \mapsto U'$ sur tout $X_\etale$.
Pour le voir, si $U$ est un objet de $X_\etale$, notons que
$$
0 \to \mathcal{I}|_{U_{Zar}} \to \mathcal{A}|_{U_{Zar}} \to
\mathcal{O}_X|_{U_{Zar}} \to 0
$$
et que $\mathcal{I}|_{U_{Zar}}$ est un faisceau quasi-cohérent sur $U$, voir
Descente,
Proposition \ref{descent-proposition-equivalence-quasi-coherent-functorial}.
Par conséquent, d'après
Compléments sur les morphismes, Lemme \ref{more-morphisms-lemma-first-order-thickening},
nous obtenons un épaississement du premier ordre de schémas $U \subset U'$ tel que
$\mathcal{O}_{U'}$ soit isomorphe à $\mathcal{A}|_{U_{Zar}}$. Il est clair que
cette construction est compatible avec la construction affine ci-dessus.

\medskip\noindent
Choisissons une présentation $X = U/R$, voir
Espaces, Définition \ref{spaces-definition-presentation},
de sorte que $s, t : R \to U$ définissent une relation d'équivalence étale.
En appliquant le foncteur ci-dessus, nous obtenons une relation d'équivalence
étale $s', t' : R' \to U'$ dans les schémas. Considérons l'espace algébrique
$X' = U'/R'$ (voir
Espaces, Théorème \ref{spaces-theorem-presentation}).
Le morphisme $X = U/R \to U'/R' = X'$ est un épaississement du premier ordre.
Considérons $\mathcal{O}_{X'}$ comme un faisceau sur $X_\etale$.
Par construction, nous avons un isomorphisme
$$
\gamma :
\mathcal{O}_{X'}|_{U_\etale}
\longrightarrow
\mathcal{A}|_{U_\etale}
$$
tel que $s^{-1}\gamma$ coïncide avec $t^{-1}\gamma$ sur $R_\etale$.
Par conséquent, d'après
Propriétés des espaces, Lemme \ref{spaces-properties-lemma-descent-sheaf},
cela implique que $\gamma$ provient d'un unique isomorphisme
$\mathcal{O}_{X'} \to \mathcal{A}$, comme voulu.

\medskip\noindent
Pour traiter le cas d'un espace algébrique de base quelconque $B$, nous
construisons d'abord $X'$ comme espace algébrique sur $\mathbf{Z}$, ainsi que ci-dessus.
Nous utilisons ensuite l'isomorphisme $\mathcal{O}_{X'} \to \mathcal{A}$ pour
définir $f^{-1}\mathcal{O}_B \to \mathcal{O}_{X'}$. D'après le
Lemme \ref{lemma-first-order-thickening-maps},
cela définit un morphisme $X' \to B$ compatible avec le morphisme donné
$X \to B$, ce qui achève la démonstration.
\end{proof}

\begin{lemma}
\label{lemma-base-change-thickening}
Soit $S$ un schéma. Soit $Y \subset Y'$ un épaississement d'espaces algébriques
sur $S$. Soit $X' \to Y'$ un morphisme et posons $X = Y \times_{Y'} X'$.
Alors $(X \subset X') \to (Y \subset Y')$
est un morphisme d'épaississements. Si $Y \subset Y'$ est un épaississement du premier
ordre (respectivement d'ordre fini), alors $X \subset X'$ est un épaississement du premier
ordre (respectivement d'ordre fini).
\end{lemma}

\begin{proof}
Omis.
\end{proof}

\begin{lemma}
\label{lemma-composition-thickening}
Soit $S$ un schéma. Si $X \subset X'$ et $X' \subset X''$ sont des
épaississements d'espaces algébriques sur $S$, alors $X \subset X''$ en est un également.
\end{lemma}

\begin{proof}
Omis.
\end{proof}

\begin{lemma}
\label{lemma-descending-property-thickening}
La propriété d'être un épaississement est locale pour la topologie fpqc.
Il en va de même pour les épaississements du premier ordre.
\end{lemma}

\begin{proof}
L'énoncé signifie ceci : soit $S$ un schéma et soit
$X \to X'$ un morphisme d'espaces algébriques sur $S$.
Soit $\{g_i : X'_i \to X'\}$ un recouvrement fpqc d'espaces algébriques
tel que le changement de base $X_i \to X'_i$ soit un épaississement pour tout $i$.
Alors $X \to X'$ est un épaississement. Comme les morphismes $g_i$ sont conjointement
surjectifs, nous concluons que $X \to X'$ est surjectif. D'après
Descente sur les espaces, Lemme
\ref{spaces-descent-lemma-descending-property-closed-immersion},
nous concluons que $X \to X'$ est une immersion fermée.
Ainsi $X \to X'$ est un épaississement. Nous omettons la démonstration dans le
cas des épaississements du premier ordre.
\end{proof}








\section{Morphismes d'épaississements}
\label{section-morphisms-thickenings}

\noindent
Si $(f, f') : (X \subset X') \to (Y \subset Y')$ est un morphisme
d'épaississements d'espaces algébriques, les propriétés de $f$
se transmettent souvent à $f'$. Il existe plusieurs variantes.

\begin{lemma}
\label{lemma-thicken-property-morphisms}
Soit $S$ un schéma. Soit $(f, f') : (X \subset X') \to (Y \subset Y')$
un morphisme d'épaississements d'espaces algébriques sur $S$. Alors
\begin{enumerate}
\item $f$ est un morphisme affine si et seulement si $f'$ est un morphisme affine,
\item $f$ est un morphisme surjectif si et seulement si $f'$ est un morphisme surjectif,
\item $f$ est quasi-compact si et seulement si $f'$ est quasi-compact,
\item $f$ est universellement fermé si et seulement si $f'$ est universellement fermé,
\item $f$ est entier si et seulement si $f'$ est entier,
\item $f$ est (quasi-)séparé si et seulement si $f'$ est (quasi-)séparé,
\item $f$ est universellement injectif si et seulement si $f'$ est universellement injectif,
\item $f$ est universellement ouvert si et seulement si $f'$ est universellement ouvert,
\item $f$ est représentable si et seulement si $f'$ est représentable, et
\item ajouter d'autres propriétés ici.
\end{enumerate}
\end{lemma}

\begin{proof}
Observons que $Y \to Y'$ et $X \to X'$ sont des morphismes entiers et des
homéomorphismes universels. Cela entraîne immédiatement les assertions
(2), (3), (4), (7) et (8).
L'assertion (1) découle de
Limites d'espaces, Proposition \ref{spaces-limits-proposition-affine},
qui nous dit qu'il existe une correspondance bijective entre les
schémas affines étales sur $X$ et sur $X'$, et entre les schémas affines
étales sur $Y$ et sur $Y'$.
L'assertion (5) découle de (1) et (4) par
Morphismes d'espaces, Lemme
\ref{spaces-morphisms-lemma-integral-universally-closed}.
Enfin, notons que
$$
X \times_Y X = X \times_{Y'} X \to X \times_{Y'} X' \to X' \times_{Y'} X'
$$
est un épaississement (les deux flèches sont des épaississements par le
Lemme \ref{lemma-base-change-thickening}).
Ainsi, en appliquant (3) et (4) au morphisme
$(X \subset X') \to (X \times_Y X \to X' \times_{Y'} X')$
nous obtenons (6). Enfin, l'assertion (9) découle du fait qu'un espace
algébrique qui est un épaississement d'un schéma est encore un schéma, voir le
Lemme \ref{lemma-thickening-scheme}.
\end{proof}

\begin{lemma}
\label{lemma-thicken-property-morphisms-cartesian}
Soit $S$ un schéma. Soit $(f, f') : (X \subset X') \to (Y \subset Y')$ un
morphisme d'épaississements d'espaces algébriques sur $S$ tel que
$X = Y \times_{Y'} X'$. Si $X \subset X'$ est un épaississement d'ordre fini, alors
\begin{enumerate}
\item $f$ est une immersion fermée si et seulement si $f'$ est une immersion fermée,
\item $f$ est localement de type fini si et seulement si $f'$ est
localement de type fini,
\item $f$ est localement quasi-fini si et seulement si $f'$ est localement
quasi-fini,
\item $f$ est localement de type fini de dimension relative $d$ si et
seulement si $f'$ est localement de type fini de dimension relative $d$,
\item $\Omega_{X/Y} = 0$ si et seulement si $\Omega_{X'/Y'} = 0$,
\item $f$ est non ramifié si et seulement si $f'$ est non ramifié,
\item $f$ est propre si et seulement si $f'$ est propre,
\item $f$ est un morphisme fini si et seulement si $f'$ est un morphisme fini,
\item $f$ est un monomorphisme si et seulement si $f'$ est un monomorphisme,
\item $f$ est une immersion si et seulement si $f'$ est une immersion, et
\item ajouter d'autres propriétés ici.
\end{enumerate}
\end{lemma}

\begin{proof}
Choisissons un schéma $V'$ et un morphisme étale surjectif $V' \to Y'$.
Choisissons un schéma $U'$ et un morphisme étale surjectif
$U' \to X' \times_{Y'} V'$. Posons $V = Y \times_{Y'} V'$ et
$U = X \times_{X'} U'$. Pour les propriétés des morphismes qui sont locales
pour la topologie étale, nous pouvons alors nous ramener au morphisme d'épaississements de schémas
$(U \subset U') \to (V \subset V')$ et appliquer Compléments sur les morphismes, Lemme
\ref{more-morphisms-lemma-thicken-property-morphisms-cartesian}.
Cela démontre (2), (3), (4), (5) et (6).

\medskip\noindent
Les propriétés des morphismes figurant en (1), (7), (8), (9) et (10) sont stables
par changement de base ; par conséquent, si $f'$ possède la propriété $\mathcal{P}$, alors
$f$ la possède également. Voir
Espaces, Lemme \ref{spaces-lemma-base-change-immersions},
et
Morphismes d'espaces, Lemmes
\ref{spaces-morphisms-lemma-base-change-proper},
\ref{spaces-morphisms-lemma-base-change-integral}, et
\ref{spaces-morphisms-lemma-base-change-monomorphism}.

\medskip\noindent
Le sens intéressant de (1), (7), (8), (9) et (10) consiste à supposer
que $f$ possède la propriété et à en déduire que $f'$ la possède aussi.
Par récurrence sur l'ordre de l'épaississement, nous pouvons
supposer que $Y \subset Y'$ est un épaississement du premier ordre, voir la
discussion ci-dessus sur les épaississements d'ordre fini.

\medskip\noindent
Démonstration de (1). Choisissons un schéma $V'$ et un morphisme étale surjectif
$V' \to Y'$. Posons $V = Y \times_{Y'} V'$, $U' = X' \times_{Y'} V'$ et
$U = X \times_Y V$. Alors $U \to V$ est une immersion fermée, ce qui
implique que $U$ est un schéma, et entraîne à son tour que $U'$ est
un schéma (Lemme \ref{lemma-thickening-scheme}). Nous pouvons donc appliquer
le lemme dans le cas des schémas
(Compléments sur les morphismes, Lemme
\ref{more-morphisms-lemma-thicken-property-morphisms-cartesian})
à $(U \subset U') \to (V \subset V')$ pour conclure.

\medskip\noindent
Démonstration de (7). Elle résulte de la combinaison de (2) avec les
résultats du Lemme \ref{lemma-thicken-property-morphisms}
et du fait qu'être propre équivaut à être quasi-compact $+$
séparé $+$ localement de type fini $+$ universellement fermé.

\medskip\noindent
Démonstration de (8). Elle résulte de la combinaison de (2) avec les
résultats du Lemme \ref{lemma-thicken-property-morphisms}
et du fait qu'être fini équivaut à être entier $+$ localement
de type fini (Morphismes, Lemme \ref{morphisms-lemma-finite-integral}).

\medskip\noindent
Démonstration de (9). Comme $f$ est un monomorphisme, nous avons $X = X \times_Y X$.
Nous pouvons appliquer les résultats déjà démontrés au morphisme
d'épaississements $(X \subset X') \to (X \times_Y X \subset X' \times_{Y'} X')$.
Nous concluons par (1) que $X' \to X' \times_{Y'} X'$ est une immersion fermée.
En fait, c'est un épaississement du premier ordre, car l'idéal qui définit
l'immersion fermée $X' \to X' \times_{Y'} X'$ est contenu dans l'image inverse
de l'idéal $\mathcal{I} \subset \mathcal{O}_{Y'}$ qui définit $Y$ dans $Y'$.
En effet, $X = X \times_Y X = (X' \times_{Y'} X') \times_{Y'} Y$ est contenu
dans $X'$. Le faisceau conormal de l'immersion fermée
$\Delta : X' \to X' \times_{Y'} X'$ est égal à $\Omega_{X'/Y'}$
(c'est l'analogue de
Morphismes, Lemme \ref{morphisms-lemma-differentials-diagonal},
pour les espaces algébriques, et cela découle soit d'une localisation étale,
soit de la combinaison des
Lemmes \ref{lemma-differentials-relative-immersion-section} et
\ref{lemma-differential-product} ; certains détails sont omis).
Il suffit donc de montrer que $\Omega_{X'/Y'} = 0$, ce qui découle de (5)
et de l'énoncé correspondant pour $X/Y$.

\medskip\noindent
Démonstration de (10). Si $f : X \to Y$ est une immersion, elle se factorise sous la forme
$X \to V \to Y$, où $V \to Y$ est un sous-espace ouvert et $X \to V$ une
immersion fermée, voir
Morphismes d'espaces, Remarque \ref{spaces-morphisms-remark-immersion}.
Soit $V' \subset Y'$ le sous-espace ouvert dont
l'espace topologique sous-jacent $|V'|$ est égal à $|V| \subset |Y| = |Y'|$.
Alors $X' \to Y'$ se factorise par $V'$, et nous concluons par (1) que $X' \to V'$
est une immersion fermée. Cela achève la démonstration.
\end{proof}

\noindent
Le lemme suivant est une variante du précédent. Au lieu de supposer
que les épaississements en jeu sont d'ordre fini (ce qui permet de transférer
de $f$ à $f'$ la propriété d'être localement de type fini),
nous supposons plutôt que $f$ et $f'$ sont tous deux localement
de type fini.

\begin{lemma}
\label{lemma-properties-that-extend-over-thickenings}
Soit $S$ un schéma. Soit $(f, f') : (X \subset X') \to (Y \to Y')$ un
morphisme d'épaississements d'espaces algébriques sur $S$. Supposons que $f$ et $f'$
soient localement de type fini et que $X = Y \times_{Y'} X'$. Alors
\begin{enumerate}
\item $f$ est localement quasi-fini si et seulement si $f'$ est localement quasi-fini,
\item $f$ est fini si et seulement si $f'$ est fini,
\item $f$ est une immersion fermée si et seulement si $f'$ est une immersion fermée,
\item $\Omega_{X/Y} = 0$ si et seulement si $\Omega_{X'/Y'} = 0$,
\item $f$ est non ramifié si et seulement si $f'$ est non ramifié,
\item $f$ est un monomorphisme si et seulement si $f'$ est un monomorphisme,
\item $f$ est une immersion si et seulement si $f'$ est une immersion,
\item $f$ est propre si et seulement si $f'$ est propre, et
\item ajouter d'autres propriétés ici.
\end{enumerate}
\end{lemma}

\begin{proof}
Choisissons un schéma $V'$ et un morphisme étale surjectif $V' \to Y'$.
Choisissons un schéma $U'$ et un morphisme étale surjectif
$U' \to X' \times_{Y'} V'$. Posons $V = Y \times_{Y'} V'$ et
$U = X \times_{X'} U'$. Pour les propriétés des morphismes qui sont locales
pour la topologie étale, nous pouvons alors nous ramener au morphisme d'épaississements de schémas
$(U \subset U') \to (V \subset V')$ et appliquer
Compléments sur les morphismes, Lemme
\ref{more-morphisms-lemma-properties-that-extend-over-thickenings}.
Cela démontre (1), (4) et (5).

\medskip\noindent
Les propriétés figurant en (2), (3), (6), (7) et (8) sont stables
par changement de base ; par conséquent, si $f'$ possède la propriété $\mathcal{P}$, alors
$f$ la possède également. Voir Espaces, Lemme \ref{spaces-lemma-base-change-immersions},
et
Morphismes d'espaces, Lemmes
\ref{spaces-morphisms-lemma-base-change-proper},
\ref{spaces-morphisms-lemma-base-change-integral}, et
\ref{spaces-morphisms-lemma-base-change-monomorphism}.
Ainsi, dans chaque cas, il suffit de démontrer que, si $f$ possède
la propriété voulue, alors $f'$ la possède aussi.

\medskip\noindent
Le cas (2) découle du cas (5) du Lemme \ref{lemma-thicken-property-morphisms}
et du fait que les morphismes finis sont exactement
les morphismes entiers qui sont localement de type fini
(Morphismes d'espaces, Lemme \ref{spaces-morphisms-lemma-finite-integral}).

\medskip\noindent
Cas (3). Cela découle immédiatement de
Limites d'espaces, Lemme
\ref{spaces-limits-lemma-check-closed-infinitesimally}.

\medskip\noindent
Démonstration de (6). Comme $f$ est un monomorphisme, nous avons $X = X \times_Y X$.
Nous pouvons appliquer les résultats déjà démontrés au morphisme d'épaississements
$(X \subset X') \to (X \times_Y X \subset X' \times_{Y'} X')$.
Nous concluons par (3) que $\Delta_{X'/Y'} : X' \to X' \times_{Y'} X'$
est une immersion fermée. En fait, $\Delta_{X'/Y'}$ induit une bijection
$|X'| \to |X' \times_{Y'} X'|$, donc $\Delta_{X'/Y'}$ est un épaississement.
D'autre part, $\Delta_{X'/Y'}$ est localement de présentation finie d'après
Morphismes d'espaces, Lemme
\ref{spaces-morphisms-lemma-diagonal-morphism-finite-type}.
Autrement dit, $\Delta_{X'/Y'}(X')$ est défini par
un faisceau quasi-cohérent d'idéaux
$\mathcal{J} \subset \mathcal{O}_{X' \times_{Y'} X'}$ de type fini.
Comme $\Omega_{X'/Y'} = 0$ par (5), nous voyons que
le faisceau conormal de $X' \to X' \times_{Y'} X'$ est nul.
(Le faisceau conormal de l'immersion fermée $\Delta_{X'/Y'}$ est égal à
$\Omega_{X'/Y'}$ ; c'est l'analogue de
Morphismes, Lemme \ref{morphisms-lemma-differentials-diagonal},
pour les espaces algébriques, et cela découle soit d'une localisation étale,
soit de la combinaison des
Lemmes \ref{lemma-differentials-relative-immersion-section} et
\ref{lemma-differential-product} ; certains détails sont omis.)
Autrement dit, $\mathcal{J}/\mathcal{J}^2 = 0$.
Cela implique que $\Delta_{X'/Y'}$ est un isomorphisme, par exemple
par Algèbre, Lemme \ref{algebra-lemma-ideal-is-squared-union-connected}.

\medskip\noindent
Démonstration de (7). Si $f : X \to Y$ est une immersion, elle se factorise sous la forme
$X \to V \to Y$, où $V \to Y$ est un sous-espace ouvert et $X \to V$ une
immersion fermée, voir
Morphismes d'espaces, Remarque \ref{spaces-morphisms-remark-immersion}.
Soit $V' \subset Y'$ le sous-espace ouvert dont
l'espace topologique sous-jacent $|V'|$ est égal à $|V| \subset |Y| = |Y'|$.
Alors $X' \to Y'$ se factorise par $V'$, et nous concluons par (3) que $X' \to V'$
est une immersion fermée.

\medskip\noindent
Le cas (8) découle du Lemme \ref{lemma-thicken-property-morphisms}
et de la définition des morphismes propres comme morphismes quasi-compacts,
universellement fermés et séparés qui sont localement de type fini.
\end{proof}










\section{Groupes de Picard des épaississements}
\label{section-picard-group-thickening}

\noindent
Quelques éléments sur les groupes de Picard des épaississements.

\begin{lemma}
\label{lemma-picard-group-first-order-thickening}
Soit $S$ un schéma. Soit $X \subset X'$ un épaississement du premier ordre
d'espaces algébriques sur $S$, dont le faisceau d'idéaux est $\mathcal{I}$.
Il existe alors une suite exacte canonique
$$
\xymatrix{
0 \ar[r] &
H^0(X, \mathcal{I}) \ar[r] &
H^0(X', \mathcal{O}_{X'}^*) \ar[r] &
H^0(X, \mathcal{O}^*_X) \ar `r[d] `d[l] `l[llld] `d[dll] [dll] \\
& H^1(X, \mathcal{I}) \ar[r] &
\Pic(X') \ar[r] &
\Pic(X) \ar `r[d] `d[l] `l[llld] `d[dll] [dll] \\
& H^2(X, \mathcal{I}) \ar[r] & \ldots \ar[r] & \ldots
}
$$
de groupes abéliens.
\end{lemma}

\begin{proof}
Rappelons que $X_\etale = X'_\etale$, voir le
Lemme \ref{lemma-thickening-equivalence} et, plus généralement, la
discussion de la section \ref{section-thickenings}.
La suite du lemme est la suite exacte longue de cohomologie
associée à la suite exacte courte de faisceaux de groupes abéliens
$$
0 \to \mathcal{I} \to \mathcal{O}_{X'}^* \to \mathcal{O}_X^* \to 0
$$
sur $X_\etale$, où le premier morphisme envoie une section locale $f$ de $\mathcal{I}$
sur la section inversible $1 + f$ de $\mathcal{O}_{X'}$.
Nous utilisons aussi l'identification du groupe de Picard d'un
site annelé au premier groupe de cohomologie du faisceau
des éléments inversibles, voir
Cohomologie sur les sites, Lemme \ref{sites-cohomology-lemma-h1-invertible}.
\end{proof}








\section{Voisinages infinitésimaux}
\label{section-first-order-infinitesimal-neighbourhood}

\noindent
Une construction naturelle d'épaississements d'ordre fini est la suivante.
Supposons que $i : Z \to X$ soit une immersion d'espaces algébriques. Choisissons un
sous-espace ouvert $U \subset X$ tel que $i$ identifie $Z$ au sous-espace
fermé $Z \subset U$ (voir
Morphismes d'espaces, Remarque \ref{spaces-morphisms-remark-immersion}).
Soit $\mathcal{I} \subset \mathcal{O}_U$ le
faisceau quasi-cohérent d'idéaux qui définit $Z$ dans $U$, voir
Morphismes d'espaces,
Lemme \ref{spaces-morphisms-lemma-closed-immersion-ideals}.
Pour $n \geq 1$, nous pouvons considérer le sous-espace fermé $Z_n \subset U$
défini par le faisceau quasi-cohérent d'idéaux $\mathcal{I}^{n + 1}$.

\begin{definition}
\label{definition-first-order-infinitesimal-neighbourhood}
Soit $i : Z \to X$ une immersion d'espaces algébriques.
\begin{enumerate}
\item Le {\it voisinage infinitésimal du premier ordre} de $Z$ dans $X$ est
l'épaississement du premier ordre $Z \subset Z_1$ sur $X$ décrit ci-dessus.
\item Le {\it voisinage infinitésimal d'ordre $n$} de $Z$ dans $X$ est
l'épaississement d'ordre $n$ $Z \subset Z_n$ sur $X$ décrit ci-dessus.
\end{enumerate}
\end{definition}

\noindent
Cet épaississement possède la propriété universelle suivante (ce qui dissipera
toute crainte que la construction ci-dessus ne dépende du choix de l'ouvert
$U$).

\begin{lemma}
\label{lemma-first-order-infinitesimal-neighbourhood}
Soit $i : Z \to X$ une immersion d'espaces algébriques.
\begin{enumerate}
\item Le voisinage infinitésimal du premier ordre $Z'$ de $Z$ dans $X$
possède la propriété universelle suivante :
pour tout diagramme commutatif
$$
\xymatrix{
Z \ar[d]_i & T \ar[l]^a \ar[d] \\
X & T' \ar[l]_b
}
$$
où $T \subset T'$ est un épaississement du premier ordre sur $X$, il existe
un unique morphisme $(a', a) : (T \subset T') \to (Z \subset Z')$
d'épaississements sur $X$.
\item Pour $n \geq 1$, le voisinage infinitésimal d'ordre $n$
$Z_n$ de $Z$ dans $X$ possède la propriété universelle suivante :
pour tout diagramme commutatif
$$
\xymatrix{
Z \ar[d]_i & T \ar[l]^a \ar[d] \\
X & T' \ar[l]_b
}
$$
où $T \subset T'$ est un épaississement d'ordre $n$ sur $X$, il existe
un unique morphisme $(a', a) : (T \subset T') \to (Z \subset Z_n)$
d'épaississements sur $X$.
\end{enumerate}
\end{lemma}

\begin{proof}
Nous ne démontrerons que (1).
Soit $U \subset X$ le sous-espace ouvert utilisé dans la construction de $Z'$,
c'est-à-dire un ouvert tel que $Z$ soit identifié à un sous-espace fermé de $U$
défini par le faisceau quasi-cohérent d'idéaux $\mathcal{I}$.
Comme $|T| = |T'|$, nous voyons que $|b|(|T'|) \subset |U|$. Nous pouvons donc
considérer $b$ comme un morphisme à valeurs dans $U$, voir
Propriétés des espaces,
Lemme \ref{spaces-properties-lemma-factor-through-open-subspace}.
Soit $\mathcal{J} \subset \mathcal{O}_{T'}$
le faisceau quasi-cohérent d'idéaux de carré nul qui définit $T$.
La commutativité du diagramme donne $b|_T = i \circ a$, où
$i : Z \to U$ est l'immersion fermée. Nous en déduisons que
$b^\sharp(b^{-1}\mathcal{I}) \subset \mathcal{J}$ d'après
Morphismes d'espaces,
Lemme \ref{spaces-morphisms-lemma-closed-immersion-ideals}.
Comme $T'$ est un épaississement du premier ordre de $T$, nous avons $\mathcal{J}^2 = 0$,
donc $b^\sharp(b^{-1}(\mathcal{I}^2)) = 0$. D'après
Morphismes d'espaces, Lemme \ref{spaces-morphisms-lemma-closed-immersion-ideals},
il s'ensuit que $b$ se factorise par $Z'$. En notant $a' : T' \to Z'$
cette factorisation, nous concluons.
\end{proof}

\begin{lemma}
\label{lemma-infinitesimal-neighbourhood-conormal}
Soit $i : Z \to X$ une immersion d'espaces algébriques.
Soit $Z \subset Z'$ le voisinage infinitésimal du premier ordre
de $Z$ dans $X$. Alors le diagramme
$$
\xymatrix{
Z \ar[r] \ar[d] & Z' \ar[d] \\
Z \ar[r] & X
}
$$
induit un morphisme de faisceaux conormaux
$\mathcal{C}_{Z/X} \to \mathcal{C}_{Z/Z'}$ d'après le
Lemme \ref{lemma-conormal-functorial}.
Ce morphisme est un isomorphisme.
\end{lemma}

\begin{proof}
Cela résulte immédiatement de la construction de $Z'$ ci-dessus.
\end{proof}





\section{Transformations formellement lisses, étales et non ramifiées}
\sectionmark{Transformations lisses, étales et non ramifiées}
\label{section-formally-smooth-etale-unramified}

\noindent
Rappelons qu'un homomorphisme d'anneaux $R \to A$ est dit
{\it formellement lisse}, resp.\ {\it formellement étale},
resp.\ {\it formellement non ramifié}
(voir Algèbre, Définition \ref{algebra-definition-formally-smooth},
resp.\ Définition \ref{algebra-definition-formally-etale},
resp.\ Définition \ref{algebra-definition-formally-unramified})
si, pour tout diagramme commutatif de flèches pleines
$$
\xymatrix{
A \ar[r] \ar@{-->}[rd] & B/I \\
R \ar[r] \ar[u] & B \ar[u]
}
$$
où $I \subset B$ est un idéal de carré nul, il
existe une flèche pointillée, resp.\ il en existe une unique, resp.\ il en existe
au plus une, qui rend le diagramme commutatif. Cela motive
l'analogue suivant pour les morphismes d'espaces algébriques et, plus
généralement, pour les foncteurs.

\begin{definition}
\label{definition-formally-smooth-etale-unramified}
Soit $S$ un schéma.
Soit $a : F \to G$ une transformation de foncteurs
$F, G : (\Sch/S)_{fppf}^{opp} \to \textit{Ens}$.
Considérons les diagrammes commutatifs de flèches pleines de la forme
$$
\xymatrix{
F \ar[d]_a & T \ar[d]^i \ar[l] \\
G & T' \ar[l] \ar@{-->}[lu]
}
$$
où $T$ et $T'$ sont des schémas affines et $i$ une immersion fermée
définie par un idéal de carré nul.
\begin{enumerate}
\item On dit que $a$ est {\it formellement lisse} si, pour tout
diagramme de flèches pleines comme ci-dessus, il existe une flèche pointillée qui rend le diagramme
commutatif\footnote{Ce n'est là qu'une des définitions possibles.
Une condition légèrement plus faible consisterait à exiger que
la flèche pointillée existe localement pour la topologie fppf sur $T'$. Cette notion plus faible
possède, en un certain sens, de meilleures propriétés formelles.}.
\item On dit que $a$ est {\it formellement étale} si, pour tout
diagramme de flèches pleines comme ci-dessus, il existe exactement une flèche pointillée qui rend le diagramme
commutatif.
\item On dit que $a$ est {\it formellement non ramifiée} si, pour tout
diagramme de flèches pleines comme ci-dessus, il existe au plus une flèche pointillée qui rend le diagramme
commutatif.
\end{enumerate}
\end{definition}

\begin{lemma}
\label{lemma-formally-etale-is-combination}
Soit $S$ un schéma.
Soit $a : F \to G$ une transformation de foncteurs
$F, G : (\Sch/S)_{fppf}^{opp} \to \textit{Ens}$.
Alors $a$ est formellement étale si et seulement si $a$ est à la fois formellement
lisse et formellement non ramifiée.
\end{lemma}

\begin{proof}
Cela résulte formellement de la définition.
\end{proof}

\begin{lemma}
\label{lemma-composition-formally-smooth-etale-unramified}
Composition.
\begin{enumerate}
\item Un composé de transformations de foncteurs formellement lisses est formellement
lisse.
\item Un composé de transformations de foncteurs formellement étales est formellement
étale.
\item Un composé de transformations de foncteurs formellement non ramifiées est
formellement non ramifié.
\end{enumerate}
\end{lemma}

\begin{proof}
C'est formel.
\end{proof}

\begin{lemma}
\label{lemma-base-change-formally-smooth-etale-unramified}
Soit $S$ un schéma appartenant à $\Sch_{fppf}$.
Soient $F, G, H : (\Sch/S)_{fppf}^{opp} \to \textit{Ens}$.
Soient $a : F \to G$, $b : H \to G$ des transformations de foncteurs.
Considérons le diagramme de produit fibré
$$
\xymatrix{
H \times_{b, G, a} F \ar[r]_-{b'} \ar[d]_{a'} & F \ar[d]^a \\
H \ar[r]^b & G
}
$$
\begin{enumerate}
\item Si $a$ est formellement lisse, alors le changement de base $a'$ est
formellement lisse.
\item Si $a$ est formellement étale, alors le changement de base $a'$ est
formellement étale.
\item Si $a$ est formellement non ramifiée, alors le changement de base $a'$ est
formellement non ramifiée.
\end{enumerate}
\end{lemma}

\begin{proof}
C'est formel.
\end{proof}

\begin{lemma}
\label{lemma-representable-property-formally-property}
Soit $S$ un schéma.
Soient $F, G : (\Sch/S)_{fppf}^{opp} \to \textit{Ens}$.
Soit $a : F \to G$ une transformation de foncteurs représentable.
\begin{enumerate}
\item Si $a$ est lisse, alors $a$ est formellement lisse.
\item Si $a$ est étale, alors $a$ est formellement étale.
\item Si $a$ est non ramifiée, alors $a$ est formellement non ramifiée.
\end{enumerate}
\end{lemma}

\begin{proof}
Considérons un diagramme commutatif de flèches pleines
$$
\xymatrix{
F \ar[d]_a & T \ar[d]^i \ar[l] \\
G & T' \ar[l] \ar@{-->}[lu]
}
$$
comme dans la
Définition \ref{definition-formally-smooth-etale-unramified}.
Alors $F \times_G T'$ est un schéma lisse (resp.\ étale, resp.\ non ramifié)
sur $T'$. Par conséquent, d'après
Compléments sur les morphismes, Lemme \ref{more-morphisms-lemma-smooth-formally-smooth}
(resp.\ Lemme \ref{more-morphisms-lemma-etale-formally-etale},
resp.\ Lemme \ref{more-morphisms-lemma-unramified-formally-unramified}),
nous pouvons tracer (resp.\ tracer de manière unique, resp.\ tracer d'au plus
une manière) la flèche pointillée du diagramme
$$
\xymatrix{
F \times_G T' \ar[d] & T \ar[d]^i \ar[l] \\
T' & T' \ar[l] \ar@{-->}[lu]
}
$$
et nous obtenons donc aussi l'assertion correspondante dans le premier diagramme.
\end{proof}

\begin{lemma}
\label{lemma-etale-on-top}
Soit $S$ un schéma appartenant à $\Sch_{fppf}$.
Soient $F, G, H : (\Sch/S)_{fppf}^{opp} \to \textit{Ens}$.
Soient $a : F \to G$, $b : G \to H$ des transformations de foncteurs.
Supposons que $a$ soit représentable, surjective et étale.
\begin{enumerate}
\item Si $b$ est formellement lisse, alors $b \circ a$ est formellement lisse.
\item Si $b$ est formellement étale, alors $b \circ a$ est formellement étale.
\item Si $b$ est formellement non ramifiée, alors $b \circ a$ est formellement non ramifiée.
\end{enumerate}
Réciproquement, considérons un diagramme commutatif de flèches pleines
$$
\xymatrix{
G \ar[d]_b & T \ar[d]^i \ar[l] \\
H & T' \ar[l] \ar@{-->}[lu]
}
$$
où $T'$ est un schéma affine sur $S$
et $i : T \to T'$ une immersion fermée définie par un idéal de carré nul.
\begin{enumerate}
\item[(4)] Si $b \circ a$ est formellement lisse, alors, pour tout $t \in T$,
il existe un morphisme étale de schémas affines $U' \to T'$ et un morphisme
$U' \to G$ tels que
$$
\xymatrix{
G \ar[d]_b & T \ar[l] & T \times_{T'} U' \ar[d] \ar[l]\\
H & T' \ar[l] & U' \ar[llu] \ar[l]
}
$$
le diagramme soit commutatif et $t$ appartienne à l'image de $U' \to T'$.
\item[(5)] Si $b \circ a$ est formellement non ramifiée, alors il existe au plus
une flèche pointillée dans le diagramme ci-dessus, c'est-à-dire que $b$ est formellement non ramifiée.
\item[(6)] Si $b \circ a$ est formellement étale, alors il existe exactement une
flèche pointillée dans le diagramme ci-dessus, c'est-à-dire que $b$ est formellement étale.
\end{enumerate}
\end{lemma}

\begin{proof}
Supposons que $b$ soit formellement lisse (resp.\ formellement étale,
resp.\ formellement non ramifiée). Puisqu'un morphisme étale est à la fois
lisse et non ramifié, nous voyons que $a$ est représentable et lisse
(resp.\ étale, resp. non ramifiée). Les assertions (1), (2) et (3)
résultent donc de la combinaison du
Lemme \ref{lemma-representable-property-formally-property}
et du
Lemme \ref{lemma-composition-formally-smooth-etale-unramified}.

\medskip\noindent
Supposons que $b \circ a$ soit formellement lisse. Considérons un diagramme
comme dans l'énoncé du lemme. Posons $W = F \times_G T$.
Par hypothèse, $W$ est un schéma étale surjectif sur $T$. D'après
Morphismes étales, Théorème \ref{etale-theorem-remarkable-equivalence},
il existe un schéma $W'$ étale sur $T'$ tel que $W = T \times_{T'} W'$.
Choisissons un sous-schéma ouvert affine $U' \subset W'$ tel que $t$ appartienne à
l'image de $U' \to T'$. Puisque $b \circ a$ est formellement
lisse, il existe un morphisme $U' \to F$ tel que
$$
\xymatrix{
F \ar[d]_{b \circ a} & W \ar[l] & T \times_{T'} U' \ar[d] \ar[l]\\
H & T' \ar[l] & U' \ar[llu] \ar[l]
}
$$
le diagramme soit commutatif. Le composé $U' \to F \to G$ fournit un
morphisme comme dans l'assertion (5) du lemme.

\medskip\noindent
Supposons que $f, g : T' \to G$ soient deux flèches pointillées qui complètent le
diagramme du lemme. Posons $W = F \times_G T$.
Par hypothèse, $W$ est un schéma étale surjectif sur $T$. D'après
Morphismes étales, Théorème \ref{etale-theorem-remarkable-equivalence},
il existe un schéma $W'$ étale sur $T'$ tel que $W = T \times_{T'} W'$.
Comme $a$ est formellement étale, les composés
$$
W' \to T' \xrightarrow{f} G
\quad\text{et}\quad
W' \to T' \xrightarrow{g} G
$$
se relèvent en des morphismes $f', g' : W' \to F$ (on relève sur les ouverts affines et l'on recolle par
unicité). Si maintenant $b \circ a : F \to H$ est formellement non ramifiée, alors
$f' = g'$ et donc $f = g$, puisque $W' \to T'$ est un recouvrement étale. Cela démontre
l'assertion (6) du lemme.

\medskip\noindent
Supposons que $b \circ a$ soit formellement étale. D'après l'assertion (4), nous
pouvons trouver, localement pour la topologie étale sur $T'$, une flèche pointillée qui complète le diagramme,
et cette flèche est unique d'après l'assertion (5). Nous pouvons donc recoller les
solutions locales pour obtenir l'assertion (6). Certains détails sont omis.
\end{proof}

\begin{remark}
\label{remark-tempting}
On pourrait être tenté de penser que, dans la situation du
Lemme \ref{lemma-etale-on-top},
nous avons
« $b$ formellement lisse » $\Leftrightarrow$ « $b \circ a$ formellement lisse ».
Cependant, cela est vraisemblablement faux en général.
\end{remark}

\begin{lemma}
\label{lemma-formally-permanence}
Soit $S$ un schéma.
Soient $F, G, H : (\Sch/S)_{fppf}^{opp} \to \textit{Ens}$.
Soient $a : F \to G$, $b : G \to H$ des transformations de foncteurs.
Supposons que $b$ soit formellement non ramifiée.
\begin{enumerate}
\item Si $b \circ a$ est formellement non ramifiée, alors $a$ est formellement non ramifiée.
\item Si $b \circ a$ est formellement étale, alors $a$ est formellement étale.
\item Si $b \circ a$ est formellement lisse, alors $a$ est formellement lisse.
\end{enumerate}
\end{lemma}

\begin{proof}
Soit $T \subset T'$ une immersion fermée de schémas affines définie par un idéal
de carré nul. Soient $g' : T' \to G$ et $f : T \to F$ tels que
$g'|_T = a \circ f$. Comme $b$ est formellement non ramifiée, il existe une correspondance
bijective entre
$$
\{f' : T' \to F \mid f = f'|_T\text{ et }a \circ f' = g'\}
$$
et
$$
\{f' : T' \to F \mid f = f'|_T\text{ et }b \circ a \circ f' = b \circ g'\}.
$$
Le lemme en résulte formellement.
\end{proof}








\section{Morphismes formellement non ramifiés}
\label{section-formally-unramified}

\noindent
Dans cette section, nous précisons ce que signifie, pour un morphisme d'espaces algébriques,
d'être formellement non ramifié.

\begin{definition}
\label{definition-formally-unramified}
Soit $S$ un schéma. Un morphisme $f : X \to Y$ d'espaces algébriques sur $S$
est dit {\it formellement non ramifié} s'il est formellement non ramifié en tant que
transformation de foncteurs au sens de la
Définition \ref{definition-formally-smooth-etale-unramified}.
\end{definition}

\noindent
Nous ne répéterons pas les résultats démontrés dans le cadre plus général des
transformations de foncteurs formellement non ramifiées à la
section \ref{section-formally-smooth-etale-unramified}.
On peut caractériser cette propriété par l'annulation du
module des différentielles relatives, voir le
Lemme \ref{lemma-characterize-formally-unramified}.

\begin{lemma}
\label{lemma-formally-unramified}
Soit $S$ un schéma. Soit $f : X \to Y$ un morphisme d'espaces algébriques sur
$S$. Les assertions suivantes sont équivalentes :
\begin{enumerate}
\item $f$ est formellement non ramifié,
\item pour tout diagramme
$$
\xymatrix{
U \ar[d] \ar[r]_\psi & V \ar[d] \\
X \ar[r]^f & Y
}
$$
où $U$ et $V$ sont des schémas et les flèches verticales sont étales,
le morphisme de schémas $\psi$ est formellement non ramifié (au sens de
Compléments sur les morphismes,
Définition \ref{more-morphisms-definition-formally-unramified}), et
\item pour un tel diagramme dont les flèches verticales sont surjectives, le morphisme
$\psi$ est formellement non ramifié.
\end{enumerate}
\end{lemma}

\begin{proof}
Supposons $f$ formellement non ramifié. D'après le
Lemme \ref{lemma-representable-property-formally-property},
les morphismes $U \to X$ et $V \to Y$ sont formellement non ramifiés. D'après le
Lemme \ref{lemma-composition-formally-smooth-etale-unramified},
le composé $U \to Y$ est donc formellement non ramifié. Il résulte alors du
Lemme \ref{lemma-formally-permanence}
que $U \to V$ est formellement non ramifié. Ainsi, (1) implique (2), et (2)
implique trivialement (3).

\medskip\noindent
Supposons donné un diagramme comme en (3). D'après le
Lemme \ref{lemma-representable-property-formally-property},
le morphisme $V \to Y$ est formellement non ramifié. D'après le
Lemme \ref{lemma-composition-formally-smooth-etale-unramified},
le composé $U \to Y$ est donc formellement non ramifié. Il résulte alors du
Lemme \ref{lemma-etale-on-top}
que $X \to Y$ est formellement non ramifié, c'est-à-dire que (1) est vérifiée.
\end{proof}

\begin{lemma}
\label{lemma-formally-unramified-not-affine}
Soit $S$ un schéma.
Si $f : X \to Y$ est un morphisme formellement non ramifié d'espaces algébriques
sur $S$, alors, pour tout diagramme commutatif de flèches pleines
$$
\xymatrix{
X \ar[d]_f & T \ar[d]^i \ar[l] \\
Y & T' \ar[l] \ar@{-->}[lu]
}
$$
où $T \subset T'$ est un épaississement du premier ordre d'espaces algébriques
sur $S$, il existe au plus une flèche pointillée qui rend le diagramme commutatif.
Autrement dit, dans la
Définition \ref{definition-formally-unramified},
on peut supprimer la condition que $T$ soit un schéma affine.
\end{lemma}

\begin{proof}
Cela tient au fait qu'il existe un morphisme étale surjectif
$U' \to T'$, où $U'$ est une réunion disjointe de schémas affines (voir
Propriétés des espaces, Lemme
\ref{spaces-properties-lemma-cover-by-union-affines}),
et qu'un morphisme $T' \to X$ est déterminé par sa restriction à $U'$.
\end{proof}

\begin{lemma}
\label{lemma-composition-formally-unramified}
Un composé de morphismes formellement non ramifiés est formellement non ramifié.
\end{lemma}

\begin{proof}
C'est formel.
\end{proof}

\begin{lemma}
\label{lemma-base-change-formally-unramified}
Un changement de base d'un morphisme formellement non ramifié est formellement non ramifié.
\end{lemma}

\begin{proof}
C'est formel.
\end{proof}

\begin{lemma}
\label{lemma-characterize-formally-unramified}
Soit $S$ un schéma. Soit $f : X \to Y$ un morphisme d'espaces algébriques sur
$S$. Les assertions suivantes sont équivalentes :
\begin{enumerate}
\item $f$ est formellement non ramifié, et
\item $\Omega_{X/Y} = 0$.
\end{enumerate}
\end{lemma}

\begin{proof}
Cela résulte de la combinaison du
Lemme \ref{lemma-formally-unramified},
de Compléments sur les morphismes,
Lemme \ref{more-morphisms-lemma-formally-unramified-differentials},
et du
Lemme \ref{lemma-localize-differentials}.
\end{proof}

\begin{lemma}
\label{lemma-unramified-formally-unramified}
Soit $S$ un schéma.
Soit $f : X \to Y$ un morphisme d'espaces algébriques sur $S$.
Les assertions suivantes sont équivalentes :
\begin{enumerate}
\item le morphisme $f$ est non ramifié,
\item le morphisme $f$ est localement de type fini et $\Omega_{X/Y} = 0$, et
\item le morphisme $f$ est localement de type fini et formellement non ramifié.
\end{enumerate}
\end{lemma}

\begin{proof}
Choisissons un diagramme
$$
\xymatrix{
U \ar[d] \ar[r]_\psi & V \ar[d] \\
X \ar[r]^f & Y
}
$$
où $U$ et $V$ sont des schémas et les flèches verticales sont étales et
surjectives. Nous avons alors
\begin{align*}
f\text{ non ramifié}
& \Leftrightarrow
\psi\text{ non ramifié} \\
& \Leftrightarrow
\psi\text{ localement de type fini et }\Omega_{U/V} = 0 \\
& \Leftrightarrow
f\text{ localement de type fini et }\Omega_{X/Y} = 0 \\
& \Leftrightarrow
f\text{ localement de type fini et formellement non ramifié}
\end{align*}
Nous avons utilisé ici
Morphismes, Lemme \ref{morphisms-lemma-unramified-omega-zero}, et le
Lemme \ref{lemma-characterize-formally-unramified}.
\end{proof}

\begin{lemma}
\label{lemma-universally-injective-unramified}
Soit $S$ un schéma.
Soit $f : X \to Y$ un morphisme d'espaces algébriques sur $S$.
Les assertions suivantes sont équivalentes :
\begin{enumerate}
\item $f$ est non ramifié et est un monomorphisme,
\item $f$ est non ramifié et universellement injectif,
\item $f$ est localement de type fini et est un monomorphisme,
\item $f$ est universellement injectif, localement de type fini et
formellement non ramifié.
\end{enumerate}
De plus, dans ce cas, $f$ est aussi représentable, séparé et
localement quasi-fini.
\end{lemma}

\begin{proof}
Nous avons vu au
Lemme \ref{lemma-unramified-formally-unramified}
qu'être formellement non ramifié et localement de type fini équivaut
à être non ramifié.
Ainsi, (4) est équivalente à (2).
Un monomorphisme est certainement formellement non ramifié ; ainsi, (3) implique (4).
Il est clair que (1) implique (3). Enfin, si (2) est vérifiée, alors
$\Delta : X \to X \times_Y X$ est à la fois une immersion ouverte
(Morphismes d'espaces, Lemme
\ref{spaces-morphisms-lemma-diagonal-unramified-morphism})
et surjective
(Morphismes d'espaces, Lemme
\ref{spaces-morphisms-lemma-universally-injective}) ;
c'est donc un isomorphisme, c'est-à-dire que $f$ est un monomorphisme. Nous voyons ainsi que
(2) implique (1).
Enfin, nous voyons que $f$ est représentable, séparé et localement
quasi-fini d'après
Morphismes d'espaces, Lemmes
\ref{spaces-morphisms-lemma-monomorphism-loc-finite-type-loc-quasi-finite} et
\ref{spaces-morphisms-lemma-locally-quasi-finite-separated-representable}.
\end{proof}

\begin{lemma}
\label{lemma-characterize-closed-immersion}
Soit $S$ un schéma.
Soit $f : X \to Y$ un morphisme d'espaces algébriques sur $S$.
Les assertions suivantes sont équivalentes :
\begin{enumerate}
\item $f$ est une immersion fermée,
\item $f$ est universellement fermé, non ramifié et est un monomorphisme,
\item $f$ est universellement fermé, non ramifié et universellement injectif,
\item $f$ est universellement fermé, localement de type fini et est un monomorphisme,
\item $f$ est universellement fermé, universellement injectif, localement de
type fini et formellement non ramifié.
\end{enumerate}
\end{lemma}

\begin{proof}
L'équivalence de (2) -- (5) résulte immédiatement du
Lemme \ref{lemma-universally-injective-unramified}.
De plus, si (2) -- (5) sont vérifiées, alors $f$ est représentable.
De même, si (1) est vérifiée, alors $f$ est représentable.
Le résultat découle donc du cas des schémas, voir
Morphismes étales, Lemme \ref{etale-lemma-characterize-closed-immersion}.
\end{proof}





\section{Épaississements universels du premier ordre}
\label{section-universal-thickening}

\noindent
Soit $S$ un schéma.
Soit $h : Z \to X$ un morphisme d'espaces algébriques sur $S$.
Un {\it épaississement universel du premier ordre} de $Z$ sur $X$ est un
épaississement du premier ordre $Z \subset Z'$ sur $X$ tel que, pour
tout épaississement du premier ordre $T \subset T'$
sur $X$ et tout diagramme commutatif de flèches pleines
\begin{equation}
\label{equation-universal-first-order-thickening}
\vcenter{
\xymatrix{
& Z \ar[ld] & & T \ar[rd] \ar[ll]^a \\
Z' \ar[rrd] & & & & T' \ar@{..>}[llll]_{a'} \ar[lld]^b \\
 & & X
}
}
\end{equation}
il existe une unique flèche pointillée qui rend le diagramme commutatif.
Notons que, dans cette situation, $(a, a') : (T \subset T') \to (Z \subset Z')$
est un morphisme d'épaississements sur $X$. Ainsi, s'il existe un épaississement universel
du premier ordre, celui-ci est unique à isomorphisme unique près.
En général, un épaississement universel du premier ordre
n'existe pas, mais il existe lorsque $h$ est formellement non ramifié.
Avant de le démontrer, montrons qu'un épaississement universel du premier ordre
dans la catégorie des schémas l'est également dans la
catégorie des espaces algébriques.

\begin{lemma}
\label{lemma-check-universal-first-order-thickening}
Soit $S$ un schéma.
Soit $h : Z \to X$ un morphisme d'espaces algébriques sur $S$.
Soit $Z \subset Z'$ un épaississement du premier ordre sur $X$.
Les assertions suivantes sont équivalentes :
\begin{enumerate}
\item $Z \subset Z'$ est un épaississement universel du premier ordre,
\item pour tout diagramme (\ref{equation-universal-first-order-thickening})
où $T'$ est un schéma, il existe une unique flèche pointillée qui rend le diagramme commutatif, et
\item pour tout diagramme (\ref{equation-universal-first-order-thickening})
où $T'$ est un schéma affine, il existe une unique flèche pointillée qui rend le
diagramme commutatif.
\end{enumerate}
\end{lemma}

\begin{proof}
Les implications (1) $\Rightarrow$ (2) $\Rightarrow$ (3) sont formelles.
Supposons (3) et considérons un diagramme arbitraire
(\ref{equation-universal-first-order-thickening}).
Choisissons une présentation $T' = U'/R'$, voir
Espaces, Définition \ref{spaces-definition-presentation}.
Nous pouvons supposer que $U' = \coprod U'_i$ est une réunion disjointe
de schémas affines, de sorte que $R' = U' \times_{T'} U' = \coprod_{i, j} U'_i \times_T' U'_j$.
Pour chaque paire $(i, j)$, choisissons un recouvrement ouvert affine
$U'_i \times_T' U'_j = \bigcup_k R'_{ijk}$. Notons $U_i, R_{ijk}$
les produits fibrés avec $T$ au-dessus de $T'$. Alors chacun des couples
$U_i \subset U'_i$ et $R_{ijk} \subset R'_{ijk}$
est un épaississement du premier ordre de schémas affines.
Notons $a_i : U_i \to Z$, resp.\ $a_{ijk} : R_{ijk} \to Z$,
le composé de $a : T \to Z$ avec le morphisme
$U_i \to T$, resp.\ $R_{ijk} \to T$.
En appliquant (3) à $a_i : U_i \to Z$,
nous obtenons des morphismes uniques $a'_i : U'_i \to Z'$.
En appliquant (3) à $a_{ijk}$, nous voyons que les deux composés
$R'_{ijk} \to R'_i \to Z'$ et $R'_{ijk} \to R'_j \to Z'$
sont égaux. Ainsi, $a' = \coprod a'_i : U' = \coprod U'_i \to Z'$
descend au faisceau quotient $T' = U'/R'$, ce qui conclut.
\end{proof}

\begin{lemma}
\label{lemma-universal-thickening-over-formally-etale}
Soit $S$ un schéma.
Soient $Z \to Y \to X$ des morphismes d'espaces algébriques sur $S$.
Si $Z \subset Z'$ est un épaississement universel du premier ordre de
$Z$ sur $Y$ et si $Y \to X$ est formellement étale, alors $Z \subset Z'$ est
un épaississement universel du premier ordre de $Z$ sur $X$.
\end{lemma}

\begin{proof}
Cela est formel. En effet, d'après le
Lemme \ref{lemma-check-universal-first-order-thickening},
il suffit de considérer les diagrammes commutatifs de flèches pleines
(\ref{equation-universal-first-order-thickening})
où $T'$ est un schéma affine. Le composé
$T \to Z \to Y$ se relève de manière unique en $T' \to Y$, puisque $Y \to X$ est
supposé formellement étale. Par conséquent, le fait que
$Z \subset Z'$ soit un épaississement universel du premier ordre sur $Y$
fournit le morphisme voulu $a' : T' \to Z'$.
\end{proof}

\begin{lemma}
\label{lemma-etale-morphism-of-universal-thickenings}
Soit $S$ un schéma.
Soient $Z \to Y \to X$ des morphismes d'espaces algébriques sur $S$.
Supposons que $Z \to Y$ soit étale.
\begin{enumerate}
\item Si $Y \subset Y'$ est un épaississement universel du premier ordre de
$Y$ sur $X$, alors l'unique morphisme étale $Z' \to Y'$ tel
que $Z = Y \times_{Y'} Z'$ (voir le
Théorème \ref{theorem-topological-invariance})
est un épaississement universel du premier ordre de $Z$ sur $X$.
\item Si $Z \to Y$ est surjectif et si
$(Z \subset Z') \to (Y \subset Y')$ est un morphisme étale
d'épaississements du premier ordre sur $X$, tandis que $Z'$ est un épaississement universel du
premier ordre de $Z$ sur $X$, alors $Y'$ est un épaississement universel du
premier ordre de $Y$ sur $X$.
\end{enumerate}
\end{lemma}

\begin{proof}
Démonstration de (1). D'après le
Lemme \ref{lemma-check-universal-first-order-thickening},
il suffit de considérer les diagrammes commutatifs de flèches pleines
(\ref{equation-universal-first-order-thickening})
où $T'$ est un schéma affine. Le composé
$T \to Z \to Y$ se relève de manière unique en $T' \to Y'$, puisque $Y'$ est
l'épaississement universel du premier ordre. Le fait que
$Z' \to Y'$ soit étale implique alors (voir le
Lemme \ref{lemma-representable-property-formally-property})
que $T' \to Y'$ se relève en le
morphisme voulu $a' : T' \to Z'$.

\medskip\noindent
Démonstration de (2). Soit $T \subset T'$ un épaississement du premier ordre sur
$X$ et soit $a : T \to Y$ un morphisme. Posons $W = T \times_Y Z$
et notons $c : W \to Z$ la projection.
Soit $W' \to T'$ l'unique morphisme étale tel que
$W = T \times_{T'} W'$, voir le
Théorème \ref{theorem-topological-invariance}.
Notons que $W' \to T'$ est surjectif puisque $Z \to Y$ est surjectif.
Par hypothèse, nous obtenons un unique morphisme $c' : W' \to Z'$
sur $X$ et dont $c$ est la restriction à $W$. Par unicité, les deux restrictions
de $c'$ à $W' \times_{T'} W'$ sont égales (puisque les deux restrictions de
$c$ à $W \times_T W$ sont égales). Ainsi, $c'$ descend en un unique
morphisme $a' : T' \to Y'$, ce qui conclut.
\end{proof}

\begin{lemma}
\label{lemma-universal-thickening}
Soit $S$ un schéma.
Soit $h : Z \to X$ un morphisme formellement non ramifié d'espaces
algébriques sur $S$.
Il existe un épaississement universel du premier ordre $Z \subset Z'$ de
$Z$ sur $X$.
\end{lemma}

\begin{proof}
Choisissons un diagramme commutatif quelconque
$$
\xymatrix{
V \ar[d] \ar[r] & U \ar[d] \\
Z \ar[r] & X
}
$$
où $V$ et $U$ sont des schémas et les flèches verticales sont étales.
Notons que $V \to U$ est un morphisme formellement non ramifié de schémas, voir le
Lemme \ref{lemma-formally-unramified}.
En combinant le
Lemme \ref{lemma-check-universal-first-order-thickening}
et
Compléments sur les morphismes, Lemme \ref{more-morphisms-lemma-universal-thickening},
nous voyons qu'il existe un épaississement universel du premier ordre $V \subset V'$
de $V$ sur $U$. D'après le
Lemme \ref{lemma-universal-thickening-over-formally-etale}, partie (1),
$V'$ est un épaississement universel du premier ordre de $V$ sur $X$.

\medskip\noindent
Fixons un schéma $U$ et un morphisme étale surjectif $U \to X$.
L'argument ci-dessus montre que, pour tout $V \to Z$ étale, où $V$ est
un schéma tel que $V \to Z \to X$ se factorise par $U$, il existe un
épaississement universel du premier ordre $V \subset V'$ de $V$ sur $X$
(qui ne dépend toutefois pas de la factorisation choisie de $V \to X$
par $U$). Nous pouvons maintenant choisir $V$ tel que $V \to Z$ soit étale
surjectif (voir
Espaces, Lemme \ref{spaces-lemma-lift-morphism-presentations}).
Alors $R = V \times_Z V$ est un schéma étale sur $Z$ tel que
$R \to X$ se factorise lui aussi par $U$.
Nous obtenons donc des épaississements universels du premier ordre
$V \subset V'$ et $R \subset R'$ sur $X$.
Puisque $V \subset V'$ est un épaississement universel du premier ordre,
les deux projections $s, t : R \to V$ se relèvent en des morphismes
$s', t': R' \to V'$. D'après le
Lemme \ref{lemma-etale-morphism-of-universal-thickenings},
puisque $R'$ est l'épaississement universel du premier ordre de $R$ sur $X$,
ces morphismes sont étales.
Alors $(t', s') : R' \to V'$ est une relation d'équivalence étale,
et nous pouvons poser $Z' = V'/R'$. Puisque $V' \to Z'$ est étale surjectif
et que $v'$ est l'épaississement universel du premier ordre de $V$ sur $X$,
nous déduisons du
Lemme \ref{lemma-universal-thickening-over-formally-etale}, partie (2),
que $Z'$ est un épaississement universel du premier ordre de $Z$ sur $X$.
\end{proof}

\begin{definition}
\label{definition-universal-thickening}
Soit $S$ un schéma.
Soit $h : Z \to X$ un morphisme formellement non ramifié
d'espaces algébriques sur $S$.
\begin{enumerate}
\item L'{\it épaississement universel du premier ordre} de $Z$ sur $X$
est l'épaississement $Z \subset Z'$ construit au
Lemme \ref{lemma-universal-thickening}.
\item Le {\it faisceau conormal de $Z$ sur $X$} est le faisceau conormal
de $Z$ dans son épaississement universel du premier ordre $Z'$ sur $X$.
\end{enumerate}
Dans cette situation, nous notons souvent le faisceau conormal $\mathcal{C}_{Z/X}$.
\end{definition}

\noindent
Nous voyons ainsi qu'il existe une suite exacte courte de faisceaux
$$
0 \to \mathcal{C}_{Z/X} \to \mathcal{O}_{Z'} \to \mathcal{O}_Z \to 0
$$
sur $Z_\etale$, et $\mathcal{C}_{Z/X}$ est un $\mathcal{O}_Z$-module
quasi-cohérent. Le lemme suivant montre qu'il n'y a pas de
conflit entre cette définition et celle du cas où $Z \to X$
est une immersion.

\begin{lemma}
\label{lemma-immersion-universal-thickening}
Soit $S$ un schéma.
Soit $i : Z \to X$ une immersion d'espaces algébriques sur $S$. Alors
\begin{enumerate}
\item $i$ est formellement non ramifié,
\item l'épaississement universel du premier ordre de $Z$ sur $X$ est le voisinage
infinitésimal du premier ordre de $Z$ dans $X$ de la
Définition \ref{definition-first-order-infinitesimal-neighbourhood},
\item le faisceau conormal de $i$ au sens de la
Définition \ref{definition-conormal-sheaf}
coïncide avec le faisceau conormal de $i$ au sens de la
Définition \ref{definition-universal-thickening}.
\end{enumerate}
\end{lemma}

\begin{proof}
Une immersion d'espaces algébriques est, par définition, un morphisme représentable.
Par conséquent, d'après
Morphismes, Lemmes \ref{morphisms-lemma-open-immersion-unramified} et
\ref{morphisms-lemma-closed-immersion-unramified},
une immersion est non ramifiée (par le principe abstrait du chapitre
Espaces, Lemme
\ref{spaces-lemma-representable-transformations-property-implication}).
Elle est donc formellement non ramifiée d'après le
Lemme \ref{lemma-unramified-formally-unramified}.
Les autres assertions résultent de la combinaison des
Lemmes \ref{lemma-first-order-infinitesimal-neighbourhood} et
\ref{lemma-infinitesimal-neighbourhood-conormal}
avec les définitions.
\end{proof}

\begin{lemma}
\label{lemma-universal-thickening-unramified}
Soit $S$ un schéma.
Soit $Z \to X$ un morphisme formellement non ramifié d'espaces algébriques sur $S$.
Alors l'épaississement universel du premier ordre $Z'$ est formellement
non ramifié sur $X$.
\end{lemma}

\begin{proof}
Soit $T \subset T'$ un épaississement du premier ordre de schémas affines sur $X$.
Soit
$$
\xymatrix{
Z' \ar[d] & T \ar[l]^c \ar[d] \\
X & T' \ar[l] \ar[lu]^{a, b}
}
$$
un diagramme commutatif. Posons $T_0 = c^{-1}(Z) \subset T$ et
$T'_a = a^{-1}(Z)$ (au sens schématique).
Comme $Z'$ est un épaississement du premier ordre de $Z$, nous voyons que $T'$
est un épaississement du premier ordre de $T'_a$. De plus, puisque $c = a|_T$, nous avons
$T_0 = T \cap T'_a$ (au sens schématique). Comme $T'$ est un épaississement du premier ordre
de $T$, il s'ensuit que $T'_a$
est un épaississement du premier ordre de $T_0$. Or $a|_{T'_a}$ et $b|_{T'_a}$
sont des morphismes de $T'_a$ vers $Z'$ sur $X$ qui coïncident sur $T_0$ comme
morphismes vers $Z$. La propriété universelle de $Z'$ donne donc
$a|_{T'_a} = b|_{T'_a}$. Ainsi, $a$ et $b$ sont des morphismes depuis
l'épaississement du premier ordre $T'$ de $T'_a$ dont les restrictions à
$T'_a$ coïncident comme morphismes vers $Z$. En utilisant une nouvelle fois la propriété universelle de
$Z'$, nous obtenons $a = b$. Autrement dit, la propriété qui définit
un morphisme formellement non ramifié est vérifiée par $Z' \to X$, comme souhaité.
\end{proof}

\begin{lemma}
\label{lemma-universal-thickening-functorial}
Soit $S$ un schéma.
Considérons un diagramme commutatif d'espaces algébriques sur $S$
$$
\xymatrix{
Z \ar[r]_h \ar[d]_f & X \ar[d]^g \\
W \ar[r]^{h'} & Y
}
$$
où $h$ et $h'$ sont formellement non ramifiés. Soit $Z \subset Z'$ l'épaississement universel
du premier ordre de $Z$ sur $X$. Soit $W \subset W'$ l'épaississement universel
du premier ordre de $W$ sur $Y$. Il existe un morphisme canonique
$(f, f') : (Z, Z') \to (W, W')$ d'épaississements sur $Y$ qui s'insère dans
le diagramme commutatif suivant
$$
\xymatrix{
& & & Z' \ar[ld] \ar[d]^{f'} \\
Z \ar[rr] \ar[d]_f \ar[rrru] & & X \ar[d] & W' \ar[ld] \\
W \ar[rrru]|!{[rr];[rruu]}\hole \ar[rr] & & Y
}
$$
En particulier, le morphisme $(f, f')$ d'épaississements induit un morphisme
de faisceaux conormaux $f^*\mathcal{C}_{W/Y} \to \mathcal{C}_{Z/X}$.
\end{lemma}

\begin{proof}
La première assertion résulte immédiatement de la propriété universelle de $W'$.
Le morphisme induit sur les faisceaux conormaux est celui du
Lemme \ref{lemma-conormal-functorial}
appliqué à $(Z \subset Z') \to (W \subset W')$.
\end{proof}

\begin{lemma}
\label{lemma-universal-thickening-fibre-product}
Soit $S$ un schéma. Soit
$$
\xymatrix{
Z \ar[r]_h \ar[d]_f & X \ar[d]^g \\
W \ar[r]^{h'} & Y
}
$$
un diagramme de produit fibré d'espaces algébriques sur $S$, où
$h'$ est formellement non ramifié. Alors $h$ est formellement non ramifié et, si
$W \subset W'$ est l'épaississement universel du premier ordre de $W$ sur $Y$,
alors $Z = X \times_Y W \subset X \times_Y W'$ est l'épaississement universel
du premier ordre de $Z$ sur $X$. En particulier, le morphisme canonique
$f^*\mathcal{C}_{W/Y} \to \mathcal{C}_{Z/X}$ du
Lemme \ref{lemma-universal-thickening-functorial}
est surjectif.
\end{lemma}

\begin{proof}
Le morphisme $h$ est formellement non ramifié d'après le
Lemme \ref{lemma-base-change-formally-unramified}.
Il est clair que $X \times_Y W'$ est un épaississement du premier ordre.
On vérifie immédiatement qu'il possède la propriété universelle,
car $W'$ la possède (d'après les propriétés universelles des
produits fibrés). Voir le
Lemme \ref{lemma-conormal-functorial-flat}
pour comprendre pourquoi cela implique la surjectivité du morphisme de faisceaux conormaux.
\end{proof}

\begin{lemma}
\label{lemma-universal-thickening-fibre-product-flat}
Soit $S$ un schéma. Soit
$$
\xymatrix{
Z \ar[r]_h \ar[d]_f & X \ar[d]^g \\
W \ar[r]^{h'} & Y
}
$$
un diagramme de produit fibré d'espaces algébriques sur $S$, où
$h'$ est formellement non ramifié et $g$ plat. Dans ce cas, le morphisme correspondant
$Z' \to W'$ entre les épaississements universels du premier ordre est plat, et
$f^*\mathcal{C}_{W/Y} \to \mathcal{C}_{Z/X}$ est un isomorphisme.
\end{lemma}

\begin{proof}
La platitude est préservée par changement de base, voir
Morphismes d'espaces, Lemme \ref{spaces-morphisms-lemma-base-change-flat}.
La première assertion découle donc de la description de $W'$ au
Lemme \ref{lemma-universal-thickening-fibre-product}.
Il est clair que $X \times_Y W'$ est un épaississement du premier ordre.
On vérifie immédiatement qu'il possède la propriété universelle,
car $W'$ la possède (d'après les propriétés universelles des
produits fibrés). Voir le
Lemme \ref{lemma-conormal-functorial-flat}
pour comprendre pourquoi cela implique que le morphisme de faisceaux conormaux est un isomorphisme.
\end{proof}

\begin{lemma}
\label{lemma-universal-thickening-localize}
La formation des épaississements universels du premier ordre commute à la
localisation étale. Plus précisément, soit $h : Z \to X$ un morphisme formellement non ramifié
d'espaces algébriques sur un schéma de base $S$.
Considérons
$$
\xymatrix{
V \ar[d] \ar[r] & U \ar[d] \\
Z \ar[r] & X
}
$$
un diagramme commutatif dont les flèches verticales sont étales.
Soit $Z'$ l'épaississement universel du premier ordre de $Z$ sur $X$.
Alors $V \to U$ est formellement non ramifié, et l'épaississement universel du premier
ordre $V'$ de $V$ sur $U$ est étale sur $Z'$.
En particulier, $\mathcal{C}_{Z/X}|_V = \mathcal{C}_{V/U}$.
\end{lemma}

\begin{proof}
La première assertion résulte du
Lemme \ref{lemma-formally-unramified}.
La compatibilité des épaississements universels du premier ordre
résulte des
Lemmes \ref{lemma-universal-thickening-over-formally-etale} et
\ref{lemma-etale-morphism-of-universal-thickenings}.
\end{proof}

\begin{lemma}
\label{lemma-differentials-universally-unramified}
Soit $S$ un schéma. Soit $B$ un espace algébrique sur $S$.
Soit $h : Z \to X$ un morphisme formellement non ramifié d'espaces algébriques
sur $B$. Soit $Z \subset Z'$ l'épaississement universel du premier ordre de $Z$
sur $X$, de morphisme structural $h' : Z' \to X$. Le morphisme canonique
$$
\text{d}h' : (h')^*\Omega_{X/B} \to \Omega_{Z'/B}
$$
induit un isomorphisme
$h^*\Omega_{X/B} \to \Omega_{Z'/B} \otimes \mathcal{O}_Z$.
\end{lemma}

\begin{proof}
Le morphisme $c_{h'}$ est celui qui est défini au
Lemme \ref{lemma-functoriality-differentials}.
Si $i : Z \to Z'$ est l'immersion fermée donnée, alors
$i^*c_{h'}$ est un morphisme
$h^*\Omega_{X/S} \to \Omega_{Z'/S} \otimes \mathcal{O}_Z$.
Pour vérifier qu'il s'agit d'un isomorphisme, on se ramène au cas des schémas
par localisation étale, voir le
Lemme \ref{lemma-universal-thickening-localize}
et le
Lemme \ref{lemma-localize-differentials}.
Dans ce cas, le résultat est
Compléments sur les morphismes,
Lemme \ref{more-morphisms-lemma-differentials-universally-unramified}.
\end{proof}

\begin{lemma}
\label{lemma-universally-unramified-differentials-sequence}
Soit $S$ un schéma. Soit $B$ un espace algébrique sur $S$.
Soit $h : Z \to X$ un morphisme formellement non ramifié d'espaces
algébriques sur $B$.
Il existe une suite exacte canonique
$$
\mathcal{C}_{Z/X} \to h^*\Omega_{X/B} \to \Omega_{Z/B} \to 0.
$$
La première flèche est induite par $\text{d}_{Z'/B}$, où
$Z'$ est le voisinage universel du premier ordre de $Z$ sur $X$.
\end{lemma}

\begin{proof}
Nous savons qu'il existe une suite exacte canonique
$$
\mathcal{C}_{Z/Z'} \to
\Omega_{Z'/S} \otimes \mathcal{O}_Z \to
\Omega_{Z/S} \to 0.
$$
Cela résulte du
Lemme \ref{lemma-differentials-relative-immersion}.
Le résultat s'obtient donc en appliquant le
Lemme \ref{lemma-differentials-universally-unramified}.
\end{proof}

\begin{lemma}
\label{lemma-two-unramified-morphisms}
Soit $S$ un schéma. Soit
$$
\xymatrix{
Z \ar[r]_i \ar[rd]_j & X \ar[d] \\
& Y
}
$$
un diagramme commutatif d'espaces algébriques sur $S$,
où $i$ et $j$ sont formellement non ramifiés. Il existe alors une
suite exacte canonique
$$
\mathcal{C}_{Z/Y} \to
\mathcal{C}_{Z/X} \to
i^*\Omega_{X/Y} \to 0
$$
où la première flèche provient du
Lemme \ref{lemma-universal-thickening-functorial}
et la seconde du
Lemme \ref{lemma-universally-unramified-differentials-sequence}.
\end{lemma}

\begin{proof}
Les morphismes étant définis, la vérification de l'exactitude de la suite
se ramène au cas des schémas par localisation étale, voir le
Lemme \ref{lemma-universal-thickening-localize}
et le
Lemme \ref{lemma-localize-differentials}.
Dans ce cas, le résultat est
Compléments sur les morphismes,
Lemme \ref{more-morphisms-lemma-two-unramified-morphisms}.
\end{proof}

\begin{lemma}
\label{lemma-transitivity-conormal-unramified}
Soit $S$ un schéma.
Soient $Z \to Y \to X$ des morphismes formellement non ramifiés
d'espaces algébriques sur $S$.
\begin{enumerate}
\item Si $Z \subset Z'$ est l'épaississement universel du premier ordre de $Z$
sur $X$ et si $Y \subset Y'$ est l'épaississement universel du premier ordre de $Y$
sur $X$, alors il existe un morphisme $Z' \to Y'$ et $Y \times_{Y'} Z'$ est
l'épaississement universel du premier ordre de $Z$ sur $Y$.
\item Il existe une suite exacte canonique
$$
i^*\mathcal{C}_{Y/X} \to
\mathcal{C}_{Z/X} \to
\mathcal{C}_{Z/Y} \to 0
$$
où les morphismes proviennent du
Lemme \ref{lemma-universal-thickening-functorial}
et $i : Z \to Y$ est le premier morphisme.
\end{enumerate}
\end{lemma}

\begin{proof}
Le morphisme $h : Z' \to Y'$ de (1) provient du
Lemme \ref{lemma-universal-thickening-functorial}.
L'assertion selon laquelle $Y \times_{Y'} Z'$ est l'épaississement universel du premier ordre
de $Z$ sur $Y$ résulte immédiatement des propriétés universelles
de $Z'$ et de $Y'$. D'après le
Lemme \ref{lemma-transitivity-conormal},
nous avons une suite exacte
$$
(i')^*\mathcal{C}_{Y \times_{Y'} Z'/Z'} \to
\mathcal{C}_{Z/Z'} \to
\mathcal{C}_{Z/Y \times_{Y'} Z'} \to 0
$$
où $i' : Z \to Y \times_{Y'} Z'$ est le morphisme donné. D'après le
Lemme \ref{lemma-conormal-functorial-flat},
il existe une surjection
$h^*\mathcal{C}_{Y/Y'} \to \mathcal{C}_{Y \times_{Y'} Z'/Z'}$.
Combinée aux égalités
$\mathcal{C}_{Y/Y'} = \mathcal{C}_{Y/X}$,
$\mathcal{C}_{Z/Z'} = \mathcal{C}_{Z/X}$, et
$\mathcal{C}_{Z/Y \times_{Y'} Z'} = \mathcal{C}_{Z/Y}$,
cette surjection démontre le lemme.
\end{proof}













\section{Morphismes formellement étales}
\label{section-formally-etale}

\noindent
Dans cette section, nous explicitons ce que signifie pour un morphisme d'espaces algébriques
d'être formellement étale.

\begin{definition}
\label{definition-formally-etale}
Soit $S$ un schéma. Un morphisme $f : X \to Y$ d'espaces algébriques sur $S$
est dit {\it formellement étale} s'il est formellement étale en tant que
transformation de foncteurs au sens de la
Définition \ref{definition-formally-smooth-etale-unramified}.
\end{definition}

\noindent
Nous ne redonnerons pas les résultats démontrés dans le cadre plus général des
transformations de foncteurs formellement étales dans la
section \ref{section-formally-smooth-etale-unramified}.

\begin{lemma}
\label{lemma-formally-etale}
Soit $S$ un schéma. Soit $f : X \to Y$ un morphisme d'espaces algébriques sur
$S$. Les assertions suivantes sont équivalentes :
\begin{enumerate}
\item $f$ est formellement étale,
\item pour tout diagramme
$$
\xymatrix{
U \ar[d] \ar[r]_\psi & V \ar[d] \\
X \ar[r]^f & Y
}
$$
où $U$ et $V$ sont des schémas et où les flèches verticales sont étales,
le morphisme de schémas $\psi$ est formellement étale (au sens de
Compléments sur les morphismes,
Définition \ref{more-morphisms-definition-formally-etale}), et
\item il existe un tel diagramme dont les flèches verticales sont surjectives et dans lequel
$\psi$ est formellement étale.
\end{enumerate}
\end{lemma}

\begin{proof}
Supposons $f$ formellement étale. D'après le
Lemme \ref{lemma-representable-property-formally-property},
les morphismes $U \to X$ et $V \to Y$ sont formellement étales. Ainsi, d'après le
Lemme \ref{lemma-composition-formally-smooth-etale-unramified},
la composée $U \to Y$ est formellement étale. Il résulte alors du
Lemme \ref{lemma-formally-permanence}
que $U \to V$ est formellement étale. Ainsi, (1) implique (2), et (2)
implique trivialement (3).

\medskip\noindent
Supposons donné un diagramme comme en (3). D'après le
Lemme \ref{lemma-representable-property-formally-property},
le morphisme $V \to Y$ est formellement étale. Ainsi, d'après le
Lemme \ref{lemma-composition-formally-smooth-etale-unramified},
la composée $U \to Y$ est formellement étale. Il résulte alors du
Lemme \ref{lemma-etale-on-top}
que $X \to Y$ est formellement étale, c'est-à-dire que (1) est vérifiée.
\end{proof}

\begin{lemma}
\label{lemma-formally-etale-not-affine}
Soit $S$ un schéma.
Soit $f : X \to Y$ un morphisme formellement étale d'espaces algébriques sur $S$.
Alors, pour tout diagramme commutatif de flèches pleines
$$
\xymatrix{
X \ar[d]_f & T \ar[d]^i \ar[l]_a \\
Y & T' \ar[l] \ar@{-->}[lu]
}
$$
où $T \subset T'$ est un épaississement du premier ordre d'espaces algébriques
sur $Y$, il existe une unique flèche pointillée qui rend le diagramme commutatif.
Autrement dit, dans la
Définition \ref{definition-formally-etale},
on peut supprimer la condition que $T$ soit affine.
\end{lemma}

\begin{proof}
Soit $U' \to T'$ un morphisme étale surjectif, où $U' = \coprod U'_i$
est une somme disjointe de schémas affines. Posons
$U_i = T \times_{T'} U'_i$. On obtient alors des morphismes
$a'_i : U'_i \to X$ tels que $a'_i|_{U_i}$ soit égal à la composée
$U_i \to T \to X$. Par unicité (voir le
Lemme \ref{lemma-formally-unramified-not-affine}),
on voit que $a'_i$ et $a'_j$ coïncident sur le produit fibré
$U'_i \times_{T'} U'_j$. Par suite, $\coprod a'_i : U' \to X$
descend en un unique morphisme $a' : T' \to X$.
\end{proof}

\begin{lemma}
\label{lemma-composition-formally-etale}
La composée de morphismes formellement étales est formellement étale.
\end{lemma}

\begin{proof}
C'est formel.
\end{proof}

\begin{lemma}
\label{lemma-base-change-formally-etale}
Un changement de base d'un morphisme formellement étale est formellement étale.
\end{lemma}

\begin{proof}
C'est formel.
\end{proof}

\begin{lemma}
\label{lemma-characterize-formally-etale}
Soit $S$ un schéma.
Soit $f : X \to Y$ un morphisme d'espaces algébriques sur $S$.
Les assertions suivantes sont équivalentes :
\begin{enumerate}
\item $f$ est formellement étale,
\item $f$ est formellement non ramifié et l'épaississement universel du premier ordre
de $X$ sur $Y$ est égal à $X$,
\item $f$ est formellement non ramifié et $\mathcal{C}_{X/Y} = 0$, et
\item $\Omega_{X/Y} = 0$ et $\mathcal{C}_{X/Y} = 0$.
\end{enumerate}
\end{lemma}

\begin{proof}
En réalité, la dernière assertion n'a de sens que parce que $\Omega_{X/Y} = 0$
implique que $\mathcal{C}_{X/Y}$ est défini par le
Lemme \ref{lemma-characterize-formally-unramified}
et la
Définition \ref{definition-universal-thickening}.
Cela montre aussi que (3) et (4) sont équivalentes.

\medskip\noindent
Chacune des hypothèses (1), (2) et (3) implique que $f$ est formellement
non ramifié. On peut donc supposer $f$ formellement non ramifié. L'équivalence
de (1), (2) et (3) résulte de la propriété universelle de l'épaississement
universel du premier ordre $X'$ de $X$ sur $S$ et du fait que
$X = X' \Leftrightarrow \mathcal{C}_{X/Y} = 0$, puisque,
par définition, $\mathcal{C}_{X/Y} = \mathcal{C}_{X/X'}$
est le faisceau d'idéaux de $X$ dans $X'$.
\end{proof}

\begin{lemma}
\label{lemma-unramified-flat-formally-etale}
Un morphisme plat et non ramifié est formellement étale.
\end{lemma}

\begin{proof}
Cela résulte du cas des schémas, voir
Compléments sur les morphismes,
Lemme \ref{more-morphisms-lemma-unramified-flat-formally-etale},
et de la localisation étale, voir les
Lemmes \ref{lemma-formally-unramified} et \ref{lemma-formally-etale},
ainsi que
Morphismes d'espaces, Lemme \ref{spaces-morphisms-lemma-flat-local}.
\end{proof}

\begin{lemma}
\label{lemma-etale-formally-etale}
Soit $S$ un schéma.
Soit $f : X \to Y$ un morphisme d'espaces algébriques sur $S$.
Les assertions suivantes sont équivalentes :
\begin{enumerate}
\item le morphisme $f$ est étale, et
\item le morphisme $f$ est localement de présentation finie et
formellement étale.
\end{enumerate}
\end{lemma}

\begin{proof}
Cela résulte du cas des schémas, voir
Compléments sur les morphismes,
Lemme \ref{more-morphisms-lemma-etale-formally-etale},
et de la localisation étale, voir le
Lemme \ref{lemma-formally-etale},
ainsi que
Morphismes d'espaces,
Lemmes \ref{spaces-morphisms-lemma-finite-presentation-local} et
\ref{spaces-morphisms-lemma-etale-local}.
\end{proof}















\section{Déformations infinitésimales de morphismes}
\label{section-action-by-derivations}

\noindent
Dans cette section, nous expliquons comment une dérivation peut servir à
déformer infinitésimalement un morphisme. Dans toute la section, nous utilisons le fait qu'un
faisceau sur un épaississement $X'$ de $X$ peut être considéré comme un faisceau sur $X$, voir les
Équations (\ref{equation-equivalence-etale-spaces}) et
(\ref{equation-fundamental-equivalence}).

\begin{lemma}
\label{lemma-difference-derivation}
Soit $S$ un schéma. Soit $B$ un espace algébrique sur $S$.
Soient $X \subset X'$ et $Y \subset Y'$ deux épaississements du premier ordre
d'espaces algébriques sur $B$.
Soient $(a, a'), (b, b') : (X \subset X') \to (Y \subset Y')$
deux morphismes d'épaississements sur $B$. Supposons que
\begin{enumerate}
\item $a = b$, et
\item les deux morphismes $a^*\mathcal{C}_{Y/Y'} \to \mathcal{C}_{X/X'}$
(Lemme \ref{lemma-conormal-functorial})
soient égaux.
\end{enumerate}
Alors le morphisme $(a')^\sharp - (b')^\sharp$ se factorise sous la forme
$$
\mathcal{O}_{Y'} \to \mathcal{O}_Y \xrightarrow{D}
a_*\mathcal{C}_{X/X'} \to a_*\mathcal{O}_{X'}
$$
où $D$ est une $\mathcal{O}_B$-dérivation.
\end{lemma}

\begin{proof}
Au lieu de travailler sur $Y$, travaillons sur $X$. L'avantage est que le foncteur d'image inverse
$a^{-1}$ est exact. En utilisant (1) et (2), on obtient un diagramme commutatif
à lignes exactes
$$
\xymatrix{
0 \ar[r] &
\mathcal{C}_{X/X'} \ar[r] &
\mathcal{O}_{X'} \ar[r] &
\mathcal{O}_X \ar[r] & 0 \\
0 \ar[r] &
a^{-1}\mathcal{C}_{Y/Y'} \ar[r] \ar[u] &
a^{-1}\mathcal{O}_{Y'}
\ar[r] \ar@<1ex>[u]^{(a')^\sharp} \ar@<-1ex>[u]_{(b')^\sharp} &
a^{-1}\mathcal{O}_Y \ar[r] \ar[u] & 0
}
$$
Or il est général que, dans une telle situation, la différence des
morphismes de $\mathcal{O}_B$-algèbres $(a')^\sharp$ et $(b')^\sharp$ est une
$\mathcal{O}_B$-dérivation de $a^{-1}\mathcal{O}_Y$ dans $\mathcal{C}_{X/X'}$.
Par adjonction entre les foncteurs $a^{-1}$ et $a_*$, cela revient
à une $\mathcal{O}_B$-dérivation de
$\mathcal{O}_Y$ dans $a_*\mathcal{C}_{X/X'}$. Nous omettons certains détails.
\end{proof}

\noindent
Notons que, dans la situation du lemme ci-dessus, on peut écrire
$D$ sous la forme
\begin{equation}
\label{equation-D}
D = \text{d}_{Y/B} \circ \theta
\end{equation}
où $\theta$ est une application $\mathcal{O}_Y$-linéaire
$\theta : \Omega_{Y/B} \to a_*\mathcal{C}_{X/X'}$.
Bien entendu, par adjonction encore, on peut alors considérer $\theta$ comme une
application $\mathcal{O}_X$-linéaire
$\theta : a^*\Omega_{Y/B} \to \mathcal{C}_{X/X'}$.

\begin{lemma}
\label{lemma-action-by-derivations}
Soit $S$ un schéma. Soit $B$ un espace algébrique sur $S$.
Soit $(a, a') : (X \subset X') \to (Y \subset Y')$
un morphisme d'épaississements du premier ordre sur $B$.
Soit
$$
\theta : a^*\Omega_{Y/B} \to \mathcal{C}_{X/X'}
$$
une application $\mathcal{O}_X$-linéaire. Il existe alors un unique morphisme d'épaississements
$(b, b') : (X \subset X') \to (Y \subset Y')$ tel que
(1) et (2) du
Lemme \ref{lemma-difference-derivation}
soient vérifiées et que la dérivation $D$ et $\theta$ soient reliées par
l'Équation (\ref{equation-D}).
\end{lemma}

\begin{proof}
Considérons le morphisme
$$
\alpha = (a')^\sharp + D : a^{-1}\mathcal{O}_{Y'} \to \mathcal{O}_{X'}
$$
où $D$ est comme dans l'Équation (\ref{equation-D}). Puisque $D$ est une
$\mathcal{O}_B$-dérivation, il en résulte que $\alpha$ est un morphisme de
faisceaux de $\mathcal{O}_B$-algèbres. Par construction, on a
$i_X^\sharp \circ \alpha = a^\sharp \circ i_Y^\sharp$, où
$i_X : X \to X'$ et $i_Y : Y \to Y'$ sont les immersions fermées données. D'après le
Lemme \ref{lemma-first-order-thickening-maps},
on obtient un unique morphisme
$(a, b') : (X  \subset X') \to (Y \subset Y')$ d'épaississements
sur $B$ tel que $\alpha = (b')^\sharp$. En posant $b = a$,
le résultat est démontré.
\end{proof}

\begin{remark}
\label{remark-action-by-derivations}
Mêmes hypothèses et notations qu'au Lemme \ref{lemma-action-by-derivations}.
L'action d'une section locale $\theta$ sur $a'$ est parfois notée
$\theta \cdot a'$. Notons que cela signifie seulement que
$(a')^\sharp$ et $(\theta \cdot a')^\sharp$ diffèrent d'une dérivation
$D$ reliée à $\theta$ par l'Équation (\ref{equation-D}).
\end{remark}

\begin{lemma}
\label{lemma-sheaf}
Soit $S$ un schéma. Soit $B$ un espace algébrique sur $S$.
Soient $X \subset X'$ et $Y \subset Y'$ des épaississements du premier ordre
sur $B$. Supposons donnés un morphisme $a : X \to Y$ et un morphisme
$A : a^*\mathcal{C}_{Y/Y'} \to \mathcal{C}_{X/X'}$ de
$\mathcal{O}_X$-modules. Pour un objet $U'$ de
$(X')_{spaces, \etale}$ avec $U = X \times_{X'} U'$
considérons les morphismes $a' : U' \to Y'$ tels que
\begin{enumerate}
\item $a'$ soit un morphisme sur $B$,
\item $a'|_U = a|_U$, et
\item le morphisme induit
$a^*\mathcal{C}_{Y/Y'}|_U \to \mathcal{C}_{X/X'}|_U$
soit la restriction de $A$ à $U$.
\end{enumerate}
Alors la règle
\begin{equation}
\label{equation-sheaf}
U' \mapsto
\{a' : U' \to Y'\text{ tels que (1), (2), (3) soient satisfaites.}\}
\end{equation}
définit un faisceau d'ensembles sur $(X')_{spaces, \etale}$.
\end{lemma}

\begin{proof}
Notons $\mathcal{F}$ la règle du lemme.
Le morphisme de restriction $\mathcal{F}(U') \to \mathcal{F}(V')$ associé à
$V' \subset U' \subset X'$
pour la règle $\mathcal{F}$ est bien le morphisme de restriction $a' \mapsto a'|_{V'}$.
Avec cette définition, il est clair que $\mathcal{F}$ est un
faisceau, car les morphismes d'espaces algébriques vérifient la descente étale, voir
Descente sur les espaces,
Lemme \ref{spaces-descent-lemma-fpqc-universal-effective-epimorphisms}.
\end{proof}

\begin{lemma}
\label{lemma-action-sheaf}
Mêmes notations et hypothèses qu'au Lemme \ref{lemma-sheaf}.
Nous identifions les faisceaux sur $X$ et $X'$ au moyen de
(\ref{equation-equivalence-etale-spaces}).
Le faisceau
$$
\SheafHom_{\mathcal{O}_X}(a^*\Omega_{Y/B}, \mathcal{C}_{X/X'})
$$
opère sur le faisceau (\ref{equation-sheaf}). En outre, cette action
est simplement transitive pour tout objet $U'$ de $(X')_{spaces, \etale}$
au-dessus duquel le faisceau (\ref{equation-sheaf}) possède une section.
\end{lemma}

\begin{proof}
C'est la combinaison des
Lemmes \ref{lemma-difference-derivation},
\ref{lemma-action-by-derivations},
et \ref{lemma-sheaf}.
\end{proof}

\begin{remark}
\label{remark-special-case}
Un cas particulier des
Lemmes \ref{lemma-difference-derivation},
\ref{lemma-action-by-derivations},
\ref{lemma-sheaf}, et
\ref{lemma-action-sheaf}
est celui où $Y = Y'$. Dans ce cas, le morphisme $A$ est toujours nul.
Le faisceau du
Lemme \ref{lemma-sheaf}
est simplement donné par la règle
$$
U' \mapsto
\{a' : U' \to Y\text{ au-dessus de }B\text{ tel que } a'|_U = a|_U\}
$$
sur lequel opère le faisceau
$\SheafHom_{\mathcal{O}_X}(a^*\Omega_{Y/B}, \mathcal{C}_{X/X'})$.
\end{remark}

\begin{remark}
\label{remark-another-special-case}
Un autre cas particulier des
Lemmes \ref{lemma-difference-derivation},
\ref{lemma-action-by-derivations},
\ref{lemma-sheaf}, et
\ref{lemma-action-sheaf}
est celui où $B$ est lui-même l'espace ambiant d'un épaississement $Z \subset Z' = B$
et où $Y = Z \times_{Z'} Y'$. Représentons la situation par
$$
\xymatrix{
(X \subset X') \ar@{..>}[rr]_{(a, ?)} \ar[rd]_{(g, g')} & &
(Y \subset Y') \ar[ld]^{(h, h')} \\
& (Z \subset Z')
}
$$
Dans ce cas, le morphisme $A : a^*\mathcal{C}_{Y/Y'} \to \mathcal{C}_{X/X'}$
est déterminé par $a$ : le morphisme
$h^*\mathcal{C}_{Z/Z'} \to \mathcal{C}_{Y/Y'}$ est surjectif (car nous avons
supposé $Y = Z \times_{Z'} Y'$) ; par conséquent, son image inverse
$g^*\mathcal{C}_{Z/Z'} = a^*h^*\mathcal{C}_{Z/Z'} \to
a^*\mathcal{C}_{Y/Y'}$ est surjective, et la composée
$g^*\mathcal{C}_{Z/Z'} \to a^*\mathcal{C}_{Y/Y'} \to \mathcal{C}_{X/X'}$
doit être le morphisme canonique induit par $g'$. Ainsi, le faisceau du
Lemme \ref{lemma-sheaf}
est simplement donné par la règle
$$
U' \mapsto
\{a' : U' \to Y'\text{ au-dessus de }Z'\text{ tel que } a'|_U = a|_U\}
$$
sur lequel opère le faisceau
$\SheafHom_{\mathcal{O}_X}(a^*\Omega_{Y/Z}, \mathcal{C}_{X/X'})$.
\end{remark}

\begin{lemma}
\label{lemma-action-by-derivations-etale-localization}
Soit $S$ un schéma. Considérons un diagramme commutatif d'épaississements
du premier ordre
$$
\vcenter{
\xymatrix{
(T_2 \subset T_2') \ar[d]_{(h, h')} \ar[rr]_{(a_2, a_2')} & &
(X_2 \subset X_2') \ar[d]^{(f, f')} \\
(T_1 \subset T_1') \ar[rr]^{(a_1, a_1')} & &
(X_1 \subset X_1')
}
}
\quad
\begin{matrix}
\text{et un} \\
\text{diagramme commutatif}
\end{matrix}
\quad
\vcenter{
\xymatrix{
X_2' \ar[r] \ar[d] & B_2 \ar[d] \\
X_1' \ar[r] & B_1
}
}
$$
d'espaces algébriques sur $S$,
où $X_2 \to X_1$ et $B_2 \to B_1$ sont étales.
Pour toute application $\mathcal{O}_{T_1}$-linéaire
$\theta_1 : a_1^*\Omega_{X_1/B_1} \to \mathcal{C}_{T_1/T'_1}$, soit
$\theta_2$ la composée
$$
\xymatrix{
a_2^*\Omega_{X_2/B_2} \ar@{=}[r] &
h^*a_1^*\Omega_{X_1/B_1} \ar[r]^-{h^*\theta_1} &
h^*\mathcal{C}_{T_1/T'_1} \ar[r] &
\mathcal{C}_{T_2/T'_2}
}
$$
(le signe d'égalité est expliqué dans la démonstration). Alors le diagramme
$$
\xymatrix{
T_2' \ar[rr]_{\theta_2 \cdot a_2'} \ar[d] & & X'_2 \ar[d] \\
T_1' \ar[rr]^{\theta_1 \cdot a_1'} & & X'_1
}
$$
est commutatif, où les actions $\theta_2 \cdot a_2'$ et $\theta_1 \cdot a_1'$
sont celles de la Remarque \ref{remark-action-by-derivations}.
\end{lemma}

\begin{proof}
Le signe d'égalité provient de l'identification
$f^*\Omega_{X_1/S_1} = \Omega_{X_2/S_2}$ que l'on obtient
parce que la construction du faisceau des différentielles est
compatible avec la localisation étale (tant à la source qu'au but), voir le
Lemme \ref{lemma-localize-differentials}.
En effet, on a ainsi
$a_2^*\Omega_{X_2/S_2} = a_2^*f^*\Omega_{X_1/S_1} =
h^*a_1^*\Omega_{X_1/S_1}$ puisque $f \circ a_2 = a_1 \circ h$.
Cela étant, la commutativité du diagramme se vérifie
localement pour la topologie étale. On peut donc supposer que $T'_i$, $X'_i$,
$B_2$ et $B_1$ sont des schémas ; dans ce cas, le lemme
résulte de
Compléments sur les morphismes, Lemme
\ref{more-morphisms-lemma-action-by-derivations-etale-localization}.
Autre démonstration : d'après le Lemme \ref{lemma-first-order-thickening-maps},
il suffit de montrer qu'un certain diagramme de faisceaux
d'anneaux sur $X_1'$ est commutatif ; on raisonne alors exactement
comme dans la démonstration du lemme précité de
Compléments sur les morphismes, Lemme
\ref{more-morphisms-lemma-action-by-derivations-etale-localization}
pour voir que c'est bien le cas.
\end{proof}










\section{Déformations infinitésimales d'espaces algébriques}
\label{section-deform}

\noindent
Le lemme élémentaire suivant est souvent un outil commode pour vérifier si
une déformation infinitésimale d'un morphisme est plate.

\begin{lemma}
\label{lemma-deform}
Soit $S$ un schéma. Soit $(f, f') : (X \subset X') \to (Y \subset Y')$ un
morphisme d'épaississements du premier ordre d'espaces algébriques sur $S$. Supposons que
$f$ soit plat. Alors les assertions suivantes sont équivalentes :
\begin{enumerate}
\item $f'$ est plat et $X = Y \times_{Y'} X'$, et
\item le morphisme canonique $f^*\mathcal{C}_{Y/Y'} \to \mathcal{C}_{X/X'}$
est un isomorphisme.
\end{enumerate}
\end{lemma}

\begin{proof}
Choisissons un schéma $V'$ et un morphisme étale surjectif $V' \to Y'$.
Choisissons un schéma $U'$ et un morphisme étale surjectif
$U' \to X' \times_{Y'} V'$. Posons $U = X \times_{X'} U'$ et
$V = Y \times_{Y'} V'$. D'après notre définition d'un morphisme plat
d'espaces algébriques, le morphisme induit $g : U \to V$ est un morphisme plat
de schémas, et $f'$ est plat si et seulement si le morphisme correspondant
$g' : U' \to V'$ est plat. De plus, $X = Y \times_{Y'} X'$ si et seulement
si $U = V \times_{V'} V'$. Enfin, le morphisme
$f^*\mathcal{C}_{Y/Y'} \to \mathcal{C}_{X/X'}$
est un isomorphisme si et seulement si
$g^*\mathcal{C}_{V/V'} \to \mathcal{C}_{U/U'}$ est un isomorphisme.
Le lemme résulte donc de son analogue pour les morphismes de schémas, voir
Compléments sur les morphismes, Lemme \ref{more-morphisms-lemma-deform}.
\end{proof}

\noindent
Le lemme suivant est la version « nilpotente » du
« critère de platitude par fibres », voir la
section \ref{section-criterion-flat-fibres}.

\begin{lemma}
\label{lemma-flatness-morphism-thickenings}
Soit $S$ un schéma. Considérons un diagramme commutatif
$$
\xymatrix{
(X \subset X') \ar[rr]_{(f, f')} \ar[rd] & & (Y \subset Y') \ar[ld] \\
& (B \subset B')
}
$$
d'épaississements d'espaces algébriques sur $S$. Supposons que
\begin{enumerate}
\item $X'$ soit plat sur $B'$,
\item $f$ soit plat,
\item $B \subset B'$ soit un épaississement d'ordre fini, et
\item $X = B \times_{B'} X'$ et $Y = B \times_{B'} Y'$.
\end{enumerate}
Alors $f'$ est plat et $Y'$ est plat sur $B'$ en tout point de
l'image de $f'$.
\end{lemma}

\begin{proof}
Choisissons un schéma $U'$ et un morphisme étale surjectif $U' \to B'$.
Choisissons un schéma $V'$ et un morphisme étale surjectif
$V' \to U' \times_{B'} Y'$.
Choisissons un schéma $W'$ et un morphisme étale surjectif
$W' \to V' \times_{Y'} X'$. Soient $U, V, W$ les changements de base
de $U', V', W'$ par $B \to B'$. Alors la platitude de $f'$ est
équivalente à celle de $W' \to V'$, et nous savons
que $W \to V$ est plat. Nous pouvons donc appliquer le lemme
dans le cas des schémas au diagramme
$$
\xymatrix{
(W \subset W') \ar[rr] \ar[rd] & & (V \subset V') \ar[ld] \\
& (U \subset U')
}
$$
d'épaississements de schémas. Voir
Compléments sur les morphismes, Lemme
\ref{more-morphisms-lemma-flatness-morphism-thickenings}.
L'assertion relative à la platitude de $Y'/B'$ aux points de
l'image de $f'$ se démontre de la même manière.
\end{proof}

\noindent
De nombreuses propriétés des morphismes de schémas sont préservées par les
déformations plates.

\begin{lemma}
\label{lemma-deform-property}
Soit $S$ un schéma. Considérons un diagramme commutatif
$$
\xymatrix{
(X \subset X') \ar[rr]_{(f, f')} \ar[rd] & & (Y \subset Y') \ar[ld] \\
& (B \subset B')
}
$$
d'épaississements d'espaces algébriques sur $S$. Supposons que $B \subset B'$
soit un épaississement d'ordre fini, que $X'$ soit plat sur $B'$, que $X = B \times_{B'} X'$,
et que $Y = B \times_{B'} Y'$. Alors
\begin{enumerate}
\item $f$ est représentable si et seulement si $f'$ est représentable,
\label{item-representable}
\item $f$ est plat si et seulement si $f'$ est plat,
\label{item-flat}
\item $f$ est un isomorphisme si et seulement si $f'$ est un isomorphisme,
\label{item-isomorphism}
\item $f$ est une immersion ouverte si et seulement si $f'$ est une immersion ouverte,
\label{item-open-immersion}
\item $f$ est quasi-compact si et seulement si $f'$ est quasi-compact,
\label{item-quasi-compact}
\item $f$ est universellement fermé si et seulement si $f'$ est universellement fermé,
\label{item-universally-closed}
\item $f$ est (quasi-)séparé si et seulement si $f'$ est (quasi-)séparé,
\label{item-separated}
\item $f$ est un monomorphisme si et seulement si $f'$ est un monomorphisme,
\label{item-monomorphism}
\item $f$ est surjectif si et seulement si $f'$ est surjectif,
\label{item-surjective}
\item $f$ est universellement injectif si et seulement si $f'$ est universellement injectif,
\label{item-universally-injective}
\item $f$ est affine si et seulement si $f'$ est affine,
\label{item-affine}
\item
\label{item-finite-type}
$f$ est localement de type fini si et seulement si $f'$ est localement de type fini,
\item $f$ est localement quasi-fini si et seulement si $f'$ est localement quasi-fini,
\label{item-quasi-finite}
\item
\label{item-finite-presentation}
$f$ est localement de présentation finie si et seulement si $f'$ est localement de
présentation finie,
\item
\label{item-relative-dimension-d}
$f$ est localement de type fini de dimension relative $d$ si et seulement si
$f'$ est localement de type fini de dimension relative $d$,
\item $f$ est universellement ouvert si et seulement si $f'$ est universellement ouvert,
\label{item-universally-open}
\item $f$ est syntomique si et seulement si $f'$ est syntomique,
\label{item-syntomic}
\item $f$ est lisse si et seulement si $f'$ est lisse,
\label{item-smooth}
\item $f$ est non ramifié si et seulement si $f'$ est non ramifié,
\label{item-unramified}
\item $f$ est étale si et seulement si $f'$ est étale,
\label{item-etale}
\item $f$ est propre si et seulement si $f'$ est propre,
\label{item-proper}
\item $f$ est entier si et seulement si $f'$ est entier,
\label{item-integral}
\item $f$ est fini si et seulement si $f'$ est fini,
\label{item-finite}
\item
\label{item-finite-locally-free}
$f$ est fini localement libre (de rang $d$) si et seulement si $f'$
est fini localement libre (de rang $d$), et
\item ajouter d'autres cas ici.
\end{enumerate}
\end{lemma}

\begin{proof}
Le cas (\ref{item-representable}) résulte du
Lemme \ref{lemma-thicken-property-morphisms}.

\medskip\noindent
Choisissons un schéma $U'$ et un morphisme étale surjectif $U' \to B'$.
Choisissons un schéma $V'$ et un morphisme étale surjectif
$V' \to U' \times_{B'} Y'$.
Choisissons un schéma $W'$ et un morphisme étale surjectif
$W' \to V' \times_{Y'} X'$. Soient $U, V, W$ les changements de base
de $U', V', W'$ par $B \to B'$. Considérons le diagramme
$$
\xymatrix{
(W \subset W') \ar[rr] \ar[rd] & & (V \subset V') \ar[ld] \\
& (U \subset U')
}
$$
d'épaississements de schémas. Pour chacune des propriétés qui sont
locales pour la topologie étale à la source et au but, le résultat découle immédiatement
du résultat correspondant pour les morphismes d'épaississements de schémas,
appliqué au diagramme ci-dessus. Ainsi, les cas
(\ref{item-flat}), (\ref{item-finite-type}),
(\ref{item-quasi-finite}), (\ref{item-finite-presentation}),
(\ref{item-relative-dimension-d}), (\ref{item-syntomic}),
(\ref{item-smooth}), (\ref{item-unramified}), (\ref{item-etale})
résultent des cas correspondants de
Compléments sur les morphismes, Lemme \ref{more-morphisms-lemma-deform-property}.

\medskip\noindent
Puisque $X \to X'$ et $Y \to Y'$ sont des homéomorphismes universels,
toute question portant sur la topologie des morphismes
$X \to Y$ et $X' \to Y'$ reçoit la même réponse. Ainsi,
les cas (\ref{item-quasi-compact}), (\ref{item-universally-closed}),
(\ref{item-surjective}), (\ref{item-universally-injective}), et
(\ref{item-universally-open}) sont vérifiés.

\medskip\noindent
Dans chacun des cas restants, nous démontrons seulement l'implication
$f\text{ possède }P \Rightarrow f'\text{ possède }P$, car l'autre
implication résulte du fait que $P$ est stable par
changement de base, voir
Espaces, Lemme \ref{spaces-lemma-base-change-immersions} et
Morphismes d'espaces, Lemmes
\ref{spaces-morphisms-lemma-base-change-separated},
\ref{spaces-morphisms-lemma-base-change-monomorphism},
\ref{spaces-morphisms-lemma-base-change-affine},
\ref{spaces-morphisms-lemma-base-change-proper},
\ref{spaces-morphisms-lemma-base-change-integral}, et
\ref{spaces-morphisms-lemma-base-change-finite-locally-free}.

\medskip\noindent
Le cas (\ref{item-open-immersion}). Supposons que $f$ soit une immersion ouverte.
Alors $f'$ est étale d'après (\ref{item-etale}) et universellement injectif
d'après (\ref{item-universally-injective}) ;
ainsi, $f'$ est une immersion ouverte, voir
Morphismes d'espaces, Lemme
\ref{spaces-morphisms-lemma-etale-universally-injective-open}.
On peut éviter d'utiliser ce lemme en
utilisant d'abord (\ref{item-representable}) pour se ramener au cas des schémas.

\medskip\noindent
Le cas (\ref{item-isomorphism}) résulte des cas
(\ref{item-open-immersion}) et (\ref{item-surjective}).

\medskip\noindent
Le cas (\ref{item-separated}). Voir le
Lemme \ref{lemma-thicken-property-morphisms}.

\medskip\noindent
Le cas (\ref{item-monomorphism}). Supposons que $f$ soit un monomorphisme.
Considérons le morphisme diagonal $\Delta_{X'/Y'} : X' \to X' \times_{Y'} X'$.
Le changement de base de $\Delta_{X'/Y'}$ par $B \to B'$ est $\Delta_{X/Y}$,
qui est un isomorphisme par hypothèse. D'après (\ref{item-isomorphism}),
nous concluons que $\Delta_{X'/Y'}$ est un isomorphisme et donc que
$f'$ est un monomorphisme.

\medskip\noindent
Le cas (\ref{item-affine}). Voir le Lemme \ref{lemma-thicken-property-morphisms}.

\medskip\noindent
Le cas (\ref{item-proper}). Voir le
Lemme \ref{lemma-thicken-property-morphisms-cartesian}.

\medskip\noindent
Le cas (\ref{item-integral}). Voir le
Lemme \ref{lemma-thicken-property-morphisms}.

\medskip\noindent
Le cas (\ref{item-finite}). Voir le
Lemme \ref{lemma-thicken-property-morphisms-cartesian}.

\medskip\noindent
Le cas (\ref{item-finite-locally-free}). Supposons $f$ fini localement libre.
D'après (\ref{item-finite}), $f'$ est fini.
D'après (\ref{item-flat}), $f'$ est plat.
D'après (\ref{item-finite-presentation}), $f'$ est localement de
présentation finie. Ainsi, $f'$ est fini localement libre d'après
Morphismes d'espaces, Lemme \ref{spaces-morphisms-lemma-finite-flat}.
\end{proof}

\noindent
Le lemme suivant est la version « localement nilpotente » du
« critère de platitude par fibres », voir la
section \ref{section-criterion-flat-fibres}.

\begin{lemma}
\label{lemma-flatness-morphism-thickenings-fp-over-ft}
Soit $S$ un schéma. Considérons un diagramme commutatif
$$
\xymatrix{
(X \subset X') \ar[rr]_{(f, f')} \ar[rd] & & (Y \subset Y') \ar[ld] \\
& (B \subset B')
}
$$
d'épaississements d'espaces algébriques sur $S$. Supposons que
\begin{enumerate}
\item $Y' \to B'$ soit localement de type fini,
\item $X' \to B'$ soit plat et localement de présentation finie,
\item $f$ soit plat, et
\item $X = B \times_{B'} X'$ et $Y = B \times_{B'} Y'$.
\end{enumerate}
Alors $f'$ est plat et, pour tout $y' \in |Y'|$ dans l'image de $|f'|$,
le morphisme $Y' \to B'$ est plat en $y'$.
\end{lemma}

\begin{proof}
Choisissons un schéma $U'$ et un morphisme étale surjectif $U' \to B'$.
Choisissons un schéma $V'$ et un morphisme étale surjectif
$V' \to U' \times_{B'} Y'$.
Choisissons un schéma $W'$ et un morphisme étale surjectif
$W' \to V' \times_{Y'} X'$. Soient $U, V, W$ les changements de base
de $U', V', W'$ par $B \to B'$. Alors la platitude de $f'$ est
équivalente à celle de $W' \to V'$, et nous savons
que $W \to V$ est plat. Nous pouvons donc appliquer le lemme
dans le cas des schémas au diagramme
$$
\xymatrix{
(W \subset W') \ar[rr] \ar[rd] & & (V \subset V') \ar[ld] \\
& (U \subset U')
}
$$
d'épaississements de schémas. Voir
Compléments sur les morphismes, Lemme
\ref{more-morphisms-lemma-flatness-morphism-thickenings-fp-over-ft}.
L'assertion relative à la platitude de $Y'/B'$ aux points de
l'image de $f'$ se démontre de la même manière.
\end{proof}

\noindent
De nombreuses propriétés des morphismes de schémas sont préservées par les
déformations plates comme dans le lemme ci-dessus.

\begin{lemma}
\label{lemma-deform-property-fp-over-ft}
Soit $S$ un schéma. Considérons un diagramme commutatif
$$
\xymatrix{
(X \subset X') \ar[rr]_{(f, f')} \ar[rd] & & (Y \subset Y') \ar[ld] \\
& (B \subset B')
}
$$
d'épaississements d'espaces algébriques sur $S$.
Supposons que $Y' \to B'$ soit localement de type fini,
que $X' \to B'$ soit plat et localement de présentation finie,
que $X = B \times_{B'} X'$ et que $Y = B \times_{B'} Y'$. Alors
\begin{enumerate}
\item $f$ est représentable si et seulement si $f'$ est représentable,
\label{item-representable-fp-over-ft}
\item $f$ est plat si et seulement si $f'$ est plat,
\label{item-flat-fp-over-ft}
\item $f$ est un isomorphisme si et seulement si $f'$ est un isomorphisme,
\label{item-isomorphism-fp-over-ft}
\item $f$ est une immersion ouverte si et seulement si $f'$ est une immersion ouverte,
\label{item-open-immersion-fp-over-ft}
\item $f$ est quasi-compact si et seulement si $f'$ est quasi-compact,
\label{item-quasi-compact-fp-over-ft}
\item $f$ est universellement fermé si et seulement si $f'$ est universellement fermé,
\label{item-universally-closed-fp-over-ft}
\item $f$ est (quasi-)séparé si et seulement si $f'$ est (quasi-)séparé,
\label{item-separated-fp-over-ft}
\item $f$ est un monomorphisme si et seulement si $f'$ est un monomorphisme,
\label{item-monomorphism-fp-over-ft}
\item $f$ est surjectif si et seulement si $f'$ est surjectif,
\label{item-surjective-fp-over-ft}
\item $f$ est universellement injectif si et seulement si $f'$ est universellement injectif,
\label{item-universally-injective-fp-over-ft}
\item $f$ est affine si et seulement si $f'$ est affine,
\label{item-affine-fp-over-ft}
\item $f$ est localement quasi-fini si et seulement si $f'$ est localement quasi-fini,
\label{item-quasi-finite-fp-over-ft}
\item
\label{item-relative-dimension-d-fp-over-ft}
$f$ est localement de type fini de dimension relative $d$ si et seulement si
$f'$ est localement de type fini de dimension relative $d$,
\item $f$ est universellement ouvert si et seulement si $f'$ est universellement ouvert,
\label{item-universally-open-fp-over-ft}
\item $f$ est syntomique si et seulement si $f'$ est syntomique,
\label{item-syntomic-fp-over-ft}
\item $f$ est lisse si et seulement si $f'$ est lisse,
\label{item-smooth-fp-over-ft}
\item $f$ est non ramifié si et seulement si $f'$ est non ramifié,
\label{item-unramified-fp-over-ft}
\item $f$ est étale si et seulement si $f'$ est étale,
\label{item-etale-fp-over-ft}
\item $f$ est propre si et seulement si $f'$ est propre,
\label{item-proper-fp-over-ft}
\item $f$ est fini si et seulement si $f'$ est fini,
\label{item-finite-fp-over-ft}
\item
\label{item-finite-locally-free-fp-over-ft}
$f$ est fini localement libre (de rang $d$) si et seulement si $f'$
est fini localement libre (de rang $d$), et
\item ajouter d'autres cas ici.
\end{enumerate}
\end{lemma}

\begin{proof}
Le cas (\ref{item-representable-fp-over-ft}) résulte du
Lemme \ref{lemma-thicken-property-morphisms}.

\medskip\noindent
Choisissons un schéma $U'$ et un morphisme étale surjectif $U' \to B'$.
Choisissons un schéma $V'$ et un morphisme étale surjectif
$V' \to U' \times_{B'} Y'$.
Choisissons un schéma $W'$ et un morphisme étale surjectif
$W' \to V' \times_{Y'} X'$. Soient $U, V, W$ les changements de base
de $U', V', W'$ par $B \to B'$. Considérons le diagramme
$$
\xymatrix{
(W \subset W') \ar[rr] \ar[rd] & & (V \subset V') \ar[ld] \\
& (U \subset U')
}
$$
d'épaississements de schémas. Pour chacune des propriétés qui sont
locales pour la topologie étale à la source et au but, le résultat découle immédiatement
du résultat correspondant pour les morphismes d'épaississements de schémas,
appliqué au diagramme ci-dessus. Ainsi, les cas
(\ref{item-flat-fp-over-ft}),
(\ref{item-quasi-finite-fp-over-ft}),
(\ref{item-relative-dimension-d-fp-over-ft}),
(\ref{item-syntomic-fp-over-ft}),
(\ref{item-smooth-fp-over-ft}),
(\ref{item-unramified-fp-over-ft}),
(\ref{item-etale-fp-over-ft})
résultent des cas correspondants de
Compléments sur les morphismes, Lemme \ref{more-morphisms-lemma-deform-property-fp-over-ft}.

\medskip\noindent
Puisque $X \to X'$ et $Y \to Y'$ sont des homéomorphismes universels,
toute question portant sur la topologie des morphismes
$X \to Y$ et $X' \to Y'$ reçoit la même réponse. Ainsi,
les cas (\ref{item-quasi-compact-fp-over-ft}),
(\ref{item-universally-closed-fp-over-ft}),
(\ref{item-surjective-fp-over-ft}),
(\ref{item-universally-injective-fp-over-ft}), et
(\ref{item-universally-open-fp-over-ft}) sont vérifiés.

\medskip\noindent
Dans chacun des cas restants, nous démontrons seulement l'implication
$f\text{ possède }P \Rightarrow f'\text{ possède }P$, car l'autre
implication résulte du fait que $P$ est stable par
changement de base, voir
Espaces, Lemme \ref{spaces-lemma-base-change-immersions} et
Morphismes d'espaces, Lemmes
\ref{spaces-morphisms-lemma-base-change-separated},
\ref{spaces-morphisms-lemma-base-change-monomorphism},
\ref{spaces-morphisms-lemma-base-change-affine},
\ref{spaces-morphisms-lemma-base-change-proper},
\ref{spaces-morphisms-lemma-base-change-integral}, et
\ref{spaces-morphisms-lemma-base-change-finite-locally-free}.

\medskip\noindent
Le cas (\ref{item-open-immersion-fp-over-ft}).
Supposons que $f$ soit une immersion ouverte.
Alors $f'$ est étale d'après (\ref{item-etale-fp-over-ft}) et universellement injectif
d'après (\ref{item-universally-injective-fp-over-ft}) ;
ainsi, $f'$ est une immersion ouverte, voir
Morphismes d'espaces, Lemme
\ref{spaces-morphisms-lemma-etale-universally-injective-open}.
On peut éviter d'utiliser ce lemme en
utilisant d'abord (\ref{item-representable-fp-over-ft}) pour se ramener au cas des schémas.

\medskip\noindent
Le cas (\ref{item-isomorphism-fp-over-ft}) résulte des cas
(\ref{item-open-immersion-fp-over-ft}) et (\ref{item-surjective-fp-over-ft}).

\medskip\noindent
Le cas (\ref{item-separated-fp-over-ft}). Voir le
Lemme \ref{lemma-thicken-property-morphisms}.

\medskip\noindent
Le cas (\ref{item-monomorphism-fp-over-ft}). Supposons que $f$ soit un monomorphisme.
Considérons le morphisme diagonal $\Delta_{X'/Y'} : X' \to X' \times_{Y'} X'$.
Le changement de base de $\Delta_{X'/Y'}$ par $B \to B'$ est $\Delta_{X/Y}$,
qui est un isomorphisme par hypothèse. D'après (\ref{item-isomorphism-fp-over-ft}),
nous concluons que $\Delta_{X'/Y'}$ est un isomorphisme et donc que
$f'$ est un monomorphisme.

\medskip\noindent
Le cas (\ref{item-affine-fp-over-ft}).
Voir le Lemme \ref{lemma-thicken-property-morphisms}.

\medskip\noindent
Le cas (\ref{item-proper-fp-over-ft}). Voir le
Lemme \ref{lemma-properties-that-extend-over-thickenings}.

\medskip\noindent
Le cas (\ref{item-finite-fp-over-ft}). Voir le
Lemme \ref{lemma-properties-that-extend-over-thickenings}.

\medskip\noindent
Le cas (\ref{item-finite-locally-free-fp-over-ft}).
Supposons $f$ fini localement libre.
D'après (\ref{item-finite-fp-over-ft}), $f'$ est fini.
D'après (\ref{item-flat-fp-over-ft}), $f'$ est plat.
De plus, $f'$ est localement de présentation finie d'après
Morphismes d'espaces, Lemme
\ref{spaces-morphisms-lemma-finite-presentation-permanence}.
Ainsi, $f'$ est fini localement libre d'après
Morphismes d'espaces, Lemme \ref{spaces-morphisms-lemma-finite-flat}.
\end{proof}













\section{Morphismes formellement lisses}
\label{section-formally-smooth}

\noindent
Dans cette section, nous introduisons la notion de morphisme formellement lisse
$X \to Y$ d'espaces algébriques. Un tel morphisme est
caractérisé par la propriété que, pour $T$, les points de $X$ à valeurs dans ce schéma se relèvent
aux épaississements infinitésimaux de $T$, pourvu que $T$ soit affine.
Le résultat principal affirme qu'un morphisme formellement lisse et
localement de présentation finie est lisse, voir le
Lemme \ref{lemma-smooth-formally-smooth}.
Il se trouve que ce critère est souvent plus facile à utiliser que le
critère jacobien.

\begin{definition}
\label{definition-formally-smooth}
Soit $S$ un schéma. Un morphisme $f : X \to Y$ d'espaces algébriques sur $S$
est dit {\it formellement lisse} s'il est formellement lisse comme
transformation de foncteurs au sens de la
Définition \ref{definition-formally-smooth-etale-unramified}.
\end{definition}

\noindent
Dans le cas des morphismes formellement non ramifiés et formellement étales,
on pourrait omettre l'hypothèse que $T'$ soit affine, voir les
Lemmes \ref{lemma-formally-unramified-not-affine} et
\ref{lemma-formally-etale-not-affine}.
Ce n'est plus vrai dans le cas des morphismes formellement lisses.
En fait, une condition un peu plus naturelle consisterait à pouvoir
compléter la flèche pointillée localement pour la topologie étale sur $T'$.
En effet, l'analyse de la démonstration du
Lemme \ref{lemma-smooth-formally-smooth}
montre que cela serait équivalent à la définition sous sa forme
actuelle. Il est aussi vrai qu'exiger l'existence de la flèche
pointillée localement pour la topologie fppf sur $T'$ suffirait, mais cela est un peu
plus difficile à démontrer.

\medskip\noindent
Nous ne reformulerons pas les résultats démontrés dans le cadre plus général des
transformations de foncteurs formellement lisses à la
section \ref{section-formally-smooth-etale-unramified}.

\begin{lemma}
\label{lemma-composition-formally-smooth}
Une composée de morphismes formellement lisses est formellement lisse.
\end{lemma}

\begin{proof}
Démonstration omise.
\end{proof}

\begin{lemma}
\label{lemma-base-change-formally-smooth}
Un changement de base d'un morphisme formellement lisse est formellement lisse.
\end{lemma}

\begin{proof}
Démonstration omise, mais voir
Algèbre, Lemme \ref{algebra-lemma-base-change-fs}
pour la version algébrique.
\end{proof}

\begin{lemma}
\label{lemma-formally-etale-unramified-smooth}
Soit $f : X \to S$ un morphisme de schémas.
Alors $f$ est formellement étale si et seulement si
$f$ est formellement lisse et formellement non ramifié.
\end{lemma}

\begin{proof}
Démonstration omise.
\end{proof}

\noindent
Voici un lemme auxiliaire qui sera englobé par le
Lemme \ref{lemma-formally-smooth}.

\begin{lemma}
\label{lemma-helper-formally-smooth}
Soit $S$ un schéma. Soit
$$
\xymatrix{
U \ar[d] \ar[r]_\psi & V \ar[d] \\
X \ar[r]^f & Y
}
$$
un diagramme commutatif de morphismes d'espaces algébriques sur $S$.
Si les flèches verticales sont étales et si $f$ est formellement lisse, alors
$\psi$ est formellement lisse.
\end{lemma}

\begin{proof}
D'après le
Lemme \ref{lemma-representable-property-formally-property},
les morphismes $U \to X$ et $V \to Y$ sont formellement étales. D'après le
Lemme \ref{lemma-composition-formally-smooth-etale-unramified},
la composée $U \to Y$ est formellement lisse. D'après le
Lemme \ref{lemma-formally-permanence},
nous voyons que $\psi : U \to V$ est formellement lisse.
\end{proof}

\noindent
Le lemme suivant est le résultat principal de cette section.
Combiné avec la
Proposition de Limites d'espaces
\ref{spaces-limits-proposition-characterize-locally-finite-presentation},
il implique que l'on peut reconnaître si un morphisme d'espaces algébriques
$f : X \to Y$ est lisse au moyen de propriétés « simples » de la
transformation de foncteurs $X \to Y$.

\begin{lemma}[Critère infinitésimal de relèvement]
\label{lemma-smooth-formally-smooth}
Soit $S$ un schéma.
Soit $f : X \to Y$ un morphisme d'espaces algébriques sur $S$.
Les assertions suivantes sont équivalentes :
\begin{enumerate}
\item Le morphisme $f$ est lisse.
\item Le morphisme $f$ est localement de présentation finie et
formellement lisse.
\end{enumerate}
\end{lemma}

\begin{proof}
Supposons que $f : X \to S$ soit localement de présentation finie et formellement lisse.
Considérons un diagramme commutatif
$$
\xymatrix{
U \ar[d] \ar[r]_\psi & V \ar[d] \\
X \ar[r]^f & Y
}
$$
où $U$ et $V$ sont des schémas et où les flèches verticales sont étales et
surjectives. D'après le
Lemme \ref{lemma-helper-formally-smooth},
nous voyons que $\psi : U \to V$ est formellement lisse. D'après
Morphismes d'espaces, Lemme
\ref{spaces-morphisms-lemma-finite-presentation-local}
le morphisme $\psi$ est localement de présentation finie.
Par conséquent, dans le cas des schémas, le morphisme
$\psi$ est lisse, voir
Compléments sur les morphismes, Lemme \ref{more-morphisms-lemma-smooth-formally-smooth}.
Ainsi, $f$ est lisse, voir
Morphismes d'espaces, Lemme
\ref{spaces-morphisms-lemma-smooth-local}.

\medskip\noindent
Réciproquement, supposons que $f : X \to Y$ soit lisse.
Considérons un diagramme commutatif de flèches pleines
$$
\xymatrix{
X \ar[d]_f & T \ar[d]^i \ar[l]^a \\
Y & T' \ar[l] \ar@{-->}[lu]
}
$$
comme dans la Définition \ref{definition-formally-smooth}. Nous allons montrer que la
flèche pointillée existe, ce qui démontrera que $f$ est formellement lisse.
Soit $\mathcal{F}$ le faisceau d'ensembles sur $(T')_{spaces, \etale}$ du
Lemme \ref{lemma-sheaf}, dans le cas particulier examiné à la
Remarque \ref{remark-special-case}. Soit
$$
\mathcal{H} =
\SheafHom_{\mathcal{O}_T}(a^*\Omega_{X/Y}, \mathcal{C}_{T/T'})
$$
le faisceau de $\mathcal{O}_T$-modules sur $T_{spaces, \etale}$
muni de l'action $\mathcal{H} \times \mathcal{F} \to \mathcal{F}$
du Lemme \ref{lemma-action-sheaf}.
L'action $\mathcal{H} \times \mathcal{F} \to \mathcal{F}$
fait de $\mathcal{F}$ un pseudo-torseur sous $\mathcal{H}$, voir
Cohomologie sur les sites, Définition \ref{sites-cohomology-definition-torsor}.
Notre but est de montrer que $\mathcal{F}$ est un torseur trivial sous $\mathcal{H}$.
Il y a deux étapes : (I) Pour montrer que $\mathcal{F}$ est un torseur,
nous devons montrer que $\mathcal{F}$ possède localement pour la topologie étale
une section. (II) Pour montrer que $\mathcal{F}$ est le torseur trivial,
il suffit de montrer que $H^1(T_\etale, \mathcal{H}) = 0$, voir
Cohomologie sur les sites, Lemme \ref{sites-cohomology-lemma-torsors-h1}.

\medskip\noindent
Démontrons d'abord (I). Pour cela, choisissons un diagramme commutatif
$$
\xymatrix{
U \ar[d] \ar[r]_\psi & V \ar[d] \\
X \ar[r]^f & Y
}
$$
où $U$ et $V$ sont des schémas et où les flèches verticales sont étales et
surjectives. Comme $f$ est supposé lisse, nous voyons que $\psi$ est lisse et
donc formellement lisse d'après le
Lemme \ref{lemma-representable-property-formally-property}.
D'après le même lemme, le morphisme $V \to Y$ est formellement étale. Ainsi, d'après le
Lemme \ref{lemma-composition-formally-smooth-etale-unramified},
la composée $U \to Y$ est formellement lisse. L'étape (I) résulte alors du
Lemme \ref{lemma-etale-on-top}, partie (4).

\medskip\noindent
Démontrons enfin (II). D'après le
Lemme \ref{lemma-finite-presentation-differentials},
nous voyons que $\Omega_{X/S}$ est de présentation finie.
Ainsi, $a^*\Omega_{X/S}$ est de présentation finie (voir
Propriétés d'espaces,
section \ref{spaces-properties-section-properties-modules}).
Par conséquent, le faisceau
$\mathcal{H} =
\SheafHom_{\mathcal{O}_T}(a^*\Omega_{X/Y}, \mathcal{C}_{T/T'})$
est quasi-cohérent d'après
Propriétés d'espaces,
Lemme \ref{spaces-properties-lemma-properties-quasi-coherent}.
Ainsi, d'après
Descente, Proposition \ref{descent-proposition-same-cohomology-quasi-coherent}
et Cohomologie des schémas, Lemme
\ref{coherent-lemma-quasi-coherent-affine-cohomology-zero}
nous avons
$$
H^1(T_{spaces, \etale}, \mathcal{H}) =
H^1(T_\etale, \mathcal{H}) =
H^1(T, \mathcal{H}) = 0
$$
comme voulu.
\end{proof}

\noindent
Les morphismes lisses satisfont une forte propriété locale de relèvement, voir le
Lemme \ref{lemma-smooth-strong-lift}. Si, dans le lemme, nous
supposons que $T'$ est affine, alors
nous ignorons s'il est nécessaire de prendre un recouvrement étale.
Plus précisément, si nous avons un diagramme commutatif
$$
\xymatrix{
X \ar[d] & T \ar[l] \ar[d] \\
Y & T' \ar[l] \ar@{..>}[lu]
}
$$
d'espaces algébriques où $X \to Y$ est lisse
et où $T \to T'$ est un épaississement de schémas affines,
existe-t-il toujours une flèche pointillée rendant le diagramme
commutatif ? Si vous connaissez la réponse ou une référence, veuillez écrire à
\href{mailto:stacks.project@gmail.com}{stacks.project@gmail.com}.

\begin{lemma}
\label{lemma-smooth-strong-lift}
Soit $S$ un schéma. Considérons un diagramme commutatif
$$
\xymatrix{
X \ar[d] & T \ar[l] \ar[d] \\
Y & T' \ar[l]
}
$$
d'espaces algébriques sur $S$, où $X \to Y$ est lisse
et où $T \to T'$ est un épaississement. Alors il existe un
recouvrement étale $\{T'_i \to T'\}$ tel que nous
puissions trouver la flèche pointillée dans
$$
\xymatrix{
X \ar[d] & T \ar[l] \ar[d] & T \times_{T'} T'_i \ar[l] \ar[d] \\
Y & T' \ar[l] & T'_i \ar[l] \ar@{..>}[llu]
}
$$
qui rend le diagramme commutatif (pour tout $i$).
\end{lemma}

\begin{proof}
Choisissons un recouvrement étale $\{Y_i \to Y\}$ où chaque $Y_i$ est affine.
Après avoir remplacé $T'$ par le recouvrement étale induit, nous pouvons supposer
que $Y$ est affine.

\medskip\noindent
Supposons que $Y$ soit affine. Choisissons un recouvrement étale $\{X_i \to X\}$.
Il donne naissance à un recouvrement étale de $T$. Ce recouvrement étale de $T$
provient d'un recouvrement étale de $T'$
(d'après le Théorème \ref{theorem-topological-invariance}, voir la
discussion de la section \ref{section-thickenings}).
Nous pouvons donc supposer que $X$ est affine.

\medskip\noindent
Supposons que $X$ et $Y$ soient affines. Nous pouvons encore remplacer
$T'$ par un recouvrement étale et supposer que $T'$ est affine. Dans ce cas, le lemme résulte de
Algèbre, Lemme \ref{algebra-lemma-smooth-strong-lift}.
\end{proof}







\noindent
Nous poursuivons un peu pour montrer que la propriété d'être formellement lisse est locale
pour la topologie étale à la source. Pour commencer, nous montrons qu'un morphisme formellement lisse possède un
faisceau de différentielles bien comporté. La notion de module quasi-cohérent
localement projectif est définie dans
Propriétés d'espaces,
section \ref{spaces-properties-section-locally-projective}.

\begin{lemma}
\label{lemma-formally-smooth-sheaf-differentials}
Soit $S$ un schéma.
Soit $f : X \to Y$ un morphisme formellement lisse d'espaces algébriques sur $S$.
Alors $\Omega_{X/Y}$ est localement projectif sur $X$.
\end{lemma}

\begin{proof}
Choisissons un diagramme
$$
\xymatrix{
U \ar[d] \ar[r]_\psi & V \ar[d] \\
X \ar[r]^f & Y
}
$$
où $U$ et $V$ sont des schémas affines (!) et où les flèches verticales sont étales.
D'après le
Lemme \ref{lemma-helper-formally-smooth},
nous voyons que $\psi : U \to V$ est formellement lisse. Ainsi,
$\Gamma(V, \mathcal{O}_V) \to \Gamma(U, \mathcal{O}_U)$ est
un homomorphisme d'anneaux formellement lisse, voir
Compléments sur les morphismes, Lemme \ref{more-morphisms-lemma-affine-formally-smooth}.
Par conséquent, d'après
Algèbre, Lemme \ref{algebra-lemma-characterize-formally-smooth-again},
le $\Gamma(U, \mathcal{O}_U)$-module
$\Omega_{\Gamma(U, \mathcal{O}_U)/\Gamma(V, \mathcal{O}_V)}$
est projectif. Ainsi, $\Omega_{U/V}$ est localement projectif, voir
Propriétés, section \ref{properties-section-locally-projective}.
Puisque $\Omega_{X/Y}|_U = \Omega_{U/V}$, nous voyons que $\Omega_{X/Y}$ est
lui aussi localement projectif. (En effet, nous pouvons trouver un recouvrement étale de
$X$ par des $U$ affines s'insérant dans des diagrammes comme ci-dessus --- les détails sont
omis.)
\end{proof}

\begin{lemma}
\label{lemma-h1-is-zero}
Soit $T$ un schéma affine.
Soient $\mathcal{F}$ et $\mathcal{G}$ des
$\mathcal{O}_T$-modules quasi-cohérents sur $T_\etale$.
Considérons le faisceau Hom interne
$\mathcal{H} = \SheafHom_{\mathcal{O}_T}(\mathcal{F}, \mathcal{G})$
sur $T_\etale$.
Si $\mathcal{F}$ est localement projectif, alors
$H^1(T_\etale, \mathcal{H}) = 0$.
\end{lemma}

\begin{proof}
D'après la définition d'un faisceau localement projectif sur un espace algébrique (voir
Propriétés d'espaces,
Définition \ref{spaces-properties-definition-locally-projective}),
nous voyons que $\mathcal{F}_{Zar} = \mathcal{F}|_{T_{Zar}}$ est un faisceau
localement projectif sur le schéma $T$. Ainsi, $\mathcal{F}_{Zar}$ est un
facteur direct d'un $\mathcal{O}_{T_{Zar}}$-module libre. Nous
concluons alors (puisque $\mathcal{F} = (\mathcal{F}_{Zar})^a$, voir
Descente, Proposition \ref{descent-proposition-equivalence-quasi-coherent})
que $\mathcal{F}$ est un facteur direct d'un $\mathcal{O}_T$-module libre
sur $T_\etale$. Nous pouvons donc supposer que
$\mathcal{F} = \bigoplus_{i \in I} \mathcal{O}_T$ est un module libre.
Dans ce cas, $\mathcal{H} = \prod_{i \in I} \mathcal{G}$ est
un produit de modules quasi-cohérents. D'après
Cohomologie sur les sites, Lemme \ref{sites-cohomology-lemma-cohomology-products},
nous concluons que $H^1 = 0$, car la cohomologie d'un faisceau quasi-cohérent
sur un schéma affine est nulle, voir
Descente, Proposition \ref{descent-proposition-same-cohomology-quasi-coherent}
et Cohomologie des schémas, Lemme
\ref{coherent-lemma-quasi-coherent-affine-cohomology-zero}.
\end{proof}

\begin{lemma}
\label{lemma-formally-smooth}
Soit $S$ un schéma.
Soit $f : X \to Y$ un morphisme d'espaces algébriques sur
$S$. Les assertions suivantes sont équivalentes :
\begin{enumerate}
\item $f$ est formellement lisse,
\item pour tout diagramme
$$
\xymatrix{
U \ar[d] \ar[r]_\psi & V \ar[d] \\
X \ar[r]^f & Y
}
$$
où $U$ et $V$ sont des schémas et où les flèches verticales sont étales,
le morphisme de schémas $\psi$ est formellement lisse (au sens de
Compléments sur les morphismes,
Définition \ref{more-morphisms-definition-formally-unramified}), et
\item pour un tel diagramme à flèches verticales surjectives, le morphisme
$\psi$ est formellement lisse.
\end{enumerate}
\end{lemma}

\begin{proof}
Nous avons vu que (1) implique (2) et (3) dans le
Lemme \ref{lemma-helper-formally-smooth}.
Supposons (3).
La démonstration du fait que $f$ est formellement lisse est entièrement semblable à
la démonstration de (1) $\Rightarrow$ (2) du
Lemme \ref{lemma-smooth-formally-smooth}.

\medskip\noindent
Considérons un diagramme commutatif de flèches pleines
$$
\xymatrix{
X \ar[d]_f & T \ar[d]^i \ar[l]^a \\
Y & T' \ar[l] \ar@{-->}[lu]
}
$$
comme dans la Définition \ref{definition-formally-smooth}.
Nous allons montrer que la flèche pointillée existe, ce qui
démontrera que $f$ est formellement lisse.
Soit $\mathcal{F}$ le faisceau d'ensembles sur $(T')_{spaces, \etale}$ du
Lemme \ref{lemma-sheaf}, dans le cas particulier examiné à la
Remarque \ref{remark-special-case}.
Soit
$$
\mathcal{H} =
\SheafHom_{\mathcal{O}_T}(a^*\Omega_{X/Y}, \mathcal{C}_{T/T'})
$$
le faisceau de $\mathcal{O}_T$-modules sur $T_{spaces, \etale}$
muni de l'action $\mathcal{H} \times \mathcal{F} \to \mathcal{F}$ comme dans le
Lemme \ref{lemma-action-sheaf}.
L'action $\mathcal{H} \times \mathcal{F} \to \mathcal{F}$
fait de $\mathcal{F}$ un pseudo-torseur sous $\mathcal{H}$, voir
Cohomologie sur les sites, Définition \ref{sites-cohomology-definition-torsor}.
Notre but est de montrer que $\mathcal{F}$ est un torseur trivial sous $\mathcal{H}$.
Il y a deux étapes : (I) Pour montrer que $\mathcal{F}$ est un torseur,
nous devons montrer que $\mathcal{F}$ possède localement pour la topologie étale
une section. (II) Pour montrer que $\mathcal{F}$ est le torseur trivial,
il suffit de montrer que $H^1(T_\etale, \mathcal{H}) = 0$, voir
Cohomologie sur les sites, Lemme \ref{sites-cohomology-lemma-torsors-h1}.

\medskip\noindent
Démontrons d'abord (I). Pour cela, considérons un diagramme
(qui existe puisque nous supposons (3))
$$
\xymatrix{
U \ar[d] \ar[r]_\psi & V \ar[d] \\
X \ar[r]^f & Y
}
$$
où $U$ et $V$ sont des schémas, où les flèches verticales sont étales et
surjectives, et où $\psi$ est formellement lisse. D'après le
Lemme \ref{lemma-representable-property-formally-property},
le morphisme $V \to Y$ est formellement étale. Ainsi, d'après le
Lemme \ref{lemma-composition-formally-smooth-etale-unramified},
la composée $U \to Y$ est formellement lisse. L'étape (I) résulte alors du
Lemme \ref{lemma-etale-on-top}, partie (4).

\medskip\noindent
Démontrons enfin (II). D'après le
Lemme \ref{lemma-formally-smooth-sheaf-differentials},
nous voyons que $\Omega_{U/V}$ est localement projectif.
Ainsi, $\Omega_{X/Y}$ est localement projectif, voir
Descente sur les espaces,
Lemme \ref{spaces-descent-lemma-locally-projective-descends}.
Par conséquent, $a^*\Omega_{X/Y}$ est localement projectif, voir
Propriétés d'espaces, Lemme
\ref{spaces-properties-lemma-locally-projective-pullback}.
Ainsi,
$$
H^1(T_\etale, \mathcal{H}) =
H^1(T_\etale,
\SheafHom_{\mathcal{O}_T}(a^*\Omega_{X/Y}, \mathcal{C}_{T/T'}) = 0
$$
d'après le
Lemme \ref{lemma-h1-is-zero},
comme voulu.
\end{proof}

\begin{lemma}
\label{lemma-descending-property-formally-smooth}
La propriété $\mathcal{P}(f) =$ « $f$ est formellement lisse »
est locale pour la topologie fpqc sur la base.
\end{lemma}

\begin{proof}
Soit $f : X \to Y$ un morphisme d'espaces algébriques sur un schéma $S$.
Choisissons un ensemble d'indices $I$ et des diagrammes
$$
\xymatrix{
U_i \ar[d] \ar[r]_{\psi_i} & V_i \ar[d] \\
X \ar[r]^f & Y
}
$$
à flèches verticales étales, où $U_i$ et $V_i$ sont des schémas affines. De plus,
supposons que $\coprod U_i \to X$ et $\coprod V_i \to Y$ soient surjectifs, voir
Propriétés d'espaces,
Lemme \ref{spaces-properties-lemma-cover-by-union-affines}.
D'après le
Lemme \ref{lemma-formally-smooth},
nous voyons que $f$ est formellement lisse si et seulement si chacun des morphismes
$\psi_i$ est formellement lisse. Nous nous ramenons donc au cas d'un morphisme
de schémas affines. Dans ce cas, le résultat résulte de
Algèbre, Lemme \ref{algebra-lemma-descent-formally-smooth}.
Certains détails sont omis.
\end{proof}

\begin{lemma}
\label{lemma-triangle-differentials-formally-smooth}
Soit $S$ un schéma.
Soient $f : X \to Y$ et $g : Y \to Z$ des morphismes d'espaces algébriques sur $S$.
Supposons que $f$ soit formellement lisse. Alors la suite
$$
0 \to f^*\Omega_{Y/Z} \to \Omega_{X/Z} \to \Omega_{X/Y} \to 0
$$
du Lemme \ref{lemma-triangle-differentials}
est exacte courte.
\end{lemma}

\begin{proof}
Cela résulte du cas des schémas, voir
Compléments sur les morphismes,
Lemme \ref{more-morphisms-lemma-triangle-differentials-formally-smooth},
par localisation étale, voir les
Lemmes \ref{lemma-formally-smooth} et \ref{lemma-localize-differentials}.
\end{proof}

\begin{lemma}
\label{lemma-differentials-formally-unramified-formally-smooth}
Soit $S$ un schéma. Soit $B$ un espace algébrique sur $S$.
Soit $h : Z \to X$ un morphisme formellement non ramifié d'espaces algébriques
sur $B$.
Supposons que $Z$ soit formellement lisse sur $B$. Alors la
suite exacte canonique
$$
0 \to \mathcal{C}_{Z/X} \to h^*\Omega_{X/B} \to \Omega_{Z/B} \to 0
$$
du
Lemme \ref{lemma-universally-unramified-differentials-sequence}
est exacte courte.
\end{lemma}

\begin{proof}
Soit $Z \to Z'$ l'épaississement universel du premier ordre de $Z$ sur $X$.
La démonstration du
Lemme \ref{lemma-universally-unramified-differentials-sequence}
montre que notre suite s'identifie à la suite
$$
\mathcal{C}_{Z/Z'} \to \Omega_{Z'/B} \otimes \mathcal{O}_Z \to
\Omega_{Z/B} \to 0.
$$
Puisque $Z \to S$ est formellement lisse, nous pouvons, localement pour la topologie étale sur $Z'$, trouver
un inverse à gauche $Z' \to Z$ sur $B$ du morphisme d'inclusion $Z \to Z'$.
Ainsi, la suite est localement scindée pour la topologie étale, voir le
Lemme \ref{lemma-differentials-relative-immersion-section}.
\end{proof}

\begin{lemma}
\label{lemma-two-unramified-morphisms-formally-smooth}
Soit $S$ un schéma. Soit
$$
\xymatrix{
Z \ar[r]_i \ar[rd]_j & X \ar[d]^f \\
& Y
}
$$
un diagramme commutatif d'espaces algébriques sur $S$,
où $i$ et $j$ sont formellement non ramifiés et où $f$ est formellement lisse.
Alors la suite exacte canonique
$$
0 \to
\mathcal{C}_{Z/Y} \to
\mathcal{C}_{Z/X} \to
i^*\Omega_{X/Y} \to 0
$$
du
Lemme \ref{lemma-two-unramified-morphisms}
est exacte et localement scindée.
\end{lemma}

\begin{proof}
Notons $Z \to Z'$ l'épaississement universel du premier ordre de $Z$ sur $X$.
Notons $Z \to Z''$ l'épaississement universel du premier ordre de $Z$ sur $Y$.
D'après le
Lemme \ref{lemma-universally-unramified-differentials-sequence},
il existe un morphisme canonique $Z' \to Z''$ tel que nous ayons un diagramme
commutatif
$$
\xymatrix{
Z \ar[r]_{i'} \ar[rd]_{j'} & Z' \ar[r]_a \ar[d]^k & X \ar[d]^f \\
& Z'' \ar[r]^b & Y
}
$$
La suite ci-dessus s'identifie à la suite
$$
\mathcal{C}_{Z/Z''} \to
\mathcal{C}_{Z/Z'} \to
(i')^*\Omega_{Z'/Z''} \to 0
$$
par nos définitions des faisceaux conormaux aux morphismes formellement non ramifiés.
Soit $U'' \to Z''$ un morphisme étale avec $U''$ affine.
Notons $U \to Z$ et $U' \to Z'$ les schémas affines correspondants,
étales sur $Z$ et $Z'$.
Comme $f$ est formellement lisse, il existe un morphisme $h : U'' \to X$
qui coïncide avec $i$ sur $U$ et tel que $f \circ h$ soit égal à $b|_{U''}$.
Puisque $Z'$ est l'épaississement universel du premier ordre, nous obtenons un unique
morphisme $g : U'' \to Z'$ tel que $g = a \circ h$. La propriété
universelle de $Z''$ implique que $k \circ g$ est le morphisme d'inclusion
$U'' \to Z''$. Ainsi, $g$ est un inverse à gauche de $k$. La situation est
$$
\xymatrix{
U \ar[d] \ar[r] & Z' \ar[d]^k \\
U'' \ar[r] \ar[ru]^g & Z''
}
$$
Ainsi, $g$ induit un morphisme $\mathcal{C}_{Z/Z'}|_U \to \mathcal{C}_{Z/Z''}|_U$
qui est un inverse à gauche du morphisme
$\mathcal{C}_{Z/Z''} \to \mathcal{C}_{Z/Z'}$ sur $U$.
\end{proof}











\section{Lissité sur une base noethérienne}
\label{section-smooth-Noetherian}

\noindent
Cette section est l'analogue de
Compléments sur les morphismes, section \ref{more-morphisms-section-smooth-Noetherian}.

\begin{lemma}
\label{lemma-lifting-along-artinian-at-point}
Soit $S$ un schéma. Soit $f : X \to Y$ un morphisme d'espaces algébriques
sur $S$. Soit $x \in |X|$.
Supposons que $Y$ soit localement noethérien et $f$ localement de type fini.
Les assertions suivantes sont équivalentes :
\begin{enumerate}
\item $f$ est lisse en $x$,
\item pour tout diagramme commutatif de flèches pleines
$$
\xymatrix{
X \ar[d]_f & \Spec(B) \ar[d]^i \ar[l]^-\alpha \\
Y & \Spec(B') \ar[l]_-{\beta} \ar@{-->}[lu]
}
$$
où $B' \to B$ est une surjection d'anneaux locaux telle que
$\Ker(B' \to B)$ soit de carré nul, et où $\alpha$ envoie le
point fermé de $\Spec(B)$ sur $x$, il existe
une flèche pointillée qui rend le diagramme commutatif, et
\item la même condition qu'en (2), mais où $B' \to B$ parcourt les petites
extensions (voir Algèbre, Définition \ref{algebra-definition-small-extension}).
\end{enumerate}
\end{lemma}

\begin{proof}
La condition (1) signifie qu'il existe un sous-espace ouvert $X' \subset X$
tel que $X' \to Y$ soit lisse. Ainsi, (1) implique les conditions (2) et (3) d'après le
Lemme \ref{lemma-smooth-formally-smooth}. La condition (2) implique
trivialement la condition (3). Supposons (3). Choisissons un diagramme commutatif
$$
\xymatrix{
X \ar[d] & U \ar[l] \ar[d] \\
Y & V \ar[l]
}
$$
où $U$ et $V$ sont affines, où les flèches horizontales sont étales, et
tel qu'il existe un point $u \in U$ envoyé sur $x$. Considérons ensuite
un diagramme
$$
\xymatrix{
X \ar[d] & U \ar[l] \ar[d] & \Spec(B) \ar[d]^i \ar[l]^-\alpha  \\
Y & V \ar[l] & \Spec(B') \ar[l]_-{\beta}
}
$$
comme en (3), mais pour $u \in U \to V$. Soit $\gamma : \Spec(B') \to X$
la flèche fournie par notre hypothèse que (3) vaut pour $X$.
Puisque $U \to X$ est étale et
donc formellement étale (Lemme \ref{lemma-etale-formally-etale}),
le morphisme $\gamma$
possède un unique relèvement à $U$ compatible avec $\alpha$. Puisque
$V \to Y$ est étale, donc formellement étale, ce relèvement est compatible
avec $\beta$. Ainsi, (3) vaut pour $u \in U \to V$, et nous concluons
que $U \to V$ est lisse en $u$ d'après
Compléments sur les morphismes, Lemme
\ref{more-morphisms-lemma-lifting-along-artinian-at-point}.
Cela montre que $X \to Y$ est lisse en $x$, et
achève la démonstration.
\end{proof}

\noindent
Il est parfois utile de savoir qu'il suffit de vérifier le
critère de relèvement pour les petites extensions « centrées » en des points
de type fini (voir
Morphismes d'espaces, section \ref{spaces-morphisms-section-points-finite-type}).

\begin{lemma}
\label{lemma-lifting-along-artinian}
Soit $S$ un schéma. Soit $f : X \to Y$ un morphisme d'espaces algébriques
sur $S$. Supposons que $Y$ soit localement noethérien et $f$ localement de type fini.
Les assertions suivantes sont équivalentes :
\begin{enumerate}
\item $f$ est lisse,
\item pour tout diagramme commutatif de flèches pleines
$$
\xymatrix{
X \ar[d]_f & \Spec(B) \ar[d]^i \ar[l]^-\alpha \\
Y & \Spec(B') \ar[l]_-{\beta} \ar@{-->}[lu]
}
$$
où $B' \to B$ est une petite extension d'anneaux locaux artiniens
et où $\beta$ est de type fini (!), il existe une flèche pointillée qui rend
le diagramme commutatif.
\end{enumerate}
\end{lemma}

\begin{proof}
Si $f$ est lisse, alors le critère infinitésimal de relèvement
(Lemme \ref{lemma-smooth-formally-smooth}) affirme que
$f$ est formellement lisse, et (2) est vérifiée.

\medskip\noindent
Supposons que $f$ ne soit pas lisse. L'ensemble des points $x \in X$ où $f$ n'est pas lisse
forme une partie fermée $T$ de $|X|$. D'après
Morphismes d'espaces, Lemme 
\ref{spaces-morphisms-lemma-enough-finite-type-points}, il existe
un point $x \in T \subset X$ tel que $x \in X_{\text{ft-pts}}$. Choisissons un
diagramme commutatif
$$
\xymatrix{
X \ar[d] & U \ar[l] \ar[d] & u \ar@{|->}[d] \\
Y & V \ar[l] & v
}
$$
où $U$ et $V$ sont affines, où les flèches horizontales sont étales, et
tel qu'il existe un point $u \in U$ envoyé sur $x$. Alors $u$
est un point de type fini de $U$. Puisque $U \to V$ n'est pas lisse au
point $u$, d'après
Compléments sur les morphismes,
Lemme \ref{more-morphisms-lemma-lifting-along-artinian-at-point},
il existe un diagramme
$$
\xymatrix{
X \ar[d] & U \ar[l] \ar[d] & \Spec(B) \ar[d]^i \ar[l]^-\alpha  \\
Y & V \ar[l] & \Spec(B') \ar[l]_-{\beta} \ar@{-->}[lu]
}
$$
où $B' \to B$ est une petite extension d'anneaux locaux (artiniens),
telle que le corps résiduel de $B$ soit égal à $\kappa(v)$ et
que la flèche pointillée n'existe pas. Puisque $U \to V$ est
de type fini, nous voyons que $v$ est un point de type fini de $V$. D'après
Morphismes, Lemme \ref{morphisms-lemma-artinian-finite-type},
le morphisme $\beta$ est de type fini ; par conséquent, la composée
$\Spec(B) \to Y$ est elle aussi de type fini.
En raisonnant exactement comme dans la démonstration du
Lemme \ref{lemma-lifting-along-artinian-at-point}
(en utilisant que $U \to X$ et $V \to Y$ sont étales, donc
formellement étales),
nous voyons qu'il ne peut exister de flèche $\Spec(B) \to X$
s'insérant dans le rectangle extérieur du dernier diagramme affiché.
Autrement dit, (2) n'est pas vérifiée, et la démonstration est terminée.
\end{proof}

\noindent
Voici une application utile.

\begin{lemma}
\label{lemma-check-smoothness-on-infinitesimal-nbhds}
Soit $S$ un schéma. Soit $f : X \to Y$ un morphisme
d'espaces algébriques sur $S$. Supposons que $f$ soit
localement de type fini et $Y$ localement noethérien.
Soit $Z \subset Y$ un sous-espace fermé dont le voisinage infinitésimal d'ordre $n$ est
$Z_n \subset Y$. Posons $X_n = Z_n \times_Y X$.
\begin{enumerate}
\item Si $X_n \to Z_n$ est lisse pour tout $n$, alors $f$
est lisse en tout point de $f^{-1}(Z)$.
\item Si $X_n \to Z_n$ est étale pour tout $n$, alors $f$
est étale en tout point de $f^{-1}(Z)$.
\end{enumerate}
\end{lemma}

\begin{proof}
Supposons que $X_n \to Z_n$ soit lisse pour tout $n$.
Soit $x \in X$ un point situé au-dessus d'un point de $Z$.
Étant donnés une petite extension $B' \to B$ et des morphismes $\alpha$, $\beta$ comme dans le
Lemme \ref{lemma-lifting-along-artinian-at-point}, partie (3),
l'idéal maximal de $B'$ est nilpotent (puisque $B'$ est artinien),
et le morphisme $\beta$ se factorise donc par $Z_n$, tandis que $\alpha$ se factorise
par $X_n$ pour un $n$ convenable. Ainsi, la propriété de relèvement pour
$X_n \to Z_n$ s'applique et fournit la flèche pointillée souhaitée dans le diagramme.
Cela démontre (1). La partie (2) résulte de (1) et du fait qu'un morphisme
est étale si et seulement s'il est lisse de dimension relative $0$.
\end{proof}











\section{Le complexe cotangent naïf}
\label{section-netherlander}

\noindent
Cette section prolonge
Modules sur les sites, section \ref{sites-modules-section-netherlander},
qui prolonge à son tour la discussion menée dans
Algèbre, section \ref{algebra-section-netherlander}.

\begin{definition}
\label{definition-netherlander}
Soit $S$ un schéma.
Soit $f : X \to Y$ un morphisme d'espaces algébriques sur $S$.
Le {\it complexe cotangent naïf de $f$}
est le complexe défini dans Modules sur les sites, Définition
\ref{sites-modules-definition-cotangent-complex-morphism-ringed-topoi}
pour le morphisme de topos annelés $f_{small}$ entre les
petits sites étales de $X$ et $Y$, voir
Propriétés d'espaces, Lemme
\ref{spaces-properties-lemma-morphism-ringed-topoi}.
Notation : $\NL_f$ ou $\NL_{X/Y}$.
\end{definition}

\noindent
Les lemmes suivants montrent que cette définition est compatible avec celle
pour les homomorphismes d'anneaux et pour les schémas, et que $\NL_{X/Y}$ est un
objet de $D_\QCoh(\mathcal{O}_X)$.

\begin{lemma}
\label{lemma-NL-etale-localization}
Soit $S$ un schéma. Considérons un diagramme commutatif
$$
\xymatrix{
U \ar[d]_p \ar[r]_g & V \ar[d]^q \\
X \ar[r]^f & Y
}
$$
d'espaces algébriques sur $S$ où $p$ et $q$ sont étales.
Alors il existe une identification canonique
$\NL_{X/Y}|_{U_\etale} = \NL_{U/V}$ dans $D(\mathcal{O}_U)$.
\end{lemma}

\begin{proof}
La formation du complexe cotangent naïf commute à l'image inverse
(Modules sur les sites, Lemme \ref{sites-modules-lemma-pullback-NL})
et nous avons $p_{small}^{-1}\mathcal{O}_X = \mathcal{O}_U$ ainsi que
$g_{small}^{-1}\mathcal{O}_{V_\etale} =
p_{small}^{-1}f_{small}^{-1}\mathcal{O}_{Y_\etale}$
car $q_{small}^{-1}\mathcal{O}_{Y_\etale} =
\mathcal{O}_{V_\etale}$ d'après Propriétés d'espaces, Lemme
\ref{spaces-properties-lemma-etale-exact-pullback}.
En suivant les définitions, nous concluons que
$\NL_{X/Y}|_{U_\etale} = \NL_{U/V}$.
\end{proof}

\begin{lemma}
\label{lemma-NL-compare-spaces-schemes}
Soit $S$ un schéma. Soit $f : X \to Y$ un morphisme d'espaces algébriques
sur $S$. Supposons $X$ et $Y$ représentés par des schémas $X_0$ et $Y_0$.
Alors il existe une identification canonique
$\NL_{X/Y} = \epsilon^*\NL_{X_0/Y_0}$ dans $D(\mathcal{O}_X)$
où $\epsilon$ est celui de Catégories dérivées d'espaces, section
\ref{spaces-perfect-section-derived-quasi-coherent-etale}
et $\NL_{X_0/Y_0}$ est défini dans
Compléments sur les morphismes, Définition
\ref{more-morphisms-definition-netherlander}.
\end{lemma}

\begin{proof}
Soit $f_0 : X_0 \to Y_0$ le morphisme de schémas correspondant à $f$.
Il existe un morphisme canonique
$\epsilon^{-1}f_0^{-1}\mathcal{O}_{Y_0} \to f_{small}^{-1}\mathcal{O}_Y$
compatible avec
$\epsilon^\sharp : \epsilon^{-1}\mathcal{O}_{X_0} \to \mathcal{O}_X$
car il existe un diagramme commutatif
$$
\xymatrix{
X_{0, Zar} \ar[d]_{f_0} & X_\etale \ar[l]^\epsilon \ar[d]^f \\
Y_{0, Zar} & Y_\etale \ar[l]_\epsilon
}
$$
voir Catégories dérivées d'espaces, Remarque
\ref{spaces-perfect-remark-match-total-direct-images}.
Nous obtenons ainsi un morphisme canonique
$$
\epsilon^{-1}\NL_{X_0/Y_0} =
\epsilon^{-1}\NL_{\mathcal{O}_{X_0}/f_0^{-1}\mathcal{O}_{Y_0}} =
\NL_{\epsilon^{-1}\mathcal{O}_{X_0}/\epsilon^{-1}f_0^{-1}\mathcal{O}_{Y_0}}
\to
\NL_{\mathcal{O}_X/f^{-1}_{small}\mathcal{O}_Y} = \NL_{X/Y}
$$
par fonctorialité du complexe cotangent naïf.
Pour voir que le morphisme induit $\epsilon^*\NL_{X_0/Y_0} \to \NL_{X/Y}$ est un
isomorphisme dans $D(\mathcal{O}_X)$, il suffit de le vérifier sur les fibres aux points géométriques
(Propriétés d'espaces, Théorème
\ref{spaces-properties-theorem-exactness-stalks}).
Soit $\overline{x} : \Spec(k) \to X_0$
un point géométrique situé au-dessus de $x \in X_0$, et soit
$\overline{y} = f \circ \overline{x}$ situé au-dessus de $y \in Y_0$. Alors
$$
\NL_{X/Y, \overline{x}} =
\NL_{\mathcal{O}_{X, \overline{x}}/\mathcal{O}_{Y, \overline{y}}}
$$
Cela est vrai parce que le passage à la fibre en $\overline{x}$ revient
à prendre l'image inverse par $\overline{x} : \Spec(k) \to X$
et que nous pouvons appliquer Modules sur les sites, Lemme \ref{sites-modules-lemma-pullback-NL}.
D'autre part, nous avons
$$
(\epsilon^*\NL_{X_0/Y_0})_{\overline{x}} =
\NL_{X_0/Y_0, x} \otimes_{\mathcal{O}_{X_0, x}}
\mathcal{O}_{X, \overline{x}} =
\NL_{\mathcal{O}_{X_0, x}/\mathcal{O}_{Y_0, y}}
\otimes_{\mathcal{O}_{X_0, x}} \mathcal{O}_{X, \overline{x}}
$$
Certains détails sont omis (indication : utiliser le fait que la fibre d'une image inverse
est la fibre au point image, voir
Sites, Lemme \ref{sites-lemma-point-morphism-sites},
ainsi que le résultat correspondant pour les modules, voir
Modules sur les sites, Lemme \ref{sites-modules-lemma-pullback-stalk}).
Observons que $\mathcal{O}_{X, \overline{x}}$ est l'hensélisé
strict de $\mathcal{O}_{X_0, x}$, et de même
pour $\mathcal{O}_{Y, \overline{y}}$
(Propriétés d'espaces, Lemme
\ref{spaces-properties-lemma-describe-etale-local-ring}).
Le résultat découle donc de
Compléments sur l'algèbre,
Lemme \ref{more-algebra-lemma-henselization-NL}.
\end{proof}

\begin{lemma}
\label{lemma-netherlander-quasi-coherent}
Soit $S$ un schéma. Soit $f : X \to Y$ un morphisme
d'espaces algébriques sur $S$. Les faisceaux de cohomologie
du complexe $\NL_{X/Y}$ sont quasi-cohérents et nuls en dehors
des degrés $-1$, $0$ ; le faisceau $\Omega_{X/Y}$ est sa cohomologie en degré $0$.
\end{lemma}

\begin{proof}
Par construction du complexe cotangent naïf dans
Modules sur les sites, section \ref{sites-modules-section-netherlander},
nous savons que $\NL_{X/Y}$ est un complexe nul en dehors des degrés $-1$, $0$
et que sa cohomologie en degré $0$ est $\Omega_{X/Y}$ (d'après notre
construction de $\Omega_{X/Y}$ à la section \ref{section-sheaf-differentials}).
Le faisceau des différentielles est quasi-cohérent (d'après le
Lemme \ref{lemma-module-differentials-quasi-coherent}).
Pour achever la démonstration, il suffit de montrer que $H^{-1}(\NL_{X/Y})$
est quasi-cohérent. Cela résulte d'une vérification locale pour la topologie étale
(permise par le Lemme \ref{lemma-NL-etale-localization} et
Propriétés d'espaces, Lemme
\ref{spaces-properties-lemma-characterize-quasi-coherent})
qui nous ramène au cas des schémas
(Lemme \ref{lemma-NL-compare-spaces-schemes}),
puis de l'emploi du résultat dans le cas des schémas
(Compléments sur les morphismes, Lemme
\ref{more-morphisms-lemma-netherlander-quasi-coherent}).
\end{proof}

\begin{lemma}
\label{lemma-netherlander-fp}
Soit $S$ un schéma.
Soit $f : X \to Y$ un morphisme d'espaces algébriques sur $S$.
Si $f$ est localement de
présentation finie, alors $\NL_{X/Y}$ est, localement pour la topologie étale sur $X$,
quasi-isomorphe à un complexe
$$
\ldots \to 0 \to \mathcal{F}^{-1} \to \mathcal{F}^0 \to 0 \to \ldots
$$
de $\mathcal{O}_X$-modules quasi-cohérents
où $\mathcal{F}^0$ est de présentation finie
et $\mathcal{F}^{-1}$ de type fini.
\end{lemma}

\begin{proof}
La formation du complexe cotangent naïf commute à la
localisation étale d'après le Lemme \ref{lemma-NL-etale-localization}.
Nous nous ramenons ainsi au cas des schémas par le
Lemme \ref{lemma-NL-compare-spaces-schemes}.
Dans le cas des schémas, le résultat est
Compléments sur les morphismes, Lemme
\ref{more-morphisms-lemma-netherlander-fp}.
\end{proof}

\begin{lemma}
\label{lemma-NL-formally-smooth}
Soit $S$ un schéma.
Soit $f : X \to Y$ un morphisme d'espaces algébriques sur $S$.
Les assertions suivantes sont équivalentes :
\begin{enumerate}
\item $f$ est formellement lisse,
\item $H^{-1}(\NL_{X/Y}) = 0$ et le module $H^0(\NL_{X/Y}) = \Omega_{X/Y}$
est localement projectif.
\end{enumerate}
\end{lemma}

\begin{proof}
Cela résulte du
Lemme \ref{lemma-formally-smooth}, du
Lemme \ref{lemma-NL-etale-localization}, du
Lemme \ref{lemma-NL-compare-spaces-schemes}
et du cas des schémas, traité dans
Compléments sur les morphismes, Lemme \ref{more-morphisms-lemma-NL-formally-smooth}.
\end{proof}

\begin{lemma}
\label{lemma-NL-formally-etale}
Soit $f : X \to Y$ un morphisme de schémas. Les assertions suivantes sont équivalentes :
\begin{enumerate}
\item $f$ est formellement étale,
\item $H^{-1}(\NL_{X/Y}) = H^0(\NL_{X/Y}) = 0$.
\end{enumerate}
\end{lemma}

\begin{proof}
Supposons (1). Tout morphisme formellement étale est formellement lisse.
Ainsi $H^{-1}(\NL_{X/Y}) = 0$ d'après le Lemme \ref{lemma-NL-formally-smooth}.
D'autre part, un morphisme formellement étale est formellement non ramifié ;
par conséquent, $\Omega_{X/Y} = 0$ d'après le
Lemme \ref{lemma-characterize-formally-unramified}.
Réciproquement, si (2) est vérifiée, alors $f$ est formellement lisse d'après le
Lemme \ref{lemma-NL-formally-smooth}
et formellement non ramifié d'après le
Lemme \ref{lemma-characterize-formally-unramified},
et donc formellement étale d'après le
Lemme \ref{lemma-formally-etale-unramified-smooth}.
\end{proof}

\begin{lemma}
\label{lemma-NL-smooth}
Soit $f : X \to Y$ un morphisme de schémas. Les assertions suivantes sont équivalentes :
\begin{enumerate}
\item $f$ est lisse, et
\item $f$ est localement de présentation finie,
$H^{-1}(\NL_{X/Y}) = 0$, et le module $H^0(\NL_{X/Y}) = \Omega_{X/Y}$
est localement libre de type fini.
\end{enumerate}
\end{lemma}

\begin{proof}
Cela résulte du
Lemme \ref{lemma-formally-smooth}, du
Lemme \ref{lemma-NL-etale-localization}, du
Lemme \ref{lemma-NL-compare-spaces-schemes}
et du cas des schémas, traité dans
Compléments sur les morphismes, Lemme \ref{more-morphisms-lemma-NL-smooth}.
\end{proof}











\section{Ouverture du lieu de platitude}
\label{section-open-flat}

\noindent
Cette section est l'analogue de
Compléments sur les morphismes, section \ref{more-morphisms-section-open-flat}.
Notons que nous avons défini la notion de platitude pour les modules
quasi-cohérents sur les espaces algébriques dans
Morphismes d'espaces, section \ref{spaces-morphisms-section-flat-modules}.

\begin{theorem}
\label{theorem-openness-flatness}
Soit $S$ un schéma.
Soit $f : X \to Y$ un morphisme d'espaces algébriques sur $S$.
Soit $\mathcal{F}$ un faisceau quasi-cohérent sur $X$.
Supposons que $f$ soit localement de présentation finie et que
$\mathcal{F}$ soit un $\mathcal{O}_X$-module
localement de présentation finie. Alors
$$
\{x \in |X| : \mathcal{F}\text{ est plat sur }Y\text{ en }x\}
$$
est ouvert dans $|X|$.
\end{theorem}

\begin{proof}
Choisissons un diagramme commutatif
$$
\xymatrix{
U \ar[d]_p \ar[r]_\alpha &
V \ar[d]^q \\
X \ar[r]^a & Y
}
$$
où $U$, $V$ sont des schémas et $p$, $q$ sont étales et surjectifs comme dans
Espaces, Lemme \ref{spaces-lemma-lift-morphism-presentations}.
D'après
Compléments sur les morphismes, Théorème \ref{more-morphisms-theorem-openness-flatness},
l'ensemble
$U' = \{u \in |U| : p^*\mathcal{F}\text{ est plat sur }V\text{ en }u\}$
est ouvert dans $U$. D'après
Morphismes d'espaces, Définition \ref{spaces-morphisms-definition-flat-module},
l'image de $U'$ dans $|X|$ est l'ensemble
du théorème. La démonstration est donc terminée, car l'application $|U| \to |X|$ est
ouverte, voir
Propriétés d'espaces, Lemme \ref{spaces-properties-lemma-topology-points}.
\end{proof}

\begin{lemma}
\label{lemma-flat-locus-base-change}
Soit $S$ un schéma. Soit
$$
\xymatrix{
X' \ar[r]_{g'} \ar[d]_{f'} & X \ar[d]^f \\
Y' \ar[r]^g & Y
}
$$
un diagramme cartésien d'espaces algébriques sur $S$.
Soit $\mathcal{F}$ un $\mathcal{O}_X$-module quasi-cohérent.
Supposons que $g$ soit plat, que $f$ soit localement de présentation finie,
et que $\mathcal{F}$ soit localement de présentation finie.
Alors
$$
\{x' \in |X'| : (g')^*\mathcal{F}\text{ est plat sur }Y'\text{ en }x'\}
$$
est l'image inverse de la partie ouverte du
Théorème \ref{theorem-openness-flatness}
par l'application continue $|g'| : |X'| \to |X|$.
\end{lemma}

\begin{proof}
Cela découle de
Morphismes d'espaces, Lemme
\ref{spaces-morphisms-lemma-base-change-module-flat}.
\end{proof}












\section{Crit\`ere de platitude par fibres}
\label{section-criterion-flat-fibres}

\noindent
Soit $S$ un schéma. Considérons un diagramme commutatif d'espaces algébriques
sur $S$
$$
\xymatrix{
X \ar[rr]_f \ar[dr]_g & & Y \ar[dl]^h \\
& Z
}
$$
et un $\mathcal{O}_X$-module quasi-cohérent $\mathcal{F}$.
Étant donné un point $x \in |X|$, nous nous demandons si
$\mathcal{F}$ est plat sur $Y$ en $x$. Si $\mathcal{F}$ est plat
sur $Z$ en $x$, le théorème ci-dessous affirme que cette question est
étroitement liée à celle de savoir si la restriction de
$\mathcal{F}$ à la fibre de $X \to Z$ au-dessus de $g(x)$
est plate sur la fibre de $Y \to Z$ au-dessus de $g(x)$. Pour donner un sens
à cet énoncé, nous proposons le lemme préliminaire suivant.

\begin{lemma}
\label{lemma-flat-on-fibres-at-point}
Dans la situation ci-dessus, les assertions suivantes sont équivalentes :
\begin{enumerate}
\item Choisissons un point géométrique $\overline{x}$ de $X$ situé au-dessus de $x$.
Posons $\overline{y} = f \circ \overline{x}$ et
$\overline{z} = g \circ \overline{x}$. Alors le module
$\mathcal{F}_{\overline{x}}/
\mathfrak m_{\overline{z}}\mathcal{F}_{\overline{x}}$
est plat sur
$\mathcal{O}_{Y, \overline{y}}/
\mathfrak m_{\overline{z}}\mathcal{O}_{Y, \overline{y}}$.
\item Choisissons un morphisme $x : \Spec(K) \to X$ dans la classe d'équivalence de
$x$. Posons $z = g \circ x$, $X_z = \Spec(K) \times_{z, Z} X$,
$Y_z = \Spec(K) \times_{z, Z} Y$, et soit $\mathcal{F}_z$ l'image inverse
de $\mathcal{F}$ sur $X_z$. Alors $\mathcal{F}_z$ est plat en $x$ sur
$Y_z$ (au sens de Morphismes d'espaces,
Définition \ref{spaces-morphisms-definition-flat-module}).
\item Choisissons un diagramme commutatif
$$
\xymatrix{
& & & U \ar[llld]_a \ar[rr] \ar[dr] & & V \ar[llld]_>>>>>>>b \ar[dl] \\
X \ar[rr]_f \ar[dr]_g & & Y \ar[dl]^h &  & W \ar[llld]_c \\
& Z
}
$$
où $U, V, W$ sont des schémas et $a, b, c$ sont étales,
ainsi qu'un point $u \in U$ envoyé sur $x$. Soit $w \in W$ l'image de
$u$. Soit $\mathcal{F}_w$ l'image inverse de $\mathcal{F}$ sur
la fibre $U_w$ de $U \to W$ en $w$. Alors $\mathcal{F}_w$
est plat sur $V_w$ en $u$.
\end{enumerate}
\end{lemma}

\begin{proof}
Notons qu'en (2), le morphisme $x : \Spec(K) \to X$ définit un
point $K$-rationnel de $X_z$ ; l'énoncé a donc bien un sens. En outre,
la condition de (2) est indépendante du choix de $\Spec(K) \to X$
dans la classe d'équivalence de $x$ (les détails sont omis ; cela résultera aussi
des arguments ci-dessous, puisque les autres conditions ne dépendent pas
de ce choix). Notons également que nous pouvons toujours choisir un diagramme comme en
(3) de la manière suivante : choisissons d'abord
un schéma $W$ et un morphisme étale surjectif $W \to Z$, puis
un schéma $V$ et un morphisme étale surjectif $V \to W \times_Z Y$, et
enfin un schéma $U$ et un morphisme étale surjectif
$U \to V \times_Y X$. Ces choix étant faits, nous prenons pour $U \to W$
la composée $U \to V \to W$ et pouvons choisir un point $u \in U$ envoyé
sur $x$, puisque le morphisme $U \to X$ est surjectif.

\medskip\noindent
Supposons donnés à la fois un diagramme comme en (3) et un point géométrique
$\overline{x} : \Spec(k) \to X$ comme en (1). D'après
Propriétés d'espaces, Lemme
\ref{spaces-properties-lemma-geometric-lift-to-usual}
nous pouvons choisir un point géométrique $\overline{u} : \Spec(k) \to U$
au-dessus de $u$ tel que $\overline{x} = a \circ \overline{u}$.
Notons $\overline{v} : \Spec(k) \to V$ et
$\overline{w} : \Spec(k) \to W$ les points géométriques induits de
$V$ et $W$. Dans cette situation, nous savons que
$\mathcal{O}_{X, \overline{x}} = \mathcal{O}_{U, u}^{sh}$
et de même pour $Y$ et $Z$, voir
Propriétés d'espaces,
Lemme \ref{spaces-properties-lemma-describe-etale-local-ring}.
De même, nous avons
$$
\mathcal{F}_{\overline{x}} =
(a^*\mathcal{F})_u \otimes_{\mathcal{O}_{U, u}}
\mathcal{O}_{U, u}^{sh}
$$
voir
Propriétés d'espaces, Lemme \ref{spaces-properties-lemma-stalk-quasi-coherent}.
Notons que la fibre de $\mathcal{F}_w$ en $u$ est donnée par
$$
(\mathcal{F}_w)_u = (a^*\mathcal{F})_u/\mathfrak m_w(a^*\mathcal{F})_u
$$
et que l'anneau local de $V_w$ en $v$ est donné par
$$
\mathcal{O}_{V_w, v} = \mathcal{O}_{V, v}/\mathfrak m_w\mathcal{O}_{V, v}.
$$
Puisque $\mathfrak m_{\overline{z}} =
\mathfrak m_w \mathcal{O}_{Z, \overline{z}} =
\mathfrak m_w \mathcal{O}_{W, w}^{sh}$
nous voyons que
\begin{align*}
\mathcal{F}_{\overline{x}}/
\mathfrak m_{\overline{z}}\mathcal{F}_{\overline{x}} & =
(a^*\mathcal{F})_u \otimes_{\mathcal{O}_{U, u}}
\mathcal{O}_{X, \overline{x}}/
\mathfrak m_{\overline{z}}\mathcal{O}_{X, \overline{x}} \\
& =
(\mathcal{F}_w)_u \otimes_{\mathcal{O}_{U_w, u}}
\mathcal{O}_{U, u}^{sh}/\mathfrak m_w\mathcal{O}_{U, u}^{sh} \\
& = (\mathcal{F}_w)_u \otimes_{\mathcal{O}_{U_w, u}}
\mathcal{O}_{U_w, \overline{u}}^{sh} \\
& = (\mathcal{F}_w)_{\overline{u}}
\end{align*}
l'avant-dernière égalité résultant du
Lemme \ref{algebra-lemma-quotient-strict-henselization} d'Algèbre
et la dernière du
Lemme \ref{spaces-properties-lemma-stalk-quasi-coherent} de Propriétés d'espaces.
Les mêmes arguments appliqués aux faisceaux structuraux de $V$ et $Y$
montrent que
$$
\mathcal{O}_{V_w, \overline{v}}^{sh} =
\mathcal{O}_{V, v}^{sh}/\mathfrak m_w \mathcal{O}_{V, v}^{sh} =
\mathcal{O}_{Y, \overline{y}}/
\mathfrak m_{\overline{z}}\mathcal{O}_{Y, \overline{y}}.
$$
Nous pouvons maintenant appliquer
Morphismes d'espaces, Lemme \ref{spaces-morphisms-lemma-flat-at-point},
pour voir que (1) est équivalente à (3).

\medskip\noindent
Enfin, démontrons l'équivalence de (2) et (3).
Pour cela, choisissons une extension de corps $\tilde K$ de $K$
et un morphisme $\tilde x : \Spec(\tilde K) \to U$ qui
est situé au-dessus de $u$ (c'est possible car $u \times_{X, x} \Spec(K)$
est un schéma non vide). Soit $\tilde z : \Spec(\tilde K) \to U \to W$
la composée. Nous obtenons un diagramme commutatif
$$
\xymatrix{
& & & U_w \times_w \tilde z \ar[llld]_a \ar[rr] \ar[dr] & &
V_w \times_w \tilde z \ar[llld]_>>>>>>>b \ar[dl] \\
X_z \ar[rr]_f \ar[dr]_g & & Y_z \ar[dl]^h &  & \tilde z \ar[llld]_c \\
& z
}
$$
où $z = \Spec(K)$ et $w = \Spec(\kappa(w))$. Il est alors
clair que $\mathcal{F}_w$ et $\mathcal{F}_z$ ont pour images inverses le
même module sur $U_w \times_w \tilde z$. Nous obtenons ainsi un diagramme
commutatif
$$
\xymatrix{
X_z \ar[d] & U_w \times_w \tilde z \ar[l] \ar[d] \ar[r] & U_w \ar[d] \\
Y_z & V_w \times_w \tilde z \ar[l] \ar[r] & V_w
}
$$
dont les deux carrés sont cartésiens et dont les flèches horizontales du bas
sont plates : la flèche horizontale inférieure gauche est la composée
du morphisme $Y \times_Z \tilde z \to Y \times_Z z = Y_z$ (changement de base
d'un morphisme plat), du morphisme étale
$V \times_Z \tilde z \to Y \times_Z \tilde z$, et
du morphisme étale $V \times_W \tilde z \to V \times_Z \tilde z$.
Il résulte donc de
Morphismes d'espaces,
Lemme \ref{spaces-morphisms-lemma-base-change-module-flat},
que
$$
\mathcal{F}_z\text{ plat en }x\text{ sur }Y_z
\Leftrightarrow
\mathcal{F}|_{U_w \times_w \tilde z}
\text{ plat en }\tilde x\text{ sur }V_w \times_w \tilde z
\Leftrightarrow
\mathcal{F}_w\text{ plat en }u\text{ sur }V_w
$$
ce qui conclut.
\end{proof}

\begin{definition}
\label{definition-module-flat-on-fibre}
Soit $S$ un schéma. Soient $X \to Y \to Z$ des morphismes d'espaces
algébriques sur $S$. Soit $\mathcal{F}$ un $\mathcal{O}_X$-module quasi-cohérent.
Soit $x \in |X|$ un point et notons $z \in |Z|$ son image.
\begin{enumerate}
\item Nous disons que {\it la restriction de $\mathcal{F}$ à sa fibre au-dessus de $z$
est plate en $x$ sur la fibre de $Y$ au-dessus de $z$} si les conditions équivalentes du
Lemme \ref{lemma-flat-on-fibres-at-point}
sont vérifiées.
\item Nous disons que {\it la fibre de $X$ au-dessus de $z$ est plate en $x$ sur la fibre de
$Y$ au-dessus de $z$} si les conditions équivalentes du
Lemme \ref{lemma-flat-on-fibres-at-point}
sont vérifiées pour $\mathcal{F} = \mathcal{O}_X$.
\item Nous disons que {\it la fibre de $X$ au-dessus de $z$ est plate sur la fibre de $Y$
au-dessus de $z$} si, pour tout $x \in |X|$ situé au-dessus de $z$, la fibre de $X$ au-dessus de $z$
est plate en $x$ sur la fibre de $Y$ au-dessus de $z$.
\end{enumerate}
\end{definition}

\noindent
Cette définition nous permet d'énoncer la version suivante du critère.
La version noethérienne se trouve à la
section \ref{section-flat-over-Noetherian}.

\begin{theorem}
\label{theorem-criterion-flatness-fibre}
Soit $S$ un schéma.
Soient $f : X \to Y$ et $Y \to Z$ des morphismes d'espaces algébriques sur $S$.
Soit $\mathcal{F}$ un $\mathcal{O}_X$-module quasi-cohérent.
Supposons que
\begin{enumerate}
\item $X$ soit localement de présentation finie sur $Z$,
\item $\mathcal{F}$ soit un $\mathcal{O}_X$-module de présentation finie, et
\item $Y$ soit localement de type fini sur $Z$.
\end{enumerate}
Soit $x \in |X|$ et soient $y \in |Y|$ et $z \in |Z|$ les images de
$x$. Si $\mathcal{F}_{\overline{x}} \not = 0$, les assertions suivantes sont
équivalentes :
\begin{enumerate}
\item $\mathcal{F}$ est plat sur $Z$ en $x$ et
la restriction de $\mathcal{F}$ à sa fibre au-dessus de $z$
est plate en $x$ sur la fibre de $Y$ au-dessus de $z$, et
\item $Y$ est plat sur $Z$ en $y$ et $\mathcal{F}$ est
plat sur $Y$ en $x$.
\end{enumerate}
En outre, l'ensemble des points $x$ où (1) et (2) sont vérifiées est ouvert dans
$\text{Supp}(\mathcal{F})$.
\end{theorem}

\begin{proof}
Choisissons un diagramme comme dans le
Lemme \ref{lemma-flat-on-fibres-at-point}, partie (3).
Il résulte des définitions que l'on se ramène au
théorème correspondant pour les morphismes de schémas
$U \to V \to W$, le faisceau quasi-cohérent $a^*\mathcal{F}$
et le point $u \in U$. Le théorème découle donc du
résultat correspondant pour les schémas, à savoir
Compléments sur les morphismes,
Théorème \ref{more-morphisms-theorem-criterion-flatness-fibre}.
\end{proof}

\begin{lemma}
\label{lemma-morphism-between-flat}
Soit $S$ un schéma.
Soient $f : X \to Y$ et $Y \to Z$ des morphismes d'espaces algébriques sur $S$.
Supposons que
\begin{enumerate}
\item $X$ soit localement de présentation finie sur $Z$,
\item $X$ soit plat sur $Z$,
\item pour tout $z \in |Z|$, la fibre de $X$ au-dessus de $z$
soit plate sur la fibre de $Y$ au-dessus de $z$, et
\item $Y$ soit localement de type fini sur $Z$.
\end{enumerate}
Alors $f$ est plat. Si $f$ est en outre surjectif, alors $Y$ est plat sur $Z$.
\end{lemma}

\begin{proof}
C'est un cas particulier du
Théorème \ref{theorem-criterion-flatness-fibre}.
\end{proof}

\begin{lemma}
\label{lemma-base-change-criterion-flatness-fibre}
Soit $S$ un schéma. Soient $f : X \to Y$ et $Y \to Z$ des morphismes
d'espaces algébriques sur $S$. Soit $\mathcal{F}$ un
$\mathcal{O}_X$-module quasi-cohérent.
Supposons que
\begin{enumerate}
\item $X$ soit localement de présentation finie sur $Z$,
\item $\mathcal{F}$ soit un $\mathcal{O}_X$-module de présentation finie,
\item $\mathcal{F}$ soit plat sur $Z$, et
\item $Y$ soit localement de type fini sur $Z$.
\end{enumerate}
Alors l'ensemble
$$
A = \{x \in |X| : \mathcal{F} \text{ plat en }x \text{ sur }Y\}.
$$
Cet ensemble est ouvert dans $|X|$ et sa formation commute à tout changement de base :
si $Z' \to Z$ est un morphisme d'espaces algébriques et si $A'$ est l'ensemble des
points de $X' = X \times_Z Z'$ où $\mathcal{F}' = \mathcal{F} \times_Z Z'$
est plat sur $Y' = Y \times_Z Z'$, alors $A'$ est l'image inverse de
$A$ par l'application continue $|X'| \to |X|$.
\end{lemma}

\begin{proof}
Une façon de le démontrer consiste à transposer la démonstration donnée dans
Compléments sur les morphismes, Lemme \ref{more-morphisms-lemma-morphism-between-flat},
à la catégorie des espaces algébriques. Nous le démontrerons plutôt
en nous ramenant au cas des schémas. Choisissons donc un diagramme comme dans le
Lemme \ref{lemma-flat-on-fibres-at-point}, partie (3),
tel que $a$, $b$ et $c$ soient surjectifs.
Il résulte des définitions que l'on se ramène au
théorème correspondant pour les morphismes de schémas
$U \to V \to W$, le faisceau quasi-cohérent $a^*\mathcal{F}$
et le point $u \in U$. Le seul point à préciser est que,
étant donné un morphisme d'espaces algébriques $Z' \to Z$, nous choisissons un schéma
$W'$ et un morphisme étale surjectif $W' \to W \times_Z Z'$.
Posons alors $U' = W' \times_W U$ et $V' = W' \times_W V$.
Notons $a', b', c'$ les morphismes de $U', V', W'$ dans
$X', Y', Z'$. Dans ce cas, $A$, resp.\ $A'$ sont les images des parties
ouvertes de $U$, resp.\ $U'$ associées à
$a^*\mathcal{F}$, resp.\ $(a')^*\mathcal{F}'$.
Cela ramène bien le lemme à
Compléments sur les morphismes, Lemme \ref{more-morphisms-lemma-morphism-between-flat}.
\end{proof}

\begin{lemma}
\label{lemma-base-change-flatness-fibres}
Soit $S$ un schéma.
Soient $f : X \to Y$ et $Y \to Z$ des morphismes d'espaces algébriques sur $S$.
Supposons que
\begin{enumerate}
\item $X$ soit localement de présentation finie sur $Z$,
\item $X$ soit plat sur $Z$, et
\item $Y$ soit localement de type fini sur $Z$.
\end{enumerate}
Alors l'ensemble
$$
\{x \in |X| : X\text{ plat en }x \text{ sur }Y\}.
$$
Cet ensemble est ouvert dans $|X|$ et sa formation commute à tout changement de base
$Z' \to Z$.
\end{lemma}

\begin{proof}
C'est un cas particulier du
Lemme \ref{lemma-base-change-criterion-flatness-fibre}.
\end{proof}

\begin{lemma}
\label{lemma-flat-and-free-at-point-fibre}
Soit $S$ un schéma. Soit $f : X \to Y$ un morphisme d'espaces algébriques
sur $S$ qui est localement de présentation finie.
Soit $\mathcal{F}$ un $\mathcal{O}_X$-module de présentation finie.
Soit $x \in |X|$, d'image $y \in |Y|$. Si $\mathcal{F}$ est plat en $x$
sur $Y$, les assertions suivantes sont équivalentes :
\begin{enumerate}
\item $(\mathcal{F}_{\overline{y}})_{\overline{x}}$ est un module plat sur
$\mathcal{O}_{X_{\overline{y}}, \overline{x}}$,
\item $(\mathcal{F}_{\overline{y}})_{\overline{x}}$ est un module libre sur
$\mathcal{O}_{X_{\overline{y}}, \overline{x}}$,
\item $\mathcal{F}_{\overline{y}}$ est libre de type fini sur un
voisinage étale de $\overline{x}$ dans $X_{\overline{y}}$, et
\item $\mathcal{F}$ est libre de type fini sur un voisinage étale de $x$ dans $X$.
\end{enumerate}
Ici, $\overline{x}$ est un point géométrique de $X$ situé au-dessus de $x$
et $\overline{y} = f \circ \overline{x}$.
\end{lemma}

\begin{proof}
Choisissons un diagramme commutatif
$$
\xymatrix{
U \ar[d] \ar[r] & V \ar[d] \\
X \ar[r] & Y
}
$$
où $U$ et $V$ sont des schémas et les flèches verticales sont étales,
et tel qu'il existe un point $u \in U$ envoyé sur $x$. Soit $v \in V$
l'image de $u$. En appliquant le Lemme \ref{lemma-flat-on-fibres-at-point}
à $\text{id} : X \to X$ sur $Y$, nous voyons que (1) devient
la condition « $\mathcal{F}|_{U_v}$ est plat sur $U_v$ en $u$ ».
Autrement dit, (1) équivaut à ce que $(\mathcal{F}|_{U_v})_u$
soit un $\mathcal{O}_{U_v, u}$-module plat.
D'après le cas des schémas (Compléments sur les morphismes, Lemme
\ref{more-morphisms-lemma-flat-and-free-at-point-fibre}),
nous voyons que cela implique que
$\mathcal{F}|_U$ est libre de type fini sur un voisinage ouvert
de $u$. Nous voyons ainsi que (1) implique (4).
Les implications (4) $\Rightarrow$ (3) et
(2) $\Rightarrow$ (1) sont immédiates.
Pour l'implication (3) $\Rightarrow$ (2), utilisons
la description des anneaux locaux et des fibres dans
Propriétés d'espaces, Lemmes
\ref{spaces-properties-lemma-describe-etale-local-ring} et
\ref{spaces-properties-lemma-stalk-quasi-coherent}.
\end{proof}

\begin{lemma}
\label{lemma-finite-free-open}
Soit $S$ un schéma. Soit $f : X \to Y$ un morphisme d'espaces algébriques
sur $S$ qui est localement de présentation finie.
Soit $\mathcal{F}$ un $\mathcal{O}_X$-module de présentation finie
plat sur $Y$. Alors l'ensemble
$$
\{x \in |X| : \mathcal{F}\text{ libre dans un voisinage étale de }x\}
$$
est ouvert dans $|X|$ et sa formation commute à tout changement de base
$Y' \to Y$.
\end{lemma}

\begin{proof}
L'ouverture est immédiate. Soit $Y' \to Y$ un morphisme d'espaces algébriques,
posons $X' = Y' \times_Y X$,
et soit $x' \in |X'|$ un point situé au-dessus de $x \in |X|$.
D'après le Lemme \ref{lemma-flat-and-free-at-point-fibre},
nous voyons que $x$ appartient à notre ensemble si et seulement si
$(\mathcal{F}_{\overline{y}})_{\overline{x}}$ est un module plat sur
$\mathcal{O}_{X_{\overline{y}}, \overline{x}}$.
De même, $x'$ appartient à l'analogue de notre ensemble pour l'image inverse
$\mathcal{F}'$ de $\mathcal{F}$ sur $X'$ si et seulement si
$(\mathcal{F}'_{\overline{y}'})_{\overline{x}'}$ est un module plat sur
$\mathcal{O}_{X'_{\overline{y}'}, \overline{x}'}$
(avec les notations évidentes). Ces deux assertions sont équivalentes
d'après le Lemme \ref{lemma-flat-on-fibres-at-point} appliqué au
morphisme $\text{id} : X \to X$ sur $Y$.
L'assertion sur le changement de base est donc démontrée.
\end{proof}







\section{Platitude sur une base noethérienne}
\label{section-flat-over-Noetherian}

\noindent
Voici le « critère de platitude par fibres » dans le cas noethérien.

\begin{theorem}
\label{theorem-criterion-flatness-fibre-Noetherian}
Soit $S$ un schéma.
Soient $f : X \to Y$ et $Y \to Z$ des morphismes d'espaces algébriques sur $S$.
Soit $\mathcal{F}$ un $\mathcal{O}_X$-module quasi-cohérent.
Supposons que
\begin{enumerate}
\item $X$, $Y$, $Z$ soient localement noethériens, et
\item $\mathcal{F}$ soit un $\mathcal{O}_X$-module cohérent.
\end{enumerate}
Soit $x \in |X|$ et soient $y \in |Y|$ et $z \in |Z|$ les images de
$x$. Si $\mathcal{F}_{\overline{x}} \not = 0$, les assertions suivantes sont
équivalentes :
\begin{enumerate}
\item $\mathcal{F}$ est plat sur $Z$ en $x$ et
la restriction de $\mathcal{F}$ à sa fibre au-dessus de $z$
est plate en $x$ sur la fibre de $Y$ au-dessus de $z$, et
\item $Y$ est plat sur $Z$ en $y$ et $\mathcal{F}$ est
plat sur $Y$ en $x$.
\end{enumerate}
\end{theorem}

\begin{proof}
Choisissons un diagramme comme dans le
Lemme \ref{lemma-flat-on-fibres-at-point}, partie (3).
Il résulte des définitions que l'on se ramène au
théorème correspondant pour les morphismes de schémas
$U \to V \to W$, le faisceau quasi-cohérent $a^*\mathcal{F}$
et le point $u \in U$. Le théorème découle donc du
résultat correspondant pour les schémas, à savoir
Compléments sur les morphismes,
Théorème \ref{more-morphisms-theorem-criterion-flatness-fibre-Noetherian}.
\end{proof}

\begin{lemma}
\label{lemma-morphism-between-flat-Noetherian}
Soit $S$ un schéma.
Soient $f : X \to Y$ et $Y \to Z$ des morphismes d'espaces algébriques sur $S$.
Supposons que
\begin{enumerate}
\item $X$, $Y$, $Z$ soient localement noethériens,
\item $X$ soit plat sur $Z$,
\item pour tout $z \in |Z|$, la fibre de $X$ au-dessus de $z$
soit plate sur la fibre de $Y$ au-dessus de $z$.
\end{enumerate}
Alors $f$ est plat. Si $f$ est en outre surjectif, alors $Y$ est plat sur $Z$.
\end{lemma}

\begin{proof}
C'est un cas particulier du
Théorème \ref{theorem-criterion-flatness-fibre-Noetherian}.
\end{proof}

\noindent
Tout comme pour vérifier la lissité, lorsque la base est noethérienne, il suffit
de vérifier la platitude sur les anneaux artiniens. Voici un énoncé type.

\begin{lemma}
\label{lemma-flatness-over-Noetherian-ring}
Soit $A$ un anneau noethérien. Soit $I \subset A$ un idéal.
Soit $X$ un espace algébrique localement de présentation finie sur
$S = \Spec(A)$. Pour $n \geq 1$, posons $S_n = \Spec(A/I^n)$ et
$X_n = S_n \times_S X$. Soit $\mathcal{F}$ un $\mathcal{O}_X$-module cohérent.
Si, pour tout $n \geq 1$, l'image inverse $\mathcal{F}_n$ de $\mathcal{F}$ sur $X$
est plate sur $S_n$, alors le lieu (ouvert) où $\mathcal{F}$
est plate sur $X$ contient l'image inverse de $V(I)$ par $X \to S$.
\end{lemma}

\begin{proof}
Le lieu où $\mathcal{F}$ est plat sur $S$ est ouvert dans $|X|$ d'après le
Théorème \ref{theorem-openness-flatness}.
L'énoncé est invariant lorsqu'on remplace $X$ par les membres d'un
recouvrement étale ; nous pouvons donc supposer que $X$ est un schéma affine.
Dans ce cas, le résultat découle immédiatement du
Lemme \ref{algebra-lemma-flat-module-powers} d'Algèbre.
Certains détails sont omis.
\end{proof}







\section{Retour sur la normalisation}
\label{section-normalization}

\noindent
La normalisation commute au changement de base lisse.

\begin{lemma}
\label{lemma-integral-closure-smooth-pullback}
Soit $S$ un schéma. Soit $f : Y \to X$ un morphisme lisse
d'espaces algébriques sur $S$. Soit $\mathcal{A}$ un faisceau quasi-cohérent
de $\mathcal{O}_X$-algèbres. La clôture intégrale
de $\mathcal{O}_Y$ dans $f^*\mathcal{A}$ est égale à $f^*\mathcal{A}'$,
où $\mathcal{A}' \subset \mathcal{A}$ est la clôture intégrale de
$\mathcal{O}_X$ dans $\mathcal{A}$.
\end{lemma}

\begin{proof}
Par notre construction de la clôture intégrale, voir
Morphismes d'espaces, Définition
\ref{spaces-morphisms-definition-integral-closure},
on se ramène immédiatement au cas où $X$ et $Y$ sont affines.
Dans ce cas, le résultat est le
Lemme \ref{algebra-lemma-integral-closure-commutes-smooth} d'Algèbre.
\end{proof}

\begin{lemma}[La normalisation commute au changement de base lisse]
\label{lemma-normalization-smooth-localization}
Soit $S$ un schéma. Soit
$$
\xymatrix{
Y_2 \ar[r] \ar[d] & Y_1 \ar[d]^f \\
X_2 \ar[r]^\varphi & X_1
}
$$
un carré cartésien d'espaces algébriques sur $S$. Supposons que $f$ soit quasi-compact
et quasi-séparé, et que $\varphi$ soit lisse.
Soit $Y_i \to X_i' \to X_i$ la normalisation de $X_i$ dans $Y_i$.
Alors $X_2' \cong X_2 \times_{X_1} X_1'$.
\end{lemma}

\begin{proof}
La factorisation $Y_1 \to X_1' \to X_1$, après changement de base à $X_2$,
devient une factorisation $Y_2 \to X_2 \times_{X_1} X_1' \to X_1$ et le morphisme
$X_2 \times_{X_1} X_1' \to X_1$ est entier
(Morphismes d'espaces, Lemme \ref{spaces-morphisms-lemma-base-change-integral}).
Nous obtenons donc un morphisme
$h : X_2' \to X_2 \times_{X_1} X_1'$ par la propriété universelle de
Morphismes d'espaces, Lemme
\ref{spaces-morphisms-lemma-characterize-normalization}.
Notons que $X_2'$ est le spectre relatif de la clôture intégrale
de $\mathcal{O}_{X_2}$ dans $f_{2, *}\mathcal{O}_{Y_2}$.
Si $\mathcal{A}' \subset f_{1, *}\mathcal{O}_{Y_1}$ désigne la clôture
intégrale de $\mathcal{O}_{X_2}$, alors $X_2 \times_{X_1} X_1'$ est le
spectre relatif de $\varphi^*\mathcal{A}'$, puisque la formation du
spectre relatif commute à tout changement de base. D'après
Cohomologie des espaces, Lemme
\ref{spaces-cohomology-lemma-flat-base-change-cohomology}
nous savons que $f_{2, *}\mathcal{O}_{Y_2} = \varphi^*f_{1, *}\mathcal{O}_{Y_1}$.
Le résultat découle donc du
Lemme \ref{lemma-integral-closure-smooth-pullback}.
\end{proof}









\section{Morphismes de Cohen-Macaulay}
\label{section-CM}

\noindent
Ceci est l'analogue de Compléments sur les morphismes, section
\ref{more-morphisms-section-CM}.

\begin{lemma}
\label{lemma-CM-local-ring-fibre}
La propriété suivante des morphismes de germes de schémas
\begin{align*}
& \mathcal{P}((X, x) \to (S, s)) = \\
& \text{l'anneau local }
\mathcal{O}_{X_s, x}
\text{ de la fibre est noethérien et Cohen-Macaulay}
\end{align*}
est de nature locale pour la topologie étale sur la source et le but (Descente, Définition
\ref{descent-definition-local-source-target-at-point}).
\end{lemma}

\begin{proof}
Étant donné un diagramme comme dans
Descente, Définition \ref{descent-definition-local-source-target-at-point},
nous obtenons un morphisme étale entre les fibres
$U'_{v'} \to U_v$ qui envoie $u'$ sur $u$ ; voir
Descente, Lemme \ref{descent-lemma-etale-on-fiber}.
Les hensélisés stricts des anneaux locaux
$\mathcal{O}_{U'_{v'}, u'}$ et $\mathcal{O}_{U_v, u}$
sont donc les mêmes. Nous concluons par le
Lemme \ref{more-algebra-lemma-henselization-CM} de Compléments sur l'algèbre.
\end{proof}

\begin{definition}
\label{definition-CM}
Soit $S$ un schéma.
Soit $f : X \to Y$ un morphisme d'espaces algébriques sur $S$.
Supposons que les fibres de $f$ soient localement noethériennes
(Diviseurs sur les espaces, Définition
\ref{spaces-divisors-definition-locally-Noetherian-fibre}).
\begin{enumerate}
\item Soient $x \in |X|$ et $y = f(x)$. Nous disons que $f$ est
{\it Cohen-Macaulay en $x$} si $f$ est plat en $x$ et si
les conditions équivalentes de
Morphismes d'espaces, Lemme
\ref{spaces-morphisms-lemma-local-source-target-at-point}
sont satisfaites pour la propriété $\mathcal{P}$ décrite dans le
Lemme \ref{lemma-CM-local-ring-fibre}.
\item Nous disons que $f$ est un {\it morphisme de Cohen-Macaulay} si $f$ est
Cohen-Macaulay en tout point de $X$.
\end{enumerate}
\end{definition}

\noindent
Voici une reformulation.

\begin{lemma}
\label{lemma-CM}
Soit $S$ un schéma. Soit $f : X \to Y$ un morphisme d'espaces algébriques
sur $S$. Supposons que les fibres de $f$ soient localement noethériennes.
Les conditions suivantes sont équivalentes :
\begin{enumerate}
\item $f$ est Cohen-Macaulay,
\item $f$ est plat et il existe un morphisme étale surjectif $V \to Y$
où $V$ est un schéma, tel que les fibres de $X_V \to V$
soient des espaces algébriques de Cohen-Macaulay, et
\item $f$ est plat et, pour tout morphisme étale $V \to Y$
où $V$ est un schéma, les fibres de $X_V \to V$
sont des espaces algébriques de Cohen-Macaulay.
\end{enumerate}
Étant donné $x \in |X|$ d'image $y \in |Y|$, les conditions suivantes sont
équivalentes :
\begin{enumerate}
\item[(a)] $f$ est Cohen-Macaulay en $x$, et
\item[(b)] $\mathcal{O}_{Y, \overline{y}} \to \mathcal{O}_{X, \overline{x}}$
est plat et
$\mathcal{O}_{X, \overline{x}}/
\mathfrak m_{\overline{y}}\mathcal{O}_{X, \overline{x}}$ est Cohen-Macaulay.
\end{enumerate}
\end{lemma}

\begin{proof}
Étant donné un morphisme étale $V \to Y$ où $V$ est un schéma,
choisissons un schéma $U$ et un morphisme étale surjectif $U \to X \times_Y V$.
Considérons le diagramme commutatif
$$
\xymatrix{
U \ar[d] \ar[r] & V \ar[d] \\
X \ar[r] & Y
}
$$
Soit $u \in U$ d'images $x \in |X|$, $y \in |Y|$ et $v \in V$.
Alors $f$ est Cohen-Macaulay en $x$ si et seulement si $U \to V$ l'est
en $u$ (par définition). De plus, le morphisme
$U_v \to X_v = (X_V)_v$ est étale surjectif. Par conséquent, le schéma $U_v$ est
Cohen-Macaulay si et seulement si l'espace algébrique $X_v$ l'est.
L'équivalence de (1), (2) et (3) découle donc de
l'équivalence correspondante pour les morphismes de
schémas ; voir Compléments sur les morphismes, Lemme \ref{more-morphisms-lemma-CM} ;
l'argument est formel.

\medskip\noindent
Démonstration de l'équivalence de (a) et (b). L'équivalence correspondante
pour la platitude est donnée par Morphismes d'espaces, Lemme
\ref{spaces-morphisms-lemma-flat-at-point-etale-local-rings}.
Nous pouvons donc supposer $f$ plat en $x$ pour démontrer l'équivalence.
Considérons un diagramme et $x, y, u, v$ comme ci-dessus. Alors
$\mathcal{O}_{Y, \overline{y}} \to \mathcal{O}_{X, \overline{x}}$
est égal à l'homomorphisme
$\mathcal{O}_{V, v}^{sh} \to \mathcal{O}_{U, u}^{sh}$
entre les hensélisés stricts des anneaux locaux ; voir
Propriétés des espaces, Lemme
\ref{spaces-properties-lemma-describe-etale-local-ring}.
Nous avons donc
$$
\mathcal{O}_{X, \overline{x}}/
\mathfrak m_{\overline{y}}\mathcal{O}_{X, \overline{x}} =
(\mathcal{O}_{U, u}/\mathfrak m_v \mathcal{O}_{U, u})^{sh}
$$
d'après Algèbre, Lemme \ref{algebra-lemma-quotient-strict-henselization}.
Il reste donc à montrer que l'anneau local noethérien
$\mathcal{O}_{U, u}/\mathfrak m_v \mathcal{O}_{U, u}$
est Cohen-Macaulay si et seulement si son hensélisé strict l'est.
C'est Compléments sur l'algèbre, Lemme \ref{more-algebra-lemma-henselization-CM}.
\end{proof}

\begin{lemma}
\label{lemma-composition-CM}
Soit $S$ un schéma.
Soient $f : X \to Y$ et $g : Y \to Z$ des morphismes d'espaces algébriques
sur $S$. Supposons que les
fibres de $f$, $g$ et $g \circ f$ soient localement noethériennes.
Soit $x \in |X|$ d'images $y \in |Y|$ et $z \in |Z|$.
\begin{enumerate}
\item Si $f$ est Cohen-Macaulay en $x$ et $g$ est Cohen-Macaulay
en $f(x)$, alors $g \circ f$ est Cohen-Macaulay en $x$.
\item Si $f$ et $g$ sont Cohen-Macaulay, alors $g \circ f$ l'est.
\item Si $g \circ f$ est Cohen-Macaulay en $x$ et $f$ est plat en $x$,
alors $f$ est Cohen-Macaulay en $x$ et $g$ l'est en $f(x)$.
\item Si $f \circ g$ est Cohen-Macaulay et $f$ est plat, alors
$f$ est Cohen-Macaulay et $g$ l'est en tout point de
l'image de $f$.
\end{enumerate}
\end{lemma}

\begin{proof}
Le résultat correspondant pour les schémas, appliqué localement pour la
topologie étale, donne l'assertion ; voir
Compléments sur les morphismes, Lemme \ref{more-morphisms-lemma-composition-CM}.
On peut aussi utiliser l'équivalence de (a) et (b) du
Lemme \ref{lemma-CM}. Nous considérons alors l'homomorphisme local
d'anneaux locaux noethériens
$$
\mathcal{O}_{Y, \overline{y}}/
\mathfrak m_{\overline{z}}\mathcal{O}_{Y, \overline{y}}
\longrightarrow
\mathcal{O}_{X, \overline{x}}/
\mathfrak m_{\overline{z}}\mathcal{O}_{X, \overline{x}}
$$
dont la fibre est
$$
\mathcal{O}_{X, \overline{x}}/
\mathfrak m_{\overline{y}}\mathcal{O}_{X, \overline{x}}
$$
et nous appliquons Algèbre, Lemme \ref{algebra-lemma-CM-goes-up}.
\end{proof}

\begin{lemma}
\label{lemma-flat-morphism-from-CM}
Soit $S$ un schéma.
Soit $f : X \to Y$ un morphisme plat entre espaces algébriques localement noethériens
sur $S$.
Si $X$ est Cohen-Macaulay, alors $f$ est Cohen-Macaulay et
$\mathcal{O}_{Y, f(\overline{x})}$ est Cohen-Macaulay pour tout $x \in |X|$.
\end{lemma}

\begin{proof}
En traduisant l'énoncé en algèbre à l'aide du Lemme \ref{lemma-CM}
(comparer avec la démonstration du
Lemme \ref{lemma-composition-CM}), le résultat découle du
Lemme \ref{algebra-lemma-CM-goes-up} d'Algèbre.
\end{proof}

\begin{lemma}
\label{lemma-base-change-CM}
Soit $S$ un schéma.
Soit $f : X \to Y$ un morphisme d'espaces algébriques sur $S$.
Supposons que les fibres de $f$ soient localement noethériennes.
Supposons que $Y' \to Y$ soit localement de type fini. Soit $f' : X' \to Y'$
le changement de base de $f$.
Soit $x' \in |X'|$ un point d'image $x \in |X|$.
\begin{enumerate}
\item Si $f$ est Cohen-Macaulay en $x$, alors
$f' : X' \to Y'$ est Cohen-Macaulay en $x'$.
\item Si $f$ est plat en $x$ et $f'$ est Cohen-Macaulay en $x'$, alors $f$
est Cohen-Macaulay en $x$.
\item Si $Y' \to Y$ est plat en $f'(x')$ et $f'$ est Cohen-Macaulay en
$x'$, alors $f$ est Cohen-Macaulay en $x$.
\end{enumerate}
\end{lemma}

\begin{proof}
Notons $y \in |Y|$ et $y' \in |Y'|$ les images de $x'$.
Choisissons un morphisme étale surjectif $V \to Y$ où $V$ est un schéma.
Choisissons un morphisme étale surjectif $U \to X \times_Y V$ où
$U$ est un schéma.
Choisissons un morphisme étale surjectif $V' \to Y' \times_Y V$
où $V'$ est un schéma.
Alors $U' = U \times_V V'$ est un schéma muni
d'un morphisme étale surjectif $U' \to X'$.
Choisissons $u' \in U'$ qui s'envoie sur $x'$. Notons $u \in U$ l'image de $u'$.
Le lemme découle alors du lemme appliqué à
$U \to V$, à son changement de base $U' \to V'$ et aux
points $u'$ et $u$ (cela découle des définitions).
Il découle donc du cas des schémas ; voir
Compléments sur les morphismes, Lemme \ref{more-morphisms-lemma-base-change-CM}.
\end{proof}

\begin{lemma}
\label{lemma-flat-finite-presentation-CM-open}
Soit $S$ un schéma.
Soit $f : X \to Y$ un morphisme d'espaces algébriques sur $S$
qui est plat et localement de présentation finie. Posons
$$
W = \{x \in |X| : f\text{ est Cohen-Macaulay en }x\}
$$
Alors $W$ est ouvert dans $|X|$ et la formation de $W$
commute à tout changement de base de $f$ :
pour tout morphisme $g : Y' \to Y$, considérons
le changement de base $f' : X' \to Y'$ de $f$ et la
projection $g' : X' \to X$. Alors l'ensemble correspondant
$W'$ pour le morphisme $f'$ est donné par $W' = (g')^{-1}(W)$.
\end{lemma}

\begin{proof}
Choisissons un diagramme commutatif
$$
\xymatrix{
U \ar[d] \ar[r] & V \ar[d] \\
X \ar[r] & Y
}
$$
dont les flèches verticales sont étales et où $U$ et $V$ sont des schémas.
Soit $u \in U$ d'image $x \in |X|$.
Alors $f$ est Cohen-Macaulay en $x$ si et seulement si $U \to V$ l'est
en $u$ (par définition).
Nous nous ramenons ainsi au cas du morphisme $U \to V$.
Voir Compléments sur les morphismes, Lemme
\ref{more-morphisms-lemma-flat-finite-presentation-CM-open}.
\end{proof}

\begin{lemma}
\label{lemma-lfp-CM-relative-dimension}
Soient $S$ un schéma et $f : X \to Y$ un
morphisme d'espaces algébriques sur $S$. Supposons que
$f$ soit localement de présentation finie et Cohen-Macaulay.
Il existe alors des sous-schémas ouverts et fermés $X_d \subset X$
tels que $X = \coprod_{d \geq 0} X_d$ et que $f|_{X_d} : X_d \to Y$
soit de dimension relative $d$.
\end{lemma}

\begin{proof}
Choisissons un diagramme commutatif
$$
\xymatrix{
U \ar[d] \ar[r] & V \ar[d] \\
X \ar[r] & Y
}
$$
dont les flèches verticales sont étales et où $U$ et $V$ sont des schémas.
Alors $U \to V$ est localement de présentation finie et Cohen-Macaulay
(cela résulte immédiatement de nos définitions).
Nous avons donc une décomposition $U = \coprod_{d \geq 0} U_d$
en sous-schémas ouverts et fermés telle que $f|_{U_d} : U_d \to V$
soit de dimension relative $d$ ; voir Morphismes, Lemme
\ref{morphisms-lemma-flat-finite-presentation-CM-fibres-relative-dimension}.
Soit $u \in U$ d'image $x \in |X|$. Alors
$f$ est de dimension relative $d$ en $x$ si et seulement si
$U \to V$ est de dimension relative $d$ en $u$
(cela découle de nos définitions).
Nous voyons ainsi que $U_d$ est l'image inverse
d'une partie $X_d \subset |X|$ qui est nécessairement
ouverte et fermée. Notons $X_d$ le sous-espace algébrique
ouvert et fermé correspondant de $X$ ; le lemme est démontré.
\end{proof}







\section{Morphismes de Gorenstein}
\label{section-gorenstein}

\noindent
Ceci est l'analogue de Dualité pour les schémas, section
\ref{duality-section-gorenstein-morphisms}.

\begin{lemma}
\label{lemma-gorenstein-local-ring-fibre}
La propriété suivante des morphismes de germes de schémas
\begin{align*}
& \mathcal{P}((X, x) \to (S, s)) = \\
& \text{l'anneau local }
\mathcal{O}_{X_s, x}
\text{ de la fibre est noethérien et de Gorenstein}
\end{align*}
est de nature locale pour la topologie étale sur la source et le but (Descente, Définition
\ref{descent-definition-local-source-target-at-point}).
\end{lemma}

\begin{proof}
Étant donné un diagramme comme dans
Descente, Définition \ref{descent-definition-local-source-target-at-point},
nous obtenons un morphisme étale entre les fibres
$U'_{v'} \to U_v$ qui envoie $u'$ sur $u$ ; voir
Descente, Lemme \ref{descent-lemma-etale-on-fiber}.
Ainsi, $\mathcal{O}_{U_v, u} \to \mathcal{O}_{U'_{v'}, u'}$
est la localisation d'un homomorphisme étale d'anneaux. Par conséquent,
le premier anneau est noethérien si et seulement si le second l'est ; voir
Compléments sur l'algèbre, Lemme \ref{more-algebra-lemma-Noetherian-etale-extension}.
Puisque $\mathcal{O}_{U'_{v'}, u'}/\mathfrak m_u \mathcal{O}_{U'_{v'}, u'}
= \kappa(u')$ (Algèbre, Lemme \ref{algebra-lemma-etale-at-prime})
est un anneau de Gorenstein, nous voyons que
$\mathcal{O}_{U_v, u}$ est de Gorenstein si et seulement si
$\mathcal{O}_{U'_{v'}, u'}$ l'est, d'après
Complexes dualisants, Lemme \ref{dualizing-lemma-flat-under-gorenstein}.
\end{proof}

\begin{definition}
\label{definition-gorenstein}
Soit $S$ un schéma.
Soit $f : X \to Y$ un morphisme d'espaces algébriques sur $S$.
Supposons que les fibres de $f$ soient localement noethériennes
(Diviseurs sur les espaces, Définition
\ref{spaces-divisors-definition-locally-Noetherian-fibre}).
\begin{enumerate}
\item Soient $x \in |X|$ et $y = f(x)$. Nous disons que $f$ est
{\it de Gorenstein en $x$} si $f$ est plat en $x$ et si
les conditions équivalentes de
Morphismes d'espaces, Lemme
\ref{spaces-morphisms-lemma-local-source-target-at-point}
sont satisfaites pour la propriété $\mathcal{P}$ décrite dans le
Lemme \ref{lemma-gorenstein-local-ring-fibre}.
\item Nous disons que $f$ est un {\it morphisme de Gorenstein} si $f$ est
de Gorenstein en tout point de $X$.
\end{enumerate}
\end{definition}

\noindent
Voici une reformulation.

\begin{lemma}
\label{lemma-gorenstein}
Soit $S$ un schéma. Soit $f : X \to Y$ un morphisme d'espaces algébriques
sur $S$. Supposons que les fibres de $f$ soient localement noethériennes.
Les conditions suivantes sont équivalentes :
\begin{enumerate}
\item $f$ est de Gorenstein,
\item $f$ est plat et il existe un morphisme étale surjectif $V \to Y$
où $V$ est un schéma, tel que les fibres de $X_V \to V$
soient des espaces algébriques de Gorenstein, et
\item $f$ est plat et, pour tout morphisme étale $V \to Y$
où $V$ est un schéma, les fibres de $X_V \to V$
sont des espaces algébriques de Gorenstein.
\end{enumerate}
Étant donné $x \in |X|$ d'image $y \in |Y|$, les conditions suivantes sont
équivalentes :
\begin{enumerate}
\item[(a)] $f$ est de Gorenstein en $x$, et
\item[(b)] $\mathcal{O}_{Y, \overline{y}} \to \mathcal{O}_{X, \overline{x}}$
est plat et
$\mathcal{O}_{X, \overline{x}}/
\mathfrak m_{\overline{y}}\mathcal{O}_{X, \overline{x}}$ est de Gorenstein.
\end{enumerate}
\end{lemma}

\begin{proof}
Étant donné un morphisme étale $V \to Y$ où $V$ est un schéma,
choisissons un schéma $U$ et un morphisme étale surjectif $U \to X \times_Y V$.
Considérons le diagramme commutatif
$$
\xymatrix{
U \ar[d] \ar[r] & V \ar[d] \\
X \ar[r] & Y
}
$$
Soit $u \in U$ d'images $x \in |X|$, $y \in |Y|$ et $v \in V$.
Alors $f$ est de Gorenstein en $x$ si et seulement si $U \to V$ l'est
en $u$ (par définition). De plus, le morphisme
$U_v \to X_v = (X_V)_v$ est étale surjectif. Par conséquent, le schéma $U_v$ est
de Gorenstein si et seulement si l'espace algébrique $X_v$ l'est.
L'équivalence de (1), (2) et (3) découle donc de
l'équivalence correspondante pour les morphismes de
schémas ; voir Dualité pour les schémas, Lemme \ref{duality-lemma-gorenstein} ;
l'argument est formel.

\medskip\noindent
Démonstration de l'équivalence de (a) et (b). L'équivalence correspondante
pour la platitude est donnée par Morphismes d'espaces, Lemme
\ref{spaces-morphisms-lemma-flat-at-point-etale-local-rings}.
Nous pouvons donc supposer $f$ plat en $x$ pour démontrer l'équivalence.
Considérons un diagramme et $x, y, u, v$ comme ci-dessus. Alors
$\mathcal{O}_{Y, \overline{y}} \to \mathcal{O}_{X, \overline{x}}$
est égal à l'homomorphisme
$\mathcal{O}_{V, v}^{sh} \to \mathcal{O}_{U, u}^{sh}$
entre les hensélisés stricts des anneaux locaux ; voir
Propriétés des espaces, Lemme
\ref{spaces-properties-lemma-describe-etale-local-ring}.
Nous avons donc
$$
\mathcal{O}_{X, \overline{x}}/
\mathfrak m_{\overline{y}}\mathcal{O}_{X, \overline{x}} =
(\mathcal{O}_{U, u}/\mathfrak m_v \mathcal{O}_{U, u})^{sh}
$$
d'après Algèbre, Lemme \ref{algebra-lemma-quotient-strict-henselization}.
Il reste donc à montrer que l'anneau local noethérien
$\mathcal{O}_{U, u}/\mathfrak m_v \mathcal{O}_{U, u}$
est de Gorenstein si et seulement si son hensélisé strict l'est.
Cela découle immédiatement de
Complexes dualisants, Lemme
\ref{dualizing-lemma-completion-henselization-dualizing}
et de la définition d'un anneau local de Gorenstein comme un
anneau local noethérien qui est un complexe dualisant sur lui-même.
\end{proof}

\begin{lemma}
\label{lemma-composition-gorenstein}
Soit $S$ un schéma.
Soient $f : X \to Y$ et $g : Y \to Z$ des morphismes d'espaces algébriques
sur $S$. Supposons que les
fibres de $f$, $g$ et $g \circ f$ soient localement noethériennes.
Soit $x \in |X|$ d'images $y \in |Y|$ et $z \in |Z|$.
\begin{enumerate}
\item Si $f$ est de Gorenstein en $x$ et $g$ est de Gorenstein
en $f(x)$, alors $g \circ f$ est de Gorenstein en $x$.
\item Si $f$ et $g$ sont de Gorenstein, alors $g \circ f$ l'est.
\item Si $g \circ f$ est de Gorenstein en $x$ et $f$ est plat en $x$,
alors $f$ est de Gorenstein en $x$ et $g$ l'est en $f(x)$.
\item Si $f \circ g$ est de Gorenstein et $f$ est plat, alors
$f$ est de Gorenstein et $g$ l'est en tout point de
l'image de $f$.
\end{enumerate}
\end{lemma}

\begin{proof}
Le résultat correspondant pour les schémas, appliqué localement pour la
topologie étale, donne l'assertion ; voir Dualité pour les schémas, Lemme
\ref{duality-lemma-composition-gorenstein}.
On peut aussi utiliser l'équivalence de (a) et (b) du
Lemme \ref{lemma-gorenstein}. Nous considérons alors l'homomorphisme local
d'anneaux locaux noethériens
$$
\mathcal{O}_{Y, \overline{y}}/
\mathfrak m_{\overline{z}}\mathcal{O}_{Y, \overline{y}}
\longrightarrow
\mathcal{O}_{X, \overline{x}}/
\mathfrak m_{\overline{z}}\mathcal{O}_{X, \overline{x}}
$$
dont la fibre est
$$
\mathcal{O}_{X, \overline{x}}/
\mathfrak m_{\overline{y}}\mathcal{O}_{X, \overline{x}}
$$
et nous appliquons Complexes dualisants, Lemme
\ref{dualizing-lemma-flat-under-gorenstein}.
\end{proof}

\begin{lemma}
\label{lemma-flat-morphism-from-gorenstein}
Soit $S$ un schéma.
Soit $f : X \to Y$ un morphisme plat entre espaces algébriques localement noethériens
sur $S$.
Si $X$ est de Gorenstein, alors $f$ est de Gorenstein et
$\mathcal{O}_{Y, f(\overline{x})}$ est de Gorenstein pour tout $x \in |X|$.
\end{lemma}

\begin{proof}
En traduisant l'énoncé en algèbre à l'aide du Lemme \ref{lemma-gorenstein}
(comparer avec la démonstration du
Lemme \ref{lemma-composition-gorenstein}), le résultat découle du
Lemme \ref{dualizing-lemma-flat-under-gorenstein} de Complexes dualisants.
\end{proof}

\begin{lemma}
\label{lemma-base-change-gorenstein}
Soit $S$ un schéma.
Soit $f : X \to Y$ un morphisme d'espaces algébriques sur $S$.
Supposons que les fibres de $f$ soient localement noethériennes.
Supposons que $Y' \to Y$ soit localement de type fini. Soit $f' : X' \to Y'$
le changement de base de $f$.
Soit $x' \in |X'|$ un point d'image $x \in |X|$.
\begin{enumerate}
\item Si $f$ est de Gorenstein en $x$, alors
$f' : X' \to Y'$ est de Gorenstein en $x'$.
\item Si $f$ est plat en $x$ et $f'$ est de Gorenstein en $x'$, alors $f$
est de Gorenstein en $x$.
\item Si $Y' \to Y$ est plat en $f'(x')$ et $f'$ est de Gorenstein en
$x'$, alors $f$ est de Gorenstein en $x$.
\end{enumerate}
\end{lemma}

\begin{proof}
Notons $y \in |Y|$ et $y' \in |Y'|$ les images de $x'$.
Choisissons un morphisme étale surjectif $V \to Y$ où $V$ est un schéma.
Choisissons un morphisme étale surjectif $U \to X \times_Y V$ où
$U$ est un schéma.
Choisissons un morphisme étale surjectif $V' \to Y' \times_Y V$
où $V'$ est un schéma.
Alors $U' = U \times_V V'$ est un schéma muni
d'un morphisme étale surjectif $U' \to X'$.
Choisissons $u' \in U'$ qui s'envoie sur $x'$. Notons $u \in U$ l'image de $u'$.
Le lemme découle alors du lemme appliqué à
$U \to V$, à son changement de base $U' \to V'$ et aux
points $u'$ et $u$ (cela découle des définitions).
Il découle donc du cas des schémas ; voir
Dualité pour les schémas, Lemme \ref{duality-lemma-base-change-gorenstein}.
\end{proof}

\begin{lemma}
\label{lemma-flat-finite-presentation-gorenstein-open}
Soit $S$ un schéma.
Soit $f : X \to Y$ un morphisme d'espaces algébriques sur $S$
qui est plat et localement de présentation finie. Posons
$$
W = \{x \in |X| : f\text{ est de Gorenstein en }x\}
$$
Alors $W$ est ouvert dans $|X|$ et la formation de $W$
commute à tout changement de base de $f$ :
pour tout morphisme $g : Y' \to Y$, considérons
le changement de base $f' : X' \to Y'$ de $f$ et la
projection $g' : X' \to X$. Alors l'ensemble correspondant
$W'$ pour le morphisme $f'$ est donné par $W' = (g')^{-1}(W)$.
\end{lemma}

\begin{proof}
Choisissons un diagramme commutatif
$$
\xymatrix{
U \ar[d] \ar[r] & V \ar[d] \\
X \ar[r] & Y
}
$$
Soit $u \in U$ d'image $x \in |X|$.
Alors $f$ est de Gorenstein en $x$ si et seulement si $U \to V$ l'est
en $u$ (par définition).
Nous nous ramenons ainsi au cas du morphisme $U \to V$
de schémas. L'ouverture est démontrée dans
Dualité pour les schémas, Lemme
\ref{duality-lemma-flat-finite-presentation-characterize-gorenstein}
et la compatibilité au changement de base dans
Dualité pour les schémas, Lemme \ref{duality-lemma-flat-lft-base-change-gorenstein}.
\end{proof}










\section{Tranchage des morphismes de Cohen-Macaulay}
\label{section-slice}

\noindent
Soit $S$ un schéma. Soit $X$ un espace algébrique sur $S$.
Soient $f_1, \ldots, f_r \in \Gamma(X, \mathcal{O}_X)$. Dans ce cas, nous
notons $V(f_1, \ldots, f_r)$ le {\it sous-espace fermé de $X$ défini par
$f_1, \ldots, f_r$}. Plus précisément, on peut définir $V(f_1, \ldots, f_r)$
comme le sous-espace fermé de $X$ correspondant au faisceau quasi-cohérent
d'idéaux engendré par $f_1, \ldots, f_r$ ; voir
Morphismes d'espaces, Lemme
\ref{spaces-morphisms-lemma-closed-immersion-ideals}.
On peut aussi choisir une présentation $X = U/R$ et considérer le
sous-schéma fermé $Z \subset U$ défini par $f_1|U, \ldots, f_r|_U$.
Il est clair que $Z$ est un sous-schéma fermé invariant par $R$ (voir
Groupoïdes, Définition \ref{groupoids-definition-invariant-open})
et l'on peut poser $V(f_1, \ldots, f_r) = Z/R_Z$.

\begin{lemma}
\label{lemma-slice}
Soit $S$ un schéma. Considérons un diagramme cartésien
$$
\xymatrix{
X \ar[d] & F \ar[l]^p \ar[d] \\
Y & \Spec(k) \ar[l]
}
$$
où $X \to Y$ est un morphisme d'espaces algébriques sur $S$
qui est plat et localement de présentation finie, et où
$k$ est un corps sur $S$. Soient $f_1, \ldots, f_r \in \Gamma(X, \mathcal{O}_X)$
et $z \in |F|$ tels que les images de $f_1, \ldots, f_r$ forment une suite régulière
dans l'anneau local $\mathcal{O}_{F, \overline{z}}$.
Alors, quitte à remplacer $X$ par un sous-espace ouvert contenant $p(z)$, le morphisme
$$
V(f_1, \ldots, f_r) \longrightarrow Y
$$
est plat et localement de présentation finie.
\end{lemma}

\begin{proof}
Posons $Z = V(f_1, \ldots, f_r)$. Il est clair que $Z \to X$ est localement de
présentation finie ; la composée $Z \to Y$ est donc localement de
présentation finie ; voir
Morphismes d'espaces,
Lemme \ref{spaces-morphisms-lemma-composition-finite-presentation}.
Il suffit donc de montrer que $Z \to Y$ est plat dans un voisinage de $p(z)$.
Soit $k'/k$ une extension de corps. Alors le morphisme de
$F' = F \times_{\Spec(k)} \Spec(k')$ vers
$F$ est surjectif et plat ; on peut donc trouver un point $z' \in |F'|$ qui s'envoie sur $z$
et l'homomorphisme local
$\mathcal{O}_{F, \overline{z}} \to \mathcal{O}_{F', \overline{z}'}$ est
plat ; voir
Morphismes d'espaces,
Lemme \ref{spaces-morphisms-lemma-flat-at-point-etale-local-rings}.
Par conséquent, les images de $f_1, \ldots, f_r$ dans
$\mathcal{O}_{F', \overline{z}'}$ forment encore une suite régulière ; voir
Algèbre, Lemme \ref{algebra-lemma-flat-increases-depth}.
Ainsi, au cours de la démonstration, nous pouvons remplacer $k$ par une extension de corps.
En particulier, nous pouvons supposer que $z \in |F|$ provient d'une section
$z : \Spec(k) \to F$ du morphisme structural $F \to \Spec(k)$.

\medskip\noindent
Choisissons un schéma $V$ et un morphisme étale surjectif
$V \to Y$. Choisissons un schéma $U$ et un morphisme étale surjectif
$U \to X \times_Y V$. Après avoir éventuellement agrandi encore $k$, nous pouvons
supposer que $\Spec(k) \to F \to X$ se factorise par $U$ (puisque
$U \to X$ est surjectif). Soit
$u : \Spec(k) \to U$ une telle factorisation et notons $v \in V$
l'image de $u$. Notons que les morphismes
$$
U_v \times_{\Spec(\kappa(v))} \Spec(k) =
U \times_V \Spec(k) \to U \times_Y \Spec(k) \to F
$$
sont étales (le premier comme changement de base de $V \to V \times_Y V$ et
le second comme changement de base de $U \to X$). De plus, par construction,
le point $u : \Spec(k) \to U$ définit un point de l'espace tout à gauche
qui s'envoie sur $z$ à droite. Par conséquent, les éléments
$f_1, \ldots, f_r$ ont pour images une suite régulière dans l'anneau local
situé à droite de l'homomorphisme suivant
$$
\mathcal{O}_{U_v, u}
\longrightarrow
\mathcal{O}_{U_v \times_{\Spec(\kappa(v)} \Spec(k), \overline{u}}
=
\mathcal{O}_{U \times_V \Spec(k), \overline{u}}.
$$
Mais, puisque la flèche affichée est plate (en combinant
Compléments sur la platitude, Lemme \ref{flat-lemma-flat-up-down-henselization}
et
Morphismes d'espaces,
Lemme \ref{spaces-morphisms-lemma-flat-at-point-etale-local-rings}),
nous déduisons de
Algèbre, Lemme \ref{algebra-lemma-flat-increases-depth},
que les images de $f_1, \ldots, f_r$ forment une suite régulière dans
$\mathcal{O}_{U_v, u}$. D'après
Compléments sur les morphismes, Lemme \ref{more-morphisms-lemma-slice-given-elements},
nous concluons que le morphisme de schémas
$$
V(f_1, \ldots, f_r) \times_X U = V(f_1|_U, \ldots, f_r|_U) \to V
$$
est plat dans un voisinage ouvert $U'$ de $u$. Soit $X' \subset X$
le sous-espace ouvert correspondant à l'image de
$|U'| \to |X|$ (voir
Propriétés des espaces, Lemmes
\ref{spaces-properties-lemma-topology-points} et
\ref{spaces-properties-lemma-open-subspaces}).
Nous concluons que $V(f_1, \ldots, f_r) \cap X' \to Y$ est plat
(voir
Morphismes d'espaces, Définition \ref{spaces-morphisms-definition-flat}),
car
nous avons le diagramme commutatif
$$
\xymatrix{
V(f_1, \ldots, f_r) \times_X U' \ar[d]_a \ar[r] & V \ar[d]^b \\
V(f_1, \ldots, f_r) \cap X' \ar[r] & Y
}
$$
où $a, b$ sont étales et $a$ est surjectif.
\end{proof}






\section{Fibres réduites}
\label{section-reduced}

\noindent
Cette section est l'analogue de
Compléments sur les morphismes, section \ref{more-morphisms-section-reduced}.

\begin{lemma}
\label{lemma-geometrically-reduced-fibre}
Soit $S$ un schéma. Soit $f : X \to Y$ un
morphisme d'espaces algébriques sur $S$. Soit $y \in |Y|$.
Les conditions suivantes sont équivalentes :
\begin{enumerate}
\item pour un certain morphisme $\Spec(k) \to Y$ dans la classe d'équivalence
de $y$, l'espace algébrique $X_k$ est géométriquement réduit sur $k$ ;
\item pour tout morphisme $\Spec(k) \to Y$ dans la classe d'équivalence
de $y$, l'espace algébrique $X_k$ est géométriquement réduit sur $k$ ;
\item pour tout morphisme $\Spec(k) \to Y$ dans la classe d'équivalence
de $y$, l'espace algébrique $X_k$ est réduit.
\end{enumerate}
\end{lemma}

\begin{proof}
Cela résulte immédiatement du Lemme
\ref{spaces-over-fields-lemma-geometrically-reduced-upstairs} d'Espaces sur les corps
et de la définition de la relation d'équivalence qui définit $|X|$,
donnée dans
Propriétés des espaces, section \ref{spaces-properties-section-points}.
\end{proof}

\begin{definition}
\label{definition-geometrically-reduced-fibre}
Soit $S$ un schéma. Soit $f : X \to Y$ un
morphisme d'espaces algébriques sur $S$. Soit $y \in |Y|$.
Nous disons que {\it la fibre de $f : X \to Y$ en $y$ est géométriquement réduite}
si les conditions équivalentes du
Lemme \ref{lemma-geometrically-reduced-fibre} sont satisfaites.
\end{definition}

\noindent
Voici les lemmes obligatoires.

\begin{lemma}
\label{lemma-base-change-fibres-geometrically-reduced}
Soit $S$ un schéma. Soient $f : X \to Y$ et $g : Y' \to Y$
des morphismes d'espaces algébriques sur $S$. Notons
$f' : X' \to Y'$ le changement de base de $f$ par $g$. Alors
\begin{align*}
\{y' \in |Y'| :
\text{la fibre de }f' : X' \to Y'\text{ en }y'
\text{ est géométriquement réduite}\} \\
= g^{-1}(\{y \in |Y| :
\text{la fibre de }f : X \to Y\text{ en }y
\text{ est géométriquement réduite}\}).
\end{align*}
\end{lemma}

\begin{proof}
Pour $y' \in |Y'|$, choisissons un morphisme $\Spec(k) \to Y'$
dans la classe d'équivalence de $y'$. Alors $g(y')$ est
représenté par la composée $\Spec(k) \to Y' \to Y$.
Ainsi $X' \times_{Y'} \Spec(k) = X \times_Y \Spec(k)$
et le résultat découle de la définition.
\end{proof}

\begin{lemma}
\label{lemma-geometrically-reduced-constructible}
Soit $S$ un schéma. Soit $f : X \to Y$ un morphisme d'espaces algébriques
sur $S$ qui est quasi-compact et
localement de présentation finie. Alors l'ensemble
$$
E = \{y \in |Y| : \text{la fibre de }f : X \to Y\text{ en }y
\text{ est géométriquement réduite}\}
$$
est localement constructible pour la topologie étale.
\end{lemma}

\begin{proof}
Choisissons un schéma affine $V$ et un morphisme étale $V  \to Y$.
L'énoncé signifie que l'image inverse de $E$
dans $|V|$ est constructible. D'après le
Lemme \ref{lemma-base-change-fibres-geometrically-reduced},
nous pouvons remplacer $Y$ par $V$, c'est-à-dire supposer que $Y$
est un schéma affine. Alors $X$ est quasi-compact. Choisissons un
schéma affine $U$ et un morphisme étale surjectif $U \to X$.
Pour un morphisme $\Spec(k) \to Y$, le morphisme entre les fibres
$U_k \to X_k$ est étale surjectif. Ainsi $U_k$ est géométriquement
réduit sur $k$ si et seulement si $X_k$ est géométriquement réduit
sur $k$ ; voir Espaces sur les corps, Lemme
\ref{spaces-over-fields-lemma-geometrically-reduced-etale-local}.
Ainsi, l'ensemble $E$ associé à $X \to Y$ est le même que l'ensemble $E$
associé à $U \to Y$.
Nous voyons de cette manière que le lemme découle du cas
des schémas ; voir Compléments sur les morphismes, Lemme
\ref{more-morphisms-lemma-geometrically-reduced-constructible}.
\end{proof}

\begin{lemma}
\label{lemma-proper-flat-over-dvr-reduced-fibre}
Soit $X$ un espace algébrique sur un anneau de valuation discrète $R$
dont le morphisme structural $X \to \Spec(R)$ est propre et plat.
Si la fibre spéciale est réduite, alors
$X$ et la fibre générique $X_\eta$ sont tous deux réduits.
\end{lemma}

\begin{proof}
Choisissons un morphisme étale $U \to X$ où $U$ est un schéma affine.
Alors $U$ est de type fini sur $R$. Soit $u \in U$ un point de la fibre spéciale.
L'anneau local $A = \mathcal{O}_{U, u}$ est essentiellement de type
fini sur $R$, donc noethérien. Soit $\pi \in R$ une uniformisante.
Comme $X$ est plat sur $R$, nous voyons que $\pi \in \mathfrak m_A$
est un non diviseur de zéro dans $A$ et, comme la fibre spéciale de $X$
est réduite, l'anneau $A/\pi A$ est réduit.
Si $a \in A$, $a \not = 0$, alors il existe un $n \geq 0$ et un élément
$a' \in A$ tels que $a = \pi^n a'$ et $a' \not \in \pi A$.
Cela résulte du théorème d'intersection de Krull
(Algèbre, Lemme \ref{algebra-lemma-intersect-powers-ideal-module-zero}).
Si $a$ est nilpotent, $a'$ l'est aussi, car $\pi$ est un non diviseur de zéro.
Mais $a'$ s'envoie sur un élément non nul de l'anneau réduit $A/\pi A$,
ce qui est impossible. Ainsi $A$ est réduit. Il s'ensuit qu'il
existe un voisinage ouvert de $u$ dans $U$ qui est réduit
(petit détail omis ; utiliser le fait que $U$ est noethérien).
On peut donc trouver un morphisme étale $U \to X$ où $U$ est un schéma
réduit, tel que tout point de la fibre spéciale de $X$ appartienne
à l'image. Comme $X$ est propre sur $R$, il s'ensuit que
$U \to X$ est surjectif. Ainsi $X$ est réduit. Comme la fibre générique de
$U \to \Spec(R)$ est également réduite (sur les morceaux affines,
elle se calcule par localisation), nous concluons qu'il en est de même
de la fibre générique.
\end{proof}

\begin{lemma}
\label{lemma-geometrically-reduced-open}
Soit $S$ un schéma. Soit $f : X \to Y$ un morphisme d'espaces
algébriques sur $S$. Si $f$ est plat, propre et de présentation finie,
alors l'ensemble
$$
E = \{y \in |Y| : \text{la fibre de }f : X \to Y\text{ en }y
\text{ est géométriquement réduite}\}
$$
est ouvert dans $|Y|$.
\end{lemma}

\begin{proof}
Par le Lemme \ref{lemma-base-change-fibres-geometrically-reduced},
la formation de $E$ commute au changement de base. Pour vérifier qu'une partie
de $|Y|$ est ouverte, nous pouvons remplacer $Y$ par les membres d'un
recouvrement étale. Nous pouvons donc supposer $Y$ affine.
Alors $Y$ est une limite projective filtrante de schémas
affines de type fini sur $\mathbf{Z}$.
Par conséquent, nous pouvons supposer que $X \to Y$ est le
changement de base de $X_0 \to Y_0$, où $Y_0$ est le spectre d'une algèbre
de type fini sur $\mathbf{Z}$ et $X_0 \to Y_0$ est plat et propre.
Voir Limites d'espaces, Lemmes
\ref{spaces-limits-lemma-descend-finite-presentation},
\ref{spaces-limits-lemma-descend-flat}, et
\ref{spaces-limits-lemma-eventually-proper}. Puisque la formation de
$E$ commute au changement de base (voir ci-dessus),
nous pouvons supposer la base noethérienne.

\medskip\noindent
Supposons $Y$ noethérien. Cet ensemble est constructible d'après le
Lemme \ref{lemma-geometrically-reduced-constructible}.
Il suffit donc de montrer que l'ensemble est stable par générisation
(Topologie, Lemme \ref{topology-lemma-characterize-closed-Noetherian}). D'après
Propriétés, Lemme \ref{properties-lemma-locally-Noetherian-specialization-dvr},
nous nous ramenons au cas où $Y = \Spec(R)$, où $R$ est un anneau de
valuation discrète et où la fibre fermée $X_y$ est géométriquement
réduite. Il reste à montrer que la fibre générique $X_\eta$ est géométriquement réduite.

\medskip\noindent
Sinon, il existe une extension finie $L$ du corps des
fractions de $R$ telle que $X_L$ ne soit pas réduit ; voir
Espaces sur les corps, Lemmes
\ref{spaces-over-fields-lemma-perfect-reduced} (caractéristique zéro) et
\ref{spaces-over-fields-lemma-geometrically-reduced-positive-characteristic}
(caractéristique positive). Il existe un anneau de valuation discrète
$R' \subset L$ de corps des fractions $L$ qui domine $R$ ; voir
Algèbre, Lemme \ref{algebra-lemma-integral-closure-Dedekind}.
Après avoir remplacé $R$ par $R'$, nous nous ramenons au
Lemme \ref{lemma-proper-flat-over-dvr-reduced-fibre}.
\end{proof}







\section{Composantes connexes des fibres}
\label{section-connected}

\noindent
Cette section est l'analogue de
Compléments sur les morphismes, section \ref{more-morphisms-section-connected}.

\begin{lemma}
\label{lemma-base-change-fibres-nr-geometrically-connected-components}
Soit $S$ un schéma.
Soit $f : X \to Y$ un morphisme d'espaces algébriques sur $S$. Soit
$$
n_{X/Y} : |Y| \to \{0, 1, 2, 3, \ldots, \infty\}
$$
la fonction qui associe à $y \in Y$ le nombre de composantes connexes
de $X_k$, où $\Spec(k) \to Y$ appartient à la classe d'équivalence
de $y$ et où $k$ est algébriquement clos.
Cette fonction est bien définie et, si $g : Y' \to Y$ est un morphisme, alors
$$
n_{X'/Y'} = n_{X/Y} \circ g
$$
où $X' \to Y'$ est le changement de base de $f$.
\end{lemma}

\begin{proof}
Supposons que $y' \in Y'$ ait pour image $y \in Y$. Soit $\Spec(k') \to Y'$
un morphisme dans la classe d'équivalence de $y'$, avec $k'$ algébriquement clos.
On peut alors choisir un diagramme commutatif
$$
\xymatrix{
\Spec(K) \ar[r] \ar[rd] &
\Spec(k') \ar[r] & Y' \ar[d] \\
& \Spec(k) \ar[r] & Y
}
$$
où $K$ est un corps algébriquement clos.
Le résultat vient de ce que les morphismes de schémas
$$
\xymatrix{
X'_{k'} & (X'_{k'})_K = (X_k)_K \ar[l] \ar[r] & X_k
}
$$
induisent des bijections entre les composantes connexes ; voir
Espaces sur les corps, Lemme
\ref{spaces-over-fields-lemma-separably-closed-field-connected-components}.
Pour en déduire que la fonction est bien définie, prendre $Y' = Y$.
\end{proof}







\section{Dimension des fibres}
\label{section-dimension}

\noindent
Cette section est l'analogue de
Compléments sur les morphismes, section \ref{more-morphisms-section-dimension}.

\begin{lemma}
\label{lemma-dimension-fibre}
Soit $S$ un schéma. Soit $f : X \to Y$ un morphisme de type fini
d'espaces algébriques sur $S$. Soit $y \in |Y|$. Les grandeurs suivantes
sont égales :
\begin{enumerate}
\item $d = -\infty$ si $y$ n'appartient pas à l'image de $|f|$ et,
sinon, le plus petit entier $d$ tel que $f$ soit de dimension relative $\leq d$
en tout $x \in |X|$ s'envoyant sur $y$ ;
\item la dimension de l'espace algébrique $X_k = \Spec(k) \times_Y X$
pour tout morphisme $\Spec(k) \to Y$ dans la classe d'équivalence qui définit $y$.
\end{enumerate}
\end{lemma}

\begin{proof}
Pour interpréter cet énoncé, il faut consulter
Morphismes d'espaces, Définition
\ref{spaces-morphisms-definition-dimension-fibre},
Propriétés des espaces,
Définition \ref{spaces-properties-definition-dimension},
Propriétés des espaces,
Définition \ref{spaces-properties-definition-dimension-at-point}.
Nous allons montrer que les nombres de (1) et (2) sont égaux pour
un morphisme fixé $\Spec(k) \to Y$.
Choisissons un morphisme étale $V \to Y$, où $V$ est un schéma
affine, ainsi qu'un point $v \in V$ s'envoyant sur $y$.
Puisque $V \times_Y \Spec(k) \to \Spec(k)$ est étale surjectif
(d'après Propriétés des espaces, Lemme \ref{spaces-properties-lemma-points-cartesian}),
on peut trouver une extension séparable finie $k'/k$
(d'après Morphismes, Lemme \ref{morphisms-lemma-etale-over-field})
et un diagramme commutatif
$$
\xymatrix{
\Spec(k') \ar[r] \ar[d] & V \ar[d] \\
\Spec(k) \ar[r] & Y
}
$$
Nous pouvons remplacer $X \to Y$ par $V \times_Y X \to V$ et
$X_k$ par $X_{k'} = \Spec(k') \times_V (V \times_Y X)$,
car cela ne change pas les dimensions considérées d'après
Propriétés des espaces, Lemme
\ref{spaces-properties-lemma-dimension-decent-invariant-under-etale}
et Morphismes d'espaces, Lemme
\ref{spaces-morphisms-lemma-dimension-fibre-after-base-change}.
Ainsi, nous pouvons supposer que $Y$ est un schéma affine.
Dans ce cas, nous pouvons supposer que $k = \kappa(y)$,
car les dimensions de $X_{\kappa(y)}$ et de $X_k$
sont égales d'après le Lemme
\ref{spaces-morphisms-lemma-dimension-fibre-after-base-change} de Morphismes d'espaces
et le fait que, pour un espace algébrique $Z$ sur un corps $K$, la
dimension relative de $Z$ en un point $z \in |Z|$
est égale à $\dim_z(Z)$ par définition.
Supposons $Y$ affine et $k = \kappa(y)$. Alors,
comme $X$ est quasi-compact, nous pouvons choisir un schéma affine $U$ et
un morphisme étale surjectif $U \to X$.
Alors $\dim(X_k) = \dim(U_k) = \max \dim_u(U_k)$
est égal, par définition, au nombre donné en (1).
\end{proof}

\begin{lemma}
\label{lemma-base-change-dimension-fibres}
Soit $S$ un schéma. Soit $f : X \to Y$ un morphisme de type fini
d'espaces algébriques sur $S$. Soit
$$
n_{X/Y} : |Y| \to \{-\infty, 0, 1, 2, 3, \ldots\}
$$
la fonction qui associe à $y \in |Y|$
l'entier étudié au Lemme \ref{lemma-dimension-fibre}.
Si $g : Y' \to Y$ est un morphisme, alors
$$
n_{X'/Y'} = n_{X/Y} \circ |g|
$$
où $X' \to Y'$ est le changement de base de $f$.
\end{lemma}

\begin{proof}
Cela résulte immédiatement du
Lemme \ref{lemma-dimension-fibre}.
\end{proof}

\begin{lemma}
\label{lemma-dimension-fibres-flat}
Soit $S$ un schéma.
Soit $f : X \to Y$ un morphisme plat et de présentation finie
d'espaces algébriques sur $S$. Soit
$n_{X/Y}$ la fonction sur $Y$ donnant la dimension des fibres de $f$
introduite au Lemme \ref{lemma-base-change-dimension-fibres}.
Alors $n_{X/Y}$ est semi-continue inférieurement.
\end{lemma}

\begin{proof}
Soit $V \to Y$ un morphisme étale surjectif, où $V$ est un schéma.
Si nous pouvons montrer que la composée $n_{X/Y} \circ |g|$
est semi-continue inférieurement, alors le lemme en découle, car l'application $|g|$ est ouverte.
Ainsi, nous pouvons supposer que $Y$ est un schéma.
En raisonnant localement, nous pouvons supposer que $V$ est un schéma affine.
On peut alors choisir un schéma affine $U$ et un morphisme
étale surjectif $U \to X$. Alors $n_{X/Y} = n_{U/Y}$.
Ainsi, nous pouvons supposer que $X$ et $Y$ sont tous deux des schémas.
Dans ce cas, le lemme résulte de
Compléments sur les morphismes, Lemme \ref{more-morphisms-lemma-dimension-fibres-flat}.
\end{proof}

\begin{lemma}
\label{lemma-dimension-fibres-proper}
Soit $S$ un schéma.
Soit $f : X \to Y$ un morphisme propre d'espaces algébriques sur $S$. Soit
$n_{X/Y}$ la fonction sur $Y$ donnant la dimension des fibres de $f$
introduite au Lemme \ref{lemma-base-change-dimension-fibres}.
Alors $n_{X/Y}$ est semi-continue supérieurement.
\end{lemma}

\begin{proof}
Soit $Z_d = \{x \in |X| :
\text{la fibre de }f\text{ en }x\text{ est de dimension }> d\}$.
Alors $Z_d$ est une partie fermée de $|X|$ d'après
Morphismes d'espaces, Lemme
\ref{spaces-morphisms-lemma-openness-bounded-dimension-fibres}.
Puisque $f$ est propre, $f(Z_d)$ est fermé dans $|Y|$.
Comme $y \in f(Z_d) \Leftrightarrow n_{X/Y}(y) > d$,
le lemme est démontré.
\end{proof}

\begin{lemma}
\label{lemma-dimension-fibres-proper-flat}
Soit $S$ un schéma. Soit $f : X \to Y$ un morphisme d'espaces algébriques
sur $S$, propre, plat et de présentation finie.
Soit $n_{X/Y}$ la fonction sur $Y$ donnant la dimension des fibres de $f$
introduite au Lemme \ref{lemma-base-change-dimension-fibres}.
Alors $n_{X/Y}$ est localement constante.
\end{lemma}

\begin{proof}
Conséquence immédiate des
Lemmes \ref{lemma-dimension-fibres-flat} et
\ref{lemma-dimension-fibres-proper}.
\end{proof}





\section{Espaces algébriques caténaires}
\label{section-catenary}

\noindent
Cette section poursuit la discussion commencée dans
Espaces décents, section \ref{decent-spaces-section-catenary}.
Le lemme qui suit sera utilisé dans la démonstration du suivant.

\begin{lemma}
\label{lemma-construct-glueing}
Soit $S$ un schéma. Soit $f : X \to Y$ un morphisme entier
d'espaces algébriques sur $S$.
Soit $y \in |Y|$ un point représentable par une immersion fermée
$y : \Spec(k) \to Y$. Il existe alors
une factorisation $X \to X' \to Y$ de $f$ telle que
\begin{enumerate}
\item $X' \to Y$ soit entier ;
\item $X \to X'$ soit un isomorphisme au-dessus de $X' \setminus X'_y$ ;
\item $X'_y$ possède un unique point $x'$ tel que $\kappa(x') = k$.
\end{enumerate}
De plus, si $f$ est fini et $Y$ est localement noethérien, alors
$X' \to Y$ est fini.
\end{lemma}

\begin{proof}
D'après Morphismes d'espaces, Lemme \ref{spaces-morphisms-lemma-pushforward},
les faisceaux $f_*\mathcal{O}_X$, $(X_y \to Y)_*\mathcal{O}_{X_y}$ et
$y_*\mathcal{O}_{\Spec(k)}$ sont des faisceaux quasi-cohérents
d'algèbres sur $\mathcal{O}_Y$. Considérons les applications
$$
f_*\mathcal{O}_Y \longrightarrow
(X_y \to Y)_*\mathcal{O}_{X_y} \longleftarrow
y_*\mathcal{O}_{\Spec(k)}
$$
Le produit fibré est un faisceau quasi-cohérent d'algèbres sur $\mathcal{O}_Y$,
noté $\mathcal{A}'$, et l'on peut définir $X' \to Y$ comme le spectre relatif
de $\mathcal{A}'$ sur $Y$ ; voir
Morphismes, Lemme \ref{morphisms-lemma-affine-equivalence-algebras}.
Cette construction commute à tout changement de base.
En particulier, il est clair que, sur le sous-espace ouvert
$|Y| \setminus \{y\}$, le morphisme $X \to X'$ est un isomorphisme
et que, sur $|Y| \setminus \{y\}$, le morphisme $X' \to Y$ est entier.
Il reste à démontrer les assertions dans un petit voisinage de $y$.
Choisissons un schéma affine $V = \Spec(R)$ et un morphisme
étale $\varphi : V \to Y$ tel que $y$ appartienne à l'image de
$\varphi$. Alors $V_y$ est un sous-schéma fermé de $V$,
étale sur $k$, et se compose donc d'un nombre fini de points
dont chacun a un corps résiduel séparable sur $k$
(voir Espaces décents, Remarque \ref{decent-spaces-remark-recall}).
Après avoir rétréci $V$,
nous pouvons supposer qu'il existe un unique point fermé $v = \Spec(l) \to V$
s'envoyant sur $y$, avec $l/k$ fini séparable.
On peut écrire $V \times_Y X = \Spec(C)$, où $R \to C$
est un homomorphisme d'anneaux entier. La compatibilité énoncée avec le
changement de base donne $U \times_X Y' = \Spec(C')$, où
$$
C' = C \times_{C \otimes_R l} l
$$
Comme $R \to l$ est surjectif, $C \to C \otimes_R l$
l'est aussi, et nous voyons qu'il s'agit d'un produit fibré du
genre étudié dans Compléments sur l'algèbre, Situation
\ref{more-algebra-situation-module-over-fibre-product}
(où $A, A', B, B'$ correspondent à $C \otimes_R l, C, l, C'$).
Remarquons que $C'$ est une $R$-sous-algèbre de $C$ et est donc
entière sur $R$ ; cela démontre (1).
Enfin, Compléments sur l'algèbre, Lemme
\ref{more-algebra-lemma-points-of-fibre-product}
montre que $V \times_X Y' = \Spec(C')$
possède un unique point $y''$ au-dessus de $v$, de corps résiduel $l$
(cela correspond à l'application surjective évidente $C' \to l$).
Ainsi, $X_y \times_{\Spec(k)} \Spec(l)$ possède un unique point
de corps résiduel $l$. Puisque $l/k$ est fini séparable,
cela implique que $X'_y$ possède un unique point de corps résiduel $k$, c'est-à-dire
que (3) est vérifiée.

\medskip\noindent
Pour démontrer la dernière assertion, remarquons que si $Y$ est localement noethérien,
alors $R$ est un anneau noethérien et, si $f$ est fini, $R \to C$ est
fini. Alors $C'$ est une $R$-algèbre de type fini d'après Compléments sur l'algèbre,
Lemme \ref{more-algebra-lemma-fibre-product-finite-type}.
Cela démontre que $X' \to Y$ est fini.
\end{proof}

\begin{lemma}
\label{lemma-universally-catenary-dimension-function}
Soit $S$ un schéma. Soit $B$ un espace algébrique sur $S$.
Soit $\delta : |B| \to \mathbf{Z}$ une fonction.
Supposons que $B$ soit décent, localement noethérien et
universellement caténaire, et que $\delta$ soit une fonction de dimension.
Si $X$ est un espace algébrique décent sur $B$ dont le morphisme structural
$f : X \to B$ est localement de type fini, nous définissons
$\delta_X : |X| \to \mathbf{Z}$ par la règle
$$
\delta_X(x) = \delta(f(x)) + \text{degré de transcendance de }x/f(x)
$$
(Morphismes d'espaces, Définition
\ref{spaces-morphisms-definition-dimension-fibre}).
Alors $\delta_X$ est une fonction de dimension.
\end{lemma}

\begin{proof}
La question est locale sur $B$. Nous pouvons donc supposer $B$ quasi-compact.
D'après Espaces décents, Lemme
\ref{decent-spaces-lemma-locally-Noetherian-decent-quasi-separated}
nous voyons que $B$ est quasi-séparé. D'après
Limites d'espaces, Proposition
\ref{spaces-limits-proposition-there-is-a-scheme-finite-over}
nous pouvons choisir un morphisme fini surjectif $\pi : Y \to X$,
où $Y$ est un schéma.
Affirmation : $\delta_Y$ est une fonction de dimension.

\medskip\noindent
L'affirmation implique le lemme. Pour $X \to B$ comme dans le lemme,
posons $Z = Y \times_B X$, de projections $p : Z \to Y$ et $q : Z \to X$.
Nous avons alors
$$
\delta_Z(z) = \delta_Y(p(z)) + \text{degré de transcendance de }z/p(z)
$$
et $\delta_Z(z) = \delta_X(q(z))$. Cela résulte de
Morphismes d'espaces, Lemme 
\ref{spaces-morphisms-lemma-dimension-fibre-at-a-point-additive}
et du fait que ces degrés de transcendance sont nuls
pour les morphismes finis. D'après Espaces décents, Lemme
\ref{decent-spaces-lemma-scheme-with-dimension-function}
et l'affirmation, nous trouvons que $\delta_Z$ est une fonction de dimension.
Nous en déduisons que $\delta_X$ est une fonction de dimension d'après
Espaces décents, Lemme
\ref{decent-spaces-lemma-check-dimension-function-finite-cover}.

\medskip\noindent
Démonstration de l'affirmation. Considérons une spécialisation $y \leadsto y'$,
$y \not = y'$, de points du schéma noethérien $Y$.
Alors $\delta_Y(y) > \delta_Y(y')$, car il n'existe
aucune spécialisation entre deux points d'une même fibre de $Y$
(voir Espaces décents, Lemme
\ref{decent-spaces-lemma-conditions-on-fibre-and-qf}).
En appliquant ceci à une chaîne de spécialisations, nous obtenons
$$
\delta_Y(y) - \delta_Y(y') \geq
\text{codim}(\overline{\{y'\}}, \overline{\{y\}})
$$
Il nous reste à démontrer l'égalité. D'après
Propriétés, Lemme \ref{properties-lemma-locally-Noetherian-closed-point},
nous pouvons choisir une spécialisation $y' \leadsto y_0$.
Il suffit de montrer que
$\delta_Y(y) - \delta_Y(y_0) =
\text{codim}(\overline{\{y_0\}}, \overline{\{y\}})$
car cela entraînera l'égalité à la fois pour
$y \leadsto y'$ et pour $y' \leadsto y_0$.

\medskip\noindent
Choisissons une chaîne maximale
$y = y_c \leadsto y_{c - 1} \leadsto \ldots \leadsto y_0$
de spécialisations dans $Y$.
Posons $b = \pi(y)$ et $b_0 = \pi(y_0)$.
Choisissons une chaîne maximale
$b = b_e \leadsto b_{e - 1} \leadsto \ldots \leadsto b_0$
de spécialisations dans $|B|$.
Nous devons montrer que $e = c$.
Puisque $\pi$ est fermé
(Morphismes d'espaces, Lemme \ref{spaces-morphisms-lemma-finite-proper}),
nous pouvons trouver une suite de spécialisations
$y = y'_e \leadsto y'_{e - 1} \leadsto \ldots \leadsto y'_0$
s'envoyant sur
$b = b_e \leadsto b_{e - 1} \leadsto \ldots \leadsto b_0$.
Remarquons que $y'_e \leadsto y'_{e - 1} \leadsto \ldots \leadsto y'_0$
est également une chaîne maximale.
Si $y_0 = y'_0$, alors, comme $Y$ est caténaire, nous concluons
que $e = c$, comme souhaité. Dans le paragraphe suivant, nous nous ramenons à
ce cas par une astuce et concluons de la même manière.

\medskip\noindent
Puisque $\pi$ est fermé, nous voyons que $b_0$ est un point fermé de $|B|$.
D'après Espaces décents, Lemme \ref{decent-spaces-lemma-decent-space-closed-point},
nous pouvons représenter $b_0$ par une immersion fermée $b_0 : \Spec(k) \to B$.
Par le Lemme \ref{lemma-construct-glueing},
nous pouvons trouver une factorisation
$$
Y \to Y' \to X
$$
où $\pi' : Y' \to X$ est fini et $Y \to Y'$ est un morphisme qui
envoie $y_0$ et $y'_0$ sur le même point et est un isomorphisme
hors de l'image inverse de $b_0$.
(Bien entendu, $Y'$ ne sera pas un schéma, mais cela n'affecte pas
l'argument qui suit.)
Il est clair que les chaînes maximales de spécialisations
$y_c \leadsto y_{c - 1} \leadsto \ldots \leadsto y_0$ et
$y'_e \leadsto y'_{e - 1} \leadsto \ldots \leadsto y'_0$
s'envoient sur des chaînes maximales de spécialisations dans $Y'$
ayant le même début et la même fin.
Comme $B$ est universellement caténaire, nous voyons que
$|Y'|$ est caténaire et concluons comme précédemment.
\end{proof}






\section{Localisation étale des morphismes}
\label{section-etale-localization}

\noindent
Cette section est l'analogue de
Compléments sur les morphismes, section \ref{more-morphisms-section-etale-localization}.

\begin{lemma}
\label{lemma-etale-splits-off-quasi-finite-part}
Soit $S$ un schéma.
Soit $f : X \to Y$ un morphisme d'espaces algébriques sur $S$.
Soit $y \in |Y|$. Soient $x_1, \ldots, x_n \in |X|$ des points s'envoyant sur $y$.
Supposons que
\begin{enumerate}
\item $f$ soit localement de type fini ;
\item $f$ soit séparé ;
\item $f$ soit quasi-fini en $x_1, \ldots, x_n$ ;
\item $f$ soit quasi-compact ou $Y$ soit décent.
\end{enumerate}
Il existe alors un morphisme étale $(U, u) \to (Y, y)$
d'espaces algébriques pointés et une décomposition
$$
U \times_Y X = W \amalg V
$$
en sous-espaces ouverts et fermés telle que le morphisme $V \to U$ soit fini,
que tout point de la fibre de $|V| \to |U|$ au-dessus de $u$ s'envoie sur l'un des $x_i$,
et que la fibre de $|W| \to |U|$ au-dessus de $u$ ne contienne aucun point
s'envoyant sur l'un des $x_i$.
\end{lemma}

\begin{proof}
Soit $(U, u) \to (Y, y)$ un morphisme étale d'espaces algébriques,
et considérons l'ensemble des $w \in |U \times_Y X|$ s'envoyant sur $u \in |U|$
et sur l'un des $x_i \in |X|$. D'après
Espaces décents, Lemme \ref{decent-spaces-lemma-qf-and-qc-finite-fibre}
(si $f$ est de type fini) ou
Espaces décents, Lemme \ref{decent-spaces-lemma-decent-finite-fibre}
(si $Y$ est décent), cet ensemble est fini.
Nous pouvons donc remplacer $f$ par le changement de base
$U \times_Y X \to U$ et $x_1, \ldots, x_n$ par l'ensemble de ces $w$.
En particulier, nous pouvons supposer, et nous le faisons, que $Y$ est un schéma affine ;
alors $X$ est un espace algébrique séparé.

\medskip\noindent
Choisissons un schéma affine $Z$ et un morphisme étale $Z \to X$ tels que
$x_1, \ldots, x_n$ appartiennent à l'image de $|Z| \to |X|$. Les fibres de
$|Z| \to |X|$ sont finies ; voir Propriétés des espaces, Lemme
\ref{spaces-properties-lemma-finite-fibres-presentation}
(ou la discussion plus générale dans Espaces décents, section
\ref{decent-spaces-section-reasonable-decent}).
Soit $\{z_1, \ldots, z_m\} \subset |Z|$ l'image inverse de
$\{x_1, \ldots, x_n\}$. D'après Compléments sur les morphismes, Lemme
\ref{more-morphisms-lemma-etale-splits-off-quasi-finite-part-technical}
il existe un morphisme étale $(U, u) \to (Y, y)$ tel
que $U \times_Y Z = Z_1 \amalg Z_2$, avec $Z_1 \to U$ fini et
$(Z_1)_y = \{z_1, \ldots, z_m\}$. Nous pouvons supposer que $U$ est affine,
et donc que $Z_1$ est affine lui aussi.

\medskip\noindent
Puisque $f$ est séparé, l'image $V$ de $Z_1 \to X$ est à la fois ouverte et fermée
(Morphismes d'espaces, Lemme
\ref{spaces-morphisms-lemma-universally-closed-permanence}).
Posons $W = X \setminus V$ pour obtenir une décomposition comme dans le lemme.
Pour achever la démonstration, il reste à montrer que $V \to U$ est fini.
Comme $Z_1 \to V$ est étale surjectif, $V$ est le quotient de
$Z_1$ par la relation d'équivalence étale $R = Z_1 \times_V Z_1$ ; voir
Espaces, Lemme \ref{spaces-lemma-space-presentation}.
Puisque $f$ est séparé, $V \to U$ est séparé et $R$ est fermé dans
$Z_1 \times_U Z_1$. Puisque $Z_1 \to U$ est fini,
les projections $s, t : R \to Z_1$ sont finies.
Ainsi, $V$ est un schéma affine d'après
Groupoïdes, Proposition \ref{groupoids-proposition-finite-flat-equivalence}.
D'après Morphismes, Lemme \ref{morphisms-lemma-image-proper-is-proper},
nous concluons que $V \to U$ est propre et, d'après
Morphismes, Lemme \ref{morphisms-lemma-finite-proper},
nous concluons que $V \to U$ est fini,
ce qui achève la démonstration.
\end{proof}

\begin{lemma}
\label{lemma-etale-splits-off-just-one-quasi-finite-part}
Soit $S$ un schéma.
Soit $f : X \to Y$ un morphisme d'espaces algébriques sur $S$.
Soit $x \in |X|$ d'image $y \in |Y|$. Supposons que
\begin{enumerate}
\item $f$ soit localement de type fini ;
\item $f$ soit séparé ;
\item $f$ soit quasi-fini en $x$.
\end{enumerate}
Il existe alors un morphisme étale $(U, u) \to (Y, y)$
d'espaces algébriques pointés et une décomposition
$$
U \times_Y X = W \amalg V
$$
en sous-espaces ouverts et fermés telle que le morphisme $V \to U$ soit fini
et qu'il existe un point $v \in |V|$ qui s'envoie sur $x$ dans $|X|$
et sur $u$ dans $|U|$.
\end{lemma}

\begin{proof}
Choisissons un schéma $U$, un point $u \in U$ et un morphisme étale
$U \to Y$ qui envoie $u$ sur $y$. Il existe un point $x' \in |U \times_Y X|$
qui s'envoie sur $x$ dans $|X|$ et sur $u$ dans $|U|$ (Propriétés des espaces,
Lemme \ref{spaces-properties-lemma-points-cartesian}).
Pour terminer, appliquons le Lemme \ref{lemma-etale-splits-off-quasi-finite-part}
au morphisme $U \times_Y X \to U$ et au point $x'$.
Il s'applique parce que $U$ est un schéma et que $u$ est donc représenté
par le monomorphisme $\Spec(\kappa(u)) \to U$.
\end{proof}










\section{Théorème principal de Zariski}
\label{section-structure-quasi-finite}

\noindent
Dans cette section, nous appliquons les résultats de la section précédente pour
démontrer le théorème principal de Zariski pour les morphismes d'espaces algébriques.
Cette section est l'analogue de la section
\ref{more-morphisms-section-application-etale-neighbourhoods} de Compléments sur les morphismes.

\begin{lemma}
\label{lemma-finite-type-separated}
Soit $S$ un schéma. Soit $f : X \to Y$ un morphisme d'espaces
algébriques sur $S$, de type fini et séparé.
Soit $Y'$ la normalisation de $Y$ dans $X$. Diagramme :
$$
\xymatrix{
X \ar[rd]_f \ar[rr]_{f'} & & Y' \ar[ld]^\nu \\
& Y &
}
$$
Il existe alors un sous-espace ouvert $U' \subset Y'$ tel que
\begin{enumerate}
\item $(f')^{-1}(U') \to U'$ soit un isomorphisme ;
\item $(f')^{-1}(U') \subset X$ soit l'ensemble des points où
$f$ est quasi-fini.
\end{enumerate}
\end{lemma}

\begin{proof}
D'après Morphismes d'espaces, Lemme
\ref{spaces-morphisms-lemma-locally-finite-type-quasi-finite-part}
il existe un sous-espace ouvert $U \subset X$ correspondant aux points
de $|X|$ où $f$ est quasi-fini. Nous devons démontrer que
\begin{enumerate}
\item[(a)] l'image de $|U| \to |Y'|$ est $|U'|$ pour un certain sous-espace ouvert
$U'$ de $Y'$ ;
\item[(b)] $U = f^{-1}(U')$ ;
\item[(c)] $U \to U'$ est un isomorphisme.
\end{enumerate}
Puisque la formation de $U$ commute à tout changement de base
(Morphismes d'espaces, Lemme
\ref{spaces-morphisms-lemma-locally-finite-type-quasi-finite-part}),
puisque la formation de la normalisation $Y'$ commute au changement de base
lisse (Lemme \ref{lemma-normalization-smooth-localization}),
puisque les morphismes étales sont ouverts, et
puisque la propriété « être un isomorphisme » est locale pour la topologie fpqc sur la base
(Descente sur les espaces, Lemme
\ref{spaces-descent-lemma-descending-property-isomorphism}),
il suffit de démontrer (a), (b), (c) localement pour la topologie étale sur $Y$
(nous omettons certains détails). Ainsi,
nous pouvons supposer que $Y$ est un schéma affine. Cela implique que $Y'$ est
lui aussi un schéma (affine).

\medskip\noindent
Soit $x \in |U|$. Affirmation : il existe un voisinage
ouvert $f'(x) \in V \subset Y'$ tel que $(f')^{-1}V \to V$ soit un
isomorphisme. Montrons d'abord que l'affirmation implique le lemme.
En effet, alors $(f')^{-1}V \cong V$ est un schéma (comme ouvert
de $Y'$), localement de type fini sur $Y$ (comme sous-espace ouvert de $X$),
et, pour $v \in V$, l'extension de corps résiduels
$\kappa(v)/\kappa(\nu(v))$ est algébrique (car
$V \subset Y'$ et $Y'$ est entier sur $Y$). Par conséquent, les fibres
de $V \to Y$ sont discrètes (Morphismes, Lemme
\ref{morphisms-lemma-algebraic-residue-field-extension-closed-point-fibre})
et $(f')^{-1}V \to Y$ est localement quasi-fini
(Morphismes, Lemme \ref{morphisms-lemma-locally-quasi-finite-fibres}).
Cela implique $(f')^{-1}V \subset U$ et $V \subset U'$. Comme $x$ était
arbitraire, nous voyons que (a), (b) et (c) sont vraies.

\medskip\noindent
Soit $y = f(x) \in |Y|$. Soit $(T, t) \to (Y, y)$ un morphisme étale
de schémas pointés. Notons par un indice inférieur ${}_T$ le changement de base à $T$.
Soit $z \in X_T$ un point de la fibre $X_t$ situé au-dessus de $x$.
Remarquons que $U_T \subset X_T$ est l'ensemble des points où $f_T$ est
quasi-fini ; voir Morphismes d'espaces, Lemme
\ref{spaces-morphisms-lemma-locally-finite-type-quasi-finite-part}.
Remarquons que
$$
X_T \xrightarrow{f'_T} Y'_T \xrightarrow{\nu_T} T
$$
est la normalisation de $T$ dans $X_T$ ; voir
Lemme \ref{lemma-normalization-smooth-localization}.
Supposons l'affirmation vraie pour $z \in U_T \subset X_T \to Y'_T \to T$, c'est-à-dire
supposons que nous puissions trouver un voisinage ouvert
$f'_T(z) \in V' \subset Y'_T$ tel que $(f'_T)^{-1}V' \to V'$ soit un
isomorphisme. Le morphisme $Y'_T \to Y'$ est étale ; par conséquent, l'image
$V \subset Y'$ de $V'$ est ouverte. Remarquons que $f'(x) \in V$ puisque $f'_T(z) \in V'$.
Remarquons que
$$
\xymatrix{
(f'_T)^{-1}V' \ar[r] \ar[d] & (f')^{-1}(V) \ar[d] \\
V' \ar[r] & V
}
$$
est un carré cartésien (car $Y'_T \times_{Y'} X = X_T$).
Puisque la flèche verticale gauche est un isomorphisme
et que $\{V' \to V\}$ est un recouvrement étale, nous concluons que la flèche verticale
droite est un isomorphisme d'après
Descente sur les espaces, Lemme
\ref{spaces-descent-lemma-descending-property-isomorphism}.
Autrement dit, l'affirmation est vraie pour $x \in U \subset X \to Y' \to Y$.

\medskip\noindent
D'après le paragraphe précédent, pour démontrer l'affirmation pour
$x \in |U|$, nous pouvons remplacer $Y$ par un voisinage étale $T$ de
$y = f(x)$ et $x$ par n'importe quel point situé au-dessus de $x$ dans $T \times_Y X$.
Ainsi, nous pouvons supposer qu'il existe une décomposition
$$
X = V \amalg W
$$
en sous-espaces ouverts et fermés, où $V \to Y$ est fini et $x \in V$ ;
voir le Lemme \ref{lemma-etale-splits-off-quasi-finite-part}.
Puisque $X$ est la somme disjointe de $V$ et $W$ sur $Y$ et que
$V \to Y$ est fini, nous voyons que la
normalisation de $Y$ dans $X$ est le morphisme
$$
X = V \amalg W \longrightarrow V \amalg W' \longrightarrow S
$$
où $W'$ est la normalisation de $Y$ dans $W$ ; voir
Morphismes d'espaces, Lemmes
\ref{spaces-morphisms-lemma-normalization-in-disjoint-union},
\ref{spaces-morphisms-lemma-finite-integral}, et
\ref{spaces-morphisms-lemma-normalization-in-integral}.
L'affirmation en résulte, et la démonstration est terminée.
\end{proof}

\noindent
Le lemme suivant est un doublon de
Morphismes d'espaces, Lemme
\ref{spaces-morphisms-lemma-quasi-finite-separated-quasi-affine}.
La présence de deux exemplaires du même lemme s'explique par le fait que les
démonstrations diffèrent quelque peu. La démonstration donnée ci-dessous
repose sur le théorème principal de Zariski pour les morphismes non représentables
d'espaces algébriques sous la forme établie ci-dessus, tandis que la démonstration de
Morphismes d'espaces, Lemme
\ref{spaces-morphisms-lemma-quasi-finite-separated-quasi-affine}
repose sur Morphismes d'espaces, Proposition
\ref{spaces-morphisms-proposition-locally-quasi-finite-separated-over-scheme}
pour se ramener au cas des morphismes de schémas.

\begin{lemma}
\label{lemma-quasi-finite-separated-quasi-affine}
Soit $S$ un schéma.
Soit $f : X \to Y$ un morphisme d'espaces algébriques sur $S$.
Supposons que $f$ soit quasi-fini et séparé.
Soit $Y'$ la normalisation de $Y$ dans $X$.
Diagramme :
$$
\xymatrix{
X \ar[rd]_f \ar[rr]_{f'} & & Y' \ar[ld]^\nu \\
& Y &
}
$$
Alors $f'$ est une immersion ouverte quasi-compacte et $\nu$ est entier.
En particulier, $f$ est quasi-affine.
\end{lemma}

\begin{proof}
Cela résulte du Lemme \ref{lemma-finite-type-separated}. En effet, d'après
ce lemme, il existe un sous-espace ouvert $U' \subset Y'$ tel que
$(f')^{-1}(U') = X$ (!) et que $X \to U'$ soit un isomorphisme ! Autrement
dit, $f'$ est une immersion ouverte. Remarquons que $f'$ est quasi-compact puisque
$f$ est quasi-compact et que $\nu : Y' \to Y$ est séparé
(Morphismes d'espaces, Lemme
\ref{spaces-morphisms-lemma-quasi-compact-permanence}).
Ainsi, pour tout schéma affine $Z$ et tout morphisme $Z \to Y$, le
produit fibré $Z \times_Y X$ est un sous-schéma ouvert quasi-compact
du schéma affine $Z \times_Y Y'$. Par conséquent, $f$ est quasi-affine par
définition.
\end{proof}

\begin{lemma}[Théorème principal de Zariski]
\label{lemma-quasi-finite-separated-pass-through-finite}
Soit $S$ un schéma.
Soit $f : X \to Y$ un morphisme d'espaces algébriques sur $S$.
Supposons que $f$ soit quasi-fini et séparé, et que
$Y$ soit quasi-compact et quasi-séparé. Il existe alors
une factorisation
$$
\xymatrix{
X \ar[rd]_f \ar[rr]_j & & T \ar[ld]^\pi \\
& Y &
}
$$
où $j$ est une immersion ouverte quasi-compacte et $\pi$ est fini.
\end{lemma}

\begin{proof}
Soit $X \to Y' \to Y$ comme dans la conclusion du
Lemme \ref{lemma-quasi-finite-separated-quasi-affine}.
D'après
Limites d'espaces, Lemme
\ref{spaces-limits-lemma-integral-algebra-directed-colimit-finite}
nous pouvons écrire
$\nu_*\mathcal{O}_{Y'} = \colim_{i \in I} \mathcal{A}_i$ comme une
limite inductive filtrante d'algèbres finies quasi-cohérentes sur $\mathcal{O}_X$,
$\mathcal{A}_i \subset \nu_*\mathcal{O}_{Y'}$. Alors
$\pi_i : T_i = \underline{\Spec}_Y(\mathcal{A}_i) \to Y$
est un morphisme fini pour tout $i$.
Remarquons que les morphismes de transition $T_{i'} \to T_i$ sont affines
et que $Y' = \lim T_i$.

\medskip\noindent
D'après Limites d'espaces, Lemme \ref{spaces-limits-lemma-descend-opens},
il existe un $i$ et un ouvert quasi-compact
$U_i \subset T_i$ dont l'image inverse dans $Y'$ est
$f'(X)$. Pour $i' \geq i$, notons $U_{i'}$ l'image inverse
de $U_i$ dans $T_{i'}$. Alors $X \cong f'(X) = \lim_{i' \geq i} U_{i'}$ ; voir
Limites d'espaces, Lemme
\ref{spaces-limits-lemma-directed-inverse-system-has-limit}.
D'après
Limites d'espaces, Lemme
\ref{spaces-limits-lemma-finite-type-eventually-closed}, nous voyons que
$X \to U_{i'}$ est une immersion fermée pour un certain $i' \geq i$.
(En fait, $X \cong U_{i'}$ pour $i'$ suffisamment
grand, mais nous n'en avons pas besoin.) Ainsi, $X \to T_{i'}$ est une immersion. D'après
Morphismes d'espaces, Lemme
\ref{spaces-morphisms-lemma-factor-the-other-way}
nous pouvons factoriser ce morphisme en $X \to T \to T_{i'}$, où la première flèche
est une immersion ouverte et la seconde une immersion fermée. Cela termine la démonstration.
\end{proof}

\begin{lemma}
\label{lemma-quasi-finite-separated-pass-through-finite-addendum}
Avec les notations et les hypothèses du
Lemme \ref{lemma-quasi-finite-separated-pass-through-finite},
supposons de plus que $f$ soit localement de présentation finie. Nous pouvons alors
choisir la factorisation de telle sorte que $T$ soit fini et de
présentation finie sur $Y$.
\end{lemma}

\begin{proof}
D'après Limites d'espaces, Lemme
\ref{spaces-limits-lemma-finite-in-finite-and-finite-presentation}, nous pouvons écrire
$T = \lim T_i$, où tous les $T_i$ sont finis et de présentation finie
sur $Y$ et où les morphismes de transition $T_{i'} \to T_i$ sont des
immersions fermées. D'après
Limites d'espaces, Lemme \ref{spaces-limits-lemma-descend-opens},
il existe un $i$ et un sous-schéma ouvert $U_i \subset T_i$ dont l'image
inverse dans $T$ est $X$. D'après
Limites d'espaces, Lemme
\ref{spaces-limits-lemma-finite-type-eventually-closed}
nous voyons que $X \cong U_i$ pour $i$ assez grand.
Remplacer $T$ par $T_i$ achève la démonstration.
\end{proof}












\section{Applications du théorème principal de Zariski, I}
\label{section-applications-zmt}

\noindent
Une première application est la caractérisation des morphismes
finis comme morphismes propres à fibres finies.

\begin{lemma}
\label{lemma-characterize-finite}
Soit $S$ un schéma. Soit $f : X \to Y$ un morphisme d'espaces algébriques
sur $S$. Les propriétés suivantes sont équivalentes :
\begin{enumerate}
\item $f$ est fini ;
\item $f$ est propre et localement quasi-fini ;
\item $f$ est propre et $|X_k|$ est un espace discret pour tout morphisme
$\Spec(k) \to Y$, où $k$ est un corps ;
\item $f$ est universellement fermé, séparé, localement de type fini,
et $|X_k|$ est un espace discret pour tout morphisme $\Spec(k) \to Y$,
où $k$ est un corps.
\end{enumerate}
\end{lemma}

\begin{proof}
Nous avons (1) $\Rightarrow$ (2) d'après
Morphismes d'espaces, Lemmes \ref{spaces-morphisms-lemma-finite-proper},
\ref{spaces-morphisms-lemma-finite-quasi-finite}.
Nous avons (2) $\Rightarrow$ (3) d'après
Morphismes d'espaces, Lemme
\ref{spaces-morphisms-lemma-locally-quasi-finite}.
Par définition, (3) implique (4).

\medskip\noindent
Supposons (4). Puisque $f$ est universellement fermé, il est quasi-compact
(Morphismes d'espaces, Lemme
\ref{spaces-morphisms-lemma-universally-closed-quasi-compact}).
Choisissons un point $y$ de $|Y|$. Représentons $y$ par un
morphisme $\Spec(k) \to Y$. Remarquons que $|X_k|$ est fini et discret,
car il est quasi-compact et discret. L'application $|X_k| \to |X|$ est surjective
sur la fibre de $|X| \to |Y|$ au-dessus de $y$
(Propriétés des espaces, Lemme \ref{spaces-properties-lemma-points-cartesian}).
D'après
Morphismes d'espaces, Lemme \ref{spaces-morphisms-lemma-quasi-finite-at-point},
nous voyons que $X \to Y$ est quasi-fini en tous les points de la fibre
de $|X| \to |Y|$ au-dessus de $y$.
Choisissons un voisinage étale élémentaire $(U, u) \to (Y, y)$
et une décomposition $X_U = V \amalg W$ comme dans le
Lemme \ref{lemma-etale-splits-off-quasi-finite-part},
adaptée à tous les points de $|X|$ situés au-dessus de $y$.
Remarquons que $W_u = \emptyset$, puisque nous avons utilisé tous les points
de la fibre de $|X| \to |Y|$ au-dessus de $y$.
Puisque $f$ est universellement fermé, nous voyons que
l'image de $|W|$ dans $|U|$ est un fermé qui ne contient pas $u$.
Après avoir rétréci $U$, nous pouvons supposer que $W = \emptyset$.
Autrement dit, nous voyons que $X_U = V$ est fini sur $U$.
Comme $y \in |Y|$ était arbitraire,
il existe donc une famille $\{U_i \to Y\}$
de morphismes étales dont les images recouvrent $Y$ et telle que
les changements de base $X_{U_i} \to U_i$ soient finis.
Nous concluons que $f$ est fini d'après
Morphismes d'espaces, Lemme \ref{spaces-morphisms-lemma-integral-local}.
\end{proof}

\noindent
Nous en déduisons le résultat utile suivant.

\begin{lemma}
\label{lemma-proper-finite-fibre-finite-in-neighbourhood}
Soit $S$ un schéma. Soit $f : X \to Y$ un morphisme d'espaces algébriques
sur $S$. Soit $y \in |Y|$. Supposons que
\begin{enumerate}
\item $f$ soit propre ;
\item $f$ soit quasi-fini en tout $x \in |X|$ situé au-dessus de $y$
(Espaces décents, Lemme \ref{decent-spaces-lemma-conditions-on-fibre-and-qf}).
\end{enumerate}
Il existe alors un voisinage ouvert $V \subset Y$ de $y$
tel que $f|_{f^{-1}(V)} : f^{-1}(V) \to V$ soit fini.
\end{lemma}

\begin{proof}
D'après
Morphismes d'espaces, Lemme
\ref{spaces-morphisms-lemma-locally-finite-type-quasi-finite-part},
l'ensemble des points où $f$ est quasi-fini est un ouvert $U \subset X$.
Soit $Z = X \setminus U$. Alors $y \not \in f(Z)$. Puisque $f$ est propre,
l'ensemble $f(Z) \subset Y$ est fermé. Choisissons un voisinage ouvert quelconque
$V \subset Y$ de $y$ tel que $Z \cap V = \emptyset$. Alors
$f^{-1}(V) \to V$ est localement quasi-fini et propre.
Par conséquent, $f^{-1}(V) \to V$ est fini d'après le Lemme \ref{lemma-characterize-finite}.
\end{proof}

\begin{lemma}
\label{lemma-flat-proper-family-cannot-collapse-fibre}
\begin{slogan}
La contraction d'une fibre d'une famille propre force également la contraction des fibres voisines.
\end{slogan}
Soit $S$ un schéma. Soit
$$
\xymatrix{
X \ar[rr]_h \ar[rd]_f & & Y \ar[ld]^g \\
& B
}
$$
un diagramme commutatif de morphismes d'espaces algébriques sur $S$.
Soit $b \in B$ et soit $\Spec(k) \to B$ un morphisme appartenant à la classe
d'équivalence de $b$. Supposons que
\begin{enumerate}
\item $X \to B$ soit un morphisme propre ;
\item $Y \to B$ soit séparé et localement de type fini ;
\item l'une des propriétés suivantes soit vraie :
\begin{enumerate}
\item l'image de $|X_k| \to |Y_k|$ est finie ;
\item l'image de $|f|^{-1}(\{b\})$ dans $|Y|$ est finie
et $B$ est décent.
\end{enumerate}
\end{enumerate}
Il existe alors un
sous-espace ouvert $B' \subset B$ contenant $b$ tel que $X_{B'} \to Y_{B'}$
se factorise par un sous-espace fermé $Z \subset Y_{B'}$ fini sur $B'$.
\end{lemma}

\begin{proof}
Soit $Z \subset Y$ l'image schématique de $h$ ; voir
Morphismes d'espaces, section
\ref{spaces-morphisms-section-scheme-theoretic-image}.
D'après Morphismes d'espaces, Lemme
\ref{spaces-morphisms-lemma-scheme-theoretic-image-is-proper}
le morphisme $X \to Z$ est surjectif et $Z \to B$ est propre.
Ainsi,
$$
\{x \in |X|\text{ situé au-dessus de }b\} \to
\{z \in |Z|\text{ situé au-dessus de }b\}
$$
et $|X_k| \to |Z_k|$ sont surjectifs. Nous voyons que l'une ou l'autre
des conditions (3)(a) et (3)(b) implique que $Z \to B$ est quasi-fini
en tous les points de $|Z|$ situés au-dessus de $b$, d'après
Espaces décents, Lemme \ref{decent-spaces-lemma-conditions-on-fibre-and-qf}.
Ainsi, $Z \to B$ est fini dans un voisinage ouvert de $b$ d'après le
Lemme \ref{lemma-proper-finite-fibre-finite-in-neighbourhood}.
\end{proof}







\section{Factorisation de Stein}
\label{section-stein-factorization}

\noindent
La factorisation de Stein affirme qu'un morphisme propre $f : X \to S$
tel que $f_*\mathcal{O}_X = \mathcal{O}_S$ a des fibres connexes.

\begin{lemma}
\label{lemma-stein-universally-closed}
Soit $S$ un schéma. Soit $f : X \to Y$ un morphisme universellement fermé et
quasi-séparé d'espaces algébriques sur $S$.
Il existe une factorisation
$$
\xymatrix{
X \ar[rr]_{f'} \ar[rd]_f & & Y' \ar[dl]^\pi \\
& Y &
}
$$
ayant les propriétés suivantes :
\begin{enumerate}
\item le morphisme $f'$ est universellement fermé, quasi-compact, quasi-séparé
et surjectif ;
\item le morphisme $\pi : Y' \to Y$ est entier ;
\item nous avons $f'_*\mathcal{O}_X = \mathcal{O}_{Y'}$ ;
\item nous avons $Y' = \underline{\Spec}_Y(f_*\mathcal{O}_X)$ ;
\item $Y'$ est la normalisation de $Y$ dans $X$ au sens de
Morphismes d'espaces, Définition
\ref{spaces-morphisms-definition-normalization-X-in-Y}.
\end{enumerate}
La formation de la factorisation $f = \pi \circ f'$ commute au changement
de base plat.
\end{lemma}

\begin{proof}
D'après Morphismes d'espaces, Lemme
\ref{spaces-morphisms-lemma-universally-closed-quasi-compact}
le morphisme $f$ est quasi-compact.
Nous définissons simplement $Y'$ comme la normalisation de $Y$ dans $X$ ; (5) et (2) sont donc
automatiquement vérifiés. D'après
Morphismes d'espaces, Lemme
\ref{spaces-morphisms-lemma-normalization-in-universally-closed}
nous voyons que (4) est vérifié. Le morphisme $f'$ est universellement fermé d'après
Morphismes d'espaces, Lemme
\ref{spaces-morphisms-lemma-universally-closed-permanence}.
Il est quasi-compact d'après
Morphismes d'espaces, Lemme
\ref{spaces-morphisms-lemma-quasi-compact-permanence}
et quasi-séparé d'après
Morphismes d'espaces, Lemme
\ref{spaces-morphisms-lemma-compose-after-separated}.

\medskip\noindent
Pour démontrer les assertions restantes, nous pouvons supposer que la base $Y$ est affine
(car la normalisation commute à la localisation étale).
Écrivons $Y = \Spec(R)$. Alors $Y' = \Spec(A)$, où
$A = \Gamma(X, \mathcal{O}_X)$ est une $R$-algèbre entière.
Il est donc clair que $f'_*\mathcal{O}_X$
est $\mathcal{O}_{Y'}$ (car $f'_*\mathcal{O}_X$ est quasi-cohérent
d'après Morphismes d'espaces, Lemme
\ref{spaces-morphisms-lemma-pushforward},
et est donc égal à $\widetilde{A}$). Cela démontre (3).

\medskip\noindent
Montrons que $f'$ est surjectif. Comme $f'$ est universellement fermé (voir ci-dessus),
l'image de $f'$ est un fermé
$V(I) \subset Y' = \Spec(A)$. Choisissons $h \in I$. Alors
$h|_X = f^\sharp(h)$ est une section globale du faisceau structural de
$X$ qui s'annule en tout point. Puisque $X$ est quasi-compact, cela signifie
que $h|_X$ est une section nilpotente, c'est-à-dire que $h^n|X = 0$ pour un certain $n > 0$.
Mais $A = \Gamma(X, \mathcal{O}_X)$, donc $h^n = 0$.
Autrement dit, $I$ est contenu dans le radical de Jacobson de $A$, et nous concluons
que $V(I) = Y'$, comme souhaité.
\end{proof}

\begin{lemma}
\label{lemma-stein-universally-closed-residue-fields}
Dans le Lemme \ref{lemma-stein-universally-closed}, supposons de plus que
$f$ soit localement de type fini et que $Y$ soit affine. Alors, pour $y \in Y$, la fibre
$\pi^{-1}(\{y\}) = \{y_1, \ldots, y_n\}$ est finie et les extensions de corps
$\kappa(y_i)/\kappa(y)$ sont finies.
\end{lemma}

\begin{proof}
Rappelons qu'il n'existe pas de spécialisations entre les points de $\pi^{-1}(\{y\})$ ;
voir Algèbre, Lemme \ref{algebra-lemma-integral-no-inclusion}.
Puisque $f'$ est surjectif, nous obtenons que $|X_y| \to \pi^{-1}(\{y\})$ est surjectif.
Remarquons que $X_y$ est un espace algébrique quasi-séparé de type fini
sur un corps (la quasi-compacité a été démontrée dans la preuve du
lemme cité). Ainsi, $|X_y|$ est un espace topologique noethérien
(Morphismes d'espaces, Lemme
\ref{spaces-morphisms-lemma-finite-presentation-noetherian}).
Un argument topologique (omis) montre alors que $\pi^{-1}(\{y\})$ est fini.
Pour chaque $i$, nous pouvons choisir un point de type fini $x_i \in |X_y|$
qui s'envoie sur $y_i$ (Morphismes d'espaces, Lemme
\ref{spaces-morphisms-lemma-enough-finite-type-points}).
Nous concluons que $\kappa(y_i)/\kappa(y)$ est finie :
$x_i$ peut être représenté par un morphisme $\Spec(k_i) \to X_y$
de type fini (d'après notre définition des points de type fini) ;
par conséquent, $\Spec(k_i) \to y = \Spec(\kappa(y))$ est de type fini
(comme composée de morphismes de type fini),
donc $k_i/\kappa(y)$ est finie (Morphismes, Lemme
\ref{morphisms-lemma-point-finite-type}).
\end{proof}

\noindent
Soit $f : X \to Y$ un morphisme d'espaces algébriques et soit
$\overline{y} : \Spec(k) \to Y$ un point géométrique. La
{\it fibre de $f$ au-dessus de $\overline{y}$} est alors l'espace algébrique
$X_{\overline{y}} = X \times_{Y, \overline{y}} \Spec(k)$ sur $k$.
Si $Y$ est un schéma et si $y \in Y$ est un point, nous notons comme d'habitude
$X_y = X \times_Y \Spec(\kappa(y))$ la fibre.

\begin{lemma}
\label{lemma-characterize-geometrically-connected-fibres}
Soit $S$ un schéma. Soit $f : X \to Y$ un morphisme d'espaces algébriques
sur $S$. Soit $\overline{y}$ un point géométrique de $Y$. Alors
$X_{\overline{y}}$ est connexe si et seulement si, pour tout voisinage
étale $(V, \overline{v}) \to (Y, \overline{y})$ tel que $V$
soit un schéma, le changement de base $X_V \to V$ a une fibre connexe $X_v$.
\end{lemma}

\begin{proof}
Puisque la catégorie des voisinages étales de $\overline{y}$ est
cofiltrante et contient une famille cofinale de schémas
(Propriétés des espaces, Lemme \ref{spaces-properties-lemma-cofinal-etale}),
nous pouvons remplacer $Y$ par l'un de ces voisinages et supposer que
$Y$ est un schéma. Soit $y \in Y$ le point correspondant à $\overline{y}$.
Alors $X_y$ est géométriquement connexe sur $\kappa(y)$ si et seulement si
$X_{\overline{y}}$ est connexe, et si et seulement si $(X_y)_{k'}$
est connexe pour toute extension séparable finie $k'$ de $\kappa(y)$.
Voir Espaces sur les corps, section
\ref{spaces-over-fields-section-geometrically-connected}, et en particulier le
Lemme \ref{spaces-over-fields-lemma-characterize-geometrically-disconnected}.
D'après Compléments sur les morphismes, Lemme
\ref{more-morphisms-lemma-realize-prescribed-residue-field-extension-etale}
il existe un voisinage étale affine $(V, v) \to (Y, y)$ tel que
$\kappa(s) \subset \kappa(u)$ s'identifie à $\kappa(s) \subset k'$
pour toute extension séparable finie donnée. Le lemme en résulte.
\end{proof}

\begin{theorem}[Factorisation de Stein ; cas noethérien]
\label{theorem-stein-factorization-Noetherian}
Soit $S$ un schéma. Soit $f : X \to Y$ un morphisme propre d'espaces
algébriques sur $S$, avec $Y$ localement noethérien.
Il existe une factorisation
$$
\xymatrix{
X \ar[rr]_{f'} \ar[rd]_f & & Y' \ar[dl]^\pi \\
& Y &
}
$$
ayant les propriétés suivantes :
\begin{enumerate}
\item le morphisme $f'$ est propre à fibres géométriques connexes ;
\item le morphisme $\pi : Y' \to Y$ est fini ;
\item nous avons $f'_*\mathcal{O}_X = \mathcal{O}_{Y'}$ ;
\item nous avons $Y' = \underline{\Spec}_Y(f_*\mathcal{O}_X)$ ;
\item $Y'$ est la normalisation de $Y$ dans $X$ ; voir
Morphismes, Définition \ref{morphisms-definition-normalization-X-in-Y}.
\end{enumerate}
\end{theorem}

\begin{proof}
Soit $f = \pi \circ f'$ la factorisation du
Lemme \ref{lemma-stein-universally-closed}. Remarquons qu'outre les
conclusions du Lemme \ref{lemma-stein-universally-closed}, nous
savons aussi que $f'$ est séparé
(Morphismes d'espaces, Lemme
\ref{spaces-morphisms-lemma-compose-after-separated})
et de type fini
(Morphismes d'espaces, Lemme
\ref{spaces-morphisms-lemma-permanence-finite-type}).
Ainsi, $f'$ est propre. D'après
Cohomologie des espaces, Lemme
\ref{spaces-cohomology-lemma-proper-pushforward-coherent}
nous voyons que $f_*\mathcal{O}_X$ est un $\mathcal{O}_Y$-module cohérent.
Nous voyons donc que $\pi$ est fini, c'est-à-dire que (2) est vérifié.

\medskip\noindent
Cela démontre toutes les assertions sauf la plus intéressante, à savoir que
les fibres géométriques de $f'$ sont connexes. Il ressort clairement de la
discussion précédente que nous pouvons remplacer $Y$ par $Y'$.
Alors $Y$ est localement noethérien,
$f : X \to Y$ est propre et $f_*\mathcal{O}_X = \mathcal{O}_Y$.
Soit $\overline{y}$ un point géométrique de $Y$.
Nous appliquons ici le théorème des fonctions formelles,
plus précisément Cohomologie des espaces, Lemme
\ref{spaces-cohomology-lemma-formal-functions-stalk}.
Il nous dit que
$$
\mathcal{O}^\wedge_{Y, \overline{y}} =
\lim_n H^0(X_n, \mathcal{O}_{X_n})
$$
où $X_n =
\Spec(\mathcal{O}_{Y, \overline{y}}/\mathfrak m_{\overline{y}}^n) \times_Y X$.
Remarquons que $X_1 = X_{\overline{y}} \to X_n$ est un épaississement (d'ordre fini)
et que l'espace topologique sous-jacent à $X_n$ est donc égal à celui
de $X_{\overline{y}}$. Ainsi, si $X_{\overline{y}} = T_1 \amalg T_2$
est une somme disjointe de sous-espaces ouverts et fermés non vides, alors, de même,
$X_n = T_{1, n} \amalg T_{2, n}$ pour tout $n$. Cela signifie à son tour que
$H^0(X_n, \mathcal{O}_{X_n})$ contient un idempotent $e_{1, n}$ distinct de 0 et de 1,
à savoir la fonction identiquement égale à $1$ sur $T_{1, n}$ et
identiquement égale à $0$ sur $T_{2, n}$. Il est clair que $e_{1, n + 1}$
se restreint à $e_{1, n}$ sur $X_n$. Par conséquent, $e_1 = \lim e_{1, n}$
est un idempotent de la limite distinct de 0 et de 1. Cela contredit le fait
que $\mathcal{O}^\wedge_{Y, \overline{y}}$ est un anneau local. Notre
hypothèse était donc fausse, c'est-à-dire que $X_{\overline{y}}$ est connexe,
comme souhaité.
\end{proof}

\begin{theorem}[Factorisation de Stein ; cas général]
\label{theorem-stein-factorization-general}
Soit $S$ un schéma. Soit $f : X \to Y$ un morphisme propre d'espaces
algébriques sur $S$. Il existe une factorisation
$$
\xymatrix{
X \ar[rr]_{f'} \ar[rd]_f & & Y' \ar[dl]^\pi \\
& Y &
}
$$
ayant les propriétés suivantes :
\begin{enumerate}
\item le morphisme $f'$ est propre à fibres géométriques connexes ;
\item le morphisme $\pi : Y' \to Y$ est entier ;
\item nous avons $f'_*\mathcal{O}_X = \mathcal{O}_{Y'}$ ;
\item nous avons $Y' = \underline{\Spec}_Y(f_*\mathcal{O}_X)$ ;
\item $Y'$ est la normalisation de $Y$ dans $X$ (Morphismes d'espaces, Définition
\ref{spaces-morphisms-definition-normalization-X-in-Y}).
\end{enumerate}
\end{theorem}

\begin{proof}
Nous pouvons appliquer le Lemme \ref{lemma-stein-universally-closed} pour obtenir le
morphisme $f' : X \to Y'$.
Remarquons qu'outre les
conclusions du Lemme \ref{lemma-stein-universally-closed}, nous
savons aussi que $f'$ est séparé
(Morphismes d'espaces, Lemme
\ref{spaces-morphisms-lemma-compose-after-separated})
et de type fini
(Morphismes d'espaces, Lemme
\ref{spaces-morphisms-lemma-permanence-finite-type}).
Ainsi, $f'$ est propre. Nous avons alors démontré toutes les
assertions, sauf celle
selon laquelle $f'$ a des fibres géométriques connexes.

\medskip\noindent
Il ressort clairement de la discussion que nous pouvons remplacer $Y$ par $Y'$.
Alors $f : X \to Y$ est propre et $f_*\mathcal{O}_X = \mathcal{O}_Y$.
Remarquons que ces conditions sont préservées par changement de base plat
(Morphismes d'espaces, Lemme \ref{spaces-morphisms-lemma-base-change-proper}
et
Cohomologie des espaces, Lemme
\ref{spaces-cohomology-lemma-flat-base-change-cohomology}).
Soit $\overline{y}$ un point géométrique de $Y$. D'après le
Lemme \ref{lemma-characterize-geometrically-connected-fibres}
et la remarque précédente, nous nous ramenons au cas où $Y$ est un
schéma, $y \in Y$ est un point, $f : X \to Y$ est un espace algébrique propre
sur $Y$ avec $f_*\mathcal{O}_X = \mathcal{O}_Y$,
et où nous devons montrer que la fibre $X_y$ est connexe.
En remplaçant $Y$ par un voisinage affine de $y$, nous pouvons
supposer que $Y = \Spec(R)$ est affine. Alors $f_*\mathcal{O}_X = \mathcal{O}_Y$
signifie que l'homomorphisme d'anneaux
$R \to \Gamma(X, \mathcal{O}_X)$ est bijectif.

\medskip\noindent
D'après Limites d'espaces, Lemme
\ref{spaces-limits-lemma-proper-limit-of-proper-finite-presentation-noetherian}
nous pouvons écrire $(X \to Y) = \lim (X_i \to Y_i)$, où $X_i \to Y_i$
est propre et de présentation finie et où $Y_i$ est noethérien. Pour $i$ assez
grand, $Y_i$ est affine (Limites d'espaces, Lemme
\ref{spaces-limits-lemma-limit-is-affine}).
Écrivons $Y_i = \Spec(R_i)$. Posons $R'_i = \Gamma(X_i, \mathcal{O}_{X_i})$.
Remarquons que nous avons des homomorphismes d'anneaux $R_i \to R_i' \to R$. En effet, nous avons
le premier parce que $X_i$ est un espace algébrique sur $R_i$, et le second parce que
nous avons $X \to X_i$ et $R = \Gamma(X, \mathcal{O}_X)$. Remarquons que
$R = \colim R'_i$ d'après Limites d'espaces, Lemme
\ref{spaces-limits-lemma-descend-section}.
Alors
$$
\xymatrix{
X \ar[d]  \ar[r] & X_i \ar[d] \\
Y \ar[r] & Y'_i \ar[r] & Y_i
}
$$
est commutatif, avec $Y'_i = \Spec(R'_i)$.
Soit $y'_i \in Y'_i$ l'image de $y$.
Nous avons $X_y = \lim X_{i, y'_i}$ parce que $X = \lim X_i$,
$Y = \lim Y'_i$ et $\kappa(y) = \colim \kappa(y'_i)$.
Écrivons maintenant $X_y = U \amalg V$, où $U$ et $V$ sont ouverts et fermés.
Alors $U, V$ sont les images inverses d'ouverts $U_i, V_i$ de $X_{i, y'_i}$
(Limites d'espaces, Lemme \ref{spaces-limits-lemma-descend-opens}).
D'après le Théorème \ref{theorem-stein-factorization-Noetherian}, les fibres
de $X_i \to Y'_i$ sont connexes ; ainsi, $U$ ou $V$ est vide.
Cela achève la démonstration.
\end{proof}

\noindent
Voici une application.

\begin{lemma}
\label{lemma-geometrically-connected-fibres-towards-normal}
Soit $S$ un schéma. Soit $f : X \to Y$ un morphisme d'espaces
algébriques sur $S$. Supposons que
\begin{enumerate}
\item $f$ soit propre ;
\item $Y$ soit intègre (Espaces sur les corps, Définition
\ref{spaces-over-fields-definition-integral-algebraic-space})
de point générique $\xi$ ;
\item $Y$ soit normal ;
\item $X$ soit réduit ;
\item tout point générique d'une composante irréductible de $|X|$ s'envoie sur $\xi$ ;
\item nous ayons $H^0(X_\xi, \mathcal{O}) = \kappa(\xi)$.
\end{enumerate}
Alors $f_*\mathcal{O}_X = \mathcal{O}_Y$ et $f$
a des fibres géométriquement connexes.
\end{lemma}

\begin{proof}
Appliquons le Théorème \ref{theorem-stein-factorization-general} pour obtenir une
factorisation $X \to Y' \to Y$. Il suffit de montrer que $Y' = Y$.
Il suffit de montrer que $Y' \times_Y V \to V$ est un isomorphisme,
où $V \to Y$ est un morphisme étale et $V$ un schéma affine intègre ;
voir Espaces sur les corps, Lemme
\ref{spaces-over-fields-lemma-normal-integral-cover-by-affines}.
La formation de $Y'$ commute au changement de base étale ; voir
Morphismes d'espaces, Lemme
\ref{spaces-morphisms-lemma-properties-normalization}.
Les points génériques de $X \times_Y V$ se trouvent au-dessus des points génériques de $X$
(Espaces décents, Lemme \ref{decent-spaces-lemma-decent-generic-points})
et s'envoient donc sur le point générique de $V$ d'après l'hypothèse (5). De plus, la condition
(6) est préservée par le changement de base $V \to Y$, par exemple par
changement de base plat (Cohomologie des espaces, Lemme
\ref{spaces-cohomology-lemma-flat-base-change-cohomology}).
Il suffit donc de démontrer le lemme lorsque $Y$ est un schéma affine
normal et intègre.

\medskip\noindent
Supposons que $Y$ soit un schéma affine normal et intègre. Nous allons montrer que $Y' \to Y$
est un isomorphisme en appliquant Morphismes, Lemme
\ref{morphisms-lemma-finite-birational-over-normal}.
En effet, $Y'$ est réduit parce que $X$ est réduit (Morphismes d'espaces, Lemme
\ref{spaces-morphisms-lemma-normalization-in-reduced}).
Le morphisme $Y' \to Y$ est entier d'après le théorème cité plus haut.
Puisque $Y$ est décent et que $X \to Y$ est séparé, nous voyons que
$X$ est lui aussi décent ; pour le voir, utilisons
Espaces décents, Lemmes
\ref{decent-spaces-lemma-properties-trivial-implications} et
\ref{decent-spaces-lemma-property-over-property}. D'après l'hypothèse (5),
Morphismes d'espaces, Lemme \ref{spaces-morphisms-lemma-normalization-generic},
et Espaces décents, Lemme \ref{decent-spaces-lemma-decent-generic-points},
nous voyons que tout point générique d'une composante irréductible de $|Y'|$
s'envoie sur $\xi$. D'autre part, puisque $Y'$ est le spectre relatif
de $f_*\mathcal{O}_X$, nous voyons que la fibre schématique
$Y'_\xi$ est le spectre de $H^0(X_\xi, \mathcal{O})$, qui est
égal à $\kappa(\xi)$ par hypothèse. Ainsi, $Y'$ est un schéma
intègre dont le corps des fonctions est égal à celui de $Y$.
Cela achève la démonstration.
\end{proof}

\noindent
Voici une autre application.

\begin{lemma}
\label{lemma-proper-flat-nr-geom-conn-comps-lower-semicontinuous}
Soit $S$ un schéma.
Soit $X \to Y$ un morphisme d'espaces algébriques sur $S$.
Notons $f$ ce morphisme. S'il est propre, plat et de présentation finie, alors la fonction
$n_{X/Y} : |Y| \to \mathbf{Z}$ qui compte le nombre de composantes
connexes géométriques des fibres de $f$
(Lemme \ref{lemma-base-change-fibres-nr-geometrically-connected-components})
est semi-continue inférieurement.
\end{lemma}

\begin{proof}
La question est locale pour la topologie étale sur $Y$ ; nous pouvons donc supposer, et supposons,
que $Y$ est un schéma affine. Soit $y \in Y$. Posons $n = n_{X/S}(y)$.
Remarquons que $n < \infty$, car la fibre géométrique de $X \to Y$ en $y$
est un espace algébrique propre sur un corps, donc noethérien, et possède
donc un nombre fini de composantes connexes.
Nous devons trouver un voisinage ouvert $V$ de $y$ tel que $n_{X/S}|_V \geq n$.
Soit $X \to Y' \to Y$ la factorisation de Stein du
Théorème \ref{theorem-stein-factorization-general}.
D'après le Lemme \ref{lemma-stein-universally-closed-residue-fields},
il existe un nombre fini de points
$y'_1, \ldots, y'_m \in Y'$ situés au-dessus de $y$,
et les extensions $\kappa(y'_i)/\kappa(y)$ sont finies.
Compléments sur les morphismes, Lemme \ref{more-morphisms-lemma-etale-makes-integral-split},
nous dit qu'après avoir remplacé $Y$ par un voisinage étale
de $y$, nous pouvons supposer que $Y' = V_1 \amalg \ldots \amalg V_m$ est un schéma,
avec $y'_i \in V_i$ et $\kappa(y'_i)/\kappa(y)$ purement inséparable.
Alors les espaces algébriques $X_{y_i'}$ sont géométriquement connexes sur
$\kappa(y)$, donc $m = n$. Les espaces algébriques
$X_i = (f')^{-1}(V_i)$, $i = 1, \ldots, n$,
sont plats et de présentation finie sur $Y$. L'image de $X_i \to Y$
est donc ouverte (Morphismes d'espaces, Lemme \ref{spaces-morphisms-lemma-fppf-open}).
Ainsi, dans un voisinage de $y$, nous voyons que $n_{X/Y}$ est
au moins égal à $n$.
\end{proof}

\begin{lemma}
\label{lemma-proper-flat-geom-red}
Soit $S$ un schéma.
Soit $f : X \to Y$ un morphisme d'espaces algébriques sur $S$. Supposons que
\begin{enumerate}
\item $f$ soit propre, plat et de présentation finie ;
\item les fibres géométriques de $f$ soient réduites.
\end{enumerate}
Alors la fonction $n_{X/S} : |Y| \to \mathbf{Z}$
qui compte le nombre de composantes connexes géométriques
des fibres de $f$
(Lemme \ref{lemma-base-change-fibres-nr-geometrically-connected-components})
est localement constante.
\end{lemma}

\begin{proof}
D'après le Lemme \ref{lemma-proper-flat-nr-geom-conn-comps-lower-semicontinuous},
la fonction $n_{X/Y}$ est semi-continue inférieurement.
Il suffit donc de montrer qu'elle est semi-continue supérieurement.
Pour cela, nous pouvons travailler localement pour la topologie étale sur $Y$ ; nous
pouvons donc supposer que $Y$ est un schéma affine.
Pour $y \in Y$, considérons la $\kappa(y)$-algèbre
$$
A = H^0(X_y, \mathcal{O}_{X_y})
$$
D'après Espaces sur les corps, Lemme
\ref{spaces-over-fields-lemma-proper-geometrically-reduced-global-sections}
et puisque $X_y$ est géométriquement réduit,
$A$ est un produit fini d'extensions finies séparables de $\kappa(y)$.
Ainsi, $A \otimes_{\kappa(y)} \kappa(\overline{y})$ est un produit
de $\beta_0(y) = \dim_{\kappa(y)} A$ exemplaires de $\kappa(\overline{y})$.
Par conséquent, $X_{\overline{y}}$ a $\beta_0(y)$ composantes connexes.
Autrement dit, nous avons $n_{X/S} = \beta_0$ comme fonctions sur $Y$.
Ainsi, $n_{X/Y}$ est semi-continue supérieurement d'après
Catégories dérivées des espaces, Lemme
\ref{spaces-perfect-lemma-jump-loci-geometric}.
Cela achève la démonstration.
\end{proof}

\begin{lemma}
\label{lemma-stein-factorization-etale}
Soit $S$ un schéma.
Soit $f : X \to Y$ un morphisme propre d'espaces algébriques sur $S$.
Soit $X \to Y' \to Y$ la factorisation de Stein de $f$
(Théorème \ref{theorem-stein-factorization-general}).
Si $f$ est de présentation finie, plat et à fibres
géométriquement réduites (Définition \ref{definition-geometrically-reduced-fibre}),
alors $Y' \to Y$ est fini étale.
\end{lemma}

\begin{proof}
La formation de la factorisation de Stein commute au changement de base plat ;
voir Lemme \ref{lemma-stein-universally-closed}.
Nous pouvons donc travailler localement pour la topologie étale sur $Y$ et supposer que $Y$
est un schéma affine. Alors $Y'$ est un schéma affine et $Y' \to Y$
est entier.

\medskip\noindent
Soit $y \in Y$. Soit $n$ le nombre de composantes connexes de
la fibre géométrique $X_{\overline{y}}$. Remarquons que $n < \infty$,
car la fibre géométrique de $X \to Y$ en $y$ est un espace
algébrique propre sur un corps, donc noethérien,
et possède donc un nombre fini de composantes connexes.
D'après le Lemme \ref{lemma-stein-universally-closed-residue-fields},
il existe un nombre fini de points $y'_1, \ldots, y'_m \in Y'$ situés au-dessus de $y$ ;
pour chaque $i$, nous pouvons choisir un point de type fini $x_i \in |X_y|$
qui s'envoie sur $y'_i$, et l'extension $\kappa(y'_i)/\kappa(y)$ est finie.
Ainsi, Compléments sur les morphismes,
Lemme \ref{more-morphisms-lemma-etale-makes-integral-split},
nous dit qu'après avoir remplacé $Y$ par un voisinage étale
de $y$, nous pouvons supposer que $Y' = V_1 \amalg \ldots \amalg V_m$ est un schéma,
avec $y'_i \in V_i$ et $\kappa(y'_i)/\kappa(y)$ purement inséparable.
Dans ce cas, les espaces algébriques $X_{y_i'}$
sont géométriquement connexes sur $\kappa(y)$, donc $m = n$.
Les espaces algébriques $X_i = (f')^{-1}(V_i)$, $i = 1, \ldots, n$,
sont propres, plats, de présentation finie et à fibres géométriquement
réduites sur $Y$. Il suffit de démontrer le lemme
pour chacun des morphismes $X_i \to Y$. Nous nous ramenons ainsi au cas où
$X_{\overline{y}}$ est connexe.

\medskip\noindent
Supposons que $X_{\overline{y}}$ soit connexe. D'après le
Lemme \ref{lemma-proper-flat-geom-red},
nous voyons que $X \to Y$ a des fibres géométriquement connexes
dans un voisinage de $y$. Ainsi,
nous pouvons supposer que les fibres de $X \to Y$ sont géométriquement connexes.
Alors $f_*\mathcal{O}_X = \mathcal{O}_Y$ d'après
Catégories dérivées des espaces, Lemme
\ref{spaces-perfect-lemma-proper-flat-geom-red-connected},
ce qui achève la démonstration.
\end{proof}

\noindent
La démonstration du lemme suivant utilise la factorisation de Stein pour les schémas ;
c'est pourquoi il se trouve dans cette section.

\begin{lemma}
\label{lemma-split-off-proper-part-henselian}
Soit $(A, I)$ un couple hensélien. Soit $X$ un espace algébrique
séparé et de type fini sur $A$. Posons
$X_0 = X \times_{\Spec(A)} \Spec(A/I)$.
Soit $Y \subset X_0$ un sous-espace ouvert et fermé tel que
$Y \to \Spec(A/I)$ soit propre. Il existe alors un sous-espace ouvert et fermé
$W \subset X$ qui est propre sur $A$ et tel que
$W \times_{\Spec(A)} \Spec(A/I) = Y$.
\end{lemma}

\begin{proof}
Nous noterons $T \mapsto T_0$ le changement de base par
$\Spec(A/I) \to \Spec(A)$.
D'après une forme faible du lemme de Chow (sous la forme de
Cohomologie des espaces, Lemme \ref{spaces-cohomology-lemma-weak-chow}),
il existe un morphisme propre surjectif $\varphi : X' \to X$ tel
que $X'$ admette une immersion dans $\mathbf{P}^n_A$.
Posons $Y' = \varphi^{-1}(Y)$. C'est un sous-schéma ouvert et fermé
de $X'_0$. Le lemme est vrai pour $(X', Y')$ d'après
Compléments sur les morphismes, Lemme
\ref{more-morphisms-lemma-split-off-proper-part-henselian}.
Soit $W' \subset X'$ le sous-schéma ouvert et fermé, propre
sur $A$, tel que $Y' = W'_0$.
D'après Morphismes d'espaces, Lemme
\ref{spaces-morphisms-lemma-universally-closed-permanence}
$Q_1 = \varphi(|W'|) \subset |X|$ et
$Q_2 = \varphi(|X' \setminus W'|) \subset |X|$
sont des fermés et, d'après
Morphismes d'espaces, Lemme \ref{spaces-morphisms-lemma-image-proper-is-proper},
toute structure de sous-espace fermé sur $Q_1$ est propre sur $A$.
L'image de $Q_1 \cap Q_2$ dans $\Spec(A)$ est fermée.
Puisque $(A, I)$ est hensélien, si $Q_1 \cap Q_2$ est non vide, alors nous
trouvons que $Q_1 \cap Q_2$ a un point situé au-dessus de $\Spec(A/I)$.
C'est impossible puisque $W'_0 = Y' = \varphi^{-1}(Y)$.
Nous concluons que $Q_1$ est ouvert et fermé dans $|X|$.
Soit $W \subset X$ le sous-espace ouvert et fermé
correspondant. Alors $W$ est propre sur $A$ et $W_0 = Y$.
\end{proof}










\section{Prolongement de propriétés à partir d'un ouvert}
\label{section-extending-properties}

\noindent
Dans cette section, nous rassemblons plusieurs résultats de la forme suivante : si $f : X \to Y$
est un morphisme plat d'espaces algébriques et si $f$ possède une certaine propriété au-dessus
d'un ouvert dense
de $Y$, alors $f$ possède la même propriété au-dessus de $Y$ tout entier.

\begin{lemma}
\label{lemma-flat-finite-type-finitely-presented-over-dense-open}
Soit $S$ un schéma.
Soit $f : X \to Y$ un morphisme d'espaces algébriques sur $S$.
Soit $\mathcal{F}$ un $\mathcal{O}_X$-module quasi-cohérent.
Soit $V \subset Y$ un sous-espace ouvert. Supposons que
\begin{enumerate}
\item $f$ soit localement de présentation finie,
\item $\mathcal{F}$ soit de type fini et plat sur $Y$,
\item $V \to Y$ soit quasi-compact et schématiquement dense,
\item $\mathcal{F}|_{f^{-1}V}$ soit de présentation finie.
\end{enumerate}
Alors $\mathcal{F}$ est de présentation finie.
\end{lemma}

\begin{proof}
Il suffit de montrer que l'image inverse de $\mathcal{F}$ sur un schéma surjectif
et étale sur $X$ est de présentation finie. Nous pouvons donc supposer que $X$
est un schéma. De même, nous pouvons remplacer $Y$ par un schéma surjectif et
étale sur $Y$ (l'image inverse de $V$ dans ce schéma est schématiquement
dense ; voir
Morphismes d'espaces, section
\ref{spaces-morphisms-section-scheme-theoretic-closure}).
Nous nous ramenons ainsi au cas des schémas, qui est
Compléments sur la platitude, Lemme
\ref{flat-lemma-flat-finite-type-finitely-presented-over-dense-open}.
\end{proof}

\begin{lemma}
\label{lemma-flat-finite-type-finitely-presented-over-dense-open-X}
Soit $S$ un schéma.
Soit $f : X \to Y$ un morphisme d'espaces algébriques sur $S$.
Soit $V \subset Y$ un sous-espace ouvert.
Supposons que
\begin{enumerate}
\item $f$ soit localement de type fini et plat,
\item $V \to Y$ soit quasi-compact et schématiquement dense,
\item $f|_{f^{-1}V} : f^{-1}V \to V$ soit localement de présentation finie.
\end{enumerate}
Alors $f$ est localement de présentation finie.
\end{lemma}

\begin{proof}
La démonstration est identique à celle du
Lemme \ref{lemma-flat-finite-type-finitely-presented-over-dense-open},
à ceci près qu'on utilise
Compléments sur la platitude, Lemme
\ref{flat-lemma-flat-finite-type-finitely-presented-over-dense-open-X}.
\end{proof}

\begin{lemma}
\label{lemma-flat-fp-dimension-over-dense-open}
Soit $S$ un schéma. Soit $f : X \to Y$ un morphisme d'espaces algébriques
sur $S$, plat et localement de type fini. Soit $V \subset Y$ un
sous-espace ouvert tel que $|V| \subset |Y|$ soit dense et que $X_V \to V$
soit de dimension relative $\leq d$. Supposons en outre que l'une des conditions suivantes soit vérifiée :
\begin{enumerate}
\item $f$ est localement de présentation finie, ou
\item $V \to Y$ est quasi-compact.
\end{enumerate}
Alors $f : X \to Y$ est de dimension relative $\leq d$.
\end{lemma}

\begin{proof}
Nous pouvons remplacer $Y$ par son réduit et donc supposer $Y$ réduit.
Alors $V$ est schématiquement dense dans $Y$ ; voir
Morphismes d'espaces, Lemme
\ref{spaces-morphisms-lemma-quasi-compact-immersion}.
Par définition, la propriété d'être de dimension relative $\leq d$ peut
se vérifier sur un recouvrement étale ; voir
Morphismes d'espaces, section \ref{spaces-morphisms-section-relative-dimension}.
Il suffit donc de montrer que $f$ est de dimension relative $\leq d$
après avoir remplacé $X$ par un schéma surjectif et étale sur $X$.
De même, nous pouvons remplacer $Y$ par un schéma surjectif et
étale sur $Y$. L'image inverse de $V$ dans ce schéma est schématiquement
dense ; voir
Morphismes d'espaces, section
\ref{spaces-morphisms-section-scheme-theoretic-closure}.
Comme un ouvert schématiquement dense d'un schéma est en particulier
dense, nous nous ramenons au cas des schémas, qui est
Compléments sur la platitude, Lemme
\ref{flat-lemma-flat-finite-presentation-dimension-over-dense-open}.
\end{proof}

\begin{lemma}
\label{lemma-proper-flat-finite-over-dense-open}
Soit $S$ un schéma. Soit $f : X \to Y$ un morphisme d'espaces algébriques
sur $S$, plat et propre. Soit $V \to Y$ un sous-espace ouvert
tel que $|V| \subset |Y|$ soit dense et que $X_V \to V$ soit fini. Si, de plus,
$f$ est localement de présentation finie ou $V \to Y$ est quasi-compact,
alors $f$ est fini.
\end{lemma}

\begin{proof}
D'après le Lemme \ref{lemma-flat-fp-dimension-over-dense-open},
les fibres de $f$ sont de dimension nulle.
D'après Morphismes d'espaces, Lemme
\ref{spaces-morphisms-lemma-locally-quasi-finite-rel-dimension-0}
il s'ensuit que $f$ est localement quasi-fini.
D'après Morphismes d'espaces, Lemme
\ref{spaces-morphisms-lemma-locally-quasi-finite-separated-representable}
il s'ensuit que $f$ est représentable.
La finitude de $f$ se vérifie localement pour la topologie étale sur $Y$ ;
nous pouvons donc supposer que $Y$ est un schéma. Comme $f$ est représentable,
nous nous ramenons au cas des schémas, qui est
Compléments sur la platitude, Lemme \ref{flat-lemma-proper-flat-finite-over-dense-open}.
\end{proof}

\begin{lemma}
\label{lemma-zariski}
Soit $S$ un schéma.
Soit $f : X \to Y$ un morphisme d'espaces algébriques sur $S$.
Soit $V \subset Y$ un sous-espace ouvert. Si
\begin{enumerate}
\item $f$ est séparé, localement de type fini et plat,
\item $f^{-1}(V) \to V$ est un isomorphisme, et
\item $V \to Y$ est quasi-compact et schématiquement dense,
\end{enumerate}
alors $f$ est une immersion ouverte.
\end{lemma}

\begin{proof}
En appliquant le
Lemme \ref{lemma-flat-finite-type-finitely-presented-over-dense-open-X},
nous voyons que $f$ est localement de présentation finie. En appliquant le
Lemme \ref{lemma-flat-fp-dimension-over-dense-open},
nous voyons que $f$ est de dimension relative $\leq 0$.
D'après Morphismes d'espaces, Lemme
\ref{spaces-morphisms-lemma-locally-quasi-finite-rel-dimension-0}
il s'ensuit que $f$ est localement quasi-fini.
D'après Morphismes d'espaces, Lemme
\ref{spaces-morphisms-lemma-locally-quasi-finite-separated-representable}
il s'ensuit que $f$ est représentable.
D'après Descente sur les espaces, Lemme
\ref{spaces-descent-lemma-descending-property-open-immersion}
le fait que $f$ soit une immersion ouverte se vérifie localement pour la topologie étale sur $Y$.
Nous pouvons donc supposer que $Y$ est un schéma. Comme $f$ est représentable,
nous nous ramenons au cas des schémas, qui est
Compléments sur la platitude, Lemme \ref{flat-lemma-zariski}.
\end{proof}







\section{Éclatement et platitude}
\label{section-blowup-flat}

\noindent
Au lieu de reprendre le travail de
Compléments sur la platitude, section \ref{flat-section-blowup-flat},
nous démontrons un analogue de Compléments sur la platitude, Lemme \ref{flat-lemma-push-ideal},
qui nous dit que le problème de trouver un éclatement convenable
est souvent local pour la topologie étale sur la base.

\begin{lemma}
\label{lemma-push-ideal}
Soit $S$ un schéma. Soit $X$ un espace algébrique quasi-compact et quasi-séparé
sur $S$. Soit $\varphi : W \to X$ un morphisme étale, quasi-compact
et séparé. Soit $U \subset X$ un sous-espace ouvert quasi-compact.
Soit $\mathcal{I} \subset \mathcal{O}_W$ un faisceau quasi-cohérent
d'idéaux de type fini tel que
$V(\mathcal{I}) \cap \varphi^{-1}(U) = \emptyset$.
Alors il existe un faisceau quasi-cohérent d'idéaux de type fini
$\mathcal{J} \subset \mathcal{O}_X$ tel que
\begin{enumerate}
\item $V(\mathcal{J}) \cap U = \emptyset$, et
\item $\varphi^{-1}(\mathcal{J})\mathcal{O}_W = \mathcal{I} \mathcal{I}'$
pour un certain idéal quasi-cohérent de type fini
$\mathcal{I}' \subset \mathcal{O}_W$.
\end{enumerate}
\end{lemma}

\begin{proof}
Choisissons une factorisation $W \to Y \to X$ dans laquelle $j : W \to Y$ est une
immersion ouverte quasi-compacte et $\pi : Y \to X$ est un morphisme fini
de présentation finie
(Lemme \ref{lemma-quasi-finite-separated-pass-through-finite-addendum}).
Posons $V = j(W) \cup \pi^{-1}(U) \subset Y$. Notons que
$\mathcal{I}$ sur $W \cong j(W)$ et $\mathcal{O}_{\pi^{-1}(U)}$
se recollent en un faisceau quasi-cohérent d'idéaux de type fini
$\mathcal{I}_1 \subset \mathcal{O}_V$. D'après
Limites d'espaces, Lemme \ref{spaces-limits-lemma-extend},
il existe un faisceau quasi-cohérent d'idéaux de type fini
$\mathcal{I}_2 \subset \mathcal{O}_Y$ tel que
$\mathcal{I}_2|_V = \mathcal{I}_1$. Autrement dit,
$\mathcal{I}_2 \subset \mathcal{O}_Y$
est un faisceau quasi-cohérent d'idéaux de type fini tel que
$V(\mathcal{I}_2)$ soit disjoint de $\pi^{-1}(U)$ et que
$j^{-1}\mathcal{I}_2 = \mathcal{I}$. Notons $i : Z \to Y$
l'immersion fermée correspondante, qui est de présentation finie
(Morphismes d'espaces, Lemme
\ref{spaces-morphisms-lemma-closed-immersion-finite-presentation}).
En particulier, la composée $\tau = \pi \circ i : Z \to X$ est finie
et de présentation finie
(Morphismes d'espaces, Lemmes
\ref{spaces-morphisms-lemma-composition-finite-presentation} et
\ref{spaces-morphisms-lemma-composition-integral}).

\medskip\noindent
Posons $\mathcal{F} = \tau_*\mathcal{O}_Z$, que nous considérons comme
un $\mathcal{O}_X$-module quasi-cohérent. D'après
Descente sur les espaces, Lemme
\ref{spaces-descent-lemma-finite-finitely-presented-module}
nous voyons que $\mathcal{F}$ est un $\mathcal{O}_X$-module de présentation finie.
Posons $\mathcal{J} = \text{Fit}_0(\mathcal{F})$. (Insérer ici une référence aux
idéaux de Fitting sur les topos annelés.) C'est un faisceau quasi-cohérent
d'idéaux de type fini sur $X$ (puisque $\mathcal{F}$ est de présentation finie, voir
Compléments d'algèbre, Lemme \ref{more-algebra-lemma-fitting-ideal-basics}).
La partie (1) du lemme vaut parce que $|\tau|(|Z|) \cap |U| = \emptyset$
par notre choix de $\mathcal{I}_2$ et parce que l'idéal de Fitting d'indice $0$
du module nul est égal au faisceau structural. Pour démontrer (2), notons que
$\varphi^{-1}(\mathcal{J})\mathcal{O}_W = \text{Fit}_0(\varphi^*\mathcal{F})$
car la formation des idéaux de Fitting commute au changement de base.
D'autre part, puisque $\varphi : W \to X$ est séparé et étale,
nous voyons que $(1, j) : W \to W \times_X Y$ est une immersion ouverte et fermée.
Ainsi $W \times_Y Z = V(\mathcal{I}) \amalg Z'$ pour un certain morphisme d'espaces
algébriques fini et de présentation finie $\tau' : Z' \to W$.
Nous voyons donc que
\begin{align*}
\text{Fit}_0(\varphi^*\mathcal{F}) & =
\text{Fit}_0((W \times_Y Z \to W)_*\mathcal{O}_{W \times_Y Z}) \\
& =
\text{Fit}_0(\mathcal{O}_W/\mathcal{I}) \cdot
\text{Fit}_0(\tau'_*\mathcal{O}_{Z'}) \\
& =
\mathcal{I} \cdot \text{Fit}_0(\tau'_*\mathcal{O}_{Z'})
\end{align*}
où la deuxième égalité résulte de
Compléments d'algèbre, Lemme \ref{more-algebra-lemma-fitting-ideal-basics},
transposé aux faisceaux sur les topos annelés.
En posant $\mathcal{I}' = \text{Fit}_0(\tau'_*\mathcal{O}_{Z'})$,
on achève la démonstration du lemme.
\end{proof}

\begin{theorem}
\label{theorem-flatten-module}
Soit $S$ un schéma. Soit $B$ un espace algébrique quasi-compact et quasi-séparé
sur $S$. Soit $X$ un espace algébrique sur $B$.
Soit $\mathcal{F}$ un module quasi-cohérent sur $X$.
Soit $U \subset B$ un sous-espace ouvert quasi-compact. Supposons que
\begin{enumerate}
\item $X$ soit quasi-compact,
\item $X$ soit localement de présentation finie sur $B$,
\item $\mathcal{F}$ soit un module de type fini,
\item $\mathcal{F}_U$ soit de présentation finie, et
\item $\mathcal{F}_U$ soit plat sur $U$.
\end{enumerate}
Alors il existe un éclatement $U$-admissible $B' \to B$ tel que la
transformée stricte $\mathcal{F}'$ de $\mathcal{F}$ soit un
$\mathcal{O}_{X \times_B B'}$-module de présentation finie et
plat sur $B'$.
\end{theorem}

\begin{proof}
Choisissons un schéma affine $V$ et un morphisme étale surjectif $V \to X$.
Puisque la transformée stricte commute à la localisation étale
(Diviseurs sur les espaces, Lemme \ref{spaces-divisors-lemma-strict-transform-local}),
il suffit de démontrer le résultat en remplaçant $X$ par $V$. Nous pouvons donc
supposer que $X \to B$ est représentable (en plus des
hypothèses du lemme).

\medskip\noindent
Supposons que $X \to B$ soit représentable. Choisissons un schéma affine $W$ et un
morphisme étale surjectif $\varphi : W \to B$. Notons que $X \times_B W$
est un schéma. D'après le cas des schémas
(Compléments sur la platitude, Théorème \ref{flat-theorem-flatten-module}),
nous pouvons trouver un faisceau quasi-cohérent d'idéaux de type fini
$\mathcal{I} \subset \mathcal{O}_W$ tel que
(a) $|V(\mathcal{I})| \cap |\varphi^{-1}(U)| = \emptyset$ et (b)
la transformée stricte de $\mathcal{F}|_{X \times_B W}$ relativement
à l'éclatement $W' \to W$ en $\mathcal{I}$ devienne plate sur $W'$
et soit un module de présentation finie. Choisissons un faisceau d'idéaux de type fini
$\mathcal{J} \subset \mathcal{O}_B$ comme dans le Lemme \ref{lemma-push-ideal}.
Soit $B' \to B$ l'éclatement de $\mathcal{J}$. Nous affirmons que cet
éclatement convient. En effet, il est clair que $B' \to B$ est $U$-admissible
par notre choix de l'idéal $\mathcal{J}$. De plus, le changement de base
$B' \times_B W \to W$ est l'éclatement de $W$ en
$\varphi^{-1}\mathcal{J} = \mathcal{I}\mathcal{I}'$
(compatibilité de l'éclatement avec le changement de base plat, voir
Diviseurs sur les espaces, Lemme
\ref{spaces-divisors-lemma-flat-base-change-blowing-up}).
Il existe donc une factorisation
$$
W \times_B B' \to W' \to W
$$
où le premier morphisme est lui aussi un éclatement, voir
Diviseurs sur les espaces, Lemme \ref{spaces-divisors-lemma-blowing-up-two-ideals}).
La restriction de $\mathcal{F}'$ (qui est défini sur $B' \times_B X$)
à $W \times_B B' \times_B X$ est la transformée stricte de
$\mathcal{F}|_{X \times_B W}$
(Diviseurs sur les espaces, Lemme \ref{spaces-divisors-lemma-strict-transform-local})
et est donc la transformée stricte effectuée deux fois de
$\mathcal{F}|_{X \times_B W}$ par les deux éclatements affichés ci-dessus
(Diviseurs sur les espaces, Lemme
\ref{spaces-divisors-lemma-strict-transform-composition-blowups}).
Après le premier éclatement, notre faisceau est déjà plat sur
la base et de présentation finie (par construction). Il en est donc encore ainsi 
après la deuxième transformée stricte (puisque celle-ci est une
image inverse d'après Diviseurs sur les espaces, Lemme
\ref{spaces-divisors-lemma-strict-transform-flat}).
Nous voyons ainsi que la restriction de $\mathcal{F}'$
à un recouvrement étale de $B' \times_B X$ possède les propriétés voulues,
ce qui démontre le théorème.
\end{proof}







\section{Applications}
\label{section-applications-flattening-by-blowing-up}

\noindent
Dans cette section, nous appliquons le résultat d'aplatissement par éclatement.

\begin{lemma}
\label{lemma-flat-after-blowing-up}
Soit $S$ un schéma. Soit $f : X \to B$ un morphisme d'espaces
algébriques sur $S$. Soit $U \subset B$ un sous-espace ouvert. Supposons que
\begin{enumerate}
\item $B$ soit quasi-compact et quasi-séparé,
\item $U$ soit quasi-compact,
\item $f : X \to B$ soit de type fini et quasi-séparé, et
\item $f^{-1}(U) \to U$ soit plat et localement de présentation finie.
\end{enumerate}
Alors il existe un éclatement $U$-admissible $B' \to B$ tel que
le transformé strict $X'$ de $X$ soit plat et de présentation finie
sur $B'$.
\end{lemma}

\begin{proof}
Soit $B' \to B$ un éclatement $U$-admissible. Notons que le transformé strict
de $X$ est quasi-compact et quasi-séparé sur $B'$, puisque $X$ est quasi-compact
et quasi-séparé sur $B$. Il nous suffit donc de trouver
un éclatement $U$-admissible tel que le transformé strict devienne plat et
localement de présentation finie. Nous ne pouvons pas appliquer directement le
Théorème \ref{theorem-flatten-module}, car $X$ n'est pas localement de
présentation finie sur $B$.

\medskip\noindent
Choisissons un schéma affine $V$ et un morphisme étale surjectif $V \to X$.
(C'est possible puisque $X$ est quasi-compact, étant un espace de type fini sur
l'espace quasi-compact $B$.) Il suffit alors de démontrer le résultat pour
le morphisme $V \to B$ (puisque le transformé strict commute à la localisation
étale, voir Diviseurs sur les espaces,
Lemme \ref{spaces-divisors-lemma-strict-transform-local}).
Nous pouvons donc supposer que $X \to B$ est séparé et de type fini.
Dans ce cas, nous pouvons trouver une immersion fermée $i : X \to Y$ avec $Y \to B$
séparé et de présentation finie, voir
Limites d'espaces, Proposition
\ref{spaces-limits-proposition-separated-closed-in-finite-presentation}.

\medskip\noindent
Appliquons le Théorème \ref{theorem-flatten-module} à $\mathcal{F} = i_*\mathcal{O}_X$
sur $Y/B$. Nous trouvons un éclatement $U$-admissible $B' \to B$ tel que la
transformée stricte de $\mathcal{F}$ soit plate sur $B'$ et de présentation finie.
Soit $X'$ le transformé strict de $X$ par l'éclatement $B' \to B$.
Soit $i' : X' \to Y \times_B B'$ le morphisme induit.
Puisque la formation de la transformée stricte commute à l'image directe par les morphismes
affines (Diviseurs sur les espaces, Lemme
\ref{spaces-divisors-lemma-strict-transform-affine}),
nous voyons que $i'_*\mathcal{O}_{X'}$ est plat sur $B'$ et de
présentation finie comme $\mathcal{O}_{Y \times_B B'}$-module.
Ainsi $X' \to B'$ est plat et localement de présentation finie.
Cela implique le lemme d'après nos remarques précédentes.
\end{proof}

\begin{lemma}
\label{lemma-finite-after-blowing-up}
Soit $S$ un schéma. Soit $f : X \to B$ un morphisme d'espaces
algébriques sur $S$. Soit $U \subset B$ un sous-espace ouvert. Supposons que
\begin{enumerate}
\item $B$ soit quasi-compact et quasi-séparé,
\item $U$ soit quasi-compact,
\item $f : X \to B$ soit propre, et
\item $f^{-1}(U) \to U$ soit fini localement libre.
\end{enumerate}
Alors il existe un éclatement $U$-admissible $B' \to B$ tel que
le transformé strict $X'$ de $X$ soit fini localement libre sur $B'$.
\end{lemma}

\begin{proof}
D'après le Lemme \ref{lemma-flat-after-blowing-up}, nous pouvons supposer que
$X \to B$ est plat et de présentation finie. Après avoir remplacé
$B$ par un éclatement $U$-admissible si nécessaire, nous pouvons supposer
que $U \subset B$ est schématiquement dense. Alors $f$ est
fini d'après le Lemme \ref{lemma-proper-flat-finite-over-dense-open}.
Par conséquent, $f$ est fini localement libre d'après
Morphismes d'espaces, Lemme \ref{spaces-morphisms-lemma-finite-flat}.
\end{proof}

\begin{lemma}
\label{lemma-isomorphism-after-blowing-up}
Soit $S$ un schéma. Soit $f : X \to B$ un morphisme d'espaces
algébriques sur $S$. Soit $U \subset B$ un sous-espace ouvert. Supposons que
\begin{enumerate}
\item $B$ soit quasi-compact et quasi-séparé,
\item $U$ soit quasi-compact,
\item $f : X \to B$ soit propre, et
\item $f^{-1}(U) \to U$ soit un isomorphisme.
\end{enumerate}
Alors il existe un éclatement $U$-admissible $B' \to B$ tel que
le transformé strict $X'$ de $X$ s'envoie isomorphiquement sur $B'$.
\end{lemma}

\begin{proof}
D'après le Lemme \ref{lemma-flat-after-blowing-up}, nous pouvons supposer que
$X \to B$ est plat et de présentation finie. Après avoir remplacé
$B$ par un éclatement $U$-admissible si nécessaire, nous pouvons supposer
que $U \subset B$ est schématiquement dense. Alors $f$ est
fini d'après le Lemme \ref{lemma-proper-flat-finite-over-dense-open}
et une immersion ouverte d'après le Lemme \ref{lemma-zariski}. Ainsi $f$
est une immersion ouverte dont l'image est fermée et contient
l'ouvert dense $U$ ; par conséquent, $f$ est un isomorphisme.
\end{proof}

\begin{lemma}
\label{lemma-zariski-after-blowup}
Soit $S$ un schéma. Soit $f : X \to B$ un morphisme d'espaces algébriques
sur $S$. Soit $U \subset B$ un sous-espace ouvert. Supposons que
\begin{enumerate}
\item $B$ soit quasi-compact et quasi-séparé,
\item $U$ soit quasi-compact,
\item $f$ soit de type fini,
\item $f^{-1}(U) \to U$ soit un isomorphisme.
\end{enumerate}
Alors il existe un éclatement $U$-admissible $B' \to B$ tel que $U$
soit schématiquement dense dans $B'$ et que le transformé
strict $X'$ de $X$ s'envoie isomorphiquement sur un sous-espace ouvert de $B'$.
\end{lemma}

\begin{proof}
Ce lemme généralise le
Lemme \ref{lemma-isomorphism-after-blowing-up}.
Comme la composée d'éclatements $U$-admissibles est $U$-admissible
(Diviseurs sur les espaces, Lemme
\ref{spaces-divisors-lemma-composition-admissible-blowups})
nous pouvons procéder par étapes. Choisissons un faisceau quasi-cohérent
d'idéaux de type fini $\mathcal{I} \subset \mathcal{O}_B$ tel que
$|B| \setminus |U| = |V(\mathcal{I})|$. Remplaçons $B$ par l'éclatement
de $B$ en $\mathcal{I}$ et $X$ par le transformé strict de $X$.
Après ce remplacement, $B \setminus U$ est le support d'un diviseur de
Cartier effectif $D$ (Diviseurs sur les espaces, Lemme
\ref{spaces-divisors-lemma-blowing-up-gives-effective-Cartier-divisor}).
En particulier, $U$ est schématiquement dense dans $B$
(Diviseurs sur les espaces, Lemme
\ref{spaces-divisors-lemma-complement-effective-Cartier-divisor}).
Ensuite, nous effectuons un autre éclatement $U$-admissible pour nous ramener au cas où
$X \to B$ est plat et de présentation finie, voir le
Lemme \ref{lemma-flat-after-blowing-up}. Notons que $U$ reste schématiquement
dense dans $B$. Par conséquent, $X \to B$ est une immersion ouverte d'après le
Lemme \ref{lemma-zariski}.
\end{proof}

\noindent
Le lemme suivant affirme qu'une modification peut être dominée
par un éclatement.

\begin{lemma}
\label{lemma-dominate-modification-by-blowup}
Soit $S$ un schéma. Soit $f : X \to B$ un morphisme d'espaces algébriques
sur $S$. Soit $U \subset B$ un sous-espace ouvert. Supposons que
\begin{enumerate}
\item $B$ soit quasi-compact et quasi-séparé,
\item $U$ soit quasi-compact,
\item $f : X \to B$ soit propre,
\item $f^{-1}(U) \to U$ soit un isomorphisme.
\end{enumerate}
Alors il existe un éclatement $U$-admissible $B' \to B$ qui domine $X$,
c'est-à-dire qu'il existe une factorisation $B' \to X \to B$
du morphisme d'éclatement.
\end{lemma}

\begin{proof}
D'après le Lemme \ref{lemma-isomorphism-after-blowing-up}, nous pouvons trouver un éclatement
$U$-admissible $B' \to B$ tel que le transformé strict $X'$ s'envoie isomorphiquement sur
$B'$. Nous pouvons alors prendre $B' = X' \to X$ comme factorisation.
\end{proof}

\begin{lemma}
\label{lemma-get-section-after-blowup}
Soit $S$ un schéma. Soient $X$, $Y$ des espaces algébriques sur $S$.
Soient $U \subset W \subset Y$ des sous-espaces ouverts.
Soient $f : X \to W$ et $s : U \to X$ des morphismes
tels que $f \circ s = \text{id}_U$. Supposons que
\begin{enumerate}
\item $f$ soit propre,
\item $Y$ soit quasi-compact et quasi-séparé, et
\item $U$ et $W$ soient quasi-compacts.
\end{enumerate}
Alors il existe un éclatement $U$-admissible $b : Y' \to Y$ et un morphisme
$s' : b^{-1}(W) \to X$ qui prolonge $s$ avec $f \circ s' = b|_{b^{-1}(W)}$.
\end{lemma}

\begin{proof}
Nous pouvons remplacer, et remplaçons, $X$ par l'image schématique de $s$.
Alors $X \to W$ est un isomorphisme au-dessus de $U$, voir
Morphismes d'espaces, Lemme
\ref{spaces-morphisms-lemma-scheme-theoretic-image-of-partial-section}.
D'après le Lemme \ref{lemma-dominate-modification-by-blowup},
il existe un éclatement $U$-admissible $W' \to W$ et un
prolongement $W' \to X$ de $s$.
Nous achevons la démonstration en appliquant
Diviseurs sur les espaces, Lemme \ref{spaces-divisors-lemma-extend-admissible-blowups},
pour prolonger $W' \to W$ en un éclatement $U$-admissible de $Y$.
\end{proof}















\section{Lemme de Chow}
\label{section-chow}

\noindent
Dans cette section, nous démontrons le lemme de Chow
(Lemme \ref{lemma-chow-noetherian-separated}).
Nous encourageons le lecteur à consulter
Cohomologie des espaces, section \ref{spaces-cohomology-section-weak-chow},
où se trouve une version faible du lemme de Chow, facile à démontrer et suffisante
pour de nombreuses applications.

\medskip\noindent
Comme nous n'avons pas encore défini les morphismes projectifs d'espaces algébriques,
les énoncés des lemmes (voir par exemple le
Lemme \ref{lemma-blowup-to-find-embedding}) feront intervenir
des morphismes propres représentables plutôt que des morphismes projectifs.

\begin{lemma}
\label{lemma-find-common-blowups}
Soit $S$ un schéma. Soit $Y$ un espace algébrique quasi-compact et quasi-séparé
sur $S$. Soient $U \to X_1$ et $U \to X_2$ des immersions ouvertes
d'espaces algébriques sur $Y$, et supposons que $U$, $X_1$, $X_2$ soient de type
fini et séparés sur $Y$. Alors il existe un diagramme commutatif
$$
\xymatrix{
X_1' \ar[d] \ar[r] & X & X_2' \ar[l] \ar[d] \\
X_1 & U \ar[l] \ar[lu] \ar[u] \ar[ru] \ar[r] & X_2
}
$$
d'espaces algébriques sur $Y$, où $X_i' \to X_i$ est un éclatement
$U$-admissible, $X_i' \to X$ est une immersion ouverte, et $X$ est séparé et de
type fini sur $Y$.
\end{lemma}

\begin{proof}
Dans toute la démonstration, tous les espaces algébriques seront séparés et de
type fini sur $Y$. Cela implique en particulier que ces espaces algébriques sont
quasi-compacts et quasi-séparés et que les morphismes entre eux
seront quasi-compacts et séparés. Voir
Morphismes d'espaces, sections
\ref{spaces-morphisms-section-separation-axioms} et
\ref{spaces-morphisms-section-quasi-compact}.
Nous utiliserons que, si $U \to W$ est une immersion de tels espaces sur $Y$,
alors l'image schématique $Z$ de $U$ dans $W$ est un sous-espace fermé
de $W$ et $U \to Z$ est une immersion ouverte, $U \subset Z$ est
schématiquement dense et $|U| \subset |Z|$ est dense. Voir
Morphismes d'espaces, Lemme
\ref{spaces-morphisms-lemma-quasi-compact-immersion}.

\medskip\noindent
Soit $X_{12} \subset X_1 \times_Y X_2$ l'image schématique
de $U \to X_1 \times_Y X_2$. Les projections $p_i : X_{12} \to X_i$
induisent des isomorphismes $p_i^{-1}(U) \to U$ d'après
Morphismes d'espaces, Lemme
\ref{spaces-morphisms-lemma-scheme-theoretic-image-of-partial-section}.
Choisissons un éclatement $U$-admissible $X_i^i \to X_i$ tel que
le transformé strict $X_{12}^i$ de $X_{12}$ soit isomorphe à un
sous-espace ouvert de $X_i^i$, voir le
Lemme \ref{lemma-zariski-after-blowup}.
Soit $\mathcal{I}_i \subset \mathcal{O}_{X_i}$ le faisceau quasi-cohérent
d'idéaux de type fini correspondant.
Rappelons que $X_{12}^i \to X_{12}$ est l'éclatement en
$p_i^{-1}\mathcal{I}_i \mathcal{O}_{X_{12}}$, voir
Diviseurs sur les espaces, Lemme \ref{spaces-divisors-lemma-strict-transform}.
Soit $X_{12}'$ l'éclatement de $X_{12}$ en
$p_1^{-1}\mathcal{I}_1 p_2^{-1}\mathcal{I}_2 \mathcal{O}_{X_{12}}$, voir
Diviseurs sur les espaces, Lemme \ref{spaces-divisors-lemma-blowing-up-two-ideals},
pour les conséquences de cette construction. Nous obtenons un diagramme commutatif
$$
\xymatrix{
X_{12}' \ar[d] \ar[r] & X_{12}^2 \ar[d] \\
X_{12}^1 \ar[r] & X_{12}
}
$$
où tous les morphismes sont des éclatements $U$-admissibles. Puisque
$X_{12}^i \subset X_i^i$ est un ouvert, nous pouvons choisir un éclatement $U$-admissible
$X_i' \to X_i^i$ dont la restriction est $X_{12}' \to X_{12}^i$, voir
Diviseurs sur les espaces, Lemme
\ref{spaces-divisors-lemma-extend-admissible-blowups}.
Alors $X_{12}' \subset X_i'$ est un sous-espace ouvert et le diagramme
$$
\xymatrix{
X_{12}' \ar[d] \ar[r] & X_i' \ar[d] \\
X_{12}^i \ar[r] & X_i^i
}
$$
est commutatif, ses flèches verticales étant des éclatements et ses flèches horizontales
des immersions ouvertes. Notons que $X'_{12} \to X_1' \times_Y X_2'$ est
une immersion propre (utiliser le fait que $X'_{12} \to X_{12}$ est propre,
que $X_{12} \to X_1 \times_Y X_2$ est fermé, que $X_1' \times_Y X_2' \to
X_1 \times_Y X_2$ est séparé, et appliquer Morphismes d'espaces, Lemme
\ref{spaces-morphisms-lemma-universally-closed-permanence}).
Ainsi $X'_{12} \to  X_1' \times_Y X_2'$ est une immersion fermée.
Si nous définissons $X$ en recollant $X_1'$ et $X_2'$ le long du sous-espace
ouvert commun $X_{12}'$, alors $X \to Y$ est de type fini et
séparé\footnote{En effet, nous pouvons vérifier que la diagonale
$X \to X \times_Y X$ est fermée sur les quatre parties ouvertes $X'_i \times_Y X'_j$
de $X \times_Y X$, où cela est clair.}. Comme les composées
d'éclatements $U$-admissibles sont des éclatements $U$-admissibles
(Diviseurs sur les espaces, Lemme
\ref{spaces-divisors-lemma-composition-admissible-blowups})
le lemme est démontré.
\end{proof}

\begin{lemma}
\label{lemma-blowup-to-find-embedding}
Soit $S$ un schéma. Soit $f : X \to Y$ un morphisme d'espaces algébriques
sur $S$. Soit $U \subset X$ un sous-espace ouvert. Supposons que
\begin{enumerate}
\item $U$ soit quasi-compact,
\item $Y$ soit quasi-compact et quasi-séparé,
\item il existe une immersion $U \to \mathbf{P}^n_Y$ sur $Y$,
\item $f$ soit de type fini et séparé.
\end{enumerate}
Alors il existe un diagramme commutatif
$$
\xymatrix{
& U \ar[ld] \ar[d] \ar[rd] \ar[rrd] \\
X \ar[rd] & X' \ar[l] \ar[d] \ar[r] & Z' \ar[ld] \ar[r] & Z \ar[ld] \\
& Y & \mathbf{P}^n_Y \ar[l]
}
$$
où
les flèches de source $U$ sont des immersions ouvertes,
$X' \to X$ est un éclatement $U$-admissible,
$X' \to Z'$ est une immersion ouverte,
$Z' \to Y$ est un morphisme propre et représentable d'espaces algébriques.
Plus précisément, $Z' \to Z$ est un éclatement $U$-admissible
et $Z \to \mathbf{P}^n_Y$ est une immersion fermée.
\end{lemma}

\begin{proof}
Soit $Z \subset \mathbf{P}^n_Y$ l'image schématique de
l'immersion $U \to \mathbf{P}^n_Y$. Puisque $U \to \mathbf{P}^n_Y$
est quasi-compact, nous voyons que $U \subset Z$ est un sous-espace
ouvert (schématiquement) dense
(Morphismes d'espaces, Lemme
\ref{spaces-morphisms-lemma-quasi-compact-immersion}).
Appliquons le Lemme \ref{lemma-find-common-blowups} pour trouver un diagramme
$$
\xymatrix{
X' \ar[d] \ar[r] & \overline{X}' & Z' \ar[l] \ar[d] \\
X & U \ar[l] \ar[lu] \ar[u] \ar[ru] \ar[r] & Z
}
$$
ayant les propriétés énumérées dans l'énoncé de ce lemme.
Comme $X' \to X$ et $Z' \to Z$ sont des éclatements $U$-admissibles,
nous trouvons que $U$ est un ouvert schématiquement dense de
$X'$ et de $Z'$ (voir Diviseurs sur les espaces, Lemmes
\ref{spaces-divisors-lemma-blowing-up-gives-effective-Cartier-divisor} et
\ref{spaces-divisors-lemma-complement-effective-Cartier-divisor}).
Puisque $Z' \to Z \to Y$ est propre, nous voyons que $Z' \subset \overline{X}'$
est un sous-espace fermé (voir Morphismes d'espaces, Lemme
\ref{spaces-morphisms-lemma-universally-closed-permanence}).
Il s'ensuit que $X' \subset Z'$ (au sens schématique) ; ainsi $X'$
est un sous-espace ouvert de $Z'$ (un petit détail est omis), ce qui démontre le lemme.
\end{proof}

\begin{lemma}
\label{lemma-chow-noetherian}
Soit $S$ un schéma. Soit $f : X \to Y$ un morphisme d'espaces algébriques
sur $S$. Supposons que $f$ soit séparé et de type fini et que $Y$ soit noethérien.
Alors il existe un sous-espace ouvert dense $U \subset X$ et
un diagramme commutatif
$$
\xymatrix{
& U \ar[ld] \ar[d] \ar[rd] \ar[rrd] \\
X \ar[rd] & X' \ar[l] \ar[d] \ar[r] & Z' \ar[ld] \ar[r] & Z \ar[ld] \\
& Y & \mathbf{P}^n_Y \ar[l]
}
$$
où
les flèches de source $U$ sont des immersions ouvertes,
$X' \to X$ est un éclatement $U$-admissible,
$X' \to Z'$ est une immersion ouverte,
$Z' \to Y$ est un morphisme propre et représentable d'espaces algébriques.
Plus précisément, $Z' \to Z$ est un éclatement $U$-admissible
et $Z \to \mathbf{P}^n_Y$ est une immersion fermée.
\end{lemma}

\begin{proof}
D'après Limites d'espaces, Lemme
\ref{spaces-limits-lemma-embedding-into-affine-over-qs}
il existe un sous-espace ouvert dense $U \subset X$ et une immersion
$U \to \mathbf{A}^n_Y$ sur $Y$. En composant avec l'immersion ouverte
$\mathbf{A}^n_Y \to \mathbf{P}^n_Y$, nous obtenons la situation du
Lemme \ref{lemma-blowup-to-find-embedding}, d'où le
résultat.
\end{proof}

\begin{remark}
\label{remark-chow-Noetherian}
Dans les Lemmes \ref{lemma-blowup-to-find-embedding} et \ref{lemma-chow-noetherian},
le morphisme $g : Z' \to Y$ est une composée de morphismes projectifs.
On peut présumer (d'après l'analogue pour les espaces algébriques de
Morphismes, Lemme \ref{morphisms-lemma-ample-composition})
qu'il existe un faisceau inversible $g$-ample sur $Z'$.
Si nous en avons un jour besoin, nous l'énoncerons et le démontrerons ici.
\end{remark}

\noindent
Le résultat suivant est \cite[IV, théorème 3.1]{Kn}. Notons que l'immersion
$X' \to \mathbf{P}^n_Y$ est quasi-compacte et peut donc se factoriser en
$X' \to Z' \to \mathbf{P}^n_Y$, où le premier morphisme est une
immersion ouverte et le second une immersion fermée
(Morphismes d'espaces, Lemme
\ref{spaces-morphisms-lemma-quasi-compact-immersion}).

\begin{lemma}[Lemme de Chow]
\label{lemma-chow-noetherian-separated}
\begin{reference}
\cite[IV, théorème 3.1]{Kn}
\end{reference}
Soit $S$ un schéma. Soit $f : X \to Y$ un morphisme d'espaces algébriques
sur $S$. Supposons que $f$ soit séparé et de type fini, et que $Y$ soit séparé et
noethérien. Alors il existe un diagramme commutatif
$$
\xymatrix{
X \ar[rd] & X' \ar[l] \ar[d] \ar[r] & \mathbf{P}^n_Y \ar[ld] \\
& Y
}
$$
où $X' \to X$ est un éclatement $U$-admissible pour un certain ouvert dense
$U \subset X$, et le morphisme $X' \to \mathbf{P}^n_Y$ est une immersion.
\end{lemma}

\begin{proof}
Dans ce premier paragraphe de la démonstration, nous ramenons le lemme au cas
où $Y$ est de type fini sur $\Spec(\mathbf{Z})$.
Nous pouvons remplacer, et remplaçons, le schéma de base $S$ par $\Spec(\mathbf{Z})$.
Nous pouvons écrire $Y = \lim Y_i$ comme limite projective d'espaces algébriques
séparés et de type fini sur $\Spec(\mathbf{Z})$, voir
Limites d'espaces, Proposition \ref{spaces-limits-proposition-approximate} et
Lemme \ref{spaces-limits-lemma-descend-separated}.
Pour tout $i$ assez grand, nous pouvons trouver un morphisme séparé et de type fini
$X_i \to Y_i$ tel que $X = Y \times_{Y_i} X_i$, voir
Limites d'espaces, Lemmes
\ref{spaces-limits-lemma-descend-finite-presentation} et
\ref{spaces-limits-lemma-descend-separated-morphism}.
Soient $\eta_1, \ldots, \eta_n$ les points génériques des composantes
irréductibles de $|X|$ ($X$ est noethérien en tant qu'espace algébrique séparé
et de type fini sur l'espace algébrique noethérien $Y$ ; par conséquent,
$|X|$ est un espace topologique noethérien).
D'après Limites d'espaces, Lemme \ref{spaces-limits-lemma-topology-limit},
nous trouvons que les images de $\eta_1, \ldots, \eta_n$ dans $|X_i|$
sont distinctes pour $i$ assez grand. Nous pouvons remplacer
$X_i$ par l'image schématique du morphisme (quasi-compact, et même affine)
$X \to X_i$.
Après ce remplacement, nous voyons que les images
de $\eta_1, \ldots, \eta_n$ dans $|X_i|$ sont les points génériques des
composantes irréductibles de $|X_i|$, voir
Morphismes d'espaces, Lemme
\ref{spaces-morphisms-lemma-quasi-compact-scheme-theoretic-image}.
Cela étant dit, supposons que nous puissions trouver un diagramme
$$
\xymatrix{
X_i \ar[rd] & X_i' \ar[l] \ar[d] \ar[r] & \mathbf{P}^n_{Y_i} \ar[ld] \\
& Y
}
$$
où $X_i' \to X_i$ est un éclatement $U_i$-admissible pour un certain ouvert dense
$U_i \subset X_i$, et où le morphisme $X_i' \to \mathbf{P}^n_{Y_i}$
est une immersion. Alors le transformé strict $X' \to X$ de $X$ relativement
à $X_i' \to X_i$ est un éclatement $U$-admissible, où $U \subset X$
est l'image inverse de $U_i$ dans $X$. Grâce à notre choix soigneux de
l'indice $i$, il s'ensuit que $\eta_1, \ldots, \eta_n \in |U|$ et que
$U \subset X$ est dense. De plus, $X' \to \mathbf{P}^n_Y$ est une
immersion, car $X'$ est fermé dans
$X_i' \times_{X_i} X = X_i' \times_{Y_i} Y$
qui est muni d'une immersion dans $\mathbf{P}^n_Y$. Nous sommes ainsi ramenés
à la situation du paragraphe suivant.

\medskip\noindent
Supposons que $Y$ soit séparé et de type fini sur $\Spec(\mathbf{Z})$.
Alors $X \to \Spec(\mathbf{Z})$ est lui aussi séparé et de type fini.
Nous appliquons le Lemme \ref{lemma-chow-noetherian} à $X \to \Spec(\mathbf{Z})$
pour trouver un sous-espace ouvert dense $U \subset X$ et un diagramme commutatif
$$
\xymatrix{
& U \ar[ld] \ar[d] \ar[rd] \ar[rrd] \\
X \ar[rd] & X' \ar[l] \ar[d] \ar[r] & Z' \ar[ld] \ar[r] & Z \ar[ld] \\
& \Spec(\mathbf{Z}) & \mathbf{P}^n_\mathbf{Z} \ar[l]
}
$$
ayant toutes les propriétés énumérées dans le lemme.
Notons que $Z$ possède un faisceau inversible ample, à savoir
$\mathcal{O}_{\mathbf{P}^n}(1)|_Z$. Par conséquent, $Z' \to Z$
est un morphisme H-projectif d'après Morphismes, Lemme
\ref{morphisms-lemma-projective-over-quasi-projective-is-H-projective}.
Il s'ensuit que $Z' \to \Spec(\mathbf{Z})$ est H-projectif
d'après Morphismes, Lemme \ref{morphisms-lemma-H-projective-composition}.
Il existe donc une immersion fermée
$Z' \to \mathbf{P}^m_{\Spec(\mathbf{Z})}$ pour un certain $m \geq 0$.
Il s'ensuit que le morphisme diagonal
$$
X' \to Y \times \mathbf{P}^m_\mathbf{Z} = \mathbf{P}^m_Y
$$
est une immersion (car sa composée avec la projection
sur $\mathbf{P}^m_\mathbf{Z}$ est une immersion), et nous avons gagné.
\end{proof}






\section{Variantes du lemme de Chow}
\label{section-chows-lemma}

\noindent
Dans cette section, nous démontrons plusieurs variantes du lemme de Chow qui portent
sur des morphismes entre espaces algébriques non noethériens.
Les versions noethériennes sont le
Lemme \ref{lemma-chow-noetherian}
et le
Lemme \ref{lemma-chow-noetherian-separated}.

\begin{lemma}
\label{lemma-chow-finite-type}
Soit $S$ un schéma. Soit $Y$ un espace algébrique quasi-compact et quasi-séparé
sur $S$. Soit $f : X \to Y$ un morphisme séparé de
type fini. Alors il existe un diagramme commutatif
$$
\xymatrix{
X \ar[rd] & X' \ar[l] \ar[d] \ar[r] & \overline{X}' \ar[ld] \\
& Y
}
$$
où $X' \to X$ est propre surjectif,
$X' \to \overline{X}'$ est une immersion ouverte, et
$\overline{X}' \to Y$ est un morphisme propre et représentable
d'espaces algébriques.
\end{lemma}

\begin{proof}
D'après
Limites d'espaces, Proposition
\ref{spaces-limits-proposition-separated-closed-in-finite-presentation}
nous pouvons trouver une immersion fermée $X \to X_1$ où $X_1$ est séparé
et de présentation finie sur $Y$. Si nous démontrons l'assertion
pour $X_1 \to Y$, il est clair que le résultat en découle pour $X$. Nous pouvons donc supposer que
$X$ est de présentation finie sur $Y$.

\medskip\noindent
Nous pouvons remplacer, et remplaçons, le schéma de base $S$ par $\Spec(\mathbf{Z})$.
Écrivons $Y = \lim_i Y_i$ comme limite projective d'espaces algébriques
quasi-séparés et de type fini sur $\Spec(\mathbf{Z})$, voir
Limites d'espaces,
Proposition \ref{spaces-limits-proposition-approximate}.
D'après
Limites d'espaces,
Lemme \ref{spaces-limits-lemma-descend-finite-presentation},
nous pouvons
trouver un indice $i \in I$ et un schéma $X_i \to Y_i$ de présentation finie
tel que $X = Y \times_{Y_i} X_i$.
D'après
Limites d'espaces,
Lemme \ref{spaces-limits-lemma-descend-separated-morphism},
nous pouvons supposer que $X_i \to Y_i$ est séparé.
Si nous démontrons l'assertion pour
$X_i$ sur $Y_i$, il est clair que l'assertion vaut pour $X$. Le cas
$X_i \to Y_i$ est traité par le
Lemme \ref{lemma-chow-noetherian}.
\end{proof}

\begin{lemma}
\label{lemma-chow-finite-type-separated}
Soit $S$ un schéma. Soit $f : X \to Y$ un morphisme d'espaces algébriques
sur $S$. Supposons que $f$ soit séparé et de type fini, et que $Y$ soit séparé et
quasi-compact. Alors il existe un diagramme commutatif
$$
\xymatrix{
X \ar[rd] & X' \ar[l] \ar[d] \ar[r] & \mathbf{P}^n_Y \ar[ld] \\
& Y
}
$$
où $X' \to X$ est un morphisme propre surjectif
et le morphisme $X' \to \mathbf{P}^n_Y$ est une immersion.
\end{lemma}

\begin{proof}
D'après
Limites d'espaces, Proposition
\ref{spaces-limits-proposition-separated-closed-in-finite-presentation}
nous pouvons trouver une immersion fermée $X \to X_1$ où $X_1$ est séparé
et de présentation finie sur $Y$. Si nous démontrons l'assertion
pour $X_1 \to Y$, il est clair que le résultat en découle pour $X$. Nous pouvons donc supposer que
$X$ est de présentation finie sur $Y$.

\medskip\noindent
Nous pouvons remplacer, et remplaçons, le schéma de base $S$ par $\Spec(\mathbf{Z})$.
Écrivons $Y = \lim_i Y_i$ comme limite projective d'espaces algébriques
quasi-séparés et de type fini sur $\Spec(\mathbf{Z})$, voir
Limites d'espaces,
Proposition \ref{spaces-limits-proposition-approximate}.
D'après Limites d'espaces, Lemme \ref{spaces-limits-lemma-descend-separated},
nous pouvons supposer que $Y_i$ est séparé pour tout $i$.
D'après
Limites d'espaces,
Lemme \ref{spaces-limits-lemma-descend-finite-presentation},
nous pouvons
trouver un indice $i \in I$ et un schéma $X_i \to Y_i$ de présentation finie
tel que $X = Y \times_{Y_i} X_i$.
D'après
Limites d'espaces,
Lemme \ref{spaces-limits-lemma-descend-separated-morphism},
nous pouvons supposer que $X_i \to Y_i$ est séparé.
Si nous démontrons l'assertion pour
$X_i$ sur $Y_i$, il est clair que l'assertion vaut pour $X$. Le cas
$X_i \to Y_i$ est traité par le
Lemme \ref{lemma-chow-noetherian-separated}.
\end{proof}






\section{Théorème d'existence de Grothendieck}
\label{section-existence-proper}

\noindent
Dans cette section, nous étudions le théorème d'existence de Grothendieck pour les
espaces algébriques. Au lieu de développer une théorie des « espaces algébriques formels »,
nous introduisons temporairement un peu de langage qui remplace la notion de
« module cohérent sur un espace formel adique noethérien ».

\medskip\noindent
Soit $S$ un schéma. Soit $X$ un espace algébrique noethérien sur $S$.
Soit $\mathcal{I} \subset \mathcal{O}_X$ un faisceau quasi-cohérent d'idéaux.
Nous considérerons ci-dessous des systèmes projectifs $(\mathcal{F}_n)$ de
$\mathcal{O}_X$-modules cohérents tels que
\begin{enumerate}
\item $\mathcal{F}_n$ soit annulé par $\mathcal{I}^n$, et
\item les morphismes de transition induisent des isomorphismes
$\mathcal{F}_{n + 1}/\mathcal{I}^n\mathcal{F}_{n + 1} \to \mathcal{F}_n$.
\end{enumerate}
Un morphisme $\alpha : (\mathcal{F}_n) \to (\mathcal{G}_n)$
entre de tels systèmes projectifs est simplement un système compatible de morphismes
$\alpha_n : \mathcal{F}_n \to \mathcal{G}_n$.
Désignons la catégorie de ces systèmes projectifs par
$\textit{Coh}(X, \mathcal{I})$. Nous allons développer pour
ces systèmes une théorie parallèle aux
résultats correspondants dans le cas des schémas, voir
Cohomologie des schémas, sections \ref{coherent-section-existence},
\ref{coherent-section-existence-proper},
\ref{coherent-section-existence-proper-support} et
\ref{coherent-section-algebraization}.

\medskip\noindent
Fonctorialité. Soit $f : X \to Y$ un
morphisme d'espaces algébriques noethériens sur un schéma $S$, et soit
$\mathcal{J} \subset \mathcal{O}_Y$ un faisceau quasi-cohérent
d'idéaux. Posons $\mathcal{I} = f^{-1}\mathcal{J}\mathcal{O}_X$.
Dans cette situation, il existe un foncteur
$$
f^* : \textit{Coh}(Y, \mathcal{J}) \longrightarrow \textit{Coh}(X, \mathcal{I})
$$
qui envoie $(\mathcal{G}_n)$ sur $(f^*\mathcal{G}_n)$. Comparer avec
Cohomologie des schémas, Lemme \ref{coherent-lemma-inverse-systems-pullback}.
Si $f$ est étale, nous pouvons voir ce foncteur simplement comme la restriction
du système à $X$, voir Propriétés des espaces, 
Équation \ref{spaces-properties-equation-restrict-modules}.

\medskip\noindent
Descente étale. Soit $S$ un schéma. Soit $U_0 \to X$ un morphisme étale
surjectif d'espaces algébriques noethériens. Posons
$U_1 = U_0 \times_X U_0$ et $U_2 = U_0 \times_X U_0 \times_X U_0$.
Soit $\mathcal{I} \subset \mathcal{O}_{X}$ un faisceau quasi-cohérent
d'idéaux. Posons $\mathcal{I}_i = \mathcal{I}|_{U_i}$. Dans cette situation,
nous obtenons un diagramme de catégories
$$
\xymatrix{
\textit{Coh}(X, \mathcal{I}) \ar[r] &
\textit{Coh}(U_0, \mathcal{I}_0) \ar@<0.5ex>[r] \ar@<-0.5ex>[r] &
\textit{Coh}(U_1, \mathcal{I}_1) \ar@<1ex>[r] \ar[r] \ar@<-1ex>[r] &
\textit{Coh}(U_2, \mathcal{I}_2)
}
$$
et la première flèche présente $\textit{Coh}(X, \mathcal{I})$ comme la
limite homotopique de la partie droite du diagramme. Plus précisément, étant donnée
une {\it donnée de descente}, c'est-à-dire un couple $((\mathcal{G}_n), \varphi)$ où
$(\mathcal{G}_n)$ est un objet de $\textit{Coh}(U_0, \mathcal{I}_0)$ et
$\varphi : \text{pr}_0^*(\mathcal{G}_n) \to \text{pr}_1^*(\mathcal{G}_n)$
est un isomorphisme dans $\textit{Coh}(U_1, \mathcal{I}_1)$
qui satisfait la condition de cocycle dans $\textit{Coh}(U_2, \mathcal{I}_2)$,
alors il existe un unique objet $(\mathcal{F}_n)$ de
$\textit{Coh}(X, \mathcal{I})$ dont la donnée de descente canonique associée
est isomorphe à $((\mathcal{G}_n), \varphi)$. Comparer avec
Descente sur les espaces, Définition
\ref{spaces-descent-definition-descent-datum-effective-quasi-coherent}.
La démonstration de cet énoncé résulte immédiatement de l'application de
Descente sur les espaces, Proposition
\ref{spaces-descent-proposition-fpqc-descent-quasi-coherent}
aux données de descente $(\mathcal{G}_n, \varphi_n)$ lorsque $n$ varie.

\begin{lemma}
\label{lemma-inverse-systems-abelian}
Soit $S$ un schéma. Soit $X$ un espace algébrique noethérien sur $S$ et soit
$\mathcal{I} \subset \mathcal{O}_X$ un faisceau quasi-cohérent d'idéaux.
\begin{enumerate}
\item La catégorie $\textit{Coh}(X, \mathcal{I})$ est abélienne.
\item L'exactitude dans $\textit{Coh}(X, \mathcal{I})$
peut se vérifier localement pour la topologie étale.
\item Pour tout morphisme plat $f : X' \to X$ d'espaces algébriques noethériens,
le foncteur $f^* : \textit{Coh}(X, \mathcal{I}) \to
\textit{Coh}(X', f^{-1}\mathcal{I}\mathcal{O}_{X'})$ est exact.
\end{enumerate}
\end{lemma}

\begin{proof}
Preuve de (1). Choisissons un schéma affine $U_0$ et un morphisme étale surjectif
$U_0 \to X$. Posons $U_1 = U_0 \times_X U_0$ et
$U_2 = U_0 \times_X U_0 \times_X U_0$ comme dans l'étude de la
descente étale ci-dessus. Les catégories $\textit{Coh}(U_i, \mathcal{I}_i)$
sont abéliennes
(Cohomologie des schémas, Lemme \ref{coherent-lemma-inverse-systems-abelian})
et les foncteurs image inverse sont les foncteurs exacts
$\textit{Coh}(U_0, \mathcal{I}_0) \to \textit{Coh}(U_1, \mathcal{I}_1)$
et
$\textit{Coh}(U_1, \mathcal{I}_1) \to \textit{Coh}(U_2, \mathcal{I}_2)$
(Cohomologie des schémas, Lemme \ref{coherent-lemma-inverse-systems-pullback}).
Le lemme résulte alors formellement de la description de
$\textit{Coh}(X, \mathcal{I})$ comme catégorie de données de descente.
Nous omettons certains détails ; comparer avec la démonstration de
Groupoïdes, Lemme \ref{groupoids-lemma-abelian}.

\medskip\noindent
La partie (2) résulte immédiatement de l'étude du paragraphe précédent.
Dans la situation de (3), choisissons un diagramme commutatif
$$
\xymatrix{
U' \ar[d] \ar[r] & U \ar[d] \\
X' \ar[r] & X
}
$$
où $U'$ et $U$ sont des schémas affines et les morphismes verticaux sont
étales surjectifs. Alors $U' \to U$ est un morphisme plat de schémas
noethériens (Morphismes d'espaces, Lemme \ref{spaces-morphisms-lemma-flat-local}),
de sorte que le foncteur image inverse
$\textit{Coh}(U, \mathcal{I}\mathcal{O}_U) \to
\textit{Coh}(U', \mathcal{I}\mathcal{O}_{U'})$
est exact d'après
Cohomologie des schémas, Lemme \ref{coherent-lemma-inverse-systems-pullback}.
Puisque l'exactitude dans $\textit{Coh}(X, \mathcal{O}_X)$ se vérifie
sur $U$, et de même pour $X', U'$, l'assertion en résulte.
\end{proof}

\begin{lemma}
\label{lemma-inverse-systems-surjective}
Soit $S$ un schéma. Soit $X$ un espace algébrique noethérien sur $S$
et soit $\mathcal{I} \subset \mathcal{O}_X$ un faisceau quasi-cohérent
d'idéaux. Un morphisme $(\mathcal{F}_n) \to (\mathcal{G}_n)$ est surjectif dans
$\textit{Coh}(X, \mathcal{I})$ si et seulement si
$\mathcal{F}_1 \to \mathcal{G}_1$ est surjectif.
\end{lemma}

\begin{proof}
On peut le vérifier sur un recouvrement étale affine de $X$ d'après le
Lemme \ref{lemma-inverse-systems-abelian}.
On se ramène ainsi au cas des schémas, qui est le
Lemme \ref{coherent-lemma-inverse-systems-surjective} de Cohomologie des schémas.
\end{proof}

\noindent
Soit $S$ un schéma. Soit $X$ un espace algébrique noethérien sur $S$ et
soit $\mathcal{I} \subset \mathcal{O}_X$ un faisceau quasi-cohérent d'idéaux.
Il existe un foncteur
\begin{equation}
\label{equation-completion-functor}
\textit{Coh}(\mathcal{O}_X) \longrightarrow \textit{Coh}(X, \mathcal{I}), \quad
\mathcal{F} \longmapsto \mathcal{F}^\wedge
\end{equation}
qui associe au $\mathcal{O}_X$-module cohérent $\mathcal{F}$
l'objet $\mathcal{F}^\wedge = (\mathcal{F}/\mathcal{I}^n\mathcal{F})$
de $\textit{Coh}(X, \mathcal{I})$.

\begin{lemma}
\label{lemma-exact}
Le foncteur (\ref{equation-completion-functor}) est exact.
\end{lemma}

\begin{proof}
Il suffit de le vérifier localement pour la topologie étale sur $X$, voir le
Lemme \ref{lemma-inverse-systems-abelian}.
On se ramène ainsi au cas des schémas, qui est le
Lemme \ref{coherent-lemma-exact} de Cohomologie des schémas.
\end{proof}

\begin{lemma}
\label{lemma-completion-internal-hom}
Soit $S$ un schéma. Soit $X$ un espace algébrique noethérien sur $S$ et
soit $\mathcal{I} \subset \mathcal{O}_X$ un faisceau quasi-cohérent d'idéaux.
Soient $\mathcal{F}$ et $\mathcal{G}$ des $\mathcal{O}_X$-modules cohérents. Posons
$\mathcal{H} = \SheafHom_{\mathcal{O}_X}(\mathcal{F}, \mathcal{G})$.
Alors
$$
\lim H^0(X, \mathcal{H}/\mathcal{I}^n\mathcal{H}) =
\Mor_{\textit{Coh}(X, \mathcal{I})}
(\mathcal{F}^\wedge, \mathcal{G}^\wedge).
$$
\end{lemma}

\begin{proof}
Puisque $\mathcal{H}$ est un faisceau sur $X_\etale$ et que l'on dispose de la
descente étale pour les objets de $\textit{Coh}(X, \mathcal{I})$,
il suffit de le démontrer localement pour la topologie étale.
On se ramène ainsi au cas des schémas, qui est le Lemme
\ref{coherent-lemma-completion-internal-hom} de Cohomologie des schémas.
\end{proof}

\noindent
Introduisons la situation qui nous occupera dans toute la
suite de cette section.

\begin{situation}
\label{situation-existence}
Ici, $A$ est un anneau noethérien complet pour la topologie $I$-adique.
De plus, $f : X \to \Spec(A)$ est un morphisme séparé de type fini
d'espaces algébriques et $\mathcal{I} = I\mathcal{O}_X$.
\end{situation}

\noindent
Dans cette situation, nous notons
$$
\textit{Coh}_{\text{support propre sur } A}(\mathcal{O}_X)
$$
la sous-catégorie pleine de $\textit{Coh}(\mathcal{O}_X)$
formée des $\mathcal{O}_X$-modules cohérents dont le
support est propre sur $\Spec(A)$ ou, de manière équivalente, dont le
support schématique est propre sur $\Spec(A)$, voir
Catégories dérivées des espaces, Lemme
\ref{spaces-perfect-lemma-module-support-proper-over-base}.
De même, nous notons
$$
\textit{Coh}_{\text{support propre sur } A}(X, \mathcal{I})
$$
la sous-catégorie pleine de $\textit{Coh}(X, \mathcal{I})$
formée des objets $(\mathcal{F}_n)$ tels que
le support de $\mathcal{F}_1$ soit propre sur $\Spec(A)$.
Puisque le support d'un module quotient est contenu dans le support
du module, il s'ensuit que (\ref{equation-completion-functor})
induit un foncteur
\begin{equation}
\label{equation-completion-functor-proper-over-A}
\textit{Coh}_{\text{support propre sur }A}(\mathcal{O}_X)
\longrightarrow
\textit{Coh}_{\text{support propre sur }A}(X, \mathcal{I})
\end{equation}
Notre premier résultat affirme que ce foncteur est pleinement fidèle.

\begin{lemma}
\label{lemma-fully-faithful}
Dans la Situation \ref{situation-existence}.
Soient $\mathcal{F}$ et $\mathcal{G}$ des $\mathcal{O}_X$-modules cohérents.
Supposons que l'intersection des supports de
$\mathcal{F}$ et $\mathcal{G}$ soit propre sur $\Spec(A)$. Alors l'application
$$
\Mor_{\textit{Coh}(\mathcal{O}_X)}(\mathcal{F}, \mathcal{G})
\longrightarrow
\Mor_{\textit{Coh}(X, \mathcal{I})}
(\mathcal{F}^\wedge, \mathcal{G}^\wedge)
$$
provenant de (\ref{equation-completion-functor}) est bijective.
En particulier, (\ref{equation-completion-functor-proper-over-A})
est pleinement fidèle.
\end{lemma}

\begin{proof}
Posons $\mathcal{H} = \SheafHom_{\mathcal{O}_X}(\mathcal{G}, \mathcal{F})$.
C'est un $\mathcal{O}_X$-module cohérent, car sa restriction
aux schémas étales sur $X$ est cohérente d'après
Modules, Lemme \ref{modules-lemma-internal-hom-locally-kernel-direct-sum}.
D'après le Lemme \ref{lemma-completion-internal-hom}, l'application
$$
\lim_n H^0(X, \mathcal{H}/\mathcal{I}^n\mathcal{H})
\to
\Mor_{\textit{Coh}(X, \mathcal{I})}
(\mathcal{G}^\wedge, \mathcal{F}^\wedge)
$$
est bijective. Soit $i : Z \to X$ le support schématique de
$\mathcal{H}$. Il est clair que $Z$ est un sous-espace fermé tel
que $|Z|$ soit contenu dans l'intersection des supports de $\mathcal{F}$
et $\mathcal{G}$. Par hypothèse, $Z \to \Spec(A)$ est donc propre
(voir Catégories dérivées des espaces, section
\ref{spaces-perfect-section-proper-over-base}).
Écrivons $\mathcal{H} = i_*\mathcal{H}'$ pour un certain
$\mathcal{O}_Z$-module cohérent $\mathcal{H}'$. Nous avons
$i_*(\mathcal{H}'/I^n\mathcal{H}') = \mathcal{H}/I^n\mathcal{H}$.
Nous obtenons donc
\begin{align*}
\lim_n H^0(X, \mathcal{H}/\mathcal{I}^n\mathcal{H})
& =
\lim_n H^0(Z, \mathcal{H}'/\mathcal{I}^n\mathcal{H}') \\
& =
H^0(Z, \mathcal{H}') \\
& =
H^0(X, \mathcal{H}) \\
& = 
\Mor_{\textit{Coh}(\mathcal{O}_X)}(\mathcal{F}, \mathcal{G})
\end{align*}
la deuxième égalité résultant du théorème des fonctions formelles
(Cohomologie des espaces, Lemme
\ref{spaces-cohomology-lemma-spell-out-theorem-formal-functions}).
Ceci démontre le lemme.
\end{proof}

\begin{remark}
\label{remark-inverse-systems-kernel-cokernel-annihilated-by}
Soit $S$ un schéma. Soit $X$ un espace algébrique noethérien sur $S$ et soient
$\mathcal{I}, \mathcal{K} \subset \mathcal{O}_X$ des faisceaux quasi-cohérents
d'idéaux. Soit $\alpha : (\mathcal{F}_n) \to (\mathcal{G}_n)$ un morphisme de
$\textit{Coh}(X, \mathcal{I})$.
Étant donnés un schéma affine $U = \Spec(A)$ et un morphisme étale surjectif
$U \to X$, notons $I, K \subset A$ les idéaux correspondant aux restrictions
$\mathcal{I}|_U, \mathcal{K}|_U$. Notons $\alpha_U : M \to N$ le morphisme de
$A^\wedge$-modules de type fini qui correspond à $\alpha|_U$ par le
Lemme \ref{coherent-lemma-inverse-systems-affine} de Cohomologie des schémas.
Nous affirmons que les conditions suivantes sont équivalentes :
\begin{enumerate}
\item il existe un entier $t \geq 1$ tel que
$\Ker(\alpha_n)$ et $\Coker(\alpha_n)$
soient annulés par $\mathcal{K}^t$ pour tout $n \geq 1$,
\item pour tout (ou pour un) ouvert affine $\Spec(A) = U \subset X$ comme ci-dessus,
les modules $\Ker(\alpha_U)$ et $\Coker(\alpha_U)$
soient annulés par $K^t$ pour un certain entier $t \geq 1$.
\end{enumerate}
Si ces conditions équivalentes sont satisfaites, nous dirons que $\alpha$ est un
{\it morphisme dont le noyau et le conoyau sont annulés par une puissance de
$\mathcal{K}$}. Pour l'équivalence, nous renvoyons à
Cohomologie des schémas, Remarque
\ref{coherent-remark-inverse-systems-kernel-cokernel-annihilated-by}.
\end{remark}

\begin{lemma}
\label{lemma-existence-easy}
Soit $S$ un schéma. Soit $X$ un espace algébrique noethérien sur $S$ et soit
$\mathcal{I} \subset \mathcal{O}_X$ un faisceau quasi-cohérent d'idéaux.
Soient $\mathcal{G}$ un $\mathcal{O}_X$-module cohérent, $(\mathcal{F}_n)$
un objet de $\textit{Coh}(X, \mathcal{I})$, et
$\alpha : (\mathcal{F}_n) \to \mathcal{G}^\wedge$
un morphisme dont le noyau et le conoyau sont annulés par une puissance de $\mathcal{I}$.
Alors il existe un unique triplet (à unique isomorphisme près)
$(\mathcal{F}, a, \beta)$ où
\begin{enumerate}
\item $\mathcal{F}$ est un $\mathcal{O}_X$-module cohérent,
\item $a : \mathcal{F} \to \mathcal{G}$ est un morphisme de $\mathcal{O}_X$-modules
dont le noyau et le conoyau sont annulés par une puissance de $\mathcal{I}$,
\item $\beta : (\mathcal{F}_n) \to \mathcal{F}^\wedge$ est un isomorphisme, et
\item $\alpha = a^\wedge \circ \beta$.
\end{enumerate}
\end{lemma}

\begin{proof}
L'unicité et la descente étale pour les objets de
$\textit{Coh}(X, \mathcal{I})$ et de $\textit{Coh}(\mathcal{O}_X)$
impliquent qu'il suffit de construire $(\mathcal{F}, a, \beta)$ localement
pour la topologie étale sur $X$. On se ramène ainsi au cas des schémas, qui est le
Lemme \ref{coherent-lemma-existence-easy} de Cohomologie des schémas.
\end{proof}

\begin{lemma}
\label{lemma-existence-tricky}
Dans la Situation \ref{situation-existence}. Soit $\mathcal{K} \subset \mathcal{O}_X$
un faisceau quasi-cohérent d'idéaux. Soit $X_e \subset X$ le sous-espace fermé
défini par $\mathcal{K}^e$. Soit $\mathcal{I}_e = \mathcal{I}\mathcal{O}_{X_e}$.
Soit $(\mathcal{F}_n)$ un objet de
$\textit{Coh}_{\text{support propre sur } A}(X, \mathcal{I})$.
Supposons que
\begin{enumerate}
\item le foncteur
$\textit{Coh}_{\text{support propre sur } A}(\mathcal{O}_{X_e})
\to \textit{Coh}_{\text{support propre sur } A}(X_e, \mathcal{I}_e)$
soit une équivalence pour tout $e \geq 1$, et
\item il existe un objet $\mathcal{H}$ de
$\textit{Coh}_{\text{support propre sur } A}(\mathcal{O}_X)$ et un morphisme
$\alpha : (\mathcal{F}_n) \to \mathcal{H}^\wedge$ dont le
noyau et le conoyau sont annulés par une puissance de $\mathcal{K}$.
\end{enumerate}
Alors $(\mathcal{F}_n)$ appartient à l'image essentielle de
(\ref{equation-completion-functor-proper-over-A}).
\end{lemma}

\begin{proof}
Dans cette démonstration, nous utiliserons sans autre mention que, pour une immersion
fermée $i : Z \to X$, le foncteur $i_*$ donne une équivalence entre la
catégorie des modules cohérents sur $Z$ et celle des modules cohérents sur $X$ annulés
par le faisceau d'idéaux de $Z$, voir
Cohomologie des espaces, Lemme \ref{spaces-cohomology-lemma-i-star-equivalence}.
En particulier, nous considérons
$$
\textit{Coh}_{\text{support propre sur } A}(\mathcal{O}_{X_e})
\subset
\textit{Coh}_{\text{support propre sur } A}(\mathcal{O}_X)
$$
comme la sous-catégorie pleine formée des modules annulés par
$\mathcal{K}^e$, et
$$
\textit{Coh}_{\text{support propre sur } A}(X_e, \mathcal{I}_e)
\subset
\textit{Coh}_{\text{support propre sur } A}(X, \mathcal{I})
$$
comme la sous-catégorie pleine des objets annulés par $\mathcal{K}^e$.
De plus, (1) nous dit que ces deux catégories sont équivalentes par le
foncteur (\ref{equation-completion-functor-proper-over-A}).

\medskip\noindent
En appliquant cette équivalence, nous obtenons un $\mathcal{O}_X$-module cohérent
$\mathcal{G}_e$ annulé par $\mathcal{K}^e$ et correspondant au système
$(\mathcal{F}_n/\mathcal{K}^e\mathcal{F}_n)$ de
$\textit{Coh}_{\text{support propre sur } A}(X, \mathcal{I})$. Les morphismes
$\mathcal{F}_n/\mathcal{K}^{e + 1}\mathcal{F}_n \to
\mathcal{F}_n/\mathcal{K}^e\mathcal{F}_n$ correspondent à des morphismes canoniques
$\mathcal{G}_{e + 1} \to \mathcal{G}_e$ qui induisent des isomorphismes
$\mathcal{G}_{e + 1}/\mathcal{K}^e\mathcal{G}_{e + 1} \to \mathcal{G}_e$.
Nous obtenons un objet $(\mathcal{G}_e)$ de la catégorie
$\textit{Coh}_{\text{support propre sur } A}(X, \mathcal{K})$.
Le morphisme $\alpha$ induit un système de morphismes
$$
\mathcal{F}_n/\mathcal{K}^e\mathcal{F}_n
\longrightarrow
\mathcal{H}/(\mathcal{I}^n + \mathcal{K}^e)\mathcal{H}
$$
d'où des morphismes $\mathcal{G}_e \to \mathcal{H}/\mathcal{K}^e\mathcal{H}$
(en utilisant encore l'équivalence de catégories).
Soit $t \geq 1$ un entier, qui existe par l'hypothèse (2),
tel que $\mathcal{K}^t$ annule le noyau et le conoyau de tous les morphismes
$\mathcal{F}_n \to \mathcal{H}/\mathcal{I}^n\mathcal{H}$.
Alors $\mathcal{K}^{2t}$ annule le noyau et le conoyau des morphismes
$\mathcal{F}_n/\mathcal{K}^e\mathcal{F}_n \to
\mathcal{H}/(\mathcal{I}^n + \mathcal{K}^e)\mathcal{H}$
(détails omis ; voir Cohomologie des schémas,
Remarque \ref{coherent-remark-inverse-systems-kernel-cokernel-annihilated-by}).
Nous en concluons que $\mathcal{K}^{4t}$ annule le noyau et
le conoyau des morphismes
$$
\mathcal{G}_e
\longrightarrow
\mathcal{H}/\mathcal{K}^e\mathcal{H},
$$
(détails omis ; voir Cohomologie des schémas,
Remarque \ref{coherent-remark-inverse-systems-kernel-cokernel-annihilated-by}).
Nous appliquons le Lemme \ref{lemma-existence-easy} pour obtenir un
$\mathcal{O}_X$-module cohérent $\mathcal{F}$, un morphisme
$a : \mathcal{F} \to \mathcal{H}$ et un isomorphisme
$\beta : (\mathcal{G}_e) \to (\mathcal{F}/\mathcal{K}^e\mathcal{F})$
dans $\textit{Coh}(X, \mathcal{K})$. En remontant les constructions, pour $n$ fixé,
le triplet
$(\mathcal{F}/\mathcal{I}^n\mathcal{F}, a \bmod \mathcal{I}^n, \beta
\bmod \mathcal{I}^n)$ est un triplet comme dans le lemme pour le morphisme
$\alpha_n \bmod \mathcal{K}^e :
(\mathcal{F}_n/\mathcal{K}^e\mathcal{F}_n) \to
(\mathcal{H}/(\mathcal{I}^n + \mathcal{K}^e)\mathcal{H})$
de $\textit{Coh}(X, \mathcal{K})$. L'unicité du
Lemme \ref{lemma-existence-easy}
donne donc un isomorphisme canonique
$\mathcal{F}/\mathcal{I}^n\mathcal{F} \to \mathcal{F}_n$
compatible avec tous les morphismes en présence.

\medskip\noindent
Pour achever la démonstration du lemme, il reste à montrer que le
support de $\mathcal{F}$ est propre sur $A$.
Par construction, le noyau de $a : \mathcal{F} \to \mathcal{H}$
est annulé par une puissance de $\mathcal{K}$. Par conséquent, le support de
ce noyau est contenu dans le support de $\mathcal{G}_1$. Puisque
$\mathcal{G}_1$ est un objet de
$\textit{Coh}_{\text{support propre sur } A}(\mathcal{O}_{X_1})$
nous voyons que celui-ci est propre sur $A$. Joint au fait que le
support de $\mathcal{H}$ est propre sur $A$, cela montre que le
support de $\mathcal{F}$ est propre sur $A$ d'après
Catégories dérivées des espaces, Lemme
\ref{spaces-perfect-lemma-union-closed-proper-over-base}.
\end{proof}

\begin{lemma}
\label{lemma-inverse-systems-push-pull}
Soit $S$ un schéma. Soit $f : X \to Y$ un morphisme propre
représentable d'espaces algébriques noethériens sur $S$. Soient
$\mathcal{J}, \mathcal{K} \subset \mathcal{O}_Y$
des faisceaux quasi-cohérents d'idéaux.
Supposons que $f$ soit un isomorphisme au-dessus de $V = Y \setminus V(\mathcal{K})$.
Posons $\mathcal{I} = f^{-1}\mathcal{J} \mathcal{O}_X$.
Soit $(\mathcal{G}_n)$ un objet de $\textit{Coh}(Y, \mathcal{J})$,
soit $\mathcal{F}$ un $\mathcal{O}_X$-module cohérent, et soit
$\beta : (f^*\mathcal{G}_n)  \to \mathcal{F}^\wedge$ un isomorphisme dans
$\textit{Coh}(X, \mathcal{I})$. Alors il existe un morphisme
$$
\alpha :
(\mathcal{G}_n)
\longrightarrow
(f_*\mathcal{F})^\wedge
$$
dans $\textit{Coh}(Y, \mathcal{J})$ dont le noyau et le conoyau
sont annulés par une puissance de $\mathcal{K}$.
\end{lemma}

\begin{proof}
Puisque $f$ est un morphisme propre, $f_*\mathcal{F}$ est un
$\mathcal{O}_Y$-module cohérent (Cohomologie des espaces, Lemme
\ref{spaces-cohomology-lemma-proper-pushforward-coherent}).
L'énoncé du lemme a donc un sens. Considérons les composés
$$
\gamma_n : \mathcal{G}_n \to
f_*f^*\mathcal{G}_n \to
f_*(\mathcal{F}/\mathcal{I}^n\mathcal{F}).
$$
Ici, le premier morphisme est le morphisme d'adjonction et le second est $f_*\beta_n$.
Nous affirmons qu'il existe un unique $\alpha$ comme dans le lemme
tel que les composés
$$
\mathcal{G}_n \xrightarrow{\alpha_n}
f_*\mathcal{F}/\mathcal{J}^nf_*\mathcal{F} \to
f_*(\mathcal{F}/\mathcal{I}^n\mathcal{F})
$$
soient égaux à $\gamma_n$ pour tout $n$. Par unicité et par descente étale
pour $\textit{Coh}(Y, \mathcal{J})$, il suffit de le démontrer
localement pour la topologie étale sur $Y$. Nous pouvons donc supposer que $Y$
est le spectre d'un anneau noethérien. Comme $f$ est représentable,
$X$ est lui aussi un schéma. On se ramène ainsi au cas des
schémas ; voir la démonstration de
Cohomologie des schémas, Lemme \ref{coherent-lemma-inverse-systems-push-pull}.
\end{proof}

\begin{theorem}[Théorème d'existence de Grothendieck]
\label{theorem-grothendieck-existence}
Dans la Situation \ref{situation-existence}, le foncteur
(\ref{equation-completion-functor-proper-over-A})
est une équivalence.
\end{theorem}

\begin{proof}
Nous utiliserons, sans autre mention dans la démonstration du théorème,
l'équivalence de catégories de Cohomologie des espaces, Lemme
\ref{spaces-cohomology-lemma-i-star-equivalence}.
D'après le Lemme \ref{lemma-fully-faithful}, le foncteur est pleinement fidèle.
Il reste donc à démontrer que le foncteur est essentiellement surjectif.

\medskip\noindent
Considérons l'ensemble $\Xi$ des faisceaux quasi-cohérents d'idéaux
$\mathcal{K} \subset \mathcal{O}_X$ tels que l'énoncé soit vrai
pour tout objet $(\mathcal{F}_n)$ de
$\textit{Coh}_{\text{support propre sur }A}(X, \mathcal{I})$
annulé par $\mathcal{K}$. Nous voulons montrer que $(0)$ appartient à $\Xi$.
Dans le cas contraire, puisque $X$ est noethérien, il existe un
faisceau quasi-cohérent d'idéaux maximal $\mathcal{K}$ qui n'appartient pas à $\Xi$, voir
Cohomologie des espaces, Lemme \ref{spaces-cohomology-lemma-acc-coherent}.
Après avoir remplacé $X$ par le sous-schéma fermé de $X$
correspondant à $\mathcal{K}$, nous pouvons supposer que tout
$\mathcal{K}$ non nul appartient à $\Xi$. Soit $(\mathcal{F}_n)$ un objet de
$\textit{Coh}_{\text{support propre sur }A}(X, \mathcal{I})$.
Nous allons montrer que cet objet appartient à l'image essentielle, ce qui
achèvera la démonstration du théorème.

\medskip\noindent
Appliquons le lemme de Chow (Lemme \ref{lemma-chow-noetherian-separated})
pour trouver un morphisme propre surjectif $f : Y \to X$ qui soit un isomorphisme
au-dessus d'un ouvert dense $U \subset X$, avec $Y$ H-quasi-projectif
sur $A$. Notons que $Y$ est un schéma et que $f$ est représentable.
Choisissons une immersion ouverte $j : Y \to Y'$ avec $Y'$ projectif sur $A$, voir
Morphismes, Lemme \ref{morphisms-lemma-H-quasi-projective-open-H-projective}.
Soit $T_n$ le support schématique de $\mathcal{F}_n$.
Notons que $|T_n| = |T_1|$ ; ainsi, $T_n$ est propre sur $A$ pour tout $n$
(Morphismes d'espaces, Lemme
\ref{spaces-morphisms-lemma-image-proper-is-proper}).
Alors $f^*\mathcal{F}_n$ est supporté sur le sous-schéma fermé
$f^{-1}T_n$, qui est propre sur $A$ (d'après
Morphismes d'espaces, Lemme \ref{spaces-morphisms-lemma-composition-proper},
et la propreté de $f$).
En particulier, le composé $f^{-1}T_n \to Y \to Y'$ est fermé
(Morphismes, Lemme \ref{morphisms-lemma-image-proper-scheme-closed}).
Soit $T'_n \subset Y'$ le sous-schéma fermé correspondant ;
il est contenu dans le sous-schéma ouvert $Y$ et égal à $f^{-1}T_n$
comme sous-schéma fermé de $Y$. Soit $\mathcal{F}_n'$
le $\mathcal{O}_{Y'}$-module cohérent correspondant à
$f^*\mathcal{F}_n$, considéré comme module cohérent sur $Y'$ au moyen de
l'immersion fermée $f^{-1}T_n = T'_n \subset Y'$.
Alors $(\mathcal{F}_n')$
est un objet de $\textit{Coh}(Y', I\mathcal{O}_{Y'})$.
D'après le cas projectif du théorème d'existence de Grothendieck
(Cohomologie des schémas, Lemme \ref{coherent-lemma-existence-projective}),
il existe un $\mathcal{O}_{Y'}$-module cohérent
$\mathcal{F}'$ et un isomorphisme
$(\mathcal{F}')^\wedge \cong (\mathcal{F}'_n)$ dans
$\textit{Coh}(Y', I\mathcal{O}_{Y'})$.
Soit $Z' \subset Y'$ le support schématique de $\mathcal{F}'$.
Puisque $\mathcal{F}'/I\mathcal{F}' = \mathcal{F}'_1$, nous voyons
que $Z' \cap V(I\mathcal{O}_{Y'}) = T'_1$ ensemblistement.
Le morphisme structural $p' : Y' \to \Spec(A)$ est propre ; par conséquent,
$p'(Z' \cap (Y' \setminus Y))$ est fermé dans $\Spec(A)$.
S'il n'était pas vide, il contiendrait un point de $V(I)$,
car $I$ est contenu dans le radical de Jacobson de $A$
(Algèbre, Lemme \ref{algebra-lemma-radical-completion}).
Mais nous avons vu ci-dessus que $Z' \cap (p')^{-1}V(I) = T'_1 \subset Y$ ;
nous en concluons donc que $Z' \subset Y$. Ainsi, $\mathcal{F}'|_Y$
est supporté sur un sous-schéma fermé de $Y$ propre sur $A$.

\medskip\noindent
Soit $\mathcal{K}$ le faisceau quasi-cohérent d'idéaux qui définit
le complémentaire réduit $X \setminus U$.
D'après Cohomologie des espaces, Lemme
\ref{spaces-cohomology-lemma-proper-pushforward-coherent}
le $\mathcal{O}_X$-module $\mathcal{H} = f_*\mathcal{F}'$ est cohérent,
et, d'après le Lemme \ref{lemma-inverse-systems-push-pull},
il existe un morphisme $\alpha : (\mathcal{F}_n) \to \mathcal{H}^\wedge$
dans la catégorie $\textit{Coh}_{\text{support propre sur } A}(X, \mathcal{I})$
dont le noyau et le conoyau sont annulés par une puissance de $\mathcal{K}$.
Soit $Z_0 \subset X$ le support schématique de $\mathcal{H}$.
Il est clair que $|Z_0| \subset f(|Z'|)$. Par conséquent,
$Z_0 \to \Spec(A)$ est propre
(Morphismes d'espaces, Lemme
\ref{spaces-morphisms-lemma-image-proper-is-proper}).
Ainsi, $\mathcal{H}$ est un objet de 
$\textit{Coh}_{\text{support propre sur } A}(\mathcal{O}_X)$.
Puisque chacun des faisceaux d'idéaux $\mathcal{K}^e$ est un élément de
$\Xi$, les hypothèses du Lemme \ref{lemma-existence-tricky}
sont satisfaites, ce qui permet de conclure.
\end{proof}

\begin{remark}[Explicitation du théorème d'existence de Grothendieck]
\label{remark-reformulate-existence-theorem}
Soit $A$ un anneau noethérien complet pour la topologie $I$-adique.
Écrivons $S = \Spec(A)$ et $S_n = \Spec(A/I^n)$.
Soit $X \to S$ un morphisme d'espaces algébriques séparé et
de type fini. Pour $n \geq 1$, posons $X_n = X \times_S S_n$.
Diagramme :
$$
\xymatrix{
X_1 \ar[r]_{i_1} \ar[d] & X_2 \ar[r]_{i_2} \ar[d] & X_3 \ar[r] \ar[d] &
\ldots & X \ar[d] \\
S_1 \ar[r] & S_2 \ar[r] & S_3 \ar[r] & \ldots & S
}
$$
Dans cette situation, nous considérons les systèmes $(\mathcal{F}_n, \varphi_n)$
où
\begin{enumerate}
\item $\mathcal{F}_n$ est un $\mathcal{O}_{X_n}$-module cohérent,
\item $\varphi_n : i_n^*\mathcal{F}_{n + 1} \to \mathcal{F}_n$
est un isomorphisme, et
\item $\text{Supp}(\mathcal{F}_1)$ est propre sur $S_1$.
\end{enumerate}
Le Théorème \ref{theorem-grothendieck-existence} affirme que le
foncteur de passage au complété
$$
\begin{matrix}
\text{modules cohérents sur }\mathcal{O}_X\text{ : }\mathcal{F} \\
\text{à support propre sur }A
\end{matrix}
\quad
\longrightarrow
\quad
\begin{matrix}
\text{systèmes }(\mathcal{F}_n) \\
\text{comme ci-dessus}
\end{matrix}
$$
est une équivalence de catégories. Dans le cas particulier où $X$ est
propre sur $A$, on peut omettre les conditions sur les supports.
\end{remark}







\section{Théorème d'algébrisation de Grothendieck}
\label{section-algebraization}

\noindent
Cette section est l'analogue de
Cohomologie des schémas, section \ref{coherent-section-algebraization}.
Toutefois, il y manque le résultat sur l'algébrisation des déformations
d'espaces algébriques propres munis de faisceaux inversibles amples, car un
espace algébrique propre muni d'un faisceau inversible ample est
déjà un schéma. Nous avons néanmoins un résultat d'algébrisation pour les
espaces algébriques propres de dimension relative $1$.
Notre premier résultat est une reformulation du théorème d'existence de Grothendieck
en termes de sous-schémas fermés et de morphismes finis.

\begin{lemma}
\label{lemma-algebraize-formal-closed-subscheme}
Soit $A$ un anneau noethérien complet pour la topologie $I$-adique.
Écrivons $S = \Spec(A)$ et $S_n = \Spec(A/I^n)$.
Soit $X \to S$ un morphisme d'espaces algébriques séparé
et de type fini.
Pour $n \geq 1$, posons $X_n = X \times_S S_n$.
Supposons donné un diagramme commutatif
$$
\xymatrix{
Z_1 \ar[r] \ar[d] & Z_2 \ar[r] \ar[d] & Z_3 \ar[r] \ar[d] & \ldots \\
X_1 \ar[r]^{i_1} & X_2 \ar[r]^{i_2} & X_3 \ar[r] & \ldots
}
$$
d'espaces algébriques à carrés cartésiens. Supposons que
\begin{enumerate}
\item $Z_1 \to X_1$ soit une immersion fermée, et
\item $Z_1 \to S_1$ soit propre.
\end{enumerate}
Alors il existe une immersion fermée d'espaces algébriques $Z \to X$ telle que
$Z_n = Z \times_S S_n$ pour tout $n \geq 1$. De plus, $Z$ est propre sur $S$.
\end{lemma}

\begin{proof}
Notons $j_n : Z_n \to X_n$ les morphismes verticaux.
Comme les carrés de l'énoncé sont cartésiens,
le changement de base de $j_n$ à $X_1$ est $j_1$.
Ainsi, Limites d'espaces, Lemme
\ref{spaces-limits-lemma-check-closed-infinitesimally}
montre que $j_n$ est une immersion fermée.
Posons $\mathcal{F}_n = j_{n, *}\mathcal{O}_{Z_n}$, de sorte que
$j_n^\sharp$ est une surjection $\mathcal{O}_{X_n} \to \mathcal{F}_n$.
En utilisant encore le caractère cartésien des carrés, nous voyons que
l'image inverse de $\mathcal{F}_{n + 1}$ sur $X_n$ est $\mathcal{F}_n$.
Le théorème d'existence de Grothendieck, sous la forme de la
Remarque \ref{remark-reformulate-existence-theorem},
nous dit donc qu'il existe un morphisme
$\mathcal{O}_X \to \mathcal{F}$
de $\mathcal{O}_X$-modules cohérents dont la restriction à
$X_n$ redonne $\mathcal{O}_{X_n} \to \mathcal{F}_n$.
De plus, le support de $\mathcal{F}$ est propre sur $S$.
Comme le foncteur de passage au complété est exact (Lemme \ref{lemma-exact}),
nous voyons que $\mathcal{O}_X \to \mathcal{F}$
est surjectif. Ainsi, $\mathcal{F} = \mathcal{O}_X/\mathcal{J}$
pour un certain faisceau quasi-cohérent d'idéaux $\mathcal{J}$.
Poser $Z = V(\mathcal{J})$ achève la démonstration.
\end{proof}

\begin{lemma}
\label{lemma-algebraize-formal-algebraic-space-finite-over-proper}
Soit $A$ un anneau noethérien complet pour la topologie $I$-adique.
Écrivons $S = \Spec(A)$ et $S_n = \Spec(A/I^n)$.
Soit $X \to S$ un morphisme d'espaces algébriques séparé
et de type fini.
Pour $n \geq 1$, posons $X_n = X \times_S S_n$.
Supposons donné un diagramme commutatif
$$
\xymatrix{
Y_1 \ar[r] \ar[d] & Y_2 \ar[r] \ar[d] & Y_3 \ar[r] \ar[d] & \ldots \\
X_1 \ar[r]^{i_1} & X_2 \ar[r]^{i_2} & X_3 \ar[r] & \ldots
}
$$
d'espaces algébriques à carrés cartésiens. Supposons que
\begin{enumerate}
\item $Y_1 \to X_1$ soit un morphisme fini, et
\item $Y_1 \to S_1$ soit propre.
\end{enumerate}
Alors il existe un morphisme fini d'espaces algébriques $Y \to X$ tel que
$Y_n = Y \times_S S_n$ pour tout $n \geq 1$. De plus, $Y$ est propre sur $S$.
\end{lemma}

\begin{proof}
Notons $f_n : Y_n \to X_n$ les morphismes verticaux.
Comme les carrés de l'énoncé sont cartésiens,
le changement de base de $f_n$ à $X_1$ est $f_1$.
Ainsi, le Lemme \ref{lemma-thicken-property-morphisms-cartesian}
montre que $f_n$ est un morphisme fini.
Posons $\mathcal{F}_n = f_{n, *}\mathcal{O}_{Y_n}$.
En utilisant le caractère cartésien des carrés, nous voyons que
l'image inverse de $\mathcal{F}_{n + 1}$ sur $X_n$ est $\mathcal{F}_n$.
Le théorème d'existence de Grothendieck, sous la forme de la
Remarque \ref{remark-reformulate-existence-theorem},
nous dit donc qu'il existe un $\mathcal{O}_X$-module cohérent $\mathcal{F}$
dont la restriction à $X_n$ redonne $\mathcal{F}_n$.
De plus, le support de $\mathcal{F}$ est propre sur $S$.
Comme le foncteur de passage au complété est pleinement fidèle
(Théorème \ref{theorem-grothendieck-existence}),
nous voyons que les applications de multiplication
$\mathcal{F}_n \otimes_{\mathcal{O}_{X_n}} \mathcal{F}_n \to
\mathcal{F}_n$ se recollent pour munir $\mathcal{F}$ d'une structure d'algèbre.
Poser $Y = \underline{\Spec}_X(\mathcal{F})$ achève la démonstration.
\end{proof}

\begin{lemma}
\label{lemma-algebraize-morphism}
Soit $A$ un anneau noethérien complet pour la topologie $I$-adique.
Écrivons $S = \Spec(A)$ et $S_n = \Spec(A/I^n)$. Soient $X$, $Y$ des espaces
algébriques sur $S$. Pour $n \geq 1$, posons $X_n = X \times_S S_n$ et
$Y_n = Y \times_S S_n$. Supposons donné un système compatible de
diagrammes commutatifs
$$
\xymatrix{
& & X_{n + 1} \ar[rd] \ar[rr]_{g_{n + 1}} & & Y_{n + 1} \ar[ld] \\
X_n \ar[rru] \ar[rd] \ar[rr]_{g_n} & & Y_n \ar[rru] \ar[ld] & S_{n + 1} \\
& S_n \ar[rru]
}
$$
Supposons que
\begin{enumerate}
\item $X \to S$ soit propre, et
\item $Y \to S$ soit séparé et de type fini.
\end{enumerate}
Alors il existe un unique morphisme d'espaces algébriques $g : X \to Y$
sur $S$ tel que $g_n$ soit le changement de base de $g$ à $S_n$.
\end{lemma}

\begin{proof}
Les morphismes $(1, g_n) : X_n \to X_n \times_S Y_n$ sont des immersions fermées
parce que $Y_n \to S_n$ est séparé
(Morphismes d'espaces, Lemme \ref{spaces-morphisms-lemma-section-immersion}).
Ainsi, d'après le Lemme \ref{lemma-algebraize-formal-closed-subscheme},
il existe un sous-espace fermé $Z \subset X \times_S Y$
propre sur $S$ dont le changement de base à $S_n$ redonne
$X_n \subset X_n \times_S Y_n$. La première projection $p : Z \to X$
est un morphisme propre (puisque $Z$ est propre sur $S$, voir
Morphismes d'espaces, Lemme
\ref{spaces-morphisms-lemma-universally-closed-permanence})
dont le changement de base à $S_n$ est un isomorphisme pour tout $n$.
En particulier, $p : Z \to X$ est quasi-fini sur un sous-espace ouvert
de $Z$ contenant tout point de $Z_0$, par exemple d'après
Morphismes d'espaces, Lemme
\ref{spaces-morphisms-lemma-locally-finite-type-quasi-finite-part}.
Comme $Z$ est propre sur $S$, ce voisinage ouvert est $Z$ tout entier.
Nous concluons que $p : Z \to X$ est fini par le théorème principal de Zariski
(on peut par exemple appliquer le
Lemme \ref{lemma-quasi-finite-separated-pass-through-finite}
et utiliser la propreté de $Z$ sur $X$ pour voir que l'immersion est
une immersion fermée). En appliquant l'équivalence du
Théorème \ref{theorem-grothendieck-existence},
nous voyons que $p_*\mathcal{O}_Z = \mathcal{O}_X$, puisque c'est vrai
modulo $I^n$ pour tout $n$. Ainsi, $p$ est un isomorphisme et nous obtenons
le morphisme $g$ comme composé $X \cong Z \to Y$.
Nous omettons la démonstration de l'unicité.
\end{proof}

\begin{remark}
\label{remark-weaken-separation-axioms-question}
On peut se demander si, dans le théorème d'algébrisation de Grothendieck (sous la forme
du Lemme \ref{lemma-algebraize-morphism}), on peut se contenter
d'axiomes de séparation plus faibles sur le but. Soyons plus précis.
Soient $A$, $I$, $S$, $S_n$, $X$, $Y$, $X_n$, $Y_n$ et $g_n$ comme dans
l'énoncé du Lemme \ref{lemma-algebraize-morphism}, et supposons que
\begin{enumerate}
\item $X \to S$ soit propre, et
\item $Y \to S$ soit localement de type fini.
\end{enumerate}
Existe-t-il un morphisme d'espaces algébriques $g : X \to Y$
sur $S$ tel que $g_n$ soit le changement de base de $g$ à $S_n$ ?
Nous ne connaissons pas la réponse en général ; si vous la connaissez, veuillez écrire à
\href{mailto:stacks.project@gmail.com}{stacks.project@gmail.com}.
Si $Y \to S$ est séparé, le résultat
découle du lemme (on se ramène immédiatement au cas où
$X$ est de type fini sur $S$, en choisissant un ouvert quasi-compact
contenant l'image de $g_1$). Si l'on suppose seulement que $Y \to S$
est quasi-séparé, le résultat est encore vrai. D'abord, comme précédemment,
nous pouvons supposer $Y$ quasi-compact aussi bien que quasi-séparé.
On peut alors utiliser soit \cite{Bhatt-Algebraize}, soit
\cite{Hall-Rydh-coherent} pour algébriser $(g_n)$. Plus précisément, pour appliquer
la première référence, nous utilisons
$$
D_{perf}(X) \to \lim D_{perf}(X_n)
\xrightarrow{\lim Lg_n^*}
\lim D_{perf}(Y_n) = D_{perf}(Y)
$$
où la dernière étape utilise un résultat d'existence de Grothendieck pour
la catégorie dérivée de l'espace algébrique propre $Y$ sur $R$
(comparer avec
Platitude sur les espaces, Remarque \ref{spaces-flat-remark-correct-generality}).
L'article cité montre que
cette flèche détermine un morphisme $Y \to X$ comme souhaité. Pour appliquer la
seconde référence, nous utilisons le même argument avec les modules cohérents :
$$
\textit{Coh}(\mathcal{O}_X) \to \lim \textit{Coh}(\mathcal{O}_{X_n})
\xrightarrow{\lim g_n^*}
\lim \textit{Coh}(\mathcal{O}_{Y_n}) = \textit{Coh}(\mathcal{O}_Y)
$$
où l'égalité finale est une conséquence du théorème d'existence de Grothendieck
(Théorème \ref{theorem-grothendieck-existence}).
La seconde référence nous dit que ce foncteur correspond à un
morphisme $Y \to X$ sur $R$. Si nous avons un jour besoin de cette généralisation,
nous énoncerons précisément et démontrerons soigneusement le résultat ici.
\end{remark}

\begin{lemma}
\label{lemma-formal-algebraic-space-proper-reldim-1}
Soit $(A, \mathfrak m, \kappa)$ un anneau local noethérien complet.
Posons $S = \Spec(A)$ et $S_n = \Spec(A/\mathfrak m^n)$.
Considérons un diagramme commutatif
$$
\xymatrix{
X_1 \ar[r]_{i_1} \ar[d] & X_2 \ar[r]_{i_2} \ar[d] & X_3 \ar[r] \ar[d] &
\ldots \\
S_1 \ar[r] & S_2 \ar[r] & S_3 \ar[r] & \ldots
}
$$
d'espaces algébriques à carrés cartésiens. Si $\dim(X_1) \leq 1$,
alors il existe un morphisme projectif de schémas $X \to S$
et des isomorphismes $X_n \cong X \times_S S_n$ compatibles avec $i_n$.
\end{lemma}

\begin{proof}
D'après Espaces sur les corps, Lemme
\ref{spaces-over-fields-lemma-codim-1-point-in-schematic-locus}
l'espace algébrique $X_1$ est un schéma. Ainsi, $X_1$
est un schéma propre de dimension $\leq 1$ sur $\kappa$.
D'après Variétés, Lemme \ref{varieties-lemma-dim-1-proper-projective},
nous voyons que $X_1$ est H-projectif sur $\kappa$.
Soit $\mathcal{L}_1$ un faisceau inversible ample sur $X_1$.

\medskip\noindent
Nous allons montrer que $\mathcal{L}_1$ se relève en un système compatible
$\{\mathcal{L}_n\}$ de faisceaux inversibles sur $\{X_n\}$.
Remarquons que $X_n$ est aussi un schéma d'après le Lemme \ref{lemma-thickening-scheme}.
Rappelons que $X_1 \to X_n$ induit des homéomorphismes entre les espaces
topologiques sous-jacents. Dans la suite de la démonstration, nous ne distinguons pas
les faisceaux sur $X_n$ des faisceaux sur $X_1$.
Supposons donné un relèvement $\mathcal{L}_n$ sur $X_n$. Considérons
la suite exacte
$$
1 \to
(1 + \mathfrak m^n\mathcal{O}_{X_{n + 1}})^* \to
\mathcal{O}_{X_{n + 1}}^* \to \mathcal{O}_{X_n}^* \to 1
$$
de faisceaux sur $X_{n + 1}$. La classe de $\mathcal{L}_n$ dans
$H^1(X_n, \mathcal{O}_{X_n}^*)$ (voir
Cohomologie, Lemme \ref{cohomology-lemma-h1-invertible})
se relève en un élément de $H^1(X_{n + 1}, \mathcal{O}_{X_{n + 1}}^*)$
si et seulement si l'obstruction dans
$H^2(X_{n + 1}, (1 + \mathfrak m^n\mathcal{O}_{X_{n + 1}})^*)$
est nulle. Comme $X_1$ est un schéma noethérien de dimension $\leq 1$,
ce groupe de cohomologie est nul (Cohomologie, Proposition
\ref{cohomology-proposition-vanishing-Noetherian}).

\medskip\noindent
D'après le théorème d'algébrisation de Grothendieck
(Cohomologie des schémas, Théorème \ref{coherent-theorem-algebraization}),
nous obtenons un morphisme projectif de schémas $X \to S = \Spec(A)$
et un système compatible d'isomorphismes $X_n = S_n \times_S X$.
\end{proof}

\begin{lemma}
\label{lemma-projective-over-complete}
Soit $(A, \mathfrak m, \kappa)$ un anneau local noethérien complet.
Soit $X$ un espace algébrique sur $\Spec(A)$.
Si $X \to \Spec(A)$ est propre et $\dim(X_\kappa) \leq 1$, alors
$X$ est un schéma projectif sur $A$.
\end{lemma}

\begin{proof}
Posons $X_n = X \times_{\Spec(A)} \Spec(A/\mathfrak m^n)$.
D'après le Lemme \ref{lemma-formal-algebraic-space-proper-reldim-1},
il existe un morphisme projectif $Y \to \Spec(A)$
et des isomorphismes compatibles
$Y \times_{\Spec(A)} \Spec(A/\mathfrak m^n) \cong
X \times_{\Spec(A)} \Spec(A/\mathfrak m^n)$.
D'après le Lemme \ref{lemma-algebraize-morphism},
nous voyons que $X \cong Y$, ce qui achève la démonstration.
\end{proof}










\section{Immersions régulières}
\label{section-regular-immersions}

\noindent
Cette section est l'analogue de
Diviseurs, section \ref{divisors-section-regular-immersions}
pour les morphismes d'espaces algébriques. Le lecteur est invité à lire d'abord
dans cette section ce qui concerne les immersions régulières de schémas.

\medskip\noindent
Dans
Diviseurs, section \ref{divisors-section-regular-immersions},
nous avons défini quatre types d'immersions régulières pour les morphismes de schémas.
Parmi eux, seuls trois sont (à notre connaissance) locaux sur le but pour
la topologie étale ; comme d'habitude, les immersions régulières au sens usuel ne le sont pas.
C'est pourquoi, pour les morphismes d'espaces algébriques, nous ne pouvons pas à proprement parler définir
les immersions régulières. (Ce genre de désagrément a conduit Grothendieck
et son école à remplacer la notion initiale d'immersion régulière par celle
d'immersion Koszul-régulière, voir
\cite[Exposé VII, définition 1.4]{SGA6}.)
Mais nous pouvons définir les immersions Koszul-régulières, $H_1$-régulières et quasi-régulières.
Remarquons aussi que, puisque les immersions Koszul-régulières ne sont pas préservées par
un changement de base arbitraire, nous ne pouvons pas employer la méthode de
Morphismes d'espaces, section \ref{spaces-morphisms-section-representable}
pour les définir. De même, comme les immersions Koszul-régulières ne sont pas locales
sur la source pour la topologie étale, nous ne pouvons pas non plus utiliser
Morphismes d'espaces, Lemme \ref{spaces-morphisms-lemma-local-source-target}
pour les définir. Nous remplaçons plutôt ce lemme par le
suivant.

\begin{lemma}
\label{lemma-representable-etale-local-target}
Soit $\mathcal{P}$ une propriété des morphismes de schémas qui est locale
sur le but pour la topologie étale. Soit $S$ un schéma.
Soit $f : X \to Y$ un morphisme représentable d'espaces algébriques sur $S$.
Considérons les diagrammes commutatifs
$$
\xymatrix{
X \times_Y V \ar[d] \ar[r] & V \ar[d] \\
X \ar[r]^f & Y
}
$$
où $V$ est un schéma et $V \to Y$ est étale.
Les conditions suivantes sont équivalentes :
\begin{enumerate}
\item pour tout diagramme comme ci-dessus, la projection $X \times_Y V \to V$
possède la propriété $\mathcal{P}$, et
\item pour un diagramme comme ci-dessus où $V \to Y$ est surjectif,
la projection $X \times_Y V \to V$ possède la propriété $\mathcal{P}$.
\end{enumerate}
Si $X$ et $Y$ sont représentables, ces conditions équivalent aussi à ce que
$f$ (considéré comme morphisme de schémas) possède la propriété $\mathcal{P}$.
\end{lemma}

\begin{proof}
Démontrons l'équivalence de (1) et (2).
L'implication (1) $\Rightarrow$ (2) est immédiate.
Supposons que
$$
\xymatrix{
X \times_Y V \ar[d] \ar[r] & V \ar[d] \\
X \ar[r]^f & Y
}
\quad\quad
\xymatrix{
X \times_Y V' \ar[d] \ar[r] & V' \ar[d] \\
X \ar[r]^f & Y
}
$$
soient deux diagrammes comme dans le lemme. Supposons que $V \to Y$ soit
surjectif et que $X \times_Y V \to V$ possède la propriété $\mathcal{P}$.
Pour montrer que (2) implique (1), nous devons prouver que
$X \times_Y V' \to V'$ possède $\mathcal{P}$. À cette
fin, considérons le diagramme
$$
\xymatrix{
X \times_Y V \ar[d] &
(X \times_Y V) \times_X (X \times_Y V') \ar[l] \ar[d] \ar[r] &
X \times_Y V' \ar[d] \\
V &
V \times_Y V' \ar[l] \ar[r] &
V'
}
$$
D'après notre hypothèse que $\mathcal{P}$ est locale sur la source pour la topologie étale,
nous voyons que $\mathcal{P}$ est préservée par changement de base étale, voir
Descente, Lemme \ref{descent-lemma-pullback-property-local-target}.
Ainsi, si la flèche verticale de gauche possède $\mathcal{P}$, il en est de même
de la flèche verticale du milieu. Puisque $U \times_X U' \to U'$ est surjectif
et étale (et définit donc un recouvrement étale de $U'$),
il en résulte (car $\mathcal{P}$ est supposée locale sur le but pour la topologie
étale) que la flèche verticale de gauche possède $\mathcal{P}$.

\medskip\noindent
Si $X$ et $Y$ sont représentables, nous pouvons prendre
$\text{id}_Y : Y \to Y$ comme recouvrement étale pour voir que la
dernière assertion du lemme est vraie.
\end{proof}

\noindent
Remarquons qu'« être une immersion Koszul-régulière (resp.\ $H_1$-régulière,
resp.\ quasi-régulière) » est une propriété des morphismes de schémas qui est locale
sur le but pour la topologie fpqc, voir
Descente, Lemme \ref{descent-lemma-descending-property-regular-immersion}.
La définition suivante a donc un sens.

\begin{definition}
\label{definition-regular-immersion}
Soit $S$ un schéma. Soit $i : X \to Y$ un morphisme d'espaces
algébriques sur $S$.
\begin{enumerate}
\item Nous disons que $i$ est une {\it immersion Koszul-régulière} si $i$ est représentable
et si les conditions équivalentes du
Lemme \ref{lemma-representable-etale-local-target}
sont satisfaites avec $\mathcal{P}(f) =$ « $f$ est une immersion Koszul-régulière ».
\item Nous disons que $i$ est une {\it immersion $H_1$-régulière} si $i$ est représentable
et si les conditions équivalentes du
Lemme \ref{lemma-representable-etale-local-target}
sont satisfaites avec $\mathcal{P}(f) =$ « $f$ est une immersion $H_1$-régulière ».
\item Nous disons que $i$ est une {\it immersion quasi-régulière} si $i$ est représentable
et si les conditions équivalentes du
Lemme \ref{lemma-representable-etale-local-target}
sont satisfaites avec $\mathcal{P}(f) =$ « $f$ est une immersion quasi-régulière ».
\end{enumerate}
\end{definition}

\begin{lemma}
\label{lemma-regular-quasi-regular-immersion}
Soit $S$ un schéma.
Soit $i : Z \to X$ une immersion d'espaces algébriques sur $S$.
Nous avons les implications suivantes :
$i$ est une immersion Koszul-régulière $\Rightarrow$
$i$ est une immersion $H_1$-régulière $\Rightarrow$
$i$ est une immersion quasi-régulière.
\end{lemma}

\begin{proof}
Par la définition, ce lemme se réduit immédiatement à
Diviseurs, Lemme \ref{divisors-lemma-regular-quasi-regular-immersion}.
\end{proof}

\begin{lemma}
\label{lemma-regular-immersion-noetherian}
Soit $S$ un schéma.
Soit $i : Z \to X$ une immersion d'espaces algébriques sur $S$.
Supposons que $X$ soit localement noethérien. Alors
$i$ est une immersion Koszul-régulière $\Leftrightarrow$
$i$ est une immersion $H_1$-régulière $\Leftrightarrow$
$i$ est une immersion quasi-régulière.
\end{lemma}

\begin{proof}
Par la Définition \ref{definition-regular-immersion}
(et la définition d'un espace algébrique localement noethérien
dans Propriétés des espaces, section
\ref{spaces-properties-section-types-properties})
ceci se ramène immédiatement au cas des schémas, qui est
Diviseurs, Lemme \ref{divisors-lemma-regular-immersion-noetherian}.
\end{proof}

\begin{lemma}
\label{lemma-flat-base-change-regular-immersion}
\begin{slogan}
Les immersions régulières sont stables par changement de base plat.
\end{slogan}
Soit $S$ un schéma. Soit $i : Z \to X$ une immersion Koszul-régulière,
$H_1$-régulière ou quasi-régulière d'espaces algébriques sur $S$.
Soit $X' \to X$ un morphisme plat d'espaces algébriques sur $S$.
Alors le changement de base $i' : Z \times_X X' \to X'$ est une immersion
Koszul-régulière, $H_1$-régulière ou quasi-régulière.
\end{lemma}

\begin{proof}
Par la Définition \ref{definition-regular-immersion}
(et la définition d'un morphisme plat d'espaces algébriques
dans Morphismes d'espaces, section
\ref{spaces-morphisms-section-flat})
ce lemme se réduit au cas des schémas, voir
Diviseurs, Lemme \ref{divisors-lemma-flat-base-change-regular-immersion}.
\end{proof}

\begin{lemma}
\label{lemma-quasi-regular-immersion}
Soit $S$ un schéma. Soit $i : Z \to X$ une immersion d'espaces algébriques
sur $S$. Alors $i$ est une immersion quasi-régulière si et seulement si les
conditions suivantes sont satisfaites :
\begin{enumerate}
\item $i$ est localement de présentation finie,
\item le faisceau conormal $\mathcal{C}_{Z/X}$ est localement libre de type fini, et
\item l'application (\ref{equation-conormal-algebra-quotient}) est un isomorphisme.
\end{enumerate}
\end{lemma}

\begin{proof}
Cela découle du cas des schémas
(Diviseurs, Lemme \ref{divisors-lemma-quasi-regular-immersion})
par localisation étale (utiliser la Définition \ref{definition-regular-immersion}
et le
Lemme \ref{lemma-etale-conormal-algebra}).
\end{proof}

\begin{lemma}
\label{lemma-transitivity-conormal-quasi-regular}
Soit $S$ un schéma. Soient $Z \to Y \to X$ des immersions d'espaces algébriques
sur $S$. Supposons que $Z \to Y$ soit une immersion $H_1$-régulière. Alors la suite canonique
du Lemme \ref{lemma-transitivity-conormal}
$$
0 \to i^*\mathcal{C}_{Y/X} \to
\mathcal{C}_{Z/X} \to
\mathcal{C}_{Z/Y} \to 0
$$
est exacte et localement scindée (pour la topologie étale).
\end{lemma}

\begin{proof}
Comme $\mathcal{C}_{Z/Y}$ est localement libre de type fini (voir le
Lemme \ref{lemma-quasi-regular-immersion}
et le
Lemme \ref{lemma-regular-quasi-regular-immersion}),
il suffit de prouver que la suite est exacte.
Il suffit de montrer que la première application est injective,
puisque la suite est déjà exacte à droite en général.
Après localisation étale sur $X$, cela se réduit au cas
des schémas, voir
Diviseurs, Lemme \ref{divisors-lemma-transitivity-conormal-quasi-regular}.
\end{proof}

\noindent
Un composé d'immersions quasi-régulières peut ne pas être une immersion quasi-régulière, voir
Algèbre, Remarque \ref{algebra-remark-join-quasi-regular-sequences}.
Les autres types d'immersions régulières sont préservés par composition.

\begin{lemma}
\label{lemma-composition-regular-immersion}
Soit $S$ un schéma. Soient $i : Z \to Y$ et $j : Y \to X$ des immersions
d'espaces algébriques sur $S$.
\begin{enumerate}
\item Si $i$ et $j$ sont des immersions Koszul-régulières, il en est de même de $j \circ i$.
\item Si $i$ et $j$ sont des immersions $H_1$-régulières, il en est de même de $j \circ i$.
\item Si $i$ est une immersion $H_1$-régulière et $j$ une immersion
quasi-régulière, alors $j \circ i$ est une immersion quasi-régulière.
\end{enumerate}
\end{lemma}

\begin{proof}
Cela découle immédiatement du cas des schémas, voir
Diviseurs, Lemme \ref{divisors-lemma-composition-regular-immersion}.
\end{proof}

\begin{lemma}
\label{lemma-permanence-regular-immersion}
Soit $S$ un schéma. Soient $i : Z \to Y$ et $j : Y \to X$ des immersions
d'espaces algébriques sur $S$. Supposons que $j$ soit localement de présentation finie
et que la suite
$$
0 \to i^*\mathcal{C}_{Y/X} \to
\mathcal{C}_{Z/X} \to
\mathcal{C}_{Z/Y} \to 0
$$
du Lemme \ref{lemma-transitivity-conormal} soit exacte et localement scindée.
\begin{enumerate}
\item Si $j \circ i$ est une immersion quasi-régulière, il en est de même de $i$.
\item Si $j \circ i$ est une immersion $H_1$-régulière, il en est de même de $i$.
\item Si $j$ et $j \circ i$ sont toutes deux des immersions Koszul-régulières, il en est de même de $i$.
\end{enumerate}
\end{lemma}

\begin{proof}
Cela découle immédiatement du cas des schémas, voir
Diviseurs, Lemme \ref{divisors-lemma-permanence-regular-immersion}.
\end{proof}

\begin{lemma}
\label{lemma-extra-permanence-regular-immersion-noetherian}
Soit $S$ un schéma. Soient $i : Z \to Y$ et $j : Y \to X$ des immersions
d'espaces algébriques sur $S$. Supposons que $X$ soit localement noethérien.
Les conditions suivantes sont équivalentes :
\begin{enumerate}
\item $i$ et $j$ sont des immersions Koszul-régulières,
\item $i$ et $j \circ i$ sont des immersions Koszul-régulières,
\item $j \circ i$ est une immersion Koszul-régulière et la suite conormale
$$
0 \to i^*\mathcal{C}_{Y/X} \to
\mathcal{C}_{Z/X} \to
\mathcal{C}_{Z/Y} \to 0
$$
est exacte et localement scindée.
\end{enumerate}
\end{lemma}

\begin{proof}
Cela découle immédiatement du cas des schémas, voir Diviseurs, Lemme
\ref{divisors-lemma-extra-permanence-regular-immersion-noetherian}.
\end{proof}














\section{Pseudo-cohérence relative}
\label{section-relative-pseudo-coherence}

\noindent
Cette section est l'analogue de
Compléments sur les morphismes, section
\ref{more-morphisms-section-relative-pseudo-coherence}.
Toutefois, dans le traitement de cette matière pour les
espaces algébriques, nous avons décidé de travailler exclusivement
avec les objets de la catégorie dérivée dont les faisceaux de cohomologie
sont quasi-cohérents. Il y a deux raisons à cela :
(1) cela simplifie grandement l'exposé et
(2) nous n'avons actuellement aucun usage de la notion plus générale.

\begin{remark}
\label{remark-match-relative-pseudo-coherence}
Soit $S$ un schéma. Soit $f : X \to Y$ un morphisme entre des espaces algébriques
représentables sur $S$, localement de type fini. Soit
$f_0 : X_0 \to Y_0$ un morphisme de schémas représentant $f$
(notation maladroite, mais temporaire). Alors $f_0$ est localement de type fini.
Si $E$ est un objet de $D_\QCoh(\mathcal{O}_X)$, alors $E$
est l'image inverse d'un unique objet $E_0$ de $D_\QCoh(\mathcal{O}_{X_0})$, voir
Catégories dérivées des espaces, Lemme
\ref{spaces-perfect-lemma-derived-quasi-coherent-small-etale-site}.
Dans cette situation, l'expression « $E$ est $m$-pseudo-cohérent relativement à $Y$ »
signifiera « $E_0$ est $m$-pseudo-cohérent relativement à $Y_0$ »,
au sens défini dans Compléments sur les morphismes, section
\ref{more-morphisms-section-relative-pseudo-coherence}.
\end{remark}

\begin{lemma}
\label{lemma-qcoh-relative-pseudo-coherence-characterize}
Soit $S$ un schéma. Soit $f : X \to Y$ un morphisme d'espaces algébriques
sur $S$, localement de type fini. Soit $m \in \mathbf{Z}$.
Soit $E \in D_\QCoh(\mathcal{O}_X)$. Avec les conventions expliquées dans la
Remarque \ref{remark-match-relative-pseudo-coherence},
les conditions suivantes sont équivalentes :
\begin{enumerate}
\item pour tout diagramme commutatif
$$
\xymatrix{
U \ar[d] \ar[r] & V \ar[d] \\
X \ar[r] & Y
}
$$
où $U$, $V$ sont des schémas et les flèches verticales sont étales, le complexe
$E|_U$ est $m$-pseudo-cohérent relativement à $V$,
\item pour un diagramme commutatif comme en (1) où $U \to X$ est
surjectif, le complexe $E|_U$ est $m$-pseudo-cohérent relativement à $V$,
\item pour tout diagramme commutatif comme en (1) où $U$ et $V$ sont
affines, le complexe $R\Gamma(U, E)$ de $\mathcal{O}_X(U)$-modules
est $m$-pseudo-cohérent relativement à $\mathcal{O}_Y(V)$.
\end{enumerate}
\end{lemma}

\begin{proof}
La partie (1) implique (3) d'après Compléments sur les morphismes, Lemme
\ref{more-morphisms-lemma-qcoh-relative-pseudo-coherence-characterize}.

\medskip\noindent
Supposons (3). Choisissons un diagramme commutatif comme en (1) où $U \to X$ est surjectif.
Choisissons un recouvrement par des ouverts affines $V = \bigcup V_j$ et des recouvrements par des ouverts affines
$(U \to V)^{-1}(V_j) = \bigcup U_{ij}$. D'après (3) et Compléments sur les morphismes, Lemme
\ref{more-morphisms-lemma-qcoh-relative-pseudo-coherence-characterize}
nous voyons que $E|_U$ est $m$-pseudo-cohérent relativement à $V$.
Ainsi (3) implique (2).

\medskip\noindent
Supposons (2). Choisissons un diagramme commutatif
$$
\xymatrix{
U \ar[d] \ar[r] & V \ar[d] \\
X \ar[r] & Y
}
$$
où $U$, $V$ sont des schémas, les flèches verticales sont étales, le
morphisme $U \to X$ est surjectif et $E|_U$ est $m$-pseudo-cohérent
relativement à $V$. Supposons ensuite donné un second diagramme commutatif
$$
\xymatrix{
U' \ar[d] \ar[r] & V' \ar[d] \\
X \ar[r] & Y
}
$$
à flèches verticales étales, où $U', V'$ sont des schémas. Nous voulons montrer
que $E|_{U'}$ est $m$-pseudo-cohérent relativement à $V'$.
Le morphisme $U'' = U \times_X U' \to U'$ est surjectif et étale
et $U'' \to V'$ se factorise par $V'' = V' \times_Y V$, qui
est étale sur $V'$. Il suffit donc de montrer que $E|_{U''}$
est $m$-pseudo-cohérent relativement à $V''$, voir
Compléments sur les morphismes, Lemmes
\ref{more-morphisms-lemma-relative-pseudo-coherent-descends-fppf} et
\ref{more-morphisms-lemma-relative-pseudo-coherent-post-compose}.
En utilisant une fois encore le second lemme, il suffit de montrer que
$E|_{U''}$ est $m$-pseudo-cohérent relativement à $V$.
C'est vrai d'après Compléments sur les morphismes, Lemme
\ref{more-morphisms-lemma-pull-relative-pseudo-coherent}
et le fait qu'un morphisme étale de schémas est pseudo-cohérent d'après
Compléments sur les morphismes, Lemme
\ref{more-morphisms-lemma-flat-finite-presentation-pseudo-coherent}.
\end{proof}

\begin{definition}
\label{definition-relative-pseudo-coherence}
Soit $S$ un schéma. Soit $f : X \to Y$ un morphisme d'espaces
algébriques sur $S$, localement de type fini.
Soit $E$ un objet de $D_\QCoh(\mathcal{O}_X)$. Soit $\mathcal{F}$ un
$\mathcal{O}_X$-module quasi-cohérent. Fixons $m \in \mathbf{Z}$.
\begin{enumerate}
\item Nous disons que $E$ est {\it $m$-pseudo-cohérent relativement à $Y$}
si les conditions équivalentes du
Lemme \ref{lemma-qcoh-relative-pseudo-coherence-characterize} sont satisfaites.
\item Nous disons que $E$ est {\it pseudo-cohérent relativement à $Y$}
si $E$ est $m$-pseudo-cohérent relativement à $Y$ pour tout $m \in \mathbf{Z}$.
\item Nous disons que $\mathcal{F}$ est {\it $m$-pseudo-cohérent relativement à $Y$} si
$\mathcal{F}$, considéré comme objet de $D_\QCoh(\mathcal{O}_X)$, est
$m$-pseudo-cohérent relativement à $Y$.
\item Nous disons que $\mathcal{F}$ est {\it pseudo-cohérent relativement à $Y$} si
$\mathcal{F}$, considéré comme objet de $D_\QCoh(\mathcal{O}_X)$, est
pseudo-cohérent relativement à $Y$.
\end{enumerate}
\end{definition}

\noindent
La plupart des propriétés des complexes pseudo-cohérents relativement à une base
découleront immédiatement des propriétés correspondantes dans le cas
des schémas. Nous ajouterons ici les lemmes pertinents au besoin.

\begin{lemma}
\label{lemma-relative-pseudo-coherent-is-moot}
Soit $S$ un schéma. Soit $f : X \to Y$ un morphisme d'espaces
algébriques sur $S$. Soit $E$ un objet de $D_\QCoh(\mathcal{O}_X)$.
Si $f$ est plat et localement de présentation finie, alors
les conditions suivantes sont équivalentes :
\begin{enumerate}
\item $E$ est pseudo-cohérent relativement à $Y$, et
\item $E$ est pseudo-cohérent sur $X$.
\end{enumerate}
\end{lemma}

\begin{proof}
Par localisation étale et d'après les définitions, nous pouvons supposer
que $X$ et $Y$ sont des schémas. Pour le cas des schémas, cela découle de
Compléments sur les morphismes, Lemme
\ref{more-morphisms-lemma-check-relative-pseudo-coherence-on-charts}.
\end{proof}

















\section{Morphismes pseudo-cohérents}
\label{section-pseudo-coherent}

\noindent
Cette section est l'analogue de
Compléments sur les morphismes, section \ref{more-morphisms-section-pseudo-coherent}
pour les morphismes de schémas. Le lecteur est invité à lire d'abord
dans cette section ce qui concerne les morphismes pseudo-cohérents de schémas.

\medskip\noindent
La propriété « pseudo-cohérent » des morphismes de schémas est
de nature locale pour la topologie étale sur la source et le but. Pour le voir, utiliser
Compléments sur les morphismes,
Lemmes \ref{more-morphisms-lemma-descending-property-pseudo-coherent} et
\ref{more-morphisms-lemma-pseudo-coherent-fppf-local-source}
et
Descente, Lemme \ref{descent-lemma-etale-local-source-target}.
D'après
Morphismes d'espaces,
Lemme \ref{spaces-morphisms-lemma-local-source-target},
nous pouvons définir la notion de morphisme pseudo-cohérent d'espaces algébriques comme
suit ; elle coïncide avec la notion déjà définie dans
Compléments sur les morphismes, section \ref{more-morphisms-section-pseudo-coherent}
lorsque les espaces algébriques considérés sont représentables.

\begin{definition}
\label{definition-pseudo-coherent}
Soit $S$ un schéma.
Soit $f : X \to Y$ un morphisme d'espaces algébriques sur $S$.
\begin{enumerate}
\item Nous disons que $f$ est {\it pseudo-cohérent} si les conditions équivalentes de
Morphismes d'espaces, Lemme \ref{spaces-morphisms-lemma-local-source-target}
sont satisfaites avec $\mathcal{P} =$ « pseudo-cohérent ».
\item Soit $x \in |X|$. Nous disons que $f$ est {\it pseudo-cohérent en $x$} s'il
existe un voisinage ouvert $X' \subset X$ de $x$ tel
que $f|_{X'} : X' \to Y$ soit pseudo-cohérent.
\end{enumerate}
\end{definition}

\noindent
Attention : un changement de base d'un morphisme pseudo-cohérent n'est pas
pseudo-cohérent en général.

\begin{lemma}
\label{lemma-flat-base-change-pseudo-coherent}
Un changement de base plat d'un morphisme pseudo-cohérent est pseudo-cohérent.
\end{lemma}

\begin{proof}
Démonstration omise. Indication : utiliser la version de ce lemme pour les schémas, voir
Compléments sur les morphismes,
Lemme \ref{more-morphisms-lemma-flat-base-change-pseudo-coherent}.
\end{proof}

\begin{lemma}
\label{lemma-composition-pseudo-coherent}
Un composé de morphismes pseudo-cohérents est pseudo-cohérent.
\end{lemma}

\begin{proof}
Démonstration omise. Indication : utiliser la version de ce lemme pour les schémas, voir
Compléments sur les morphismes,
Lemme \ref{more-morphisms-lemma-composition-pseudo-coherent}.
\end{proof}

\begin{lemma}
\label{lemma-pseudo-coherent-finite-presentation}
Un morphisme pseudo-cohérent est localement de présentation finie.
\end{lemma}

\begin{proof}
Cela résulte immédiatement des définitions.
\end{proof}

\begin{lemma}
\label{lemma-flat-finite-presentation-pseudo-coherent}
Un morphisme plat qui est localement de présentation finie est pseudo-cohérent.
\end{lemma}

\begin{proof}
Démonstration omise. Indication : utiliser la version de ce lemme pour les schémas, voir
Compléments sur les morphismes,
Lemme \ref{more-morphisms-lemma-flat-finite-presentation-pseudo-coherent}.
\end{proof}

\begin{lemma}
\label{lemma-permanence-pseudo-coherent}
Soit $f : X \to Y$ un morphisme d'espaces algébriques pseudo-cohérent
au-dessus d'un espace algébrique de base $B$. Alors $f$ est pseudo-cohérent.
\end{lemma}

\begin{proof}
Démonstration omise. Indication : utiliser la version de ce lemme pour les schémas, voir
Compléments sur les morphismes,
Lemme \ref{more-morphisms-lemma-permanence-pseudo-coherent}.
\end{proof}

\begin{lemma}
\label{lemma-Noetherian-pseudo-coherent}
Soit $S$ un schéma. Soit $f : X \to Y$ un morphisme d'espaces algébriques
sur $S$. Si $Y$ est localement noethérien, alors $f$ est pseudo-cohérent si
et seulement si $f$ est localement de type fini.
\end{lemma}

\begin{proof}
Démonstration omise. Indication : utiliser la version de ce lemme pour les schémas, voir
Compléments sur les morphismes,
Lemme \ref{more-morphisms-lemma-Noetherian-pseudo-coherent}.
\end{proof}








\section{Morphismes parfaits}
\label{section-perfect}

\noindent
Cette section est l'analogue de
Compléments sur les morphismes, section \ref{more-morphisms-section-perfect}
pour les morphismes de schémas. Le lecteur est invité à lire d'abord
dans cette section ce qui concerne les morphismes parfaits de schémas.

\medskip\noindent
La propriété « parfait » des morphismes de schémas est
de nature locale pour la topologie étale sur la source et le but. Pour le voir, utiliser
Compléments sur les morphismes,
Lemmes \ref{more-morphisms-lemma-descending-property-perfect} et
\ref{more-morphisms-lemma-perfect-fppf-local-source}
et
Descente, Lemme \ref{descent-lemma-etale-local-source-target}.
D'après
Morphismes d'espaces,
Lemme \ref{spaces-morphisms-lemma-local-source-target},
nous pouvons définir la notion de morphisme parfait d'espaces algébriques comme
suit ; elle coïncide avec la notion déjà définie dans
Compléments sur les morphismes, section \ref{more-morphisms-section-perfect}
lorsque les espaces algébriques considérés sont représentables.

\begin{definition}
\label{definition-perfect}
Soit $S$ un schéma.
Soit $f : X \to Y$ un morphisme d'espaces algébriques sur $S$.
\begin{enumerate}
\item Nous disons que $f$ est {\it parfait} si les conditions équivalentes de
Morphismes d'espaces, Lemme \ref{spaces-morphisms-lemma-local-source-target}
sont satisfaites avec $\mathcal{P} =$ « parfait ».
\item Soit $x \in |X|$. Nous disons que $f$ est {\it parfait en $x$} s'il
existe un voisinage ouvert $X' \subset X$ de $x$ tel
que $f|_{X'} : X' \to Y$ soit parfait.
\end{enumerate}
\end{definition}

\noindent
Remarquons qu'un morphisme parfait est pseudo-cohérent, donc localement de
présentation finie. Attention : un changement de base d'un morphisme parfait n'est pas parfait
en général.

\begin{lemma}
\label{lemma-flat-base-change-perfect}
Un changement de base plat d'un morphisme parfait est parfait.
\end{lemma}

\begin{proof}
Démonstration omise. Indication : utiliser la version de ce lemme pour les schémas, voir
Compléments sur les morphismes,
Lemme \ref{more-morphisms-lemma-flat-base-change-perfect}.
\end{proof}

\begin{lemma}
\label{lemma-composition-perfect}
Un composé de morphismes parfaits est parfait.
\end{lemma}

\begin{proof}
Démonstration omise. Indication : utiliser la version de ce lemme pour les schémas, voir
Compléments sur les morphismes,
Lemme \ref{more-morphisms-lemma-composition-perfect}.
\end{proof}

\begin{lemma}
\label{lemma-flat-finite-presentation-perfect}
Soit $S$ un schéma. Soit $f : X \to Y$ un morphisme d'espaces algébriques sur
$S$. Les conditions suivantes sont équivalentes :
\begin{enumerate}
\item $f$ est plat et parfait, et
\item $f$ est plat et localement de présentation finie.
\end{enumerate}
\end{lemma}

\begin{proof}
Démonstration omise. Indication : utiliser la version de ce lemme pour les schémas, voir
Compléments sur les morphismes,
Lemme \ref{more-morphisms-lemma-flat-finite-presentation-perfect}.
\end{proof}

\begin{lemma}
\label{lemma-perfect-proper-perfect-direct-image}
Soit $S$ un schéma. Soit $Y$ un espace algébrique noethérien sur $S$.
Soit $f : X \to Y$ un morphisme propre parfait d'espaces algébriques.
Soit $E \in D(\mathcal{O}_X)$ un objet parfait. Alors
$Rf_*E$ est un objet parfait de $D(\mathcal{O}_Y)$.
\end{lemma}

\begin{proof}
Nous affirmons que Catégories dérivées des espaces, Lemme
\ref{spaces-perfect-lemma-perfect-direct-image} s'applique.
Les conditions (1) et (2) sont immédiates. La condition (3) est locale
sur $X$. Nous pouvons donc supposer $X$ et $Y$ affines et $E$
représenté par un complexe strictement parfait de $\mathcal{O}_X$-modules.
Il suffit donc de montrer que $\mathcal{O}_X$ est de dimension de Tor
finie en tant que faisceau de $f^{-1}\mathcal{O}_Y$-modules
sur le site étale. D'après Catégories dérivées des espaces, Lemme
\ref{spaces-perfect-lemma-tor-dimension-rel}, il suffit de
le vérifier sur le site de Zariski. Pour les morphismes de schémas de type fini,
cela équivaut à la perfection du morphisme d'après Compléments sur les morphismes,
Lemme \ref{more-morphisms-lemma-check-perfect-stalks}.
\end{proof}











\section{Morphismes localement d'intersection complète}
\label{section-lci}

\noindent
Cette section est l'analogue de
Compléments sur les morphismes, section \ref{more-morphisms-section-lci},
pour les morphismes de schémas. Le lecteur est invité à lire d'abord
dans cette section ce qui concerne les morphismes de schémas localement d'intersection complète.

\medskip\noindent
La propriété « être un morphisme localement d'intersection complète » des morphismes
de schémas est de nature locale pour la topologie étale sur la source et le but. Pour le voir, utiliser
Compléments sur les morphismes,
Lemmes \ref{more-morphisms-lemma-descending-property-lci} et
\ref{more-morphisms-lemma-lci-syntomic-local-source}
et
Descente, Lemme \ref{descent-lemma-etale-local-source-target}.
D'après
Morphismes d'espaces,
Lemme \ref{spaces-morphisms-lemma-local-source-target},
nous pouvons définir la notion de morphisme localement d'intersection complète
d'espaces algébriques comme suit ; elle coïncide avec la notion déjà
définie dans
Compléments sur les morphismes, section \ref{more-morphisms-section-lci},
lorsque les espaces algébriques considérés sont représentables.

\begin{definition}
\label{definition-lci}
Soit $S$ un schéma.
Soit $f : X \to Y$ un morphisme d'espaces algébriques sur $S$.
\begin{enumerate}
\item Nous disons que $f$ est un {\it morphisme de Koszul}, ou que $f$ est un
{\it morphisme localement d'intersection complète}, si les conditions équivalentes de
Morphismes d'espaces, Lemme \ref{spaces-morphisms-lemma-local-source-target},
sont satisfaites avec $\mathcal{P}(f) =$ « $f$ est un morphisme localement d'intersection complète ».
\item Soit $x \in |X|$. Nous disons que $f$ est {\it Koszul en $x$} s'il
existe un voisinage ouvert $X' \subset X$ de $x$ tel
que $f|_{X'} : X' \to Y$ soit un morphisme localement d'intersection complète.
\end{enumerate}
\end{definition}

\noindent
En un certain sens, la propriété qui définit un morphisme localement
d'intersection complète est le résultat du lemme suivant.

\begin{lemma}
\label{lemma-lci}
Soit $S$ un schéma.
Soit $f : X \to Y$ un morphisme localement d'intersection complète
d'espaces algébriques sur $S$.
Soit $P$ un espace algébrique lisse sur $Y$.
Soit $U \to X$ un morphisme étale d'espaces algébriques
et soit $i : U \to P$ une immersion d'espaces algébriques sur $Y$.
Diagramme :
$$
\xymatrix{
X \ar[rd] & U \ar[l] \ar[d] \ar[r]_i & P \ar[ld] \\
& Y
}
$$
Alors $i$ est une immersion Koszul-régulière d'espaces algébriques.
\end{lemma}

\begin{proof}
Choisissons un schéma $V$ et un morphisme étale surjectif $V \to Y$.
Choisissons un schéma $W$ et un morphisme étale surjectif $W \to P \times_Y V$.
Posons $U' = U \times_P W$ ; c'est un schéma étale sur $U$.
Nous devons montrer que $U' \to W$ est une immersion Koszul-régulière de
schémas, voir
Définition \ref{definition-regular-immersion}.
D'après la
Définition \ref{definition-lci}
ci-dessus, le morphisme de schémas $U' \to V$ est localement d'intersection
complète. Le résultat découle donc de
Compléments sur les morphismes, Lemme \ref{more-morphisms-lemma-lci}.
\end{proof}

\noindent
Il paraît opportun de rassembler ici quelques propriétés communes
à tous les morphismes de Koszul.

\begin{lemma}
\label{lemma-lci-properties}
Soit $S$ un schéma. Soit $f : X \to Y$ un morphisme localement d'intersection
complète d'espaces algébriques sur $S$. Alors
\begin{enumerate}
\item $f$ est localement de présentation finie,
\item $f$ est pseudo-cohérent, et
\item $f$ est parfait.
\end{enumerate}
\end{lemma}

\begin{proof}
Démonstration omise. Indication : utiliser la version de ce lemme pour les schémas, voir
Compléments sur les morphismes,
Lemme \ref{more-morphisms-lemma-lci-properties}.
\end{proof}

\noindent
Attention : un changement de base d'un morphisme de Koszul n'est pas Koszul en général.

\begin{lemma}
\label{lemma-flat-base-change-lci}
Un changement de base plat d'un morphisme localement d'intersection complète est un
morphisme localement d'intersection complète.
\end{lemma}

\begin{proof}
Démonstration omise. Indication : utiliser la version de ce lemme pour les schémas, voir
Compléments sur les morphismes,
Lemme \ref{more-morphisms-lemma-flat-base-change-lci}.
\end{proof}

\begin{lemma}
\label{lemma-composition-lci}
Un composé de morphismes localement d'intersection complète est un
morphisme localement d'intersection complète.
\end{lemma}

\begin{proof}
Démonstration omise. Indication : utiliser la version de ce lemme pour les schémas, voir
Compléments sur les morphismes,
Lemme \ref{more-morphisms-lemma-composition-lci}.
\end{proof}

\begin{lemma}
\label{lemma-flat-lci}
\begin{slogan}
Syntomique équivaut à plat et localement d'intersection complète (pour les espaces algébriques).
\end{slogan}
Soit $S$ un schéma.
Soit $f : X \to Y$ un morphisme d'espaces algébriques sur $S$.
Les conditions suivantes sont équivalentes :
\begin{enumerate}
\item $f$ est plat et localement d'intersection complète, et
\item $f$ est syntomique.
\end{enumerate}
\end{lemma}

\begin{proof}
Démonstration omise. Indication : utiliser la version de ce lemme pour les schémas, voir
Compléments sur les morphismes,
Lemme \ref{more-morphisms-lemma-flat-lci}.
\end{proof}

\begin{lemma}
\label{lemma-regular-immersion-lci}
Soit $S$ un schéma. Une immersion Koszul-régulière d'espaces algébriques
sur $S$ est un morphisme localement d'intersection complète.
\end{lemma}

\begin{proof}
Soit $i : X \to Y$ une immersion Koszul-régulière d'espaces algébriques
sur $S$. Par définition, il existe un morphisme étale surjectif
$V \to Y$, où $V$ est un schéma, tel que $X \times_Y V$ soit un schéma
et que le changement de base $X \times_Y V \to V$ soit une immersion Koszul-régulière de
schémas. D'après Compléments sur les morphismes, Lemme
\ref{more-morphisms-lemma-regular-immersion-lci},
$X \times_Y V \to V$ est un morphisme localement d'intersection complète.
La Définition \ref{definition-lci} montre alors que $i$ est un
morphisme localement d'intersection complète d'espaces algébriques.
\end{proof}

\begin{lemma}
\label{lemma-lci-permanence}
Soit $S$ un schéma. Soit
$$
\xymatrix{
X \ar[rr]_f \ar[rd] & & Y \ar[ld] \\
& Z
}
$$
un diagramme commutatif de morphismes d'espaces algébriques sur $S$.
Supposons que $Y \to Z$ soit lisse et que $X \to Z$ soit un
morphisme localement d'intersection complète.
Alors $f : X \to Y$ est un morphisme localement d'intersection complète.
\end{lemma}

\begin{proof}
Choisissons un schéma $W$ et un morphisme étale surjectif $W \to Z$.
Choisissons un schéma $V$ et un morphisme étale surjectif $V \to W \times_Z Y$.
Choisissons un schéma $U$ et un morphisme étale surjectif $U \to V \times_Y X$.
Alors $U \to W$ est un morphisme de schémas localement d'intersection complète et
$V \to W$ est un morphisme de schémas lisse. Le résultat pour les schémas
(Compléments sur les morphismes, Lemme \ref{more-morphisms-lemma-lci-permanence})
montre que $U \to V$ est un morphisme localement d'intersection complète.
Par définition, cela signifie que $f$ est un morphisme localement d'intersection complète.
\end{proof}

\begin{lemma}
\label{lemma-descending-property-lci}
La propriété $\mathcal{P}(f) =$ « $f$ est un morphisme localement
d'intersection complète » est de nature locale fpqc sur la base.
\end{lemma}

\begin{proof}
Soit $S$ un schéma. Soit $f : X \to Y$ un
morphisme d'espaces algébriques sur $S$.
Soit $\{Y_i \to Y\}$ un recouvrement fpqc
(Topologies sur les espaces, Définition
\ref{spaces-topologies-definition-fpqc-covering}).
Soit $f_i : X_i \to Y_i$ le changement de base de $f$ par $Y_i \to Y$.
Si $f$ est un morphisme localement d'intersection complète,
alors chaque $f_i$ est un morphisme localement d'intersection complète
d'après le Lemme \ref{lemma-flat-base-change-lci}.

\medskip\noindent
Réciproquement, supposons que chaque $f_i$ soit un morphisme localement d'intersection complète.
Nous pouvons remplacer le recouvrement par un raffinement (de nouveau parce que le
changement de base plat préserve la propriété d'être un
morphisme localement d'intersection complète). Nous pouvons donc supposer que
$Y_i$ est un schéma pour tout $i$, voir
Topologies sur les espaces, Lemme \ref{spaces-topologies-lemma-refine-fpqc-schemes}.
Choisissons un schéma $V$ et un morphisme étale surjectif $V \to Y$.
Choisissons un schéma $U$ et un morphisme étale surjectif
$U \to V \times_Y X$. Nous devons montrer que $U \to V$ est un
morphisme de schémas localement d'intersection complète.
D'après Topologies sur les espaces, Lemme
\ref{spaces-topologies-lemma-recognize-fpqc-covering}
la famille $\{Y_i \times_Y V \to V\}$ est un recouvrement fpqc
de schémas. Le cas des schémas
(Compléments sur les morphismes, Lemme \ref{more-morphisms-lemma-descending-property-lci})
montre qu'il suffit de prouver que le changement de base
$$
U \times_Y Y_i = U \times_V (V \times_Y Y_i) \longrightarrow V
$$
de $U \to V$ par $V \times_Y Y_i \to V$ est un
morphisme localement d'intersection complète. Nous pouvons l'écrire comme le
composé
$$
U \times_Y Y_i \longrightarrow
(V \times_Y X) \times_Y Y_i = V \times_Y X_i \longrightarrow
V \times_Y Y_i
$$
La première flèche est un morphisme étale de schémas (comme changement de base de
$U \to V \times_Y X$) et la seconde est un
morphisme de schémas localement d'intersection complète, comme changement de base plat de $f_i$.
Le résultat s'ensuit, car la propriété d'être un morphisme localement d'intersection complète
est de nature locale pour la topologie syntomique sur la source et les morphismes étales
sont syntomiques (Compléments sur les morphismes, Lemme
\ref{more-morphisms-lemma-lci-syntomic-local-source}
et Morphismes, Lemme \ref{morphisms-lemma-etale-syntomic}).
\end{proof}

\begin{lemma}
\label{lemma-lci-syntomic-local-source}
La propriété $\mathcal{P}(f) =$ « $f$ est un morphisme localement
d'intersection complète » est de nature locale pour la topologie syntomique sur la source.
\end{lemma}

\begin{proof}
Cela résulte de
Descente sur les espaces, Lemme \ref{spaces-descent-lemma-transfer-from-schemes}, et de
Compléments sur les morphismes, Lemme \ref{more-morphisms-lemma-lci-syntomic-local-source}.
\end{proof}

\begin{lemma}
\label{lemma-base-change-lci-fibres}
Soit $S$ un schéma. Considérons un diagramme commutatif
$$
\xymatrix{
X \ar[rr]_f \ar[rd]_p & & Y \ar[ld]^q \\
& Z
}
$$
d'espaces algébriques sur $S$. Supposons que $p$ et $q$
soient plats et localement de présentation finie.
Alors il existe un sous-espace ouvert $U(f) \subset X$
tel que $|U(f)| \subset |X|$ soit l'ensemble des points où $f$ est Koszul.
De plus, pour tout morphisme d'espaces algébriques $Z' \to Z$, si
$f' : X' \to Y'$ est le changement de base de $f$ par $Z' \to Z$, alors
$U(f')$ est l'image inverse de $U(f)$ par la projection $X' \to X$.
\end{lemma}

\begin{proof}
Ce lemme est l'analogue de
Compléments sur les morphismes, Lemme \ref{more-morphisms-lemma-base-change-lci-fibres},
et nous allons en fait l'en déduire. D'après la
Définition \ref{definition-lci},
l'ensemble $\{x \in |X| : f \text{ est Koszul en }x\}$ est
ouvert dans $|X|$ ; d'après
Propriétés des espaces, Lemme \ref{spaces-properties-lemma-open-subspaces},
il correspond donc à un sous-espace ouvert $U(f)$ de $X$. Il suffit par conséquent de
prouver la dernière assertion.

\medskip\noindent
Choisissons un schéma $W$ et un morphisme étale surjectif $W \to Z$.
Choisissons un schéma $V$ et un morphisme étale surjectif $V \to W \times_Z Y$.
Choisissons un schéma $U$ et un morphisme étale surjectif $U \to V \times_Y X$.
Enfin, choisissons un schéma $W'$ et un morphisme étale surjectif
$W' \to W \times_Z Z'$.
Posons $V' = W' \times_W V$ et $U' = W' \times_W U$ ; nous obtenons ainsi des
morphismes étales surjectifs $V' \to Y'$ et $U' \to X'$.
Nous utiliserons sans autre mention le fait qu'un morphisme étale d'espaces algébriques
induit une application ouverte entre les espaces topologiques associés (voir
Propriétés des espaces, Lemme
\ref{spaces-properties-lemma-etale-open}).
Notons que, par définition, $U(f)$ est l'image dans $|X|$ de l'ensemble $T$
des points de $U$ où le morphisme de schémas $U \to V$ est Koszul.
De même, $U(f')$ est l'image dans $|X'|$ de l'ensemble $T'$ des points de
$U'$ où le morphisme de schémas $U' \to V'$ est Koszul. Or, par construction, le
diagramme
$$
\xymatrix{
U' \ar[r] \ar[d] & U \ar[d] \\
V' \ar[r] & V
}
$$
est cartésien (dans la catégorie des schémas). Le
Lemme \ref{more-morphisms-lemma-base-change-lci-fibres} de Compléments sur les morphismes
cité ci-dessus montre donc que $T'$ est l'image inverse de $T$. Puisque
$|U'| \to |X'|$ est surjective, cela démontre le lemme.
\end{proof}

\begin{lemma}
\label{lemma-unramified-lci}
Soit $S$ un schéma. Soit $f : X \to Y$ un morphisme localement d'intersection
complète d'espaces algébriques sur $S$. Alors $f$ est non ramifié si et seulement
si $f$ est formellement non ramifié et, dans ce cas, le faisceau conormal
$\mathcal{C}_{X/Y}$ est localement libre de type fini sur $X$.
\end{lemma}

\begin{proof}
Cela découle du résultat correspondant pour les morphismes de schémas, voir
Compléments sur les morphismes, Lemme \ref{more-morphisms-lemma-unramified-lci},
par localisation étale, voir
Lemme \ref{lemma-universal-thickening-localize}.
(Notons que, dans la situation de ce lemme, le morphisme $V \to U$
est non ramifié et localement d'intersection complète par définition.)
\end{proof}

\begin{lemma}
\label{lemma-transitivity-conormal-lci}
Soit $S$ un schéma. Soient $Z \to Y \to X$ des morphismes formellement non ramifiés
d'espaces algébriques sur $S$. Supposons que $Z \to Y$ soit un morphisme localement
d'intersection complète. La suite exacte
$$
0 \to i^*\mathcal{C}_{Y/X} \to
\mathcal{C}_{Z/X} \to
\mathcal{C}_{Z/Y} \to 0
$$
du
Lemme \ref{lemma-transitivity-conormal}
est exacte courte.
\end{lemma}

\begin{proof}
Choisissons un schéma $U$ et un morphisme étale surjectif $U \to X$.
Choisissons un schéma $V$ et un morphisme étale surjectif $V \to U \times_X Y$.
Choisissons un schéma $W$ et un morphisme étale surjectif $W \to V \times_Y Z$.
D'après le
Lemme \ref{lemma-universal-thickening-localize},
les morphismes $W \to V$ et $V \to U$ sont formellement non ramifiés.
De plus, la suite
$i^*\mathcal{C}_{Y/X} \to \mathcal{C}_{Z/X} \to \mathcal{C}_{Z/Y} \to 0$
se restreint à la suite correspondante
$i^*\mathcal{C}_{V/U} \to \mathcal{C}_{W/U} \to \mathcal{C}_{W/V} \to 0$
pour $W \to V \to U$. Le résultat découle donc du cas des schémas
(Compléments sur les morphismes, Lemme \ref{more-morphisms-lemma-transitivity-conormal-lci}),
car, par définition, le morphisme $W \to V$ est localement d'intersection
complète.
\end{proof}









\section{Quand un morphisme est-il un isomorphisme ?}
\label{section-when-isomorphism}

\noindent
Plus généralement, on peut demander :
« Quand un morphisme possède-t-il la propriété $\mathcal{P}$ ? »
Une question plus précise est la suivante. Supposons donné un diagramme commutatif
$$
\xymatrix{
X \ar[rr]_f \ar[rd]_p & & Y \ar[ld]^q \\
& Z
}
$$
d'espaces algébriques. Existe-t-il un monomorphisme d'espaces algébriques
$W \to Z$ vérifiant les deux propriétés suivantes :
\begin{enumerate}
\item le changement de base $f_W : X_W \to Y_W$ possède la propriété $\mathcal{P}$, et
\item tout morphisme $Z' \to Z$ d'espaces algébriques se factorise par $W$ si
et seulement si le changement de base $f_{Z'} : X_{Z'} \to Y_{Z'}$ possède la propriété
$\mathcal{P}$.
\end{enumerate}
Dans de nombreux cas, si $W \to Z$ existe, c'est une immersion, une immersion ouverte
ou une immersion fermée.

\medskip\noindent
La réponse à cette question peut dépendre de propriétés auxiliaires des
morphismes $f$, $p$ et $q$. Un exemple est $\mathcal{P}(f) =$ « $f$ est plat »,
dont nous avons traité en grand détail, pour les morphismes de schémas, dans le cas $Y = S$ au
chapitre « Compléments sur la platitude », à partir de
Compléments sur la platitude, section \ref{flat-section-flattening-functors}.

\begin{lemma}
\label{lemma-where-unramified}
Considérons un diagramme commutatif
$$
\xymatrix{
X \ar[rr]_f \ar[rd]_p & & Y \ar[ld]^q \\
& Z
}
$$
d'espaces algébriques. Supposons que $p$ soit localement de type fini et fermé.
Alors il existe un sous-espace ouvert $W \subset Z$
tel qu'un morphisme $Z' \to Z$ se factorise par $W$ si et seulement si le
changement de base $f_{Z'} : X_{Z'} \to Y_{Z'}$ est non ramifié.
\end{lemma}

\begin{proof}
D'après le
Lemme \ref{spaces-morphisms-lemma-where-unramified} de Morphismes d'espaces,
il existe un sous-espace ouvert $U(f) \subset X$ qui est l'ensemble des
points où $f$ est non ramifié. De plus, la formation de $U(f)$ commute
à tout changement de base. Soit $W \subset Z$ le sous-espace ouvert
(voir
Propriétés des espaces, Lemme
\ref{spaces-properties-lemma-open-subspaces})
dont l'ensemble sous-jacent est
$$
|W| = |Z| \setminus |p|\left(|X| \setminus |U(f)|\right)
$$
c'est-à-dire que $z \in |Z|$ est un point de $W$ si et seulement si $f$ est non ramifié
en tout point de $X$ au-dessus de $z$. Cet ensemble est ouvert, car nous
avons supposé que $p$ est fermé. Puisque la formation de $U(f)$
commute à tout changement de base, on voit immédiatement (en utilisant
Propriétés des espaces, Lemme
\ref{spaces-properties-lemma-factor-through-open-subspace})
que $W$ possède la propriété universelle voulue.
\end{proof}

\begin{lemma}
\label{lemma-where-unramified-universally-injective}
Considérons un diagramme commutatif
$$
\xymatrix{
X \ar[rr]_f \ar[rd]_p & & Y \ar[ld]^q \\
& Z
}
$$
d'espaces algébriques. Supposons que
\begin{enumerate}
\item $p$ soit localement de type fini,
\item $p$ soit fermé, et
\item $p_2 : X \times_Y X \to Z$ soit fermé.
\end{enumerate}
Alors il existe un sous-espace ouvert $W \subset Z$
tel qu'un morphisme $Z' \to Z$ se factorise par $W$ si et seulement si le
changement de base $f_{Z'} : X_{Z'} \to Y_{Z'}$ est non ramifié et universellement
injectif.
\end{lemma}

\begin{proof}
Après avoir remplacé $Z$ par le sous-espace ouvert trouvé dans le
Lemme \ref{lemma-where-unramified},
on peut supposer que $f$ est déjà non ramifié ; notons que cela ne
remet en cause ni l'hypothèse (2) ni l'hypothèse (3). D'après
Morphismes d'espaces, Lemme
\ref{spaces-morphisms-lemma-diagonal-unramified-morphism},
on voit que $\Delta_{X/Y} : X \to X \times_Y X$ est une immersion ouverte.
Cela reste vrai après tout changement de base. Par conséquent, d'après
Morphismes d'espaces, Lemme
\ref{spaces-morphisms-lemma-universally-injective},
on voit que $f_{Z'}$ est universellement injectif si et seulement si
le changement de base de la diagonale $X_{Z'} \to (X \times_Y X)_{Z'}$
est un isomorphisme. Soit $W \subset Z$ le sous-espace ouvert
(voir
Propriétés des espaces, Lemme
\ref{spaces-properties-lemma-open-subspaces})
dont l'ensemble sous-jacent est
$$
|W| = |Z| \setminus
|p_2|\left(|X \times_Y X| \setminus \Im(|\Delta_{X/Y}|)\right)
$$
c'est-à-dire que $z \in |Z|$ est un point de $W$ si et seulement si la fibre de
$|X \times_Y X| \to |Z|$ au-dessus de $z$ est contenue dans l'image de
$|X| \to |X \times_Y X|$. Il ressort alors clairement de ce qui précède
que la restriction $p^{-1}(W) \to q^{-1}(W)$ de $f$ est
non ramifiée et universellement injective.

\medskip\noindent
Réciproquement, supposons que $f_{Z'}$ soit non ramifié et universellement injectif.
Pour montrer que $Z' \to Z$ se factorise par $W$, il suffit de montrer
que l'image de $|Z'| \to |Z|$ est contenue dans $|W|$, voir
Propriétés des espaces, Lemme
\ref{spaces-properties-lemma-factor-through-open-subspace}.
Il suffit donc de démontrer le résultat lorsque $Z'$ est le spectre d'un corps.
Notons $z \in |Z|$ l'image de $|Z'| \to |Z|$. La discussion ci-dessus montre
que
$$
|X_{Z'}| \longrightarrow |(X \times_Y X)_{Z'}|
$$
est surjective. D'après
Propriétés des espaces,
Lemme \ref{spaces-properties-lemma-points-cartesian},
dans le diagramme commutatif
$$
\xymatrix{
|X_{Z'}| \ar[d] \ar[r] &
|(X \times_Y X)_{Z'}| \ar[d] \\
|p|^{-1}(\{z\}) \ar[r] &
|p_2|^{-1}(\{z\})
}
$$
les flèches verticales sont surjectives. Il s'ensuit que $z \in |W|$, comme voulu.
\end{proof}

\begin{lemma}
\label{lemma-where-closed-immersion}
Considérons un diagramme commutatif
$$
\xymatrix{
X \ar[rr]_f \ar[rd]_p & & Y \ar[ld]^q \\
& Z
}
$$
d'espaces algébriques. Supposons que
\begin{enumerate}
\item $p$ soit localement de type fini,
\item $p$ soit universellement fermé, et
\item $q : Y \to Z$ soit séparé.
\end{enumerate}
Alors il existe un sous-espace ouvert $W \subset Z$
tel qu'un morphisme $Z' \to Z$ se factorise par $W$ si et seulement si le
changement de base $f_{Z'} : X_{Z'} \to Y_{Z'}$ est une immersion fermée.
\end{lemma}

\begin{proof}
Nous utiliserons la caractérisation des immersions fermées comme morphismes
universellement fermés, non ramifiés et universellement injectifs, voir
Lemme \ref{lemma-characterize-closed-immersion}.
Notons d'abord que, puisque $p$ est universellement fermé et que $q$ est
séparé, $f$ est universellement fermé, voir
Morphismes d'espaces, Lemme
\ref{spaces-morphisms-lemma-universally-closed-permanence}.
Il en résulte que tout changement de base de $f$ est universellement fermé, voir
Morphismes d'espaces, Lemme
\ref{spaces-morphisms-lemma-base-change-universally-closed}.
Ainsi, pour achever la démonstration du lemme, il suffit de prouver que
les hypothèses du
Lemme \ref{lemma-where-unramified-universally-injective}
sont satisfaites. La projection $\text{pr}_0 : X \times_Y X \to X$
est universellement fermée comme changement de base de $f$, voir
Morphismes d'espaces, Lemme
\ref{spaces-morphisms-lemma-base-change-universally-closed}.
Par suite, $X \times_Y X \to Z$ est universellement fermé comme
composé de morphismes universellement fermés (voir
Morphismes d'espaces, Lemme
\ref{spaces-morphisms-lemma-composition-universally-closed}).
Cela achève la démonstration du lemme.
\end{proof}

\begin{lemma}
\label{lemma-where-flat}
Considérons un diagramme commutatif
$$
\xymatrix{
X \ar[rr]_f \ar[rd]_p & & Y \ar[ld]^q \\
& Z
}
$$
d'espaces algébriques. Supposons que
\begin{enumerate}
\item $p$ soit localement de présentation finie,
\item $p$ soit plat,
\item $p$ soit fermé, et
\item $q$ soit localement de type fini.
\end{enumerate}
Alors il existe un sous-espace ouvert $W \subset Z$
tel qu'un morphisme $Z' \to Z$ se factorise par $W$ si et seulement si le
changement de base $f_{Z'} : X_{Z'} \to Y_{Z'}$ est plat.
\end{lemma}

\begin{proof}
D'après le Lemme \ref{lemma-base-change-flatness-fibres},
l'ensemble
$$
A = \{x \in |X| : X\text{ est plat en }x \text{ sur }Y\}.
$$
Cet ensemble est ouvert dans $|X|$ et sa formation commute à tout changement de
base. Soit $W \subset Z$ le sous-espace ouvert
(voir
Propriétés des espaces, Lemme
\ref{spaces-properties-lemma-open-subspaces})
dont l'ensemble sous-jacent est
$$
|W| = |Z| \setminus |p|\left(|X| \setminus A\right)
$$
c'est-à-dire que $z \in |Z|$ est un point de $W$ si et seulement si la fibre entière
de $|X| \to |Z|$ au-dessus de $z$ est contenue dans $A$. Cet ensemble est ouvert car
$p$ est fermé. Puisque la formation de $A$ commute à tout changement de
base, il s'ensuit que $W$ convient.
\end{proof}

\begin{lemma}
\label{lemma-where-surjective-flat}
Considérons un diagramme commutatif
$$
\xymatrix{
X \ar[rr]_f \ar[rd]_p & & Y \ar[ld]^q \\
& Z
}
$$
d'espaces algébriques. Supposons que
\begin{enumerate}
\item $p$ soit localement de présentation finie,
\item $p$ soit plat,
\item $p$ soit fermé,
\item $q$ soit localement de type fini, et
\item $q$ soit fermé.
\end{enumerate}
Alors il existe un sous-espace ouvert $W \subset Z$
tel qu'un morphisme $Z' \to Z$ se factorise par $W$ si et seulement si le
changement de base $f_{Z'} : X_{Z'} \to Y_{Z'}$ est surjectif et plat.
\end{lemma}

\begin{proof}
D'après le
Lemme \ref{lemma-where-flat},
on peut supposer que $f$ est plat.
Notons que $f$ est localement de présentation finie d'après
Morphismes d'espaces,
Lemme \ref{spaces-morphisms-lemma-finite-presentation-permanence}.
Par conséquent, $f$ est ouvert, voir
Morphismes d'espaces, Lemme \ref{spaces-morphisms-lemma-fppf-open}.
Soit $W \subset Z$ le sous-espace ouvert (voir
Propriétés des espaces, Lemme
\ref{spaces-properties-lemma-open-subspaces})
dont l'ensemble sous-jacent est
$$
|W| = |Z| \setminus |q|\left(|Y| \setminus |f|(|X|)\right).
$$
Autrement dit, pour $z \in |Z|$, on a $z \in |W|$ si et seulement
si la fibre entière de $|Y| \to |Z|$ au-dessus de $z$ est contenue dans l'image de
$|X| \to |Y|$. Puisque $q$ est fermé, cet ensemble est ouvert dans $|Z|$.
Le morphisme $X_W \to Y_W$ est surjectif par construction.
Enfin, supposons que $X_{Z'} \to Y_{Z'}$ soit surjectif.
Pour montrer que $Z' \to Z$ se factorise par $W$, il suffit de montrer
que l'image de $|Z'| \to |Z|$ est contenue dans $|W|$, voir
Propriétés des espaces, Lemme
\ref{spaces-properties-lemma-factor-through-open-subspace}.
Il suffit donc de démontrer le résultat lorsque $Z'$ est le spectre d'un corps.
Notons $z \in |Z|$ l'image de $|Z'| \to |Z|$. D'après
Propriétés des espaces,
Lemme \ref{spaces-properties-lemma-points-cartesian},
dans le diagramme commutatif
$$
\xymatrix{
|X_{Z'}| \ar[d] \ar[r] &
|Y_{Z'}| \ar[d] \\
|p|^{-1}(\{z\}) \ar[r] &
|q|^{-1}(\{z\})
}
$$
les flèches verticales sont surjectives. Il s'ensuit que $z \in |W|$, comme voulu.
\end{proof}

\begin{lemma}
\label{lemma-where-isomorphism}
Considérons un diagramme commutatif
$$
\xymatrix{
X \ar[rr]_f \ar[rd]_p & & Y \ar[ld]^q \\
& Z
}
$$
d'espaces algébriques. Supposons que
\begin{enumerate}
\item $p$ soit localement de présentation finie,
\item $p$ soit plat,
\item $p$ soit universellement fermé,
\item $q$ soit localement de type fini,
\item $q$ soit fermé, et
\item $q$ soit séparé.
\end{enumerate}
Alors il existe un sous-espace ouvert $W \subset Z$
tel qu'un morphisme $Z' \to Z$ se factorise par $W$ si et seulement si le
changement de base $f_{Z'} : X_{Z'} \to Y_{Z'}$ est un isomorphisme.
\end{lemma}

\begin{proof}
D'après le
Lemme \ref{lemma-where-surjective-flat},
il existe un sous-espace ouvert $W_1 \subset Z$ tel que
$f_{Z'}$ soit surjectif et plat si et seulement si $Z' \to Z$
se factorise par $W_1$. D'après le
Lemme \ref{lemma-where-closed-immersion},
il existe un sous-espace ouvert $W_2 \subset Z$ tel que
$f_{Z'}$ soit une immersion fermée si et seulement si $Z' \to Z$
se factorise par $W_2$. Nous affirmons que $W = W_1 \cap W_2$ convient.
Si $f_{Z'}$ est un isomorphisme, alors $Z' \to Z$
se factorise par $W$. Il suffit donc de montrer que
$f_W$ est un isomorphisme. Par construction, $f_W$ est une
immersion fermée surjective et plate. En particulier, $f_W$ est
représentable. Puisqu'une immersion fermée surjective et plate de
schémas est un isomorphisme (voir
Morphismes, Lemme \ref{morphisms-lemma-characterize-flat-closed-immersions}),
cela achève la démonstration. (Notons qu'en fait $f_W$ est localement de présentation finie,
donc ouvert, de sorte que l'on peut éviter d'utiliser ce lemme si on le souhaite.)
\end{proof}

\begin{lemma}
\label{lemma-where-lci}
Considérons un diagramme commutatif
$$
\xymatrix{
X \ar[rr]_f \ar[rd]_p & & Y \ar[ld]^q \\
& Z
}
$$
d'espaces algébriques. Supposons que
\begin{enumerate}
\item $p$ soit plat et localement de présentation finie,
\item $p$ soit fermé, et
\item $q$ soit plat et localement de présentation finie.
\end{enumerate}
Alors il existe un sous-espace ouvert $W \subset Z$
tel qu'un morphisme $Z' \to Z$ se factorise par $W$ si et seulement si le
changement de base $f_{Z'} : X_{Z'} \to Y_{Z'}$ est un morphisme localement d'intersection
complète.
\end{lemma}

\begin{proof}
D'après le Lemme \ref{lemma-base-change-lci-fibres},
il existe un sous-espace ouvert $U(f) \subset X$ qui est l'ensemble des
points où $f$ est Koszul. De plus, la formation de $U(f)$ commute
à tout changement de base. Soit $W \subset Z$ le sous-espace ouvert
(voir
Propriétés des espaces, Lemme
\ref{spaces-properties-lemma-open-subspaces})
dont l'ensemble sous-jacent est
$$
|W| = |Z| \setminus |p|\left(|X| \setminus |U(f)|\right)
$$
c'est-à-dire que $z \in |Z|$ est un point de $W$ si et seulement si $f$ est Koszul
en tout point de $X$ au-dessus de $z$. Cet ensemble est ouvert car nous
avons supposé que $p$ est fermé. Puisque la formation de $U(f)$
commute à tout changement de base, on voit immédiatement (en utilisant
Propriétés des espaces, Lemme
\ref{spaces-properties-lemma-factor-through-open-subspace})
que $W$ possède la propriété universelle voulue.
\end{proof}








\section{Suites exactes de différentielles et de faisceaux conormaux}
\label{section-exact}

\noindent
Dans cette section, nous rassemblons quelques résultats sur les suites exactes de faisceaux
conormaux et de faisceaux de différentielles. En un certain sens, ce sont toutes des
réalisations du triangle des complexes cotangents associé à des
morphismes composables d'espaces algébriques.

\medskip\noindent
Dans les suites ci-dessous, chacune des applications
est construite soit dans le
lemme \ref{lemma-functoriality-differentials}, soit dans le
lemme \ref{lemma-universal-thickening-functorial}.
Soit $S$ un schéma. Soient $g : Z \to Y$ et $f : Y \to X$ des morphismes
d'espaces algébriques sur $S$.
\begin{enumerate}
\item Il existe une suite exacte canonique
$$
g^*\Omega_{Y/X} \to \Omega_{Z/X} \to \Omega_{Z/Y} \to 0,
$$
voir le
lemme \ref{lemma-triangle-differentials}.
Si $g : Z \to Y$ est formellement lisse, cette suite est une suite
exacte courte, voir le
lemme \ref{lemma-triangle-differentials-formally-smooth}.
\item Si $g$ est formellement non ramifié, il existe une suite exacte canonique
$$
\mathcal{C}_{Z/Y} \to g^*\Omega_{Y/X} \to \Omega_{Z/X} \to 0,
$$
voir le
lemme \ref{lemma-universally-unramified-differentials-sequence}.
Si $f \circ g : Z \to X$ est formellement lisse, cette suite est une suite
exacte courte, voir le
lemme \ref{lemma-differentials-formally-unramified-formally-smooth}.
\item Si $g$ et $f \circ g$ sont formellement non ramifiés, il existe une
suite exacte canonique
$$
\mathcal{C}_{Z/X} \to \mathcal{C}_{Z/Y} \to g^*\Omega_{Y/X} \to 0,
$$
voir le
lemme \ref{lemma-two-unramified-morphisms}.
Si $f : Y \to X$ est formellement lisse, cette suite est une suite
exacte courte, voir le
lemme \ref{lemma-two-unramified-morphisms-formally-smooth}.
\item Si $g$ et $f$ sont formellement non ramifiés, il existe une suite
exacte canonique
$$
g^*\mathcal{C}_{Y/X} \to \mathcal{C}_{Z/X} \to \mathcal{C}_{Z/Y} \to 0.
$$
voir le
lemme \ref{lemma-transitivity-conormal-unramified}.
Si $g : Z \to Y$ est un morphisme localement d'intersection complète,
alors cette suite est une suite exacte courte, voir le
lemme \ref{lemma-transitivity-conormal-lci}.
\end{enumerate}










\section{Caractérisation des complexes pseudo-cohérents, II}
\label{section-characterize-pseudo-coherent}

\noindent
Dans cette section, nous étudions une caractérisation des complexes pseudo-cohérents
en termes de cohomologie. On trouvera des résultats antérieurs sur les complexes pseudo-cohérents
sur les espaces algébriques dans
Catégories dérivées des espaces, section
\ref{spaces-perfect-section-spell-out}
et dans
Catégories dérivées des espaces, section
\ref{spaces-perfect-section-pseudo-coherent-hocolim}.
L'analogue de cette section pour les schémas est Compléments sur les morphismes, section
\ref{more-morphisms-section-characterize-pseudo-coherent}.
Un outil fondamental consistera à nous ramener au cas de l'espace projectif
au moyen d'une version dérivée du lemme de Chow, voir
Lemme \ref{lemma-derived-chow}.

\begin{lemma}
\label{lemma-case-of-tor-independence}
Soit $S$ un schéma. Considérons un diagramme commutatif d'espaces algébriques
$$
\xymatrix{
Z' \ar[d] \ar[r] & Y' \ar[d] \\
X' \ar[r] & B'
}
$$
sur $S$.
Soit $B \to B'$ un morphisme. Notons $X$ et $Y$ les changements
de base de $X'$ et $Y'$ à $B$.
Supposons que $Y' \to B'$ et $Z' \to X'$ soient plats.
Alors $X \times_B Y$ et $Z'$ sont Tor-indépendants sur $X' \times_{B'} Y'$.
\end{lemma}

\begin{proof}
D'après Catégories dérivées des espaces, Lemme
\ref{spaces-perfect-lemma-tor-independent},
on peut vérifier l'indépendance pour Tor localement pour la topologie \'etale sur $X \times_B Y$
et $Z'$. Cela\footnote{Voici l'argument plus en détail.
Choisissons un morphisme \'etale surjectif $W' \to B'$
où $W'$ est un schéma. Choisissons un morphisme \'etale surjectif
$W \to B \times_{B'} W'$ où $W$ est un schéma. Choisissons un
morphisme \'etale surjectif
$U' \to X' \times_{B'} W'$ où $U'$ est un schéma. Choisissons un
morphisme \'etale surjectif $V' \to Y' \times_{B'} W'$ où $V'$ est un schéma.
Remarquons que $U' \times_{W'} V' \to X' \times_{B'} Y'$ est surjectif et
\'etale. Choisissons un morphisme \'etale surjectif
$T' \to Z' \times_{X' \times_{B'} Y'} U' \times_{W'} V'$
où $T'$ est un schéma. Notons $U$ et $V$ les changements de base de $U'$ et $V'$
à $W$. Le lemme affirme alors que $X \times_B Y$ et $Z'$
sont Tor-indépendants sur $X' \times_{B'} Y'$ comme espaces algébriques
si et seulement si $U \times_W V$ et $T'$ sont Tor-indépendants sur
$U' \times_{W'} V'$ comme schémas. Ainsi,
il suffit de démontrer le lemme pour
le carré de sommets $T', U', V', W'$ et le changement de base par $W \to W'$.
La platitude de $Y' \to B'$ et de $Z' \to X'$ entraîne celle
de $V' \to W'$ et de $T' \to U'$.}
ramène le lemme au cas des schémas,
qui est Compléments sur les morphismes, Lemme
\ref{more-morphisms-lemma-case-of-tor-independence}.
\end{proof}

\begin{lemma}[Lemme de Chow dérivé]
\label{lemma-derived-chow}
Soit $A$ un anneau. Soit $X$ un espace algébrique séparé
de présentation finie sur $A$. Soit $x \in |X|$. Alors il existe
un $n \geq 0$,
un sous-espace fermé $Z \subset X \times_A \mathbf{P}^n_A$,
un point $z \in |Z|$,
un ouvert $V \subset \mathbf{P}^n_A$, et
un objet $E$ de $D(\mathcal{O}_{X \times_A \mathbf{P}^n_A})$ tels que
\begin{enumerate}
\item $Z \to X \times_A \mathbf{P}^n_A$ est de présentation finie,
\item $c : Z \to \mathbf{P}^n_A$ est une immersion fermée au-dessus de $V$,
posons $W = c^{-1}(V)$,
\item la restriction de $b : Z \to X$ à $W$ est \'etale,
$z \in W$, et $b(z) = x$,
\item $E|_{X \times_A V} \cong
(b, c)_*\mathcal{O}_Z|_{X \times_A V}$,
\item $E$ est pseudo-cohérent et à support dans $Z$.
\end{enumerate}
\end{lemma}

\begin{proof}
Nous pouvons trouver une sous-$\mathbf{Z}$-algèbre de type fini $A' \subset A$
et un espace algébrique $X'$ séparé et de présentation finie sur $A'$
dont le changement de base à $A$ est $X$. Voir
Limites d'espaces, Lemmes
\ref{spaces-limits-lemma-descend-finite-presentation} et
\ref{spaces-limits-lemma-descend-separated-morphism}.
Soit $x' \in |X'|$ l'image de $x$.
Si nous pouvons démontrer le lemme pour $(X'/A', x')$, alors
le lemme en découle pour $(X/A, x)$.
En effet, si $n', Z', z', V', E'$ fournissent les données voulues
pour $(X'/A', x')$, nous pouvons poser
$n = n'$,
prendre pour $Z \subset X \times \mathbf{P}^n$ l'image inverse de $Z'$,
prendre pour $z \in Z$ l'unique point d'image $x$,
prendre pour $V \subset \mathbf{P}^n_A$ l'image inverse de $V'$, et
prendre pour $E$ l'image inverse dérivée de $E'$.
Remarquons que $E$ est pseudo-cohérent d'après
Cohomologie sur les sites, Lemme
\ref{sites-cohomology-lemma-pseudo-coherent-pullback}.
Il ne reste qu'à vérifier (5). Pour cela,
posons $W = c^{-1}(V)$ et $W' = (c')^{-1}(V')$
et considérons le carré cartésien
$$
\xymatrix{
W \ar[d]_{(b, c)} \ar[r] & W' \ar[d]^{(b', c')} \\
X \times_A V \ar[r] & X' \times_{A'} V'
}
$$
D'après le Lemme \ref{lemma-case-of-tor-independence}, $X \times_A V$ et $W'$
sont Tor-indépendants sur $X' \times_{A'} V'$.
Ainsi, l'image inverse dérivée de
$(b', c')_*\mathcal{O}_{W'}$ sur $X \times_A V$
est $(b, c)_*\mathcal{O}_W$ d'après
Catégories dérivées des espaces,
Lemme \ref{spaces-perfect-lemma-compare-base-change}.
On utilise aussi ici le fait que $R(b', c')_*\mathcal{O}_{Z'} = (b', c')_*\mathcal{O}_{Z'}$
parce que $(b', c')$ est une immersion fermée, et de même pour
$(b, c)_*\mathcal{O}_Z$.
Comme $E'|_{U' \times_{A'} V'} =
(b', c')_*\mathcal{O}_{W'}$, on obtient
$E|_{U \times_A V} = (b, c)_*\mathcal{O}_W$,
et (5) est vérifié.
Nous sommes ainsi ramenés à la situation décrite au
paragraphe suivant.

\medskip\noindent
Supposons que $A$ soit de type fini sur $\mathbf{Z}$.
Choisissons un morphisme \'etale $U \to X$, où $U$ est un schéma affine,
et un point $u \in U$ d'image $x$. Alors $U$ est de type fini sur $A$.
Choisissons une immersion fermée $U \to \mathbf{A}^n_A$ et notons
$j : U \to \mathbf{P}^n_A$ l'immersion obtenue en la composant
avec l'immersion ouverte $\mathbf{A}^n_A \to \mathbf{P}^n_A$.
Soit $Z$ l'adhérence schématique de
$$
(\text{id}_U, j) : U \longrightarrow X \times_A \mathbf{P}^n_A
$$
Soit $z \in Z$ l'image de $u$.
Soit $Y \subset \mathbf{P}^n_A$ l'adhérence schématique
de $j$. Il est alors clair que $Z \subset X \times_A Y$
est l'adhérence schématique de
$(\text{id}_U, j) : U \to X \times_A Y$.
Comme $X$ est séparé, le morphisme
$X \times_A Y \to Y$ est lui aussi séparé.
On voit donc que $Z \to Y$ est un isomorphisme au-dessus
du sous-schéma ouvert $j(U) \subset Y$ d'après
Morphismes d'espaces, Lemme
\ref{spaces-morphisms-lemma-scheme-theoretic-image-of-partial-section}.
Choisissons un ouvert $V \subset \mathbf{P}^n_A$ tel que $V \cap Y = j(U)$.
On voit alors que (2) est vérifié, que $W = (\text{id}_U, j)(U)$, et donc
que (3) est vérifié. La partie (1) est vérifiée parce que $A$ est noethérien.

\medskip\noindent
Comme $A$ est noethérien, $X$ et $X \times_A \mathbf{P}^n_A$
sont des espaces algébriques noethériens. Nous pouvons donc prendre $E = (b, c)_*\mathcal{O}_Z$
dans ce cas : (4) est clair et, pour (5), voir
Catégories dérivées des espaces, Lemme
\ref{spaces-perfect-lemma-identify-pseudo-coherent-noetherian}.
Ceci achève la démonstration.
\end{proof}

\begin{lemma}
\label{lemma-compute-Fourier-Mukai-for-derived-chow}
Soient $X/A$, $x \in |X|$, et
$n, Z, z, V, E$ comme dans le Lemme \ref{lemma-derived-chow}.
Pour tout $K \in D_\QCoh(\mathcal{O}_X)$, on a
$$
Rq_*(Lp^*K \otimes^\mathbf{L} E)|_V = R(W \to V)_*K|_W
$$
où $p : X \times_A \mathbf{P}^n_A \to X$ et
$q : X \times_A \mathbf{P}^n_A \to \mathbf{P}^n_A$ sont
les projections et où le morphisme $W \to V$ est
l'immersion fermée de présentation finie $c|_W : W \to V$.
\end{lemma}

\begin{proof}
Comme $W = c^{-1}(V)$ et que $c$ est une immersion fermée
au-dessus de $V$, on voit que $c|_W$ est une immersion fermée.
Elle est de présentation finie parce que $W$ et $V$ sont de présentation
finie sur $A$, voir
Morphismes d'espaces, Lemme
\ref{spaces-morphisms-lemma-finite-presentation-permanence}.
Tout d'abord, on a
$$
Rq_*(Lp^*K \otimes^\mathbf{L} E)|_V =
Rq'_*\left((Lp^*K \otimes^\mathbf{L} E)|_{X \times_A V}\right)
$$
où $q' : X \times_A V \to V$ est la projection, car
la formation de l'image directe totale commute à la localisation.
Notons $i = (b, c)|_W : W \to X \times_A V$ l'immersion fermée ainsi obtenue.
Alors
$$
Rq'_*\left((Lp^*K \otimes^\mathbf{L} E)|_{X \times_A V}\right) =
Rq'_*(Lp^*K|_{X \times_A V} \otimes^\mathbf{L} i_*\mathcal{O}_W)
$$
d'après la propriété (5). Comme $i$ est une immersion fermée, on a
$i_*\mathcal{O}_W = Ri_*\mathcal{O}_W$.
En utilisant
Catégories dérivées des espaces,
Lemme \ref{spaces-perfect-lemma-cohomology-base-change},
nous pouvons réécrire cette expression sous la forme
$$
Rq'_* Ri_* Li^* Lp^*K|_{X \times_A V} =
R(q' \circ i)_* Lb^*K|_W =
R(W \to V)_* K|_W
$$
ce qui est l'égalité cherchée. (Notons que se restreindre à $W$
et prendre l'image inverse dérivée par $W \to X$ reviennent au même
puisque $W$ est \'etale sur $X$.)
\end{proof}

\begin{lemma}
\label{lemma-characterize-pseudo-coherent}
Soit $A$ un anneau. Soit $X$ un espace algébrique séparé et
de présentation finie sur $A$. Soit $K \in D_\QCoh(\mathcal{O}_X)$.
Si $R\Gamma(X, E \otimes^\mathbf{L} K)$ est pseudo-cohérent
dans $D(A)$ pour tout $E$ pseudo-cohérent dans $D(\mathcal{O}_X)$,
alors $K$ est pseudo-cohérent relativement à $A$
(Définition \ref{definition-relative-pseudo-coherence}).
\end{lemma}

\begin{proof}
Supposons que $K \in D_\QCoh(\mathcal{O}_X)$ et que
$R\Gamma(X, E \otimes^\mathbf{L} K)$ soit pseudo-cohérent
dans $D(A)$ pour tout $E$ pseudo-cohérent dans $D(\mathcal{O}_X)$.
Soit $x \in |X|$. Nous allons montrer que $K$ est pseudo-cohérent relativement à $A$
dans un voisinage \'etale de $x$. Cela démontrera le lemme
par notre définition de la pseudo-cohérence relative.

\medskip\noindent
Choisissons $n, Z, z, V, E$ comme dans le Lemme \ref{lemma-derived-chow}.
Notons $p : X \times \mathbf{P}^n \to X$ et
$q : X \times \mathbf{P}^n \to \mathbf{P}^n_A$
les projections.
Alors, pour tout $i \in \mathbf{Z}$, on a
\begin{align*}
& R\Gamma(\mathbf{P}^n_A,
Rq_*(Lp^*K \otimes^\mathbf{L} E)
\otimes^\mathbf{L}
\mathcal{O}_{\mathbf{P}^n_A}(i)) \\
& =
R\Gamma(X \times \mathbf{P}^n,
Lp^*K \otimes^\mathbf{L} E \otimes^\mathbf{L}
Lq^*\mathcal{O}_{\mathbf{P}^n_A}(i)) \\
& =
R\Gamma(X,
K \otimes^\mathbf{L} Rp_*(E \otimes^\mathbf{L}
Lq^*\mathcal{O}_{\mathbf{P}^n_A}(i)))
\end{align*}
d'après
Catégories dérivées des espaces,
Lemme \ref{spaces-perfect-lemma-cohomology-base-change}.
D'après
Catégories dérivées des espaces, Lemme
\ref{spaces-perfect-lemma-flat-proper-pseudo-coherent-direct-image-general},
le complexe $Rp_*(E \otimes^\mathbf{L} Lq^*\mathcal{O}_{\mathbf{P}^n_A}(i))$
est pseudo-cohérent sur $X$. L'hypothèse nous dit donc que l'expression
figurant dans la formule affichée est un objet pseudo-cohérent de $D(A)$.
D'après
Catégories dérivées des schémas,
Lemme \ref{perfect-lemma-pseudo-coherent-on-projective-space},
nous en concluons que $Rq_*(Lp^*K \otimes^\mathbf{L} E)$ est
pseudo-cohérent sur $\mathbf{P}^n_A$.
D'après le Lemme \ref{lemma-compute-Fourier-Mukai-for-derived-chow},
nous avons
$$
Rq_*(Lp^*K \otimes^\mathbf{L} E)|_{X \times_A V} =
R(W \to V)_*K|_W
$$
Comme $W \to V$ est une immersion fermée dans un sous-schéma ouvert de
$\mathbf{P}^n_A$, cela signifie que $K|_W$ est pseudo-cohérent relativement à $A$,
par exemple d'après
Compléments sur les morphismes,
Lemme \ref{more-morphisms-lemma-check-relative-pseudo-coherence-on-charts}.
\end{proof}

\begin{lemma}
\label{lemma-characterize-pseudo-coh-improved}
Soit $A$ un anneau. Soit $X$ un espace algébrique séparé et
de présentation finie sur $A$. Soit $K \in D_\QCoh(\mathcal{O}_X)$. Si
$R \Gamma (X, E \otimes ^{\mathbf{L}} K)$ est
pseudo-cohérent dans $D(A)$ pour tout objet parfait
$E \in D(\mathcal{O}_X)$, alors $K$ est pseudo-cohérent
relativement à $A$.
\end{lemma}

\begin{proof}
Compte tenu du Lemme \ref{lemma-characterize-pseudo-coherent}, il suffit
de montrer que $R \Gamma (X, E \otimes ^{\mathbf{L}} K)$ est
pseudo-cohérent dans $D(A)$ pour tout objet pseudo-cohérent
$E \in D(\mathcal{O}_X)$. D'après Catégories dérivées des espaces,
Proposition \ref{spaces-perfect-proposition-detecting-bounded-above},
il en résulte que $K \in D^-_\QCoh (\mathcal{O}_X)$. Le
résultat découle alors de Catégories dérivées des espaces, Lemme
\ref{spaces-perfect-lemma-perfect-enough}.
\end{proof}












\section{Objets relativement parfaits}
\label{section-relatively-perfect}

\noindent
Dans cette section, nous introduisons une notion provenant de \cite{lieblich-complexes}.
Cette notion a été étudiée pour les morphismes de schémas dans
Catégories dérivées de schémas, section
\ref{perfect-section-relatively-perfect}.

\begin{definition}
\label{definition-relatively-perfect}
Soit $S$ un schéma. Soit $f : X \to Y$ un morphisme d'espaces algébriques
sur $S$ qui soit plat et localement de présentation finie.
Un objet $E$ de $D(\mathcal{O}_X)$ est {\it parfait relativement à $Y$} ou
{\it $Y$-parfait} si $E$ est pseudo-cohérent
(Cohomologie sur les sites, définition
\ref{sites-cohomology-definition-pseudo-coherent}) et si
$E$ est localement de dimension de Tor finie en tant qu'objet de
$D(f^{-1}\mathcal{O}_Y)$
(Cohomologie sur les sites, définition
\ref{sites-cohomology-definition-tor-amplitude}).
\end{definition}

\noindent
Voir Catégories dérivées de schémas,
remarque \ref{perfect-remark-discuss-rel-perfect}, pour une discussion ;
signalons seulement ici que la pseudo-cohérence de $E$ est
équivalente à la pseudo-cohérence de $E$ relativement à $Y$ d'après le
lemme \ref{lemma-relative-pseudo-coherent-is-moot}.
De plus, la pseudo-cohérence de $E$ implique $E \in D_\QCoh(\mathcal{O}_X)$ ; voir
Catégories dérivées d'espaces, lemme
\ref{spaces-perfect-lemma-pseudo-coherent}.

\begin{example}
\label{example-relatively-perfect-field}
Soit $k$ un corps. Soit $X$ un espace algébrique de présentation finie
sur $k$ (en particulier, $X$ est quasi-compact). Alors un objet $E$ de
$D(\mathcal{O}_X)$ est $k$-parfait si et seulement s'il est borné et
pseudo-cohérent (par définition), c'est-à-dire si et seulement s'il appartient à
$D^b_{\textit{Coh}}(X)$ (d'après
Catégories dérivées d'espaces, lemme
\ref{spaces-perfect-lemma-identify-pseudo-coherent-noetherian}).
Ainsi, être relativement parfait ne signifie {\bf pas} ``être parfait sur les fibres''.
\end{example}

\noindent
La notion algébrique correspondante est étudiée dans
Compléments d'algèbre, section \ref{more-algebra-section-relatively-perfect}.
Nous pouvons relier la notion pour les espaces algébriques à la
notion algébrique comme suit.

\begin{lemma}
\label{lemma-affine-locally-rel-perfect}
Soit $S$ un schéma. Soit $f : X \to Y$ un morphisme d'espaces algébriques
sur $S$ qui soit plat et localement de présentation finie.
Soit $E \in D_\QCoh(\mathcal{O}_X)$. Les conditions suivantes sont équivalentes :
\begin{enumerate}
\item $E$ est $Y$-parfait,
\item pour tout diagramme commutatif
$$
\xymatrix{
U \ar[d] \ar[r]_g & V \ar[d] \\
X \ar[r]^f & Y
}
$$
où $U$ et $V$ sont des schémas et où les flèches verticales sont \'etales, le complexe
$E|_U$ est $V$-parfait au sens de Catégories dérivées de schémas,
définition \ref{perfect-definition-relatively-perfect},
\item pour un diagramme commutatif comme en (2) tel que $U \to X$ soit
surjectif, le complexe $E|_U$ est $V$-parfait au sens de
Catégories dérivées de schémas,
définition \ref{perfect-definition-relatively-perfect},
\item pour tout diagramme commutatif comme en (2) avec $U$ et $V$
affines, le complexe $R\Gamma(U, E)$ est $\mathcal{O}_Y(V)$-parfait.
\end{enumerate}
\end{lemma}

\begin{proof}
Pour donner un sens aux assertions (2), (3), (4) du lemme, remarquons que
l'objet $E|_U$ de $D_\QCoh(\mathcal{O}_U)$ correspond à
un objet $E_0$ de $D_\QCoh(\mathcal{O}_{U_0})$, où $U_0$
désigne le schéma sous-jacent à $U$ ; voir Catégories dérivées d'espaces,
lemme \ref{spaces-perfect-lemma-derived-quasi-coherent-small-etale-site}.
De plus, dans ce cas, $E_0$ est pseudo-cohérent si et seulement si
$E|_U$ est pseudo-cohérent ; voir Catégories dérivées d'espaces,
lemme \ref{spaces-perfect-lemma-descend-pseudo-coherent}.
De même, $E|_U$ est localement de dimension de Tor finie sur
$f^{-1}\mathcal{O}_Y|_U = g^{-1}\mathcal{O}_V$ si et seulement si
$E_0$ est localement de dimension de Tor finie sur $g_0^{-1}\mathcal{O}_{V_0}$
d'après Catégories dérivées d'espaces, lemme
\ref{spaces-perfect-lemma-tor-dimension-rel}.
Ici, $g_0 : U_0 \to V_0$ est le morphisme de schémas représentant
$g : U \to V$ (notations de Catégories dérivées d'espaces,
remarque \ref{spaces-perfect-remark-match-total-direct-images}).
Enfin, remarquons que ``être pseudo-cohérent'' est une propriété locale pour la topologie \'etale et
que, bien entendu, ``avoir localement une dimension de Tor finie'' est une propriété locale pour la topologie \'etale.
Nous voyons ainsi qu'il suffit de vérifier localement pour la topologie \'etale que l'objet est $Y$-parfait
et, d'après la discussion ci-dessus, que (1) implique (2) et que
(3) implique (1). Comme l'assertion (4) est équivalente
à (2) et (3) d'après
Catégories dérivées de schémas, lemme
\ref{perfect-lemma-affine-locally-rel-perfect},
la démonstration est achevée.
\end{proof}

\begin{lemma}
\label{lemma-triangulated}
Soit $S$ un schéma. Soit $f : X \to Y$ un morphisme
d'espaces algébriques sur $S$ qui soit plat et localement de présentation finie.
La sous-catégorie pleine de $D(\mathcal{O}_X)$ formée des objets $Y$-parfaits est
une sous-catégorie triangulée saturée\footnote{Catégories dérivées, définition
\ref{derived-definition-saturated}.}.
\end{lemma}

\begin{proof}
Cela résulte des lemmes de Cohomologie sur les sites
\ref{sites-cohomology-lemma-cone-pseudo-coherent},
\ref{sites-cohomology-lemma-summands-pseudo-coherent},
\ref{sites-cohomology-lemma-cone-tor-amplitude} et
\ref{sites-cohomology-lemma-summands-tor-amplitude}.
\end{proof}

\begin{lemma}
\label{lemma-perfect-relatively-perfect}
Soit $S$ un schéma.
Soit $f : X \to Y$ un morphisme d'espaces algébriques sur $S$
qui soit plat et localement de présentation finie.
Tout objet parfait de $D(\mathcal{O}_X)$ est $Y$-parfait.
Si $K, M \in D(\mathcal{O}_X)$, alors $K \otimes_{\mathcal{O}_X}^\mathbf{L} M$
est $Y$-parfait si $K$ est parfait et si $M$ est $Y$-parfait.
\end{lemma}

\begin{proof}
Réduisons-nous au cas des schémas à l'aide du
lemme \ref{lemma-affine-locally-rel-perfect},
puis appliquons
Catégories dérivées de schémas, lemme
\ref{perfect-lemma-perfect-relatively-perfect}.
\end{proof}

\begin{lemma}
\label{lemma-base-change-relatively-perfect}
Soit $S$ un schéma.
Soit $f : X \to Y$ un morphisme d'espaces algébriques sur $S$
qui soit plat et localement de présentation finie.
Soit $g : Y' \to Y$ un morphisme d'espaces algébriques sur $S$.
Posons $X' = Y' \times_Y X$ et notons $g' : X' \to X$ la projection.
Si $K \in D(\mathcal{O}_X)$ est $Y$-parfait, alors $L(g')^*K$
est $Y'$-parfait.
\end{lemma}

\begin{proof}
Réduisons-nous au cas des schémas à l'aide du
lemme \ref{lemma-affine-locally-rel-perfect},
puis appliquons
Catégories dérivées de schémas, lemme
\ref{perfect-lemma-base-change-relatively-perfect}.
\end{proof}

\begin{situation}
\label{situation-relative-descent}
Soit $S$ un schéma.
Soit $Y = \lim_{i \in I} Y_i$ la limite projective d'un système filtrant
d'espaces algébriques sur $S$
à morphismes de transition affines $g_{i'i} : Y_{i'} \to Y_i$.
Nous supposons que $Y_i$ est quasi-compact et quasi-séparé pour tout $i \in I$.
Notons $g_i : Y \to Y_i$ la projection. Fixons un élément $0 \in I$
et un morphisme plat de présentation finie $X_0 \to Y_0$.
Posons $X_i = Y_i \times_{Y_0} X_0$ et $X = Y \times_{Y_0} X_0$
et notons les morphismes de transition $f_{i'i} : X_{i'} \to X_i$
et les projections $f_i : X \to X_i$.
\end{situation}

\begin{lemma}
\label{lemma-relative-descend-homomorphisms}
Plaçons-nous dans la Situation \ref{situation-relative-descent}.
Soient $K_0$ et $L_0$ des objets de $D(\mathcal{O}_{X_0})$.
Posons $K_i = Lf_{i0}^*K_0$ et $L_i = Lf_{i0}^*L_0$ pour $i \geq 0$
et posons $K = Lf_0^*K_0$ et $L = Lf_0^*L_0$. Alors l'application
$$
\colim_{i \geq 0} \Hom_{D(\mathcal{O}_{X_i})}(K_i, L_i)
\longrightarrow
\Hom_{D(\mathcal{O}_X)}(K, L)
$$
est un isomorphisme si $K_0$ est pseudo-cohérent et si
$L_0 \in D_\QCoh(\mathcal{O}_{X_0})$ est (localement) de
dimension de Tor finie en tant qu'objet de
$D((X_0 \to Y_0)^{-1}\mathcal{O}_{Y_0})$.
\end{lemma}

\begin{proof}
Pour tout objet quasi-compact et quasi-séparé $U_0$ de
$(X_0)_{spaces, \etale}$, considérons la condition $P$ selon laquelle
$$
\colim_{i \geq 0} \Hom_{D(\mathcal{O}_{U_i})}(K_i|_{U_i}, L_i|_{U_i})
\longrightarrow
\Hom_{D(\mathcal{O}_U)}(K|_U, L|_U)
$$
est un isomorphisme, où $U = X \times_{X_0} U_0$ et
$U_i = X_i \times_{X_0} U_0$. Nous allons démontrer que $P$ vaut pour tout $U_0$.

\medskip\noindent
Supposons que $(U_0 \subset W_0, V_0 \to W_0)$ soit un carré
distingué élémentaire de $(X_0)_{spaces, \etale}$
et que $P$ vaille pour les trois objets $U_0, V_0, U_0 \times_{W_0} V_0$.
Alors $P$ vaut pour $W_0$ par Mayer--Vietoris
pour les Hom dans la catégorie dérivée ; voir Catégories dérivées d'espaces,
lemme \ref{spaces-perfect-lemma-mayer-vietoris-hom}.

\medskip\noindent
Considérons d'abord $U_0 = W_0 \times_{Y_0} X_0$, où $W_0$ est un objet
quasi-compact et quasi-séparé de $(Y_0)_{spaces, \etale}$.
D'après le principe d'induction de Catégories dérivées d'espaces,
lemme \ref{spaces-perfect-lemma-induction-principle},
appliqué à ces $W_0$ et au paragraphe précédent,
nous constatons qu'il suffit de démontrer
$P$ pour $U_0 = W_0 \times_{Y_0} X_0$ avec $W_0$ affine.
Autrement dit, nous nous sommes ramenés au cas où $Y_0$ est affine.
Ensuite, appliquons de nouveau le principe d'induction, cette fois à tous les
ouverts quasi-compacts et quasi-séparés de $X_0$, pour nous ramener au
cas où $X_0$ est également affine.

\medskip\noindent
Si $X_0$ et $Y_0$ sont affines, nous sommes ramenés au cas
des schémas, démontré dans
Catégories dérivées de schémas, lemme
\ref{perfect-lemma-relative-descend-homomorphisms}.
Le lecteur pourra utiliser
Catégories dérivées d'espaces, lemmes
\ref{spaces-perfect-lemma-pseudo-coherent},
\ref{spaces-perfect-lemma-derived-quasi-coherent-small-etale-site},
\ref{spaces-perfect-lemma-descend-pseudo-coherent} et
\ref{spaces-perfect-lemma-tor-dimension-rel}
pour traduire l'énoncé en un énoncé
ne faisant intervenir que des schémas et des catégories dérivées de modules sur des schémas.
\end{proof}

\begin{lemma}
\label{lemma-descend-relatively-perfect}
Dans la Situation \ref{situation-relative-descent}, la catégorie des
objets $Y$-parfaits de $D(\mathcal{O}_X)$ est la limite inductive des catégories
d'objets $Y_i$-parfaits de $D(\mathcal{O}_{X_i})$.
\end{lemma}

\begin{proof}
Pour tout objet quasi-compact et quasi-séparé $U_0$ de
$(X_0)_{spaces, \etale}$, considérons la condition $P$
selon laquelle le foncteur
$$
\colim_{i \geq 0} D_{Y_i\text{-parfait}}(\mathcal{O}_{U_i})
\longrightarrow
D_{Y\text{-parfait}}(\mathcal{O}_U)
$$
est une équivalence, où $U = X \times_{X_0} U_0$ et
$U_i = X_i \times_{X_0} U_0$.
Nous savons déjà que ce foncteur est pleinement fidèle
d'après le lemme \ref{lemma-relative-descend-homomorphisms}. Il suffit donc de montrer
qu'il est essentiellement surjectif.

\medskip\noindent
Supposons que $(U_0 \subset W_0, V_0 \to W_0)$ soit un carré
distingué élémentaire de $(X_0)_{spaces, \etale}$
et que $P$ vaille pour les trois objets $U_0, V_0, U_0 \times_{W_0} V_0$.
Nous affirmons que $P$ vaut pour $W_0$. Nous emploierons les notations
$U_i = X_i \times_{X_0} U_0$, $U = X \times_{X_0} U_0$,
et de même pour $V_0$ et $W_0$. Par abus de notation, nous emploierons le symbole
$f_i$ pour tous les morphismes $U \to U_i$, $V \to V_i$,
$U \times_W V \to U_i \times_{W_i} V_i$ et $W \to W_i$.
Supposons que $E$ soit un objet $Y$-parfait de $D(\mathcal{O}_W)$.
But : montrer que $E$ appartient à l'image essentielle du foncteur.
Par hypothèse,
nous pouvons trouver $i \geq 0$, un objet $Y_i$-parfait $E_{U, i}$ sur $U_i$,
un objet $Y_i$-parfait $E_{V, i}$ sur $V_i$, et
des isomorphismes $Lf_i^*E_{U, i} \to E|_U$ et $Lf_i^*E_{V, i} \to E|_V$.
Soient
$$
a : E_{U, i} \to (Rf_{i, *}E)|_{U_i}
\quad\text{et}\quad
b : E_{V, i} \to (Rf_{i, *}E)|_{V_i}
$$
les morphismes adjoints aux isomorphismes $Lf_i^*E_{U, i} \to E|_U$
et $Lf_i^*E_{V, i} \to E|_V$.
Puisque le foncteur est pleinement fidèle, quitte à augmenter $i$,
nous pouvons trouver un isomorphisme
$c : E_{U, i}|_{U_i \times_{W_i} V_i} \to E_{V, i}|_{U_i \times_{W_i} V_i}$
qui, par image inverse, donne les identifications
$$
Lf_i^*E_{U, i}|_{U \times_W V} \to E|_{U \times_W V} \to
Lf_i^*E_{V, i}|_{U \times_W V}.
$$
Appliquons Catégories dérivées d'espaces, lemme
\ref{spaces-perfect-lemma-glue},
pour obtenir un objet $E_i$ sur $W_i$ et un morphisme $d : E_i \to Rf_{i, *}E$
dont les restrictions sont les morphismes $a$ et $b$ sur $U_i$ et $V_i$.
Il est alors clair que $E_i$ est $Y_i$-parfait (car la propriété d'être
relativement parfait est locale pour la topologie \'etale) et que
$d$ est adjoint à un isomorphisme $Lf_i^*E_i \to E$.

\medskip\noindent
En reprenant exactement l'argument employé dans
la démonstration du lemme \ref{lemma-relative-descend-homomorphisms},
en utilisant le principe d'induction
(Catégories dérivées d'espaces, lemme
\ref{spaces-perfect-lemma-induction-principle}),
nous nous ramenons au cas où $X_0$ et $Y_0$ sont tous deux
affines : travaillons d'abord avec les objets quasi-compacts et quasi-séparés
de $(Y_0)_{spaces, \etale}$ pour nous ramener au cas où
$Y_0$ est affine, puis avec les objets quasi-compacts et quasi-séparés
de $(X_0)_{spaces, \etale}$ pour nous ramener au cas où $X_0$ est affine.
Dans le cas affine, le résultat découle du cas des schémas, à savoir
Catégories dérivées de schémas, lemme
\ref{perfect-lemma-descend-relatively-perfect}.
Le passage au cas des schémas se fait grâce au
lemme \ref{lemma-affine-locally-rel-perfect}.
\end{proof}

\begin{lemma}
\label{lemma-derived-pushforward-rel-perfect}
Soit $S$ un schéma.
Soit $f : X \to Y$ un morphisme d'espaces algébriques sur $S$
qui soit plat, propre et
de présentation finie. Soit $E \in D(\mathcal{O}_X)$ un objet $Y$-parfait.
Alors $Rf_*E$ est un objet parfait de $D(\mathcal{O}_Y)$
et sa formation commute à tout changement de base.
\end{lemma}

\begin{proof}
L'assertion relative au changement de base est Catégories dérivées d'espaces,
lemme \ref{spaces-perfect-lemma-base-change-tensor} (avec
$\mathcal{G}^\bullet$ égal à $\mathcal{O}_X$ placé en degré $0$).
Il suffit donc de montrer que $Rf_*E$ est un objet parfait. Nous allons nous ramener
au cas où $Y$ est affine noethérien par un argument de passage à la limite.

\medskip\noindent
La question est locale sur $Y$ pour la topologie \'etale ; nous pouvons donc supposer $Y$ affine.
Posons $Y = \Spec(R)$. Écrivons $R = \colim R_i$ comme limite inductive filtrante
d'anneaux noethériens $R_i$. D'après Limites d'espaces, lemme
\ref{spaces-limits-lemma-descend-finite-presentation},
il existe un $i$ et un espace algébrique
$X_i$ de présentation finie sur $R_i$
dont le changement de base à $R$ est $X$. D'après
Limites d'espaces, lemmes \ref{spaces-limits-lemma-eventually-proper} et
\ref{spaces-limits-lemma-descend-flat},
on peut supposer que $X_i$ est propre et plat sur $R_i$.
D'après le lemme \ref{lemma-descend-relatively-perfect},
on peut supposer qu'il existe un objet $R_i$-parfait $E_i$ de
$D(\mathcal{O}_{X_i})$ dont l'image inverse sur $X$ est $E$.
En appliquant Catégories dérivées d'espaces,
lemme \ref{spaces-perfect-lemma-perfect-direct-image},
à $X_i \to \Spec(R_i)$ et à $E_i$, et en utilisant la propriété de
changement de base déjà démontrée, on obtient le résultat.
\end{proof}

\begin{lemma}
\label{lemma-compute-ext-rel-perfect}
Soit $S$ un schéma.
Soit $f : X \to Y$ un morphisme d'espaces algébriques sur $S$.
Soient $E, K \in D(\mathcal{O}_X)$.
Supposons que
\begin{enumerate}
\item $Y$ soit quasi-compact et quasi-séparé,
\item $f$ soit propre, plat et de présentation finie,
\item $E$ soit $Y$-parfait,
\item $K$ soit pseudo-cohérent.
\end{enumerate}
Alors il existe un objet pseudo-cohérent $L \in D(\mathcal{O}_Y)$ tel que
$$
Rf_*R\SheafHom(K, E) = R\SheafHom(L, \mathcal{O}_Y)
$$
et il en va de même après tout changement de base : étant donné
$$
\vcenter{
\xymatrix{
X' \ar[r]_{g'} \ar[d]_{f'} &
X \ar[d]^f \\
Y' \ar[r]^g &
Y
}
}
\quad\quad
\begin{matrix}
\text{cartésien, alors on a } \\
Rf'_*R\SheafHom(L(g')^*K, L(g')^*E) \\
= R\SheafHom(Lg^*L, \mathcal{O}_{Y'})
\end{matrix}
$$
\end{lemma}

\begin{proof}
Puisque $Y$ est quasi-compact et quasi-séparé, il en va de même de $X$.
D'après Catégories dérivées d'espaces, lemme
\ref{spaces-perfect-lemma-pseudo-coherent-hocolim}, on peut écrire
$K = \text{hocolim} K_n$, avec $K_n$ parfait et $K_n \to K$ induisant
un isomorphisme sur les troncations $\tau_{\geq -n}$. Soit $K_n^\vee$
le complexe parfait dual
(Cohomologie sur les sites, lemme \ref{sites-cohomology-lemma-dual-perfect-complex}).
On obtient un système projectif $\ldots \to K_3^\vee \to K_2^\vee \to K_1^\vee$
d'objets parfaits. D'après le lemme \ref{lemma-perfect-relatively-perfect},
on voit que $K_n^\vee \otimes_{\mathcal{O}_X} E$ est $Y$-parfait.
On peut donc appliquer le lemme \ref{lemma-derived-pushforward-rel-perfect}
à $K_n^\vee \otimes_{\mathcal{O}_X} E$, et on obtient un système projectif
$$
\ldots \to M_3 \to M_2 \to M_1
$$
de complexes parfaits sur $Y$, avec
$$
M_n = Rf_*(K_n^\vee \otimes_{\mathcal{O}_X}^\mathbf{L} E) =
Rf_*R\SheafHom(K_n, E)
$$
De plus, la formation de ces complexes commute à tout
changement de base, à savoir $Lg^*M_n =
Rf'_*((L(g')^*K_n)^\vee \otimes_{\mathcal{O}_{X'}}^\mathbf{L} L(g')^*E) =
Rf'_*R\SheafHom(L(g')^*K_n, L(g')^*E)$.

\medskip\noindent
Comme $K_n \to K$ induit un isomorphisme sur $\tau_{\geq -n}$, on voit que
$K_n \to K_{n + 1}$ induit un isomorphisme sur $\tau_{\geq -n}$.
Il s'ensuit que $K_{n + 1}^\vee \to K_n^\vee$
induit un isomorphisme sur $\tau_{\leq n}$, puisque
$K_n^\vee = R\SheafHom(K_n, \mathcal{O}_X)$.
Supposons que $E$ soit de Tor-amplitude dans $[a, b]$ en tant que complexe
de $f^{-1}\mathcal{O}_Y$-modules. Il en va alors de même après
tout changement de base, voir
Catégories dérivées d'espaces, lemme
\ref{spaces-perfect-lemma-tor-independence-and-tor-amplitude}.
On trouve que
$K_{n + 1}^\vee \otimes_{\mathcal{O}_X} E \to
K_n^\vee \otimes_{\mathcal{O}_X} E$
induit un isomorphisme sur $\tau_{\leq n + a}$,
et il en va de même après tout changement de base.
En appliquant le foncteur dérivé à droite $Rf_*$,
on conclut que les applications $M_{n + 1} \to M_n$
induisent des isomorphismes sur $\tau_{\leq n + a}$,
et il en va de même après tout changement de base.
Choisissons un triangle distingué
$$
M_{n + 1} \to M_n \to C_n \to M_{n + 1}[1]
$$
Prenons $y \in |Y|$. Choisissons un voisinage \'etale élémentaire
$(U, u) \to (Y, y)$ ; cela est possible d'après
Espaces décents, lemme
\ref{decent-spaces-lemma-decent-space-elementary-etale-neighbourhood}.
Prenons pour $Y'$ le spectre du corps résiduel en $u$.
En prenant l'image inverse, on voit que $C_n|_U \otimes_{\mathcal{O}_U}^\mathbf{L} \kappa(u)$
n'a de cohomologie non nulle qu'en degrés $\geq n + a$. D'après
Compléments d'algèbre, lemme
\ref{more-algebra-lemma-lift-perfect-from-residue-field},
on voit que le complexe parfait $C_n|_U$ est de Tor-amplitude dans
$[n + a, m_n]$ pour un certain entier $m_n$, quitte à rétrécir $U$.
Ainsi, $C_n$ est de Tor-amplitude dans $[n + a, m_n]$ pour un certain entier $m_n$
(car $Y$ est quasi-compact).
En particulier, le complexe parfait dual $C_n^\vee$ est de Tor-amplitude dans
$[-m_n, -n - a]$.

\medskip\noindent
Soit $L_n = M_n^\vee$ le complexe parfait dual. La
conclusion de la discussion de l'alinéa précédent est que
$L_n \to L_{n + 1}$ induit des isomorphismes sur $\tau_{\geq -n - a}$.
Ainsi, $L = \text{hocolim} L_n$ est pseudo-cohérent, voir
Catégories dérivées d'espaces, lemme
\ref{spaces-perfect-lemma-pseudo-coherent-hocolim}.
Puisque l'on a
$$
R\SheafHom(K, E) = R\SheafHom(\text{hocolim} K_n, E) =
R\lim R\SheafHom(K_n, E) = R\lim K_n^\vee \otimes_{\mathcal{O}_X} E
$$
(Cohomologie sur les sites, lemme
\ref{sites-cohomology-lemma-colim-and-lim-of-duals})
et puisque $R\lim$ commute à $Rf_*$, on trouve que
$$
Rf_*R\SheafHom(K, E) = R\lim M_n = R\lim R\SheafHom(L_n, \mathcal{O}_Y) =
R\SheafHom(L, \mathcal{O}_Y)
$$
Ceci démontre la formule sur $Y$. Puisque la construction de $M_n$ est
compatible avec le changement de base, la formule reste vraie après
tout changement de base.
\end{proof}

\begin{remark}
\label{remark-compare-L}
Le lecteur aura peut-être remarqué la similitude entre le
lemme \ref{lemma-compute-ext-rel-perfect} et
Catégories dérivées d'espaces, lemme \ref{spaces-perfect-lemma-compute-ext}.
En effet, le complexe pseudo-cohérent $L$ du
lemme \ref{lemma-compute-ext-rel-perfect}
peut être caractérisé comme l'unique complexe pseudo-cohérent
sur $Y$ tel qu'il existe des isomorphismes fonctoriels
$$
\Ext^i_{\mathcal{O}_Y}(L, \mathcal{F}) \longrightarrow
\Ext^i_{\mathcal{O}_X}(K,
E \otimes_{\mathcal{O}_X}^\mathbf{L} Lf^*\mathcal{F})
$$
compatibles avec les homomorphismes cobord, lorsque $\mathcal{F}$ parcourt
$\QCoh(\mathcal{O}_Y)$. Si cela est un jour nécessaire, on
formulera ici un résultat précis et en donnera une démonstration détaillée.
\end{remark}

\begin{lemma}
\label{lemma-bounded-on-fibres}
Soit $S$ un schéma. Soit $X$ un espace algébrique sur $S$ tel que
le morphisme structural $f : X \to S$ soit plat et
localement de présentation finie. Soit $E$ un objet pseudo-cohérent
de $D(\mathcal{O}_X)$. Les conditions suivantes sont équivalentes :
\begin{enumerate}
\item $E$ est $S$-parfait, et
\item $E$ est localement borné inférieurement et, pour tout point $s \in S$,
l'objet $L(X_s \to X)^*E$ de $D(\mathcal{O}_{X_s})$
est localement borné inférieurement.
\end{enumerate}
\end{lemma}

\begin{proof}
Comme tout est local, on se ramène immédiatement au
cas où $X$ et $S$ sont affines, voir le lemme
\ref{lemma-affine-locally-rel-perfect}.
Ce cas est traité dans
Catégories dérivées de schémas, lemme \ref{perfect-lemma-bounded-on-fibres}.
\end{proof}

\begin{lemma}
\label{lemma-characterize-relatively-perfect}
Soit $A$ un anneau. Soit $X$ un espace algébrique séparé, de
présentation finie et plat sur $A$. Soit $K \in D_\QCoh(\mathcal{O}_X)$.
Si $R \Gamma (X, E \otimes^\mathbf{L} K)$ est parfait dans
$D(A)$ pour tout objet parfait $E \in D(\mathcal{O}_X)$, alors $K$ est
$\Spec(A)$-parfait.
\end{lemma}

\begin{proof}
D'après le lemme \ref{lemma-characterize-pseudo-coh-improved},
$K$ est pseudo-cohérent relativement à $A$. D'après le lemme
\ref{lemma-relative-pseudo-coherent-is-moot},
$K$ est pseudo-cohérent dans $D(\mathcal{O}_X)$.
D'après Catégories dérivées d'espaces, proposition
\ref{spaces-perfect-proposition-detecting-bounded-below},
on voit que $K$ appartient à $D^-(\mathcal{O}_X)$.
Soit $\mathfrak{p}$ un idéal premier de $A$, et notons
$i : Y \to X$ l'inclusion de la fibre schématique
au-dessus de $\mathfrak{p}$, c'est-à-dire que $Y$ est un schéma sur $\kappa(\mathfrak p)$.
D'après le lemme \ref{lemma-bounded-on-fibres},
il suffit de montrer que $Li^*(K)$ est borné inférieurement.
Soit $G \in D_{perf} (\mathcal{O}_X)$ un complexe
parfait qui engendre $D_\QCoh (\mathcal{O}_X)$,
voir Catégories dérivées d'espaces, théorème
\ref{spaces-perfect-theorem-bondal-van-den-Bergh}.
On a
\begin{align*}
R\Hom _{\mathcal{O}_Y}(Li^*(G), Li^*(K))
& =
R\Gamma(Y, Li^*(G ^\vee \otimes ^\mathbf{L} K)) \\
& =
R\Gamma(X, G^\vee \otimes ^{\mathbf{L}} K)
\otimes^\mathbf{L}_A \kappa(\mathfrak{p})
\end{align*}
La première égalité utilise le fait que $Li^*$ préserve les objets parfaits et les duaux,
ainsi que Cohomologie sur les sites, lemme
\ref{sites-cohomology-lemma-dual-perfect-complex} ; on omet
certains détails. La seconde égalité résulte de
Catégories dérivées d'espaces, lemme
\ref{spaces-perfect-lemma-compare-base-change},
puisque $X$ est plat sur $A$. Il résulte de notre hypothèse qu'il s'agit d'un
objet parfait de $D(\kappa(\mathfrak{p}))$. L'objet
$Li^*(G) \in D_{perf}(\mathcal{O}_Y)$ engendre $D_\QCoh(\mathcal{O}_Y)$ d'après
Catégories dérivées d'espaces, remarque
\ref{spaces-perfect-remark-pullback-generator}.
Par conséquent, Catégories dérivées d'espaces, proposition
\ref{spaces-perfect-proposition-detecting-bounded-below},
implique maintenant que $Li^*(K)$ est borné inférieurement, ce qui achève la démonstration.
\end{proof}














\section{Théorème du cube}
\label{section-theorem-cube}

\noindent
Cette section est l'analogue de Compléments sur les morphismes, section
\ref{more-morphisms-section-theorem-cube}.
Le lemme suivant montre que la diagonale du foncteur de Picard
est représentable par des immersions localement fermées sous
les hypothèses du lemme.

\begin{lemma}
\label{lemma-diagonal-picard-flat-proper}
Soit $S$ un schéma.
Soit $f : X \to Y$ un morphisme plat, propre et de présentation finie
d'espaces algébriques sur $S$.
Soit $\mathcal{E}$ un $\mathcal{O}_X$-module localement libre de type fini.
Pour un morphisme $g : Y' \to Y$, considérons le diagramme de changement de base
$$
\xymatrix{
X' \ar[d]_{f'} \ar[r]_{g'} & X \ar[d]^f \\
Y' \ar[r]^g & Y
}
$$
Supposons que $\mathcal{O}_{Y'} \to f'_*\mathcal{O}_{X'}$ soit un
isomorphisme pour tout $g : Y' \to Y$.
Alors il existe une immersion $j : Z \to Y$ de présentation finie
telle qu'un morphisme $g : Y' \to Y$ se factorise par $Z$ si et seulement s'il
existe un $\mathcal{O}_{Y'}$-module $\mathcal{N}$ localement libre de type fini
tel que $(f')^*\mathcal{N} \cong (g')^*\mathcal{L}$.
\end{lemma}

\begin{proof}
Soit $y : \Spec(k) \to Y$ un point à valeurs dans un corps.
Alors la fibre $X_y$ de $f$ en $y$ est connexe puisque, par hypothèse,
$H^0(X_y, \mathcal{O}_{X_y}) = k$. Ainsi, le rang de
$\mathcal{E}$ est constant sur les fibres. Puisque $f$ est ouvert
(Morphismes d'espaces, lemme \ref{spaces-morphisms-lemma-fppf-open})
et fermé, on conclut qu'il existe une décomposition $Y = \coprod Y_r$
de $Y$ en sous-espaces ouverts et fermés telle que $\mathcal{E}$
soit de rang constant $r$ sur l'image inverse de $Y_r$.
On peut donc supposer que $\mathcal{E}$ est de rang constant $r$.
On notera $\mathcal{E}^\vee = \SheafHom(\mathcal{E}, \mathcal{O}_X)$
le module dual de rang $r$.

\medskip\noindent
D'après le théorème de cohomologie et changement de base (plus précisément, d'après
Catégories dérivées d'espaces, lemme
\ref{spaces-perfect-lemma-flat-proper-perfect-direct-image-general}),
on voit que $E = Rf_*\mathcal{E}$ est un objet parfait de la
catégorie dérivée de $Y$ et que sa formation commute à
tout changement de base. De même pour $E' = Rf_*\mathcal{E}^\vee$.
Comme il n'y a jamais de cohomologie en degrés $< 0$, on voit que
$E$ et $E'$ sont (localement) de Tor-amplitude dans $[0, b]$ pour un certain $b$.
Observons que, pour tout $g : Y' \to Y$, on a
$f'_*((g')^*\mathcal{E}) = H^0(Lg^*E)$ et
$f'_*((g')^*\mathcal{E}^\vee) = H^0(Lg^*E')$.
Soient $j : Z \to Y$ et $j' : Z' \to Y$ les immersions localement
fermées construites dans Catégories dérivées d'espaces, lemme
\ref{spaces-perfect-lemma-locally-closed-where-H0-locally-free},
pour $E$ et $E'$, avec $a = 0$ et $r = r$ ; elles sont caractérisées
par le fait que $H^0(Lj^*E)$ et $H^0((j')^*E')$
sont des modules localement libres de rang $r$ dont la formation commute à l'image inverse.

\medskip\noindent
Soit $g : Y' \to Y$ un morphisme. S'il existe un module $\mathcal{N}$
comme dans le lemme, alors, en utilisant la formule de projection
Cohomologie sur les sites, lemme \ref{sites-cohomology-lemma-projection-formula},
on voit que les modules
$$
f'_*((g')^*\mathcal{E}) \cong
f'_*((f')^*\mathcal{N}) \cong
\mathcal{N} \otimes_{\mathcal{O}_{Y'}} f'_*\mathcal{O}_{X'} \cong
\mathcal{N}\quad\text{et de même }\quad
f'_*((g')^*\mathcal{E}^\vee) \cong \mathcal{N}^\vee
$$
sont localement libres de rang $r$ et restent localement libres de rang $r$
après tout changement de base ultérieur $Y'' \to Y'$.
Ainsi, dans ce cas, $g : Y' \to Y$ se factorise par $j$ et par $j'$.
On peut donc remplacer $Y$ par $Z \times_Y Z'$ et supposer que
$f_*\mathcal{E}$ et $f_*\mathcal{E}^\vee$ sont des
$\mathcal{O}_Y$-modules localement libres de rang $r$
dont la formation commute à tout changement de base.

\medskip\noindent
Dans cette situation, si $g : Y' \to Y$ est un morphisme et s'il existe un
module $\mathcal{N}$ comme dans le lemme, alors l'application (le cup-produit en degré $0$)
$$
f'_*((g')^*\mathcal{E})
\otimes_{\mathcal{O}_{Y'}}
f'_*((g')^*\mathcal{E}^\vee)
\longrightarrow \mathcal{O}_{Y'}
$$
est un accouplement parfait. Réciproquement, si cette application de cup-produit est un
accouplement parfait, on voit que, localement sur $Y'$, il existe
une base de sections
$\sigma_1, \ldots, \sigma_r$ dans $f'_*((g')^*\mathcal{L})$ et
$\tau_1, \ldots, \tau_r$ dans
$f'_*((g')^*\mathcal{E}^\vee)$ dont les produits vérifient
$\sigma_i \tau_j = \delta_{ij}$. En considérant $\sigma_i$
comme une section de $(g')^*\mathcal{L}$ sur $X'$
et $\tau_j$ comme une section de $(g')^*\mathcal{L}^\vee$ sur $X'$,
on conclut que
$$
\sigma_1, \ldots, \sigma_r :
\mathcal{O}_{X'}^{\oplus r}
\longrightarrow
(g')^*\mathcal{E}
$$
est un isomorphisme dont l'inverse est donné par
$$
\tau_1, \ldots, \tau_r :
(g')^*\mathcal{E}
\longrightarrow
\mathcal{O}_{X'}^{\oplus r}
$$
Autrement dit, on voit que
$(f')^*f'_*(g')^*\mathcal{E} \cong (g')^*\mathcal{E}$.
Mais la condition que le cup-produit soit non dégénéré
définit un sous-schéma ouvert rétrocompact (à savoir le lieu où un déterminant
convenable est non nul), ce qui achève la démonstration.
\end{proof}









\section{Descente des propriétés de finitude des complexes}
\label{section-descent-finiteness}

\noindent
Cette section est l'analogue de Compléments sur les morphismes,
section \ref{more-morphisms-section-descent-finiteness},
et de
Catégories dérivées de schémas, section
\ref{perfect-section-descent-finiteness}.

\begin{lemma}
\label{lemma-pseudo-coherent-descends-fpqc}
Soit $S$ un schéma. Soit $\{f_i : X_i \to X\}$ un recouvrement fpqc
d'espaces algébriques sur $S$. Soit $E \in D_\QCoh(\mathcal{O}_X)$.
Soit $m \in \mathbf{Z}$. Alors $E$ est $m$-pseudo-cohérent si et seulement si chaque
$Lf_i^*E$ est $m$-pseudo-cohérent.
\end{lemma}

\begin{proof}
L'image inverse préserve toujours la $m$-pseudo-cohérence, voir
Cohomologie sur les sites, lemme
\ref{sites-cohomology-lemma-pseudo-coherent-pullback}.
Il suffit donc de supposer que $Lf_i^*E$ est $m$-pseudo-cohérent
et de démontrer que $E$ est $m$-pseudo-cohérent.
On peut tout d'abord supposer que $X_i$ est un schéma pour tout $i$, voir
Topologies sur les espaces, lemme \ref{spaces-topologies-lemma-refine-fpqc-schemes}.
Ensuite, choisissons un morphisme \'etale surjectif $U \to X$, où $U$ est un schéma.
Alors $U_i = U \times_X X_i$ est un schéma et nous obtenons un recouvrement fpqc
$\{U_i \to U\}$ de schémas, voir
Topologies sur les espaces, Lemme
\ref{spaces-topologies-lemma-recognize-fpqc-covering}.
Nous savons que le résultat est vrai pour
$\{U_i \to U\}_{i \in I}$ d'après le cas des schémas, voir
Catégories dérivées des schémas, Lemme
\ref{perfect-lemma-pseudo-coherent-descends-fpqc}.
D'autre part, la restriction $E|_U$ provient
d'un objet de $D_\QCoh(\mathcal{O}_U)$ (défini au moyen de la topologie de Zariski
et du faisceau structural ``usuel'' de $U$), voir
Catégories dérivées des espaces, Lemme
\ref{spaces-perfect-lemma-derived-quasi-coherent-small-etale-site}.
Le lemme en résulte, puisque les deux notions de pseudo-cohérence
(\'etale et de Zariski) coïncident d'après
Catégories dérivées des espaces,
Lemme \ref{spaces-perfect-lemma-descend-pseudo-coherent}.
\end{proof}

\begin{lemma}
\label{lemma-tor-amplitude-descends-fppf}
Soit $S$ un schéma. Soit $\{g_i : Y_i \to Y\}$ un recouvrement fpqc par des
espaces algébriques sur $S$. Soit $f : X \to Y$ un morphisme d'espaces algébriques
et posons $X_i = Y_i \times_Y X$, avec les projections $f_i : X_i \to Y_i$
et $g'_i : X_i \to X$. Soit $E \in D_\QCoh(\mathcal{O}_X)$.
Soient $a, b \in \mathbf{Z}$. Alors les conditions suivantes sont équivalentes :
\begin{enumerate}
\item $E$ est d'amplitude de Tor dans $[a, b]$ en tant qu'objet de
$D(f^{-1}\mathcal{O}_Y)$, et
\item $L(g'_i)^*E$ est d'amplitude de Tor dans $[a, b]$ en tant qu'objet de
$D(f_i^{-1}\mathcal{O}_{Y_i})$ pour tout $i$.
\end{enumerate}
L'énoncé reste vrai si ``d'amplitude de Tor dans $[a, b]$'' est remplacé par
``de dimension de Tor localement finie''.
\end{lemma}

\begin{proof}
L'image réciproque préserve la propriété ``être d'amplitude de Tor dans $[a, b]$'' d'après
Catégories dérivées des espaces, Lemme
\ref{spaces-perfect-lemma-tor-independence-and-tor-amplitude}.
Remarquons que $Y_i$ et $X$ sont Tor-indépendants au-dessus de $Y$
puisque $Y_i \to Y$ est plat. Supposons (2) et démontrons (1).
Nous pouvons calculer la dimension de Tor sur les fibres, voir
Cohomologie sur les sites, Lemme \ref{sites-cohomology-lemma-tor-amplitude-stalk}
et Propriétés des espaces, Théorème
\ref{spaces-properties-theorem-exactness-stalks}.
Soit $\overline{x}$ un point géométrique de $X$. Choisissons
un $i$ et un point géométrique $\overline{x}_i$ de $X_i$ d'image
$\overline{x}$ dans $X$. Alors
$$
(L(g_i')^*E)_{\overline{x}_i} =
E_{\overline{x}}
\otimes_{\mathcal{O}_{X, \overline{x}}}^\mathbf{L}
\mathcal{O}_{X_, \overline{x}_i}
$$
Soient $\overline{y}_i$ dans $Y_i$ et $\overline{y}$ dans $Y$
les images de $\overline{x}_i$ et $\overline{x}$.
Puisque $X$ et $Y_i$ sont Tor-indépendants au-dessus de $Y$, nous pouvons appliquer
Compléments d'algèbre, Lemme \ref{more-algebra-lemma-base-change-comparison}
pour voir que le membre de droite de la formule affichée est égal à
$E_{\overline{x}}
\otimes_{\mathcal{O}_{Y, \overline{y}}}^\mathbf{L}
\mathcal{O}_{Y_i, \overline{y}_i}$
dans $D(\mathcal{O}_{Y_i, \overline{y}_i})$.
Puisque nous avons supposé que son amplitude de Tor est contenue dans
$[a, b]$, nous concluons que l'amplitude de Tor de
$E_{\overline{x}}$ dans $D(\mathcal{O}_{Y, \overline{y}})$
est contenue dans $[a, b]$ d'après Compléments d'algèbre, Lemme
\ref{more-algebra-lemma-flat-descent-tor-amplitude}.
Ainsi (1) en résulte.

\medskip\noindent
Un peu de topologie élémentaire donne aussi le cas
``de dimension de Tor localement finie''.
\end{proof}

\noindent
Les lemmes suivants ne trouvent pas vraiment leur place dans cette section.

\begin{lemma}
\label{lemma-thickening-pseudo-coherent}
Soit $S$ un schéma. Soit $i : X \to X'$ un épaississement d'ordre fini entre
espaces algébriques. Soit $K' \in D(\mathcal{O}_{X'})$ un objet tel que
$K = Li^*K'$ soit pseudo-cohérent. Alors $K'$ est pseudo-cohérent.
\end{lemma}

\begin{proof}
Nous démontrons d'abord que $K'$ a des faisceaux de cohomologie quasi-cohérents ; nous invitons
le lecteur à omettre cette partie.
Pour cela, nous pouvons nous ramener au cas d'un épaississement du premier ordre, voir
section \ref{section-thickenings}. Soit $\mathcal{I} \subset \mathcal{O}_{X'}$
le faisceau quasi-cohérent d'idéaux définissant $X$.
En tensorisant la suite exacte courte
$$
0 \to \mathcal{I} \to \mathcal{O}_{X'} \to i_*\mathcal{O}_X \to 0
$$
par $K'$, nous obtenons un triangle distingué
$$
K' \otimes_{\mathcal{O}_{X'}}^\mathbf{L} \mathcal{I}
\to K' \to
K' \otimes_{\mathcal{O}_{X'}}^\mathbf{L} i_*\mathcal{O}_X
\to
(K' \otimes_{\mathcal{O}_{X'}}^\mathbf{L} \mathcal{I})[1]
$$
Puisque $i_* = Ri_*$ et que nous pouvons considérer $\mathcal{I}$
comme un $\mathcal{O}_X$-module quasi-cohérent (puisque nous avons un épaississement du
premier ordre), nous pouvons récrire ceci sous la forme
$$
i_*(K \otimes_{\mathcal{O}_X}^\mathbf{L} \mathcal{I})
\to K' \to
i_*K \to
i_*(K \otimes_{\mathcal{O}_X}^\mathbf{L} \mathcal{I})[1]
$$
Utilisons Cohomologie des espaces, Lemme
\ref{spaces-cohomology-lemma-projection-formula-finite}
pour identifier les termes. Puisque $K$ appartient à
$D_\QCoh(\mathcal{O}_X)$, nous concluons que
$K'$ appartient à $D_\QCoh(\mathcal{O}_{X'})$ ; ceci utilise
Catégories dérivées des espaces, Lemmes
\ref{spaces-perfect-lemma-pseudo-coherent},
\ref{spaces-perfect-lemma-quasi-coherence-tensor-product}, et
\ref{spaces-perfect-lemma-quasi-coherence-direct-image}.

\medskip\noindent
Supposons que $K'$ appartienne à $D_\QCoh(\mathcal{O}_{X'})$.
La question est locale pour la topologie \'etale sur $X'$ ;
nous pouvons donc supposer que $X'$ est affine.
Dans ce cas, le résultat découle du cas des schémas
(Compléments sur les morphismes, Lemme
\ref{more-morphisms-lemma-thickening-pseudo-coherent}).
Le passage au langage des schémas utilise
Catégories dérivées des espaces, Lemmes
\ref{spaces-perfect-lemma-derived-quasi-coherent-small-etale-site} et
\ref{spaces-perfect-lemma-descend-pseudo-coherent} et
Remarque \ref{spaces-perfect-remark-match-total-direct-images}.
\end{proof}

\begin{lemma}
\label{lemma-thickening-relatively-perfect}
Soit $S$ un schéma. Considérons un diagramme cartésien
$$
\xymatrix{
X \ar[r]_i \ar[d]_f & X' \ar[d]^{f'} \\
Y \ar[r]^j & Y'
}
$$
d'espaces algébriques sur $S$. Supposons que $X' \to Y'$ soit plat et localement
de présentation finie, et que $Y \to Y'$ soit un épaississement d'ordre fini.
Soit $E' \in D(\mathcal{O}_{X'})$. Si $E = Li^*(E')$ est $Y$-parfait,
alors $E'$ est $Y'$-parfait.
\end{lemma}

\begin{proof}
Rappelons qu'être $Y$-parfait signifie, pour $E$, que $E$ est
pseudo-cohérent et localement de dimension de Tor finie en tant que complexe
de $f^{-1}\mathcal{O}_Y$-modules
(Définition \ref{definition-relatively-perfect}).
D'après le Lemme \ref{lemma-thickening-pseudo-coherent},
nous voyons que $E'$ est pseudo-cohérent.
En particulier, $E'$ appartient à $D_\QCoh(\mathcal{O}_{X'})$, voir
Catégories dérivées des espaces, Lemme
\ref{spaces-perfect-lemma-pseudo-coherent}.
Le Lemme \ref{lemma-affine-locally-rel-perfect}
nous ramène au cas des schémas.
Le cas des schémas est
Compléments sur les morphismes, Lemme
\ref{more-morphisms-lemma-thickening-relatively-perfect}.
\end{proof}

\begin{lemma}
\label{lemma-henselian-relatively-perfect}
Soit $(R, I)$ un couple formé d'un anneau et d'un idéal $I$
contenu dans le radical de Jacobson. Posons $S = \Spec(R)$ et $S_0 = \Spec(R/I)$.
Soit $X$ un espace algébrique sur $R$ dont le morphisme structural
$f : X \to S$ est propre, plat et de présentation finie.
Notons $X_0 = S_0 \times_S X$. Soit $E \in D(\mathcal{O}_X)$
pseudo-cohérent. Si la restriction dérivée $E_0$ de $E$
à $X_0$ est $S_0$-parfaite, alors $E$ est $S$-parfait.
\end{lemma}

\begin{proof}
Choisissons un morphisme \'etale surjectif $U \to X$ avec $U$ affine.
Choisissons une immersion fermée $U \to \mathbf{A}^d_S$.
Posons $U_0 = S_0 \times_S U$.
Le complexe $E_0|_{U_0}$ est d'amplitude de Tor
dans $[a, b]$ pour certains $a, b \in \mathbf{Z}$.
Soit $\overline{x}$ un point géométrique de $X$.
Nous allons montrer que l'amplitude de Tor de
$E_{\overline{x}}$ sur $R$ est contenue dans $[a - d, b]$.
Cela achèvera la démonstration, puisque l'amplitude de Tor peut se
lire sur les fibres d'après
Cohomologie sur les sites, Lemme \ref{sites-cohomology-lemma-tor-amplitude-stalk}
et Propriétés des espaces, Théorème
\ref{spaces-properties-theorem-exactness-stalks}.

\medskip\noindent
Soit $x \in |X|$ le point déterminé par $\overline{x}$.
Rappelons que $|X| \to |S|$ est fermé (par définition des morphismes propres).
Puisque $I$ est contenu dans le radical de Jacobson, toute partie fermée non vide
de $S$ contient un point du sous-schéma fermé $S_0$.
Nous pouvons donc trouver une spécialisation $x \leadsto x_0$ dans $|X|$
avec $x_0 \in |X_0|$. Choisissons $u_0 \in U_0$ s'envoyant sur $x_0$.
D'après Espaces décents, Lemme \ref{decent-spaces-lemma-generalizations-lift-flat}
(ou Espaces décents, Lemme \ref{decent-spaces-lemma-specialization},
qui s'applique directement aux morphismes \'etales),
nous trouvons une spécialisation $u \leadsto u_0$ dans $U$
telle que $u$ s'envoie sur $x$. Nous pouvons relever $\overline{x}$
en un point géométrique $\overline{u}$ de $U$ situé au-dessus de $u$.
Alors $E_{\overline{x}} = (E|_U)_{\overline{u}}$.

\medskip\noindent
Écrivons $U = \Spec(A)$. Alors $A$ est plate et de présentation finie
comme $R$-algèbre, et c'est un quotient d'une $R$-algèbre polynomiale en
$d$ variables. La restriction $E|_U$ correspond
(d'après Catégories dérivées des espaces, Lemmes
\ref{spaces-perfect-lemma-pseudo-coherent},
\ref{spaces-perfect-lemma-derived-quasi-coherent-small-etale-site}, et
\ref{spaces-perfect-lemma-descend-pseudo-coherent},
et
Catégories dérivées des schémas, Lemme
\ref{perfect-lemma-affine-compare-bounded} et
\ref{perfect-lemma-pseudo-coherent-affine})
à un objet pseudo-cohérent $K$ de $D(A)$.
Remarquons que $E_0$ correspond à $K \otimes_A^\mathbf{L} A/IA$.
Soient $\mathfrak q \subset \mathfrak q_0 \subset A$ les idéaux premiers
correspondant à $u \leadsto u_0$.
Alors
$$
E_{\overline{x}} =
(E|_U)_{\overline{u}} =
E_u \otimes_{\mathcal{O}_{U, u}}^\mathbf{L} \mathcal{O}_{U, \overline{u}} =
K_{\mathfrak q} \otimes_{A_\mathfrak q}^\mathbf{L} A_{\mathfrak q}^{sh}
$$
(certains détails sont omis). Puisque $A_\mathfrak q \to A_\mathfrak q^{sh}$
est plat, l'amplitude de Tor de cet objet comme $R$-module est la même que
l'amplitude de Tor de $K_\mathfrak q$ comme $R$-module
(Compléments d'algèbre, Lemme \ref{more-algebra-lemma-no-change-tor-amplitude}).
De plus, $K_{\mathfrak q}$ est un localisé de $K_{\mathfrak q_0}$.
Il suffit donc de montrer que $K_{\mathfrak q_0}$ est d'amplitude de Tor dans
$[a - d, b]$ en tant que complexe de $R$-modules.

\medskip\noindent
Soit $I \subset \mathfrak p_0 \subset R$ l'idéal premier
correspondant à $f(x_0)$. Alors nous avons
\begin{align*}
K \otimes_R^\mathbf{L} \kappa(\mathfrak p_0)
& =
(K \otimes_R^\mathbf{L} R/I) \otimes_{R/I}^\mathbf{L}
\kappa(\mathfrak p_0) \\
& =
(K \otimes_A^\mathbf{L} A/IA) \otimes_{R/I}^\mathbf{L} \kappa(\mathfrak p_0)
\end{align*}
la seconde égalité résultant de ce que $R \to A$ est plat.
Par notre choix de $a, b$, ce complexe n'a de cohomologie
que dans les degrés de l'intervalle $[a, b]$.
Nous pouvons donc finalement appliquer
Compléments d'algèbre, Lemme
\ref{more-algebra-lemma-lift-from-fibre-relatively-perfect}
à $R \to A$, $\mathfrak q_0$, $\mathfrak p_0$ et $K$
pour conclure.
\end{proof}







\section{Familles de courbes nodales}
\label{section-families-nodal}

\noindent
Cette section prolonge
Courbes algébriques, section \ref{curves-section-families-nodal}.
Nous renvoyons aussi à cette section pour notre choix
de terminologie.

\medskip\noindent
La propriété ``à singularités au plus nodales de dimension relative $1$''
des morphismes de schémas est locale pour la topologie \'etale sur la source et le but, voir
Descente, Lemme \ref{descent-lemma-etale-local-source-target}
et
Courbes algébriques, Lemmes \ref{curves-lemma-nodal-family-etale-local-source},
\ref{curves-lemma-nodal-family-fpqc-local-target}, et
\ref{curves-lemma-nodal-family-postcompose-etale}.
Elle est aussi stable par changement de base et locale pour la topologie fpqc sur le but, voir
Courbes algébriques, Lemmes \ref{curves-lemma-base-change-nodal-family} et
\ref{curves-lemma-nodal-family-fpqc-local-target}. Par conséquent, d'après
Morphismes d'espaces, Lemme \ref{spaces-morphisms-lemma-local-source-target},
nous pouvons définir comme suit la notion de morphisme à singularités au plus nodales de dimension relative
$1$ pour les espaces algébriques ; elle coïncide avec la
notion déjà existante définie dans
Morphismes d'espaces, section \ref{spaces-morphisms-section-representable},
lorsque le morphisme est représentable.

\begin{definition}
\label{definition-nodal-family}
Soit $S$ un schéma.
Soit $f : X \to Y$ un morphisme d'espaces algébriques sur $S$.
Nous disons que $f$ est {\it à singularités au plus nodales de dimension relative $1$}
si les conditions équivalentes de
Morphismes d'espaces, Lemme \ref{spaces-morphisms-lemma-local-source-target},
sont satisfaites avec $\mathcal{P} =$``à singularités au plus nodales de dimension relative $1$''.
\end{definition}

\begin{lemma}
\label{lemma-base-change-nodal}
La propriété d'être à singularités au plus nodales de dimension relative $1$
est préservée par changement de base.
\end{lemma}

\begin{proof}
Voir
Morphismes d'espaces, Remarque \ref{spaces-morphisms-remark-base-change-P}
et
Courbes algébriques, Lemme \ref{curves-lemma-base-change-nodal-family}.
\end{proof}

\begin{lemma}
\label{lemma-nodal-local}
Soit $S$ un schéma.
Soit $f : X \to Y$ un morphisme d'espaces algébriques sur $S$.
Les conditions suivantes sont équivalentes :
\begin{enumerate}
\item $f$ est à singularités au plus nodales de dimension relative $1$,
\item pour tout schéma $Z$ et tout morphisme $Z \to Y$, le morphisme
$Z \times_Y X \to Z$ est à singularités au plus nodales de dimension relative $1$,
\item pour tout schéma affine $Z$ et tout morphisme
$Z \to Y$, le morphisme $Z \times_Y X \to Z$ est
à singularités au plus nodales de dimension relative $1$,
\item il existe un schéma $V$ et un morphisme \'etale surjectif
$V \to Y$ tels que $V \times_Y X \to V$ soit
à singularités au plus nodales de dimension relative $1$,
\item il existe un schéma $U$ et un morphisme \'etale surjectif
$\varphi : U \to X$ tels que le composé $f \circ \varphi$
soit à singularités au plus nodales de dimension relative $1$,
\item pour tout diagramme commutatif
$$
\xymatrix{
U \ar[d] \ar[r] & V \ar[d] \\
X \ar[r] & Y
}
$$
où $U$, $V$ sont des schémas et les flèches verticales sont \'etales,
la flèche horizontale supérieure est à singularités au plus nodales de dimension relative $1$,
\item il existe un diagramme commutatif
$$
\xymatrix{
U \ar[d] \ar[r] & V \ar[d] \\
X \ar[r] & Y
}
$$
où $U$, $V$ sont des schémas, les flèches verticales sont \'etales et
$U \to X$ est surjectif, tel que la flèche horizontale supérieure soit
à singularités au plus nodales de dimension relative $1$, et
\item il existe des recouvrements de Zariski $Y = \bigcup_{i \in I} Y_i$
et $f^{-1}(Y_i) = \bigcup X_{ij}$ tels que
chaque morphisme $X_{ij} \to Y_i$ soit
à singularités au plus nodales de dimension relative $1$.
\end{enumerate}
\end{lemma}

\begin{proof}
La démonstration est omise.
\end{proof}

\noindent
Le lemme suivant nous dit que nous pouvons vérifier fibre par fibre si un morphisme
est à singularités au plus nodales de dimension relative $1$.

\begin{lemma}
\label{lemma-locus-where-nodal}
Soit $S$ un schéma.
Soit $f : X \to Y$ un morphisme d'espaces algébriques sur $S$,
plat et localement de présentation finie. Il existe alors
un plus grand sous-espace ouvert $X' \subset X$ tel que $f|_{X'} : X' \to Y$
soit à singularités au plus nodales de dimension relative $1$. De plus, la formation
de $X'$ commute à tout changement de base.
\end{lemma}

\begin{proof}
Choisissons un diagramme commutatif
$$
\xymatrix{
U \ar[d] \ar[r]_h & V \ar[d] \\
X \ar[r]^f & Y
}
$$
où $U$, $V$ sont des schémas, les flèches verticales sont \'etales et
$U \to X$ est surjectif. D'après le lemme pour le cas des schémas
(Courbes algébriques, Lemme \ref{curves-lemma-locus-where-nodal}),
nous trouvons un plus grand sous-schéma ouvert $U' \subset U$
tel que $h|_{U'} : U' \to V$ soit à singularités au plus nodales de dimension relative $1$
et que la formation de $U'$ commute au changement de base.
Soit $X' \subset X$ le sous-espace ouvert dont les points correspondent
à l'ouvert $\Im(|U'| \to |X|)$.
Le Lemme \ref{lemma-nodal-local} montre que $X' \to Y$ est
à singularités au plus nodales de dimension relative $1$ et que $X'$ est le
plus grand sous-espace ouvert possédant cette propriété (ceci implique aussi
que $U'$ est l'image réciproque de $X'$ dans $U$, mais nous n'en avons
pas besoin). Puisque la même assertion vaut après changement de base,
la démonstration est achevée.
\end{proof}




\section{La propriété de résolution}
\label{section-resolution-property}

\noindent
Nous poursuivons la discussion de Catégories dérivées des espaces, section
\ref{spaces-perfect-section-resolution-property}.

\begin{situation}
\label{situation-descend-resolution-property}
Soit $S$ un schéma. Soit $X$ un espace algébrique quasi-compact et quasi-séparé
sur $S$. Soit $V \to X$ un morphisme \'etale surjectif,
où $V$ est un schéma affine (un tel morphisme existe d'après
Propriétés des espaces, Lemme
\ref{spaces-properties-lemma-quasi-compact-affine-cover}).
Choisissons un diagramme commutatif
$$
\xymatrix{
V \ar[rd]_\varphi \ar[rr]_j & & Y \ar[ld]^\pi \\
& X
}
$$
où $j$ est une immersion ouverte et $\pi$ un morphisme fini
d'espaces algébriques (un tel diagramme existe d'après
le Lemme \ref{lemma-quasi-finite-separated-pass-through-finite}).
Soit $\mathcal{I} \subset \mathcal{O}_Y$ un faisceau quasi-cohérent d'idéaux
de type fini sur $Y$ tel que
$V(\mathcal{I}) = Y \setminus j(V)$ (un tel faisceau d'idéaux existe
d'après Limites d'espaces, Lemme
\ref{spaces-limits-lemma-quasi-coherent-finite-type-ideals}).
\end{situation}

\begin{lemma}
\label{lemma-surjection-in-Noetherian-case}
Dans la Situation \ref{situation-descend-resolution-property}, supposons que
$X$ soit noethérien. Alors, pour tout $\mathcal{O}_X$-module cohérent
$\mathcal{F}$, il existe un entier $r \geq 0$,
des entiers $n_1, \ldots, n_r \geq 0$ et une surjection
$$
\bigoplus\nolimits_{i = 1, \ldots, r} \pi_*(\mathcal{I}^{n_i})
\longrightarrow \mathcal{F}
$$
de $\mathcal{O}_X$-modules.
\end{lemma}

\begin{proof}
Notons $\omega_{Y/X}$ le $\mathcal{O}_Y$-module cohérent tel qu'il existe
un isomorphisme
$$
\pi_*\omega_{Y/X} \cong
\SheafHom_{\mathcal{O}_X}(\pi_*\mathcal{O}_Y, \mathcal{O}_X)
$$
de $\pi_*\mathcal{O}_Y$-modules, voir Morphismes d'espaces, lemme
\ref{spaces-morphisms-lemma-affine-equivalence-modules} et
Descente sur les espaces, lemme \ref{spaces-descent-lemma-finite-over-finite-module}.
Le morphisme canonique $\mathcal{O}_X \to \pi_*\mathcal{O}_Y$
induit un morphisme canonique
$$
\text{Tr}_\pi : \pi_*\omega_{Y/X} \longrightarrow \mathcal{O}_X
$$
Comme $V$ est noethérien affine, nous pouvons choisir des sections
$$
s_1, \ldots, s_r \in
\Gamma(V, \pi^*\mathcal{F} \otimes_{\mathcal{O}_Y} \omega_{Y/X})
$$
qui engendrent le module cohérent
$\pi^*\mathcal{F} \otimes_{\mathcal{O}_X} \omega_{Y/X}$
sur $V$. D'après
Cohomologie des espaces, lemme \ref{spaces-cohomology-lemma-homs-over-open},
nous pouvons choisir des entiers $n_i \geq 0$ tels que $s_i$
se prolonge en un morphisme
$s_i' : \mathcal{I}^{n_i} \to
\pi^*\mathcal{F} \otimes_{\mathcal{O}_Y} \omega_{Y/X}$.
En prenant les images directes sur $X$, nous obtenons des morphismes
$$
\sigma_i :
\pi_*\mathcal{I}^{n_i}
\xrightarrow{\pi_*s'_i}
\pi_*(\pi^*\mathcal{F} \otimes_{\mathcal{O}_Y} \omega_{Y/X}) =
\mathcal{F} \otimes_{\mathcal{O}_X} \pi_*\omega_{Y/X}
\xrightarrow{\text{Tr}_\pi} \mathcal{F}
$$
où l'égalité résulte de
Cohomologie des espaces, lemme \ref{spaces-cohomology-lemma-finite-rings}.
Pour achever la démonstration, nous allons montrer que la somme de ces morphismes est surjective.

\medskip\noindent
Soit $x \in |X|$ un point de $X$. Soit $v \in |V|$ un point d'image $x$.
Nous pouvons choisir un voisinage \'etale $(U, u) \to (X, x)$ tel que
$$
U \times_X Y = W \coprod W'
$$
(somme disjointe d'espaces algébriques), où $W \to U$ est un
isomorphisme et où l'unique point $w \in W$ au-dessus de
$u$ s'envoie sur $v$ dans $V \subset Y$. Cela résulte du
lemme \ref{lemma-etale-splits-off-just-one-quasi-finite-part} et de
Morphismes \'etales, lemme \ref{etale-lemma-etale-etale-local}.
Après avoir encore rétréci $U$ si nécessaire, nous pouvons supposer que la projection envoie $W$ dans
$V \subset Y$. Comme la formation de $\omega_{Y/X}$
commute à la localisation \'etale, nous voyons que
$$
\pi_*\omega_{Y/X}|_U = (\pi|_W)_*\omega_{W/U} \oplus (\pi|_{W'})_*\omega_{W'/U}
$$
Nous avons $(\pi|_W)_*\omega_{W/U} = \mathcal{O}_U$, et cet isomorphisme
est donné par l'application trace $\text{Tr}_\pi|_U$ restreinte
au premier facteur de la décomposition ci-dessus. Comme $W$ s'envoie
dans $V$, nous voyons que $\mathcal{I}^{n_i}|_W = \mathcal{O}_W$.
Par conséquent
$$
\pi_*(\mathcal{I}^{n_i})|_U = \mathcal{O}_U \oplus
(W' \to U)_*(\mathcal{I}^{n_i}|_{W'})
$$
Une chasse au diagramme montre (détails omis) que $\sigma_i|_U$
sur le facteur $\mathcal{O}_U$ est le morphisme $\mathcal{O}_U \to \mathcal{F}$
correspondant à la section
$$
s_i|_W  \in
\Gamma(W,\pi^*\mathcal{F} \otimes_{\mathcal{O}_Y} \omega_{Y/X})
=
\Gamma(W, \mathcal{F}|_W \otimes_{\mathcal{O}_W} \omega_{W/U})
=
\Gamma(U, \mathcal{F})
$$
Comme les sections $s_i$ engendrent le module
$\pi^*\mathcal{F} \otimes_{\mathcal{O}_Y} \omega_{Y/X}$
sur $V$ et comme $W$ s'envoie dans $V$, nous concluons
que la restriction de $\bigoplus \sigma_i$ à $U$ est surjective.
Comme $x$ était arbitraire, la démonstration est achevée.
\end{proof}

\begin{lemma}
\label{lemma-resolution-property-in-Noetherian-case}
Dans la situation \ref{situation-descend-resolution-property}, supposons
$X$ noethérien. Alors $X$ possède la propriété de résolution si et seulement
si $\pi_*\mathcal{I}$ est quotient d'un module localement libre de type fini
sur $\mathcal{O}_X$.
\end{lemma}

\begin{proof}
Le module $\pi_*\mathcal{I}$ est cohérent d'après
Cohomologie des espaces, lemme
\ref{spaces-cohomology-lemma-finite-pushforward-coherent}.
Par conséquent, si $X$ possède la propriété de résolution, alors $\pi_*\mathcal{I}$
est quotient d'un $\mathcal{O}_X$-module localement libre de type fini.
Réciproquement, supposons donnée une surjection $\mathcal{E} \to \pi_*\mathcal{I}$
pour un $\mathcal{O}_X$-module localement libre de type fini $\mathcal{E}$.
Remarquons que, pour tout $n \geq 1$, il existe une surjection
$$
\pi_*\mathcal{I} \otimes_{\mathcal{O}_X} \pi_*\mathcal{I}^n
\longrightarrow
\pi_*\mathcal{I}^{n + 1}
$$
Ainsi, $\mathcal{E}^{\otimes n}$ s'envoie surjectivement sur
$\pi_*\mathcal{I}^n$ pour tout $n \geq 1$. Nous concluons que $X$
possède la propriété de résolution en combinant ceci avec le
résultat du lemme \ref{lemma-surjection-in-Noetherian-case}.
\end{proof}

\begin{lemma}
\label{lemma-resolution-property-one-sheaf}
Dans la situation \ref{situation-descend-resolution-property},
l'espace algébrique $X$ possède la propriété de résolution
si et seulement si $\pi_*\mathcal{I}$ est quotient d'un
$\mathcal{O}_X$-module localement libre de type fini.
\end{lemma}

\begin{proof}
L'image directe $\pi_*\mathcal{G}$ d'un module
quasi-cohérent de type fini sur $\mathcal{O}_Y$, noté $\mathcal{G}$, est un module
quasi-cohérent de type fini sur $\mathcal{O}_X$ d'après
Descente sur les espaces, lemme
\ref{spaces-descent-lemma-finite-over-finite-module}.
En particulier, si $X$ possède la propriété de résolution, alors $\pi_*\mathcal{I}$
est quotient d'un $\mathcal{O}_X$-module localement libre de type fini d'après
Catégories dérivées des espaces, définition
\ref{spaces-perfect-definition-resolution-property}.

\medskip\noindent
Supposons que nous ayons une surjection $\mathcal{E} \to \pi_*\mathcal{I}$
pour un $\mathcal{O}_X$-module localement libre de type fini $\mathcal{E}$.
Dans la suite de la démonstration, nous montrons que $X$ possède la propriété de résolution
en nous ramenant au cas noethérien traité dans le
lemme \ref{lemma-resolution-property-in-Noetherian-case}.
Nous suggérons au lecteur d'omettre la suite de la démonstration.

\medskip\noindent
Une première réduction consiste à considérer $X$ comme un espace algébrique sur
$\Spec(\mathbf{Z})$, voir
Espaces, définition \ref{spaces-definition-base-change}.
(Cela n'affecte ni les conditions ni la conclusion du lemme.)

\medskip\noindent
D'après Limites d'espaces, lemme
\ref{spaces-limits-lemma-finite-in-finite-and-finite-presentation},
nous pouvons écrire $Y = \lim Y_i$, où $Y_i$ est fini et de
présentation finie sur $X$ et où les morphismes de transition sont des
immersions fermées. Considérons le sous-espace fermé $Z = V(\mathcal{I})$ de $Y$.
Comme $\mathcal{I}$ est de type fini, le morphisme $Z \to Y$ est de
présentation finie. Nous pouvons donc trouver un $i$ et un morphisme
$Z_i \to Y_i$ de présentation finie dont le changement de base à $Y$
est $Z \to Y$, voir Limites d'espaces, lemme
\ref{spaces-limits-lemma-descend-finite-presentation}.
Pour $i' \geq i$, notons $Z_{i'} = Z_i \times_{Y_i} Y_{i'}$.
Après avoir augmenté $i$, nous pouvons supposer que $Z_i \to Y_i$ est une
immersion fermée (de présentation finie), voir Limites d'espaces, lemme
\ref{spaces-limits-lemma-descend-closed-immersion}. Notons
$\mathcal{I}_i \subset \mathcal{O}_{Y_i}$ le faisceau d'idéaux
de $Z_i$ et $\pi_i : Y_i \to X$ le morphisme structural.
De même pour $i' \geq i$. Comme $Z = \lim_{i' \geq i} Z_{i'}$, nous avons
$$
\pi_*\mathcal{I} = \colim \pi_{i', *}\mathcal{I}_{i'}
$$
Les morphismes de transition du système sont tous surjectifs,
comme il résulte de la surjectivité des morphismes
$\pi_{i, *}\mathcal{O}_{Y_i} \to \pi_{i', *}\mathcal{O}_{Y_{i'}}$
et du fait que $Z_{i'} = Z_i \times_{Y_i} Y_{i'}$.
D'après Cohomologie des espaces, lemme
\ref{spaces-cohomology-lemma-finite-presentation-quasi-compact-colimit},
pour un certain $i' \geq i$, le morphisme $\mathcal{E} \to \pi_*\mathcal{I}$
se relève en un morphisme $\mathcal{E} \to \pi_{i', *}\mathcal{I}_{i'}$.
Après avoir augmenté $i'$, ce morphisme
$\mathcal{E} \to \pi_{i', *}\mathcal{I}_{i'}$
devient surjectif (sinon la limite inductive des conoyaux,
dont les morphismes de transition sont surjectifs, serait non nulle).
Cela nous ramène au cas examiné au paragraphe suivant.

\medskip\noindent
Supposons que $X$ soit un espace algébrique sur $\mathbf{Z}$ et que
$Y \to X$ soit de présentation finie. Par approximation noethérienne
absolue, nous pouvons écrire $X = \lim X_i$ comme une limite projective filtrante,
où chaque $X_i$ est un espace algébrique quasi-séparé de type fini
sur $\mathbf{Z}$ et les morphismes de transition sont affines, voir
Limites d'espaces, proposition \ref{spaces-limits-proposition-approximate}.
Comme $\pi : Y \to X$ est de présentation finie, nous pouvons trouver un $i$
et un morphisme $\pi_i : Y_i \to X_i$ de présentation finie
dont le changement de base à $X$ est $\pi$, voir
Limites d'espaces, lemme
\ref{spaces-limits-lemma-descend-finite-presentation}.
Après avoir augmenté $i$, nous pouvons supposer que $\pi_i$ est fini, voir
Limites d'espaces, lemme
\ref{spaces-limits-lemma-descend-finite}.
Ensuite, nous pouvons supposer qu'il existe un module localement
libre de type fini sur $\mathcal{O}_{X_i}$, noté $\mathcal{E}_i$,
dont l'image inverse sur $X$ est $\mathcal{E}$, voir
Limites d'espaces, lemme \ref{spaces-limits-lemma-descend-invertible-modules}.
Nous pouvons aussi supposer qu'il existe un morphisme
$\mathcal{E}_i \to \pi_{i, *}\mathcal{O}_{Y_i}$
dont l'image inverse sur $X$ est la composée
$\mathcal{E} \to \pi_*\mathcal{I} \to \pi_*\mathcal{O}_Y$, voir
Limites d'espaces, lemme
\ref{spaces-limits-lemma-descend-modules-finite-presentation}.
Le conoyau
$$
\mathcal{E}_i \to \pi_{i, *}\mathcal{O}_{Y_i} \to \mathcal{Q}_i \to 0
$$
est un $\mathcal{O}_{Y_i}$-module cohérent dont l'image inverse sur $X$
est le conoyau (de présentation finie) $\mathcal{Q}$ du morphisme
$\mathcal{E} \to \pi_*\mathcal{O}_Y$. Autrement dit, nous avons
$\mathcal{Q} = \pi_*(\mathcal{O}_Y/\mathcal{I})$. Considérons
le morphisme
$$
\mathcal{E}_i
\otimes_{\mathcal{O}_{X_i}}
\pi_{i, *} \mathcal{O}_{Y_i}
\longrightarrow
\pi_{i, *}\mathcal{O}_{Y_i}
\otimes_{\mathcal{O}_{X_i}}
\pi_{i, *} \mathcal{O}_{Y_i}
\to
\pi_{i, *} \mathcal{O}_{Y_i}
\to
\mathcal{Q}_i
$$
où la seconde flèche est donnée par la structure d'algèbre
sur $\pi_{i, *}\mathcal{O}_{Y_i}$. L'image inverse de ce morphisme
sur $Y$ est nulle, car l'image de $\mathcal{E} \to \pi_*\mathcal{O}_Y$
est l'idéal $\pi_*\mathcal{I}$. Par conséquent, d'après
Limites d'espaces, lemme
\ref{spaces-limits-lemma-descend-modules-finite-presentation},
après avoir augmenté $i$, nous pouvons supposer que la composée affichée
est nulle. Cela signifie exactement que l'image de
$\mathcal{E}_i \to \pi_{i, *}\mathcal{O}_{Y_i}$
est de la forme $\pi_{i, *}\mathcal{I}_i$
pour un faisceau d'idéaux cohérent $\mathcal{I}_i \subset \mathcal{O}_{Y_i}$.
Comme $\mathcal{E}_i \to \pi_{i, *}\mathcal{O}_{Y_i}$
s'identifie après changement de base à $\mathcal{E} \to \pi_*\mathcal{O}_Y$, nous
voyons que l'image inverse de $\mathcal{I}_i$ sur $Y$ engendre $\mathcal{I}$.
Notons $V_i \subset Y_i$ le sous-espace ouvert dont le complémentaire
est $V(\mathcal{I}_i) \subset Y_i$. Alors $V$ est l'image inverse
de $V_i$ d'après les remarques ci-dessus. Après avoir augmenté $i$, nous pouvons supposer
que $V_i$ est affine et que $\pi_i|_{V_i} : V_i \to X_i$ est 
\'etale, voir
Limites d'espaces, lemmes
\ref{spaces-limits-lemma-limit-is-affine} et
\ref{spaces-limits-lemma-descend-etale}.
Ceci étant établi, nous pouvons appliquer le
lemme \ref{lemma-resolution-property-in-Noetherian-case}
pour conclure que $X_i$ possède la
propriété de résolution. Comme $X \to X_i$ est affine,
nous concluons que $X$ possède lui aussi la propriété de résolution
d'après Catégories dérivées des espaces, lemme
\ref{spaces-perfect-lemma-resolution-property-goes-up-affine}.
\end{proof}

\begin{lemma}
\label{lemma-resolution-property-descends}
Soit $S$ un schéma.
Soit $X = \lim X_i$ la limite d'un système projectif d'espaces algébriques quasi-compacts
et quasi-séparés sur $S$, dont les morphismes de transition sont affines.
Alors $X$ possède la propriété de résolution si et seulement si $X_i$ possède
la propriété de résolution pour un certain $i$.
\end{lemma}

\begin{proof}
Si $X_i$ possède la propriété de résolution, alors $X$ la possède d'après
Catégories dérivées des espaces, lemme
\ref{spaces-perfect-lemma-resolution-property-goes-up-affine}.
Supposons que $X$ possède la propriété de résolution.
Choisissons $i \in I$. Nous pouvons choisir un schéma affine $V_i$ et un morphisme
\'etale surjectif $V_i \to X_i$ 
(Propriétés des espaces, lemme
\ref{spaces-properties-lemma-quasi-compact-affine-cover}).
Nous pouvons choisir une immersion $j : V_i \to Y_i$
où $Y_i$ est fini et de présentation finie sur $X_i$
(lemme \ref{lemma-quasi-finite-separated-pass-through-finite-addendum}).
Nous pouvons choisir un idéal quasi-cohérent de type fini
$\mathcal{I}_i \subset \mathcal{O}_{Y_i}$
tel que $V_i = Y_i \setminus V(\mathcal{I}_i)$
(Limites d'espaces, lemme
\ref{spaces-limits-lemma-quasi-coherent-finite-type-ideals}).
Notons $V \to Y \to X$ les changements de base de $V_i \to Y_i \to X_i$
à $X$. Notons $\mathcal{I} \subset \mathcal{O}_Y$ l'image inverse
de l'idéal $\mathcal{I}_i$. D'après l'implication facile du
lemme \ref{lemma-resolution-property-one-sheaf},
il existe un $\mathcal{O}_X$-module localement libre de type fini
$\mathcal{E}$ et une surjection $\mathcal{E} \to \pi_*\mathcal{I}$.
Remarquons que, puisque $\pi_i : Y_i \to X_i$ est fini et de présentation finie,
le morphisme $\pi : Y \to X$ est lui aussi fini et de présentation finie,
et les $\mathcal{O}_{X_i}$-modules
$\pi_{i, *}\mathcal{O}_{Y_i}$ et
$\pi_{i, *}(\mathcal{O}_{Y_i}/\mathcal{I}_i)$
sont de présentation finie et donnent après changement de base à $X$
$\pi_*\mathcal{O}_Y$ et
$\pi_*(\mathcal{O}_Y/\mathcal{I})$. Ainsi, d'après Limites d'espaces, lemme
\ref{spaces-limits-lemma-descend-modules-finite-presentation},
après avoir augmenté $i$, nous pouvons trouver un module localement libre
de type fini sur $\mathcal{O}_{X_i}$, noté $\mathcal{E}_i$,
et un morphisme $\mathcal{E}_i \to \pi_{i, *}\mathcal{O}_{Y_i}$
dont le changement de base à $X$ redonne la composée
$\mathcal{E} \to \pi_*\mathcal{I} \to \pi_*\mathcal{O}_Y$.
Les images inverses des $\mathcal{O}_{X_i}$-modules de présentation finie
$\Coker(\mathcal{E}_i \to \pi_{i, *}\mathcal{O}_{Y_i})$
et $\pi_{i, *}(\mathcal{O}_{Y_i}/\mathcal{I}_i)$
sur $X$ coïncident comme quotients de $\pi_*\mathcal{O}_Y$.
Ainsi, d'après Limites d'espaces, lemme
\ref{spaces-limits-lemma-descend-modules-finite-presentation},
nous pouvons supposer qu'ils coïncident, autrement dit que
l'image de $\mathcal{E}_i \to \pi_{i, *}\mathcal{O}_{X_i}$
est égale à $\pi_{i, *}\mathcal{I}_i$.
Nous concluons alors que $X_i$ possède la propriété de résolution d'après le
lemme \ref{lemma-resolution-property-one-sheaf}.
\end{proof}

\begin{lemma}
\label{lemma-resolution-property-affine-diagonal}
Soit $S$ un schéma.
Soit $X$ un espace algébrique quasi-compact et quasi-séparé qui possède
la propriété de résolution. Alors $X$ a une diagonale affine sur $\mathbf{Z}$
(au sens de Propriétés des espaces, définition
\ref{spaces-properties-definition-separated}).
\end{lemma}

\begin{proof}
Nous pourrions le démontrer comme dans le cas des schémas, mais nous allons plutôt
déduire le lemme du cas des schémas.
Tout d'abord, nous pouvons et voulons supposer $S = \Spec(\mathbf{Z})$.
Ensuite, choisissons un schéma $Y$ et un morphisme entier surjectif
$f : Y \to X$, voir
Espaces décents, lemme \ref{decent-spaces-lemma-there-is-a-scheme-integral-over}.
Alors $f$ est affine, donc $Y$ possède la propriété de résolution d'après
Catégories dérivées des espaces, lemme
\ref{spaces-perfect-lemma-resolution-property-goes-up-affine}.
Ainsi, d'après le cas des schémas, le schéma $Y$ a une diagonale affine, voir
Catégories dérivées des schémas, lemme
\ref{perfect-lemma-resolution-property-affine-diagonal}.
Considérons ensuite le diagramme commutatif
$$
\xymatrix{
Y \ar[d] \ar[rr]_{\Delta_Y} & & Y \times_{\mathbf{Z}} Y \ar[d] \\
X \ar[rr]^{\Delta_X} & & X \times_{\mathbf{Z}} X
}
$$
Remarquons que la flèche verticale de droite est entière, donc en particulier
affine. Soit $W \to X \times_{\mathbf{Z}} X$ un morphisme avec $W$ affine.
Nous voyons alors que
$$
Y \times_{X \times_{\mathbf{Z}} X} W =
Y \times_{\Delta_Y, Y \times_{\mathbf{Z}} Y}
(Y \times_{\mathbf{Z}} Y) \times_{X \times_{\mathbf{Z}} X} W
$$
est affine. D'autre part, $Y \to X$ est entier et surjectif,
donc
$$
Y \times_{X \times_{\mathbf{Z}} X} W
\longrightarrow
X \times_{X \times_{\mathbf{Z}} X} W
$$
est entier et surjectif comme changement de base de $Y \to X$ à $W$.
Nous concluons que le but de cette flèche est affine d'après
Limites d'espaces, proposition \ref{spaces-limits-proposition-affine}.
Il s'ensuit que $\Delta_X$ est affine, comme souhaité.
\end{proof}







\section{Éclatements et propriété de résolution}
\label{section-resolution-property-by-blowing-up}

\noindent
Nous démontrons que la propriété de résolution est vérifiée après un éclatement.

\begin{lemma}
\label{lemma-blow-up-resolution-property}
Soit $S$ un schéma. Soit $X$ un espace algébrique quasi-compact et quasi-séparé
sur $S$. Supposons que $|X|$ ait un nombre fini de composantes
irréductibles. Il existe un ouvert dense quasi-compact $U \subset X$
et un éclatement $U$-admissible $X' \to X$ tels que l'espace algébrique
$X'$ possède la propriété de résolution.
\end{lemma}

\begin{proof}
D'après Limites d'espaces, lemme
\ref{spaces-limits-lemma-there-is-a-scheme-finite-over-refined},
il existe un morphisme surjectif, fini et de présentation finie
$f : Y \to X$, où $Y$ est un schéma, et un ouvert dense quasi-compact
$U \subset X$ tel que $f^{-1}(U) \to U$ soit fini \'etale.
D'après Compléments sur les morphismes, lemme \ref{more-morphisms-lemma-blow-up-ample-family},
il existe un ouvert dense quasi-compact $V \subset Y$ et un éclatement $V$-admissible
$Y' \to Y$ tel que $Y'$ possède une famille ample de modules
inversibles. Après avoir rétréci $U$, nous pouvons supposer que $f^{-1}(U) \subset V$
(détails omis). Ainsi $f' : Y' \to X$ est fini \'etale au-dessus de $U$
et, en particulier, le morphisme $(f')^{-1}(U) \to U$ est fini localement libre.
D'après le lemme \ref{lemma-finite-after-blowing-up}, il existe un éclatement $U$-admissible
$X' \to X$ tel que le transformé strict $Y''$ de $Y'$
soit fini localement libre sur $X'$. Considérons le diagramme
$$
\xymatrix{
Y'' \ar[d] \ar[r]_g &
Y' \ar[r] &
Y \ar[d] \\
X' \ar[rr] & & X
}
$$
Comme $g : Y'' \to Y'$ est un éclatement dont le centre est contenu
(Diviseurs sur les espaces, lemme \ref{spaces-divisors-lemma-strict-transform})
dans l'image inverse du centre de $X' \to X$, nous voyons que
$g : Y'' \to Y'$ est projectif et qu'il existe un module inversible 
$g$-ample sur $Y''$. Par conséquent, d'après
Compléments sur les morphismes, lemme \ref{more-morphisms-lemma-ample-family-ample-relative},
nous voyons que $Y''$ possède une famille ample de modules inversibles.
Ainsi $Y''$ possède la propriété de résolution, voir
Catégories dérivées des schémas, lemme
\ref{perfect-lemma-resolution-property-ample-family}.
Nous concluons que $X'$ possède la propriété de résolution
d'après
Catégories dérivées des espaces, lemme
\ref{spaces-perfect-lemma-resolution-property-goes-down-finite-flat}.
\end{proof}

\begin{lemma}
\label{lemma-blow-up-resolution-property-filtered}
Soit $S$ un schéma. Soit $X$ un espace algébrique quasi-compact et quasi-séparé
sur $S$. Il existe un $t \geq 0$ et des sous-espaces
fermés
$$
X \supset Z_0 \supset Z_1 \supset \ldots \supset Z_t = \emptyset
$$
tels que $Z_i \to X$ soit de présentation finie,
$Z_0 \subset X$ soit un épaississement et que, pour chaque $i = 0, \ldots t - 1$,
il existe un éclatement $(Z_i \setminus Z_{i - 1})$-admissible
$Z'_i \to Z_i$ tel que $Z'_i$ possède la propriété de résolution.
\end{lemma}

\begin{proof}
Dans ce paragraphe, nous utilisons l'approximation noethérienne absolue pour nous ramener au
cas des espaces algébriques de présentation finie sur $\Spec(\mathbf{Z})$.
Nous pouvons considérer $X$ comme un espace algébrique sur $\Spec(\mathbf{Z})$, voir
Espaces, définition \ref{spaces-definition-base-change} et
Propriétés des espaces, définition \ref{spaces-properties-definition-separated}.
Nous pouvons donc appliquer Limites d'espaces, proposition
\ref{spaces-limits-proposition-approximate}.
Il s'ensuit que nous pouvons trouver un morphisme affine $X \to X_0$
avec $X_0$ de présentation finie sur $\mathbf{Z}$.
Si nous pouvons démontrer le lemme pour $X_0$, nous pouvons alors prendre les images inverses
de la stratification et des centres des éclatements sur $X$
et obtenir le résultat pour $X$; nous utilisons ici le fait que la propriété de résolution
monte le long des morphismes affines (Catégories dérivées des espaces,
lemme \ref{spaces-perfect-lemma-resolution-property-goes-up-affine}) et
que le transformé strict d'un morphisme affine est affine -- détails omis.
Cela nous ramène au cas examiné au paragraphe suivant.

\medskip\noindent
Supposons que $X$ soit de présentation finie sur $\mathbf{Z}$.
Alors $X$ est noethérien et $|X|$ est un espace topologique
noethérien (ayant un nombre fini de composantes irréductibles) de dimension finie.
Nous pouvons donc raisonner par récurrence sur $\dim(|X|)$.
D'après le lemme \ref{lemma-blow-up-resolution-property},
il existe un ouvert dense $U \subset X$ et un éclatement $U$-admissible
$X' \to X$ tel que $X'$ possède la propriété de résolution.
Posons $Z_0 = X$ et soit $Z_1 \subset X$ le sous-espace fermé réduit
tel que $|Z_1| = |X| \setminus |U|$.
Par récurrence, nous trouvons un entier $t \geq 0$ et une filtration
$$
Z_1 \supset Z_{1, 0} \supset Z_{1, 1} \supset \ldots
\supset Z_{1, t} = \emptyset
$$
par des sous-espaces fermés, où $Z_{1, 0} \to Z_1$ est un épaississement
et il existe des éclatements $(Z_{1, i} \setminus Z_{1, i + 1})$-admissibles
$Z'_{1, i} \to Z_{1, i}$ tels que $Z'_{1, i}$
possède la propriété de résolution. Comme $Z_1$ est réduit, nous avons $Z_1 = Z_{1, 0}$.
Nous pouvons donc poser $Z_i = Z_{1, i - 1}$ et $Z'_i = Z'_{1, i - 1}$
pour $i \geq 1$, et le lemme est démontré.
\end{proof}










\begin{multicols}{2}[\section{Autres chapitres}]
\noindent
Préliminaires
\begin{enumerate}
\item \hyperref[introduction-section-phantom]{Introduction}
\item \hyperref[conventions-section-phantom]{Conventions}
\item \hyperref[sets-section-phantom]{Théorie des ensembles}
\item \hyperref[categories-section-phantom]{Catégories}
\item \hyperref[topology-section-phantom]{Topologie}
\item \hyperref[sheaves-section-phantom]{Faisceaux sur les espaces}
\item \hyperref[sites-section-phantom]{Sites et faisceaux}
\item \hyperref[stacks-section-phantom]{Champs}
\item \hyperref[fields-section-phantom]{Corps}
\item \hyperref[algebra-section-phantom]{Algèbre commutative}
\item \hyperref[brauer-section-phantom]{Groupes de Brauer}
\item \hyperref[homology-section-phantom]{Algèbre homologique}
\item \hyperref[derived-section-phantom]{Catégories dérivées}
\item \hyperref[simplicial-section-phantom]{Méthodes simpliciales}
\item \hyperref[more-algebra-section-phantom]{Compléments d'algèbre}
\item \hyperref[smoothing-section-phantom]{Lissage des morphismes d'anneaux}
\item \hyperref[modules-section-phantom]{Faisceaux de modules}
\item \hyperref[sites-modules-section-phantom]{Modules sur les sites}
\item \hyperref[injectives-section-phantom]{Objets injectifs}
\item \hyperref[cohomology-section-phantom]{Cohomologie des faisceaux}
\item \hyperref[sites-cohomology-section-phantom]{Cohomologie sur les sites}
\item \hyperref[dga-section-phantom]{Algèbre différentielle graduée}
\item \hyperref[dpa-section-phantom]{Algèbre à puissances divisées}
\item \hyperref[sdga-section-phantom]{Faisceaux différentiels gradués}
\item \hyperref[hypercovering-section-phantom]{Hyperrecouvrements}
\end{enumerate}
Schémas
\begin{enumerate}
\setcounter{enumi}{25}
\item \hyperref[schemes-section-phantom]{Schémas}
\item \hyperref[constructions-section-phantom]{Constructions de schémas}
\item \hyperref[properties-section-phantom]{Propriétés des schémas}
\item \hyperref[morphisms-section-phantom]{Morphismes de schémas}
\item \hyperref[coherent-section-phantom]{Cohomologie des schémas}
\item \hyperref[divisors-section-phantom]{Diviseurs}
\item \hyperref[limits-section-phantom]{Limites de schémas}
\item \hyperref[varieties-section-phantom]{Variétés}
\item \hyperref[topologies-section-phantom]{Topologies sur les schémas}
\item \hyperref[descent-section-phantom]{Descente}
\item \hyperref[perfect-section-phantom]{Catégories dérivées de schémas}
\item \hyperref[more-morphisms-section-phantom]{Compléments sur les morphismes}
\item \hyperref[flat-section-phantom]{Compléments sur la platitude}
\item \hyperref[groupoids-section-phantom]{Schémas en groupoïdes}
\item \hyperref[more-groupoids-section-phantom]{Compléments sur les schémas en groupoïdes}
\item \hyperref[etale-section-phantom]{Morphismes étales de schémas}
\end{enumerate}
Thèmes de théorie des schémas
\begin{enumerate}
\setcounter{enumi}{41}
\item \hyperref[chow-section-phantom]{Homologie de Chow}
\item \hyperref[intersection-section-phantom]{Théorie de l'intersection}
\item \hyperref[pic-section-phantom]{Schémas de Picard des courbes}
\item \hyperref[weil-section-phantom]{Théories cohomologiques de Weil}
\item \hyperref[adequate-section-phantom]{Modules adéquats}
\item \hyperref[dualizing-section-phantom]{Complexes dualisants}
\item \hyperref[duality-section-phantom]{Dualité pour les schémas}
\item \hyperref[discriminant-section-phantom]{Discriminants et différentes}
\item \hyperref[derham-section-phantom]{Cohomologie de de Rham}
\item \hyperref[local-cohomology-section-phantom]{Cohomologie locale}
\item \hyperref[algebraization-section-phantom]{Géométrie algébrique et formelle}
\item \hyperref[curves-section-phantom]{Courbes algébriques}
\item \hyperref[resolve-section-phantom]{Résolution des surfaces}
\item \hyperref[models-section-phantom]{Réduction semi-stable}
\item \hyperref[functors-section-phantom]{Foncteurs et morphismes}
\item \hyperref[equiv-section-phantom]{Catégories dérivées de variétés}
\item \hyperref[pione-section-phantom]{Groupes fondamentaux de schémas}
\item \hyperref[etale-cohomology-section-phantom]{Cohomologie étale}
\item \hyperref[crystalline-section-phantom]{Cohomologie cristalline}
\item \hyperref[proetale-section-phantom]{Cohomologie pro-étale}
\item \hyperref[relative-cycles-section-phantom]{Cycles relatifs}
\item \hyperref[more-etale-section-phantom]{Compléments de cohomologie étale}
\item \hyperref[trace-section-phantom]{La formule des traces}
\end{enumerate}
Espaces algébriques
\begin{enumerate}
\setcounter{enumi}{64}
\item \hyperref[spaces-section-phantom]{Espaces algébriques}
\item \hyperref[spaces-properties-section-phantom]{Propriétés des espaces algébriques}
\item \hyperref[spaces-morphisms-section-phantom]{Morphismes d'espaces algébriques}
\item \hyperref[decent-spaces-section-phantom]{Espaces algébriques décents}
\item \hyperref[spaces-cohomology-section-phantom]{Cohomologie des espaces algébriques}
\item \hyperref[spaces-limits-section-phantom]{Limites d'espaces algébriques}
\item \hyperref[spaces-divisors-section-phantom]{Diviseurs sur les espaces algébriques}
\item \hyperref[spaces-over-fields-section-phantom]{Espaces algébriques sur des corps}
\item \hyperref[spaces-topologies-section-phantom]{Topologies sur les espaces algébriques}
\item \hyperref[spaces-descent-section-phantom]{Descente et espaces algébriques}
\item \hyperref[spaces-perfect-section-phantom]{Catégories dérivées d'espaces}
\item \hyperref[spaces-more-morphisms-section-phantom]{Compléments sur les morphismes d'espaces}
\item \hyperref[spaces-flat-section-phantom]{Platitude sur les espaces algébriques}
\item \hyperref[spaces-groupoids-section-phantom]{Groupoïdes dans les espaces algébriques}
\item \hyperref[spaces-more-groupoids-section-phantom]{Compléments sur les groupoïdes dans les espaces}
\item \hyperref[bootstrap-section-phantom]{Amorçage}
\item \hyperref[spaces-pushouts-section-phantom]{Sommes amalgamées d'espaces algébriques}
\end{enumerate}
Thèmes de géométrie
\begin{enumerate}
\setcounter{enumi}{81}
\item \hyperref[spaces-chow-section-phantom]{Groupes de Chow des espaces}
\item \hyperref[groupoids-quotients-section-phantom]{Quotients de groupoïdes}
\item \hyperref[spaces-more-cohomology-section-phantom]{Compléments sur la cohomologie des espaces}
\item \hyperref[spaces-simplicial-section-phantom]{Espaces simpliciaux}
\item \hyperref[spaces-duality-section-phantom]{Dualité pour les espaces}
\item \hyperref[formal-spaces-section-phantom]{Espaces algébriques formels}
\item \hyperref[restricted-section-phantom]{Algébrisation des espaces formels}
\item \hyperref[spaces-resolve-section-phantom]{Résolution des surfaces revisitée}
\end{enumerate}
Théorie des déformations
\begin{enumerate}
\setcounter{enumi}{89}
\item \hyperref[formal-defos-section-phantom]{Théorie formelle des déformations}
\item \hyperref[defos-section-phantom]{Théorie des déformations}
\item \hyperref[cotangent-section-phantom]{Le complexe cotangent}
\item \hyperref[examples-defos-section-phantom]{Problèmes de déformation}
\end{enumerate}
Champs algébriques
\begin{enumerate}
\setcounter{enumi}{93}
\item \hyperref[algebraic-section-phantom]{Champs algébriques}
\item \hyperref[examples-stacks-section-phantom]{Exemples de champs}
\item \hyperref[stacks-sheaves-section-phantom]{Faisceaux sur les champs algébriques}
\item \hyperref[criteria-section-phantom]{Critères de représentabilité}
\item \hyperref[artin-section-phantom]{Axiomes d'Artin}
\item \hyperref[quot-section-phantom]{Espaces de Quot et de Hilbert}
\item \hyperref[stacks-properties-section-phantom]{Propriétés des champs algébriques}
\item \hyperref[stacks-morphisms-section-phantom]{Morphismes de champs algébriques}
\item \hyperref[stacks-limits-section-phantom]{Limites de champs algébriques}
\item \hyperref[stacks-cohomology-section-phantom]{Cohomologie des champs algébriques}
\item \hyperref[stacks-perfect-section-phantom]{Catégories dérivées de champs}
\item \hyperref[stacks-introduction-section-phantom]{Introduction aux champs algébriques}
\item \hyperref[stacks-more-morphisms-section-phantom]{Compléments sur les morphismes de champs}
\item \hyperref[stacks-geometry-section-phantom]{La géométrie des champs}
\end{enumerate}
Thèmes de théorie des espaces de modules
\begin{enumerate}
\setcounter{enumi}{107}
\item \hyperref[moduli-section-phantom]{Champs de modules}
\item \hyperref[moduli-curves-section-phantom]{Espaces de modules de courbes}
\end{enumerate}
Divers
\begin{enumerate}
\setcounter{enumi}{109}
\item \hyperref[examples-section-phantom]{Exemples}
\item \hyperref[exercises-section-phantom]{Exercices}
\item \hyperref[guide-section-phantom]{Guide bibliographique}
\item \hyperref[desirables-section-phantom]{Desiderata}
\item \hyperref[coding-section-phantom]{Style de programmation}
\item \hyperref[obsolete-section-phantom]{Obsolète}
\item \hyperref[fdl-section-phantom]{Licence de documentation libre GNU}
\item \hyperref[index-section-phantom]{Index généré automatiquement}
\end{enumerate}
\end{multicols}


\bibliography{my}
\bibliographystyle{amsalpha}

\end{document}
